\documentclass{article}
\usepackage[paperwidth=600bp,paperheight=160bp,margin=0bp]{geometry}
\usepackage{fontspec,amsmath,xcolor,tikz}
\setmainfont{FandolKaiSC-Regular.ttf}[Path=fonts/]
\newfontfamily\scfallback{NotoSerifSC-Regular.ttf}[Path=fonts/]
\newfontfamily\symbolfont{DejaVu Sans}[Path=fonts/,Extension=.ttf,UprightFont=DejaVuSans]
\XeTeXlinebreaklocale "zh"
\XeTeXlinebreakskip=0pt plus 1pt
\pagestyle{empty}\setlength{\parindent}{0pt}
\begin{document}
\definecolor{answer}{rgb}{0.850980,0.047059,0.070588}
\null\begin{tikzpicture}[remember picture,overlay]\begin{scope}[shift={(current page.north west)},x=1bp,y=-1bp]
\node[anchor=north west,inner sep=0pt,text=answer,align=left,text width=234.0bp] at (15,15) {\fontsize{8.5bp}{10.62bp}\selectfont 销售总额等于象棋总价与围棋总价之和。\allowbreak{}\par};
\end{scope}\end{tikzpicture}\newpage
\definecolor{answer}{rgb}{0.850980,0.047059,0.070588}
\null\begin{tikzpicture}[remember picture,overlay]\begin{scope}[shift={(current page.north west)},x=1bp,y=-1bp]
\node[anchor=north west,inner sep=0pt,text=answer,align=left,text width=234.0bp] at (15,15) {\fontsize{8.5bp}{10.62bp}\selectfont 法\ensuremath{2}：\allowbreak{}假设每副都按\ensuremath{30}元卖，\allowbreak{}共有\ensuremath{75}副棋；\allowbreak{}先全部按\ensuremath{30}元计，\allowbreak{}\ensuremath{30}副围棋每副再补\ensuremath{20}元。\allowbreak{}\par};
\end{scope}\end{tikzpicture}\newpage
\definecolor{answer}{rgb}{0.850980,0.047059,0.070588}
\null\begin{tikzpicture}[remember picture,overlay]\begin{scope}[shift={(current page.north west)},x=1bp,y=-1bp]
\node[anchor=north west,inner sep=0pt,text=answer,align=left,text width=234.0bp] at (15,15) {\fontsize{8.5bp}{10.62bp}\selectfont 相同时间里，\allowbreak{}两车路程差等于速度差乘时间。\allowbreak{}\par};
\end{scope}\end{tikzpicture}\newpage
\definecolor{answer}{rgb}{0.850980,0.047059,0.070588}
\null\begin{tikzpicture}[remember picture,overlay]\begin{scope}[shift={(current page.north west)},x=1bp,y=-1bp]
\node[anchor=north west,inner sep=0pt,text=answer,align=left,text width=229.0bp] at (15,15) {\fontsize{8.5bp}{10.62bp}\selectfont 先分别求大、\allowbreak{}小七巧板用纸量，\allowbreak{}再相加。\allowbreak{}\par};
\end{scope}\end{tikzpicture}\newpage
\definecolor{answer}{rgb}{0.850980,0.047059,0.070588}
\null\begin{tikzpicture}[remember picture,overlay]\begin{scope}[shift={(current page.north west)},x=1bp,y=-1bp]
\node[anchor=north west,inner sep=0pt,text=answer,align=left,text width=229.0bp] at (15,15) {\fontsize{8.5bp}{10.62bp}\selectfont 法\ensuremath{1} :\allowbreak{}分别计算大、\allowbreak{}小七巧板用纸量\par};
\end{scope}\end{tikzpicture}\newpage
\definecolor{answer}{rgb}{0.850980,0.047059,0.070588}
\null\begin{tikzpicture}[remember picture,overlay]\begin{scope}[shift={(current page.north west)},x=1bp,y=-1bp]
\node[anchor=north west,inner sep=0pt,text=answer,align=left,text width=229.0bp] at (15,15) {\fontsize{8.5bp}{10.62bp}\selectfont 法\ensuremath{2}：\allowbreak{}假设\ensuremath{15}副全是大的，\allowbreak{}先全部按每副\ensuremath{36}张算；\allowbreak{}\ensuremath{6}副小的每副应少\ensuremath{12}张。\allowbreak{}\par};
\end{scope}\end{tikzpicture}\newpage
\definecolor{answer}{rgb}{0.850980,0.047059,0.070588}
\null\begin{tikzpicture}[remember picture,overlay]\begin{scope}[shift={(current page.north west)},x=1bp,y=-1bp]
\node[anchor=north west,inner sep=0pt,text=answer,align=left,text width=229.0bp] at (15,15) {\fontsize{7.5bp}{9.5bp}\selectfont 析：\allowbreak{}只比较单价：\allowbreak{}帐篷用④，\allowbreak{}桌椅用⑤⑥求每套价，\allowbreak{}勾④⑤⑥。\allowbreak{}\par};
\end{scope}\end{tikzpicture}\newpage
\definecolor{answer}{rgb}{0.850980,0.047059,0.070588}
\null\begin{tikzpicture}[remember picture,overlay]\begin{scope}[shift={(current page.north west)},x=1bp,y=-1bp]
\node[anchor=north west,inner sep=0pt,text=answer,align=left,text width=229.0bp] at (15,15) {\fontsize{8.5bp}{10.62bp}\selectfont 先用桌椅总价除以套数，\allowbreak{}求出每套价格，\allowbreak{}再从帐篷单价中减去。\allowbreak{}\par};
\end{scope}\end{tikzpicture}\newpage
\definecolor{answer}{rgb}{0.850980,0.047059,0.070588}
\null\begin{tikzpicture}[remember picture,overlay]\begin{scope}[shift={(current page.north west)},x=1bp,y=-1bp]
\node[anchor=north west,inner sep=0pt,text=answer,align=left,text width=229.0bp] at (15,15) {\fontsize{8.5bp}{10.62bp}\selectfont 先求篮球与足球的价格和，\allowbreak{}再取\ensuremath{2}倍，\allowbreak{}最后加\ensuremath{2}元。\allowbreak{}\par};
\end{scope}\end{tikzpicture}\newpage
\definecolor{answer}{rgb}{0.850980,0.047059,0.070588}
\null\begin{tikzpicture}[remember picture,overlay]\begin{scope}[shift={(current page.north west)},x=1bp,y=-1bp]
\node[anchor=north west,inner sep=0pt,text=answer,align=left,text width=228.99999237060547bp] at (15,15) {\fontsize{10.0bp}{12.5bp}\selectfont 先算括号内乘、\allowbreak{}除,\allowbreak{} 再算括号外除。\allowbreak{}\par};
\end{scope}\end{tikzpicture}\newpage
\definecolor{answer}{rgb}{0.850980,0.047059,0.070588}
\null\begin{tikzpicture}[remember picture,overlay]\begin{scope}[shift={(current page.north west)},x=1bp,y=-1bp]
\node[anchor=north west,inner sep=0pt,text=answer,align=left,text width=228.99999237060547bp] at (15,15) {\fontsize{10.0bp}{12.5bp}\selectfont 先算两处括号内,\allowbreak{}再算除法。\allowbreak{}\par};
\end{scope}\end{tikzpicture}\newpage
\definecolor{answer}{rgb}{0.850980,0.047059,0.070588}
\null\begin{tikzpicture}[remember picture,overlay]\begin{scope}[shift={(current page.north west)},x=1bp,y=-1bp]
\node[anchor=north west,inner sep=0pt,text=answer,align=left,text width=229.0bp] at (15,15) {\fontsize{8.5bp}{10.62bp}\selectfont 黄芪比板蓝根多\ensuremath{350}千克；\allowbreak{}当归比黄芪的\ensuremath{2}倍少\ensuremath{100}千克。\allowbreak{}\par};
\end{scope}\end{tikzpicture}\newpage
\definecolor{answer}{rgb}{0.850980,0.047059,0.070588}
\null\begin{tikzpicture}[remember picture,overlay]\begin{scope}[shift={(current page.north west)},x=1bp,y=-1bp]
\node[anchor=north west,inner sep=0pt,text=answer,align=left,text width=234.0bp] at (15,15) {\fontsize{8.5bp}{10.62bp}\selectfont 按“单价×数量”分别求模型和徽章总价。\allowbreak{}\par};
\end{scope}\end{tikzpicture}\newpage
\definecolor{answer}{rgb}{0.850980,0.047059,0.070588}
\null\begin{tikzpicture}[remember picture,overlay]\begin{scope}[shift={(current page.north west)},x=1bp,y=-1bp]
\node[anchor=north west,inner sep=0pt,text=answer,align=left,text width=229.0bp] at (15,15) {\fontsize{8.5bp}{10.62bp}\selectfont 当归产量已经求出为\ensuremath{1600}千克，\allowbreak{}灵芝是它的\ensuremath{3}倍。\allowbreak{}\par};
\end{scope}\end{tikzpicture}\newpage
\definecolor{answer}{rgb}{0.850980,0.047059,0.070588}
\null\begin{tikzpicture}[remember picture,overlay]\begin{scope}[shift={(current page.north west)},x=1bp,y=-1bp]
\node[anchor=north west,inner sep=0pt,text=answer,align=left,text width=229.0bp] at (15,15) {\fontsize{8.5bp}{10.62bp}\selectfont 答：\allowbreak{}灵芝产量是金银花的\ensuremath{6}倍。\allowbreak{}\par};
\end{scope}\end{tikzpicture}\newpage
\definecolor{answer}{rgb}{0.850980,0.047059,0.070588}
\null\begin{tikzpicture}[remember picture,overlay]\begin{scope}[shift={(current page.north west)},x=1bp,y=-1bp]
\node[anchor=north west,inner sep=0pt,text=answer,align=left,text width=215.66666412353516bp] at (15,15) {\fontsize{10.0bp}{12.5bp}\selectfont 先小括号,\allowbreak{} 再中括号,\allowbreak{}最后相加。\allowbreak{}\par};
\end{scope}\end{tikzpicture}\newpage
\definecolor{answer}{rgb}{0.850980,0.047059,0.070588}
\null\begin{tikzpicture}[remember picture,overlay]\begin{scope}[shift={(current page.north west)},x=1bp,y=-1bp]
\node[anchor=north west,inner sep=0pt,text=answer,align=left,text width=229.0bp] at (15,15) {\fontsize{8.5bp}{10.62bp}\selectfont 美术组人数是航模组的\ensuremath{2}倍；\allowbreak{}用合唱、\allowbreak{}体操组总人数除以美术组人数。\allowbreak{}\par};
\end{scope}\end{tikzpicture}\newpage
\definecolor{answer}{rgb}{0.850980,0.047059,0.070588}
\null\begin{tikzpicture}[remember picture,overlay]\begin{scope}[shift={(current page.north west)},x=1bp,y=-1bp]
\node[anchor=north west,inner sep=0pt,text=answer,align=left,text width=215.66666412353516bp] at (15,15) {\fontsize{10.0bp}{12.5bp}\selectfont 先小括号,\allowbreak{} 再中括号,\allowbreak{}最后相除。\allowbreak{}\par};
\end{scope}\end{tikzpicture}\newpage
\definecolor{answer}{rgb}{0.850980,0.047059,0.070588}
\null\begin{tikzpicture}[remember picture,overlay]\begin{scope}[shift={(current page.north west)},x=1bp,y=-1bp]
\node[anchor=north west,inner sep=0pt,text=answer,align=left,text width=215.66666412353516bp] at (15,15) {\fontsize{10.0bp}{12.5bp}\selectfont 先小括号,\allowbreak{} 再中括号,\allowbreak{}最后相乘。\allowbreak{}\par};
\end{scope}\end{tikzpicture}\newpage
\definecolor{answer}{rgb}{0.850980,0.047059,0.070588}
\null\begin{tikzpicture}[remember picture,overlay]\begin{scope}[shift={(current page.north west)},x=1bp,y=-1bp]
\node[anchor=north west,inner sep=0pt,text=answer,align=left,text width=229.0bp] at (15,15) {\fontsize{8.5bp}{10.62bp}\selectfont 手表价格是耳机的\ensuremath{2}倍多\ensuremath{80}元；\allowbreak{}平板价格比耳机与手表的差价的\ensuremath{3}倍少\ensuremath{50}元。\allowbreak{}原题是“差价”，\allowbreak{}不能把两者总价代入。\allowbreak{}\par};
\end{scope}\end{tikzpicture}\newpage
\definecolor{answer}{rgb}{0.850980,0.047059,0.070588}
\null\begin{tikzpicture}[remember picture,overlay]\begin{scope}[shift={(current page.north west)},x=1bp,y=-1bp]
\node[anchor=north west,inner sep=0pt,text=answer,align=left,text width=229.0bp] at (15,15) {\fontsize{8.5bp}{10.62bp}\selectfont 法\ensuremath{1} :\allowbreak{}先求手表价格，\allowbreak{}再求差价\par};
\end{scope}\end{tikzpicture}\newpage
\definecolor{answer}{rgb}{0.850980,0.047059,0.070588}
\null\begin{tikzpicture}[remember picture,overlay]\begin{scope}[shift={(current page.north west)},x=1bp,y=-1bp]
\node[anchor=north west,inner sep=0pt,text=answer,align=left,text width=229.0bp] at (15,15) {\fontsize{8.5bp}{10.62bp}\selectfont 法\ensuremath{2} :\allowbreak{}将差价化为耳机价格加\ensuremath{80}元\par};
\end{scope}\end{tikzpicture}\newpage
\definecolor{answer}{rgb}{0.850980,0.047059,0.070588}
\null\begin{tikzpicture}[remember picture,overlay]\begin{scope}[shift={(current page.north west)},x=1bp,y=-1bp]
\node[anchor=north west,inner sep=0pt,text=answer,align=left,text width=109.0bp] at (15,15) {\fontsize{7.5bp}{9.5bp}\selectfont 析：\allowbreak{}①代入②，\allowbreak{}再逐层代入③④。\allowbreak{}\par};
\end{scope}\end{tikzpicture}\newpage
\definecolor{answer}{rgb}{0.850980,0.047059,0.070588}
\null\begin{tikzpicture}[remember picture,overlay]\begin{scope}[shift={(current page.north west)},x=1bp,y=-1bp]
\node[anchor=north west,inner sep=0pt,text=answer,align=left,text width=108.0bp] at (15,15) {\fontsize{7.5bp}{9.5bp}\selectfont 析：\allowbreak{}先\ensuremath{120} ÷ \ensuremath{5},\allowbreak{}再\ensuremath{35}减商，\allowbreak{}乘\ensuremath{8}后从\ensuremath{150}中减。\allowbreak{}\par};
\end{scope}\end{tikzpicture}\newpage
\definecolor{answer}{rgb}{0.850980,0.047059,0.070588}
\null\begin{tikzpicture}[remember picture,overlay]\begin{scope}[shift={(current page.north west)},x=1bp,y=-1bp]
\node[anchor=north west,inner sep=0pt,text=answer,align=left,text width=233.41253662109375bp] at (15,15) {\fontsize{8.5bp}{10.62bp}\selectfont 先求货车运量，\allowbreak{}卡车比面包车与货车运量和的\ensuremath{5}倍少\ensuremath{25}箱。\allowbreak{}\par};
\end{scope}\end{tikzpicture}\newpage
\definecolor{answer}{rgb}{0.850980,0.047059,0.070588}
\null\begin{tikzpicture}[remember picture,overlay]\begin{scope}[shift={(current page.north west)},x=1bp,y=-1bp]
\node[anchor=north west,inner sep=0pt,text=answer,align=left,text width=229.0bp] at (15,15) {\fontsize{8.5bp}{10.62bp}\selectfont 法\ensuremath{2} :\allowbreak{}以面包车运量为基准补差\par};
\end{scope}\end{tikzpicture}\newpage
\definecolor{answer}{rgb}{0.850980,0.047059,0.070588}
\null\begin{tikzpicture}[remember picture,overlay]\begin{scope}[shift={(current page.north west)},x=1bp,y=-1bp]
\node[anchor=north west,inner sep=0pt,text=answer,align=left,text width=231.0bp] at (15,15) {\fontsize{10.0bp}{12.5bp}\selectfont 先求第二组,\allowbreak{} 再求前两组总数,\allowbreak{}最后乘\ensuremath{2}。\allowbreak{}也可用分步计算检验。\allowbreak{}\par};
\end{scope}\end{tikzpicture}\newpage
\definecolor{answer}{rgb}{0.850980,0.047059,0.070588}
\null\begin{tikzpicture}[remember picture,overlay]\begin{scope}[shift={(current page.north west)},x=1bp,y=-1bp]
\node[anchor=north west,inner sep=0pt,text=answer,align=left,text width=229.0bp] at (15,15) {\fontsize{8.5bp}{10.62bp}\selectfont “\ensuremath{19}米的\ensuremath{3}倍多\ensuremath{1}米”表示先求\ensuremath{3}倍，\allowbreak{}再加\ensuremath{1}米。\allowbreak{}\par};
\end{scope}\end{tikzpicture}\newpage
\definecolor{answer}{rgb}{0.850980,0.047059,0.070588}
\null\begin{tikzpicture}[remember picture,overlay]\begin{scope}[shift={(current page.north west)},x=1bp,y=-1bp]
\node[anchor=north west,inner sep=0pt,text=answer,align=left,text width=229.0bp] at (15,15) {\fontsize{8.5bp}{10.62bp}\selectfont 先求节水龙头流量，\allowbreak{}再加\ensuremath{300}千克得到消防龙头流量，\allowbreak{}最后减普通龙头流量。\allowbreak{}\par};
\end{scope}\end{tikzpicture}\newpage
\definecolor{answer}{rgb}{0.850980,0.047059,0.070588}
\null\begin{tikzpicture}[remember picture,overlay]\begin{scope}[shift={(current page.north west)},x=1bp,y=-1bp]
\node[anchor=north west,inner sep=0pt,text=answer,align=left,text width=253.0bp] at (15,15) {\fontsize{8.5bp}{10.62bp}\selectfont 第二组栽\ensuremath{70} + \ensuremath{23}棵；\allowbreak{}第三组是前两组总数的\ensuremath{2}倍。\allowbreak{}\par};
\end{scope}\end{tikzpicture}\newpage
\definecolor{answer}{rgb}{0.850980,0.047059,0.070588}
\null\begin{tikzpicture}[remember picture,overlay]\begin{scope}[shift={(current page.north west)},x=1bp,y=-1bp]
\node[anchor=north west,inner sep=0pt,text=answer,align=left,text width=229.0bp] at (15,15) {\fontsize{8.5bp}{10.62bp}\selectfont 答：\allowbreak{}消防龙头每分钟比普通龙头多流\ensuremath{295}千克水。\allowbreak{}\par};
\end{scope}\end{tikzpicture}\newpage
\definecolor{answer}{rgb}{0.850980,0.047059,0.070588}
\null\begin{tikzpicture}[remember picture,overlay]\begin{scope}[shift={(current page.north west)},x=1bp,y=-1bp]
\node[anchor=north west,inner sep=0pt,text=answer,align=left,text width=234.0bp] at (15,15) {\fontsize{7.5bp}{9.5bp}\selectfont 析：\allowbreak{}括住\ensuremath{600} \ensuremath{-}\ensuremath{300},\allowbreak{}使减法先算。\allowbreak{}\par};
\end{scope}\end{tikzpicture}\newpage
\definecolor{answer}{rgb}{0.850980,0.047059,0.070588}
\null\begin{tikzpicture}[remember picture,overlay]\begin{scope}[shift={(current page.north west)},x=1bp,y=-1bp]
\node[anchor=north west,inner sep=0pt,text=answer,align=left,text width=234.0bp] at (15,15) {\fontsize{7.5bp}{9.5bp}\selectfont 析：\allowbreak{}括住\ensuremath{600} \ensuremath{-}\ensuremath{300} ÷ \ensuremath{50},\allowbreak{}先除再减。\allowbreak{}\par};
\end{scope}\end{tikzpicture}\newpage
\definecolor{answer}{rgb}{0.850980,0.047059,0.070588}
\null\begin{tikzpicture}[remember picture,overlay]\begin{scope}[shift={(current page.north west)},x=1bp,y=-1bp]
\node[anchor=north west,inner sep=0pt,text=answer,align=left,text width=234.0bp] at (15,15) {\fontsize{7.5bp}{9.5bp}\selectfont 析：\allowbreak{}先小括号减法，\allowbreak{}再中括号除法。\allowbreak{}\par};
\end{scope}\end{tikzpicture}\newpage
\definecolor{answer}{rgb}{0.850980,0.047059,0.070588}
\null\begin{tikzpicture}[remember picture,overlay]\begin{scope}[shift={(current page.north west)},x=1bp,y=-1bp]
\node[anchor=north west,inner sep=0pt,text=answer,align=left,text width=234.0bp] at (15,15) {\fontsize{8.5bp}{10.62bp}\selectfont 法\ensuremath{2} :\allowbreak{}先用乘法分配律消去除法\par};
\end{scope}\end{tikzpicture}\newpage
\definecolor{answer}{rgb}{0.850980,0.047059,0.070588}
\null\begin{tikzpicture}[remember picture,overlay]\begin{scope}[shift={(current page.north west)},x=1bp,y=-1bp]
\node[anchor=north west,inner sep=0pt,text=answer,align=left,text width=229.0bp] at (15,15) {\fontsize{8.5bp}{10.62bp}\selectfont 法\ensuremath{2}：\allowbreak{}以每天\ensuremath{76}个为基准，\allowbreak{}七天先都按每天\ensuremath{76}个算；\allowbreak{}后\ensuremath{5}天每天多\ensuremath{7}个，\allowbreak{}合计多\ensuremath{35}个，\allowbreak{}平均分到七天，\allowbreak{}每天多\ensuremath{5}个。\allowbreak{}\par};
\end{scope}\end{tikzpicture}\newpage
\definecolor{answer}{rgb}{0.850980,0.047059,0.070588}
\null\begin{tikzpicture}[remember picture,overlay]\begin{scope}[shift={(current page.north west)},x=1bp,y=-1bp]
\node[anchor=north west,inner sep=0pt,text=answer,align=left,text width=230.25bp] at (15,15) {\fontsize{8.5bp}{10.62bp}\selectfont 先求已装质量，\allowbreak{}再求每袋质量，\allowbreak{}剩余质量除以每袋质量。\allowbreak{}\par};
\end{scope}\end{tikzpicture}\newpage
\definecolor{answer}{rgb}{0.850980,0.047059,0.070588}
\null\begin{tikzpicture}[remember picture,overlay]\begin{scope}[shift={(current page.north west)},x=1bp,y=-1bp]
\node[anchor=north west,inner sep=0pt,text=answer,align=left,text width=229.0bp] at (15,15) {\fontsize{8.5bp}{10.62bp}\selectfont 法\ensuremath{1} :\allowbreak{}先求每袋质量，\allowbreak{}再求剩余袋数\par};
\end{scope}\end{tikzpicture}\newpage
\definecolor{answer}{rgb}{0.850980,0.047059,0.070588}
\null\begin{tikzpicture}[remember picture,overlay]\begin{scope}[shift={(current page.north west)},x=1bp,y=-1bp]
\node[anchor=north west,inner sep=0pt,text=answer,align=left,text width=229.0bp] at (15,15) {\fontsize{8.5bp}{10.62bp}\selectfont 法\ensuremath{2} :\allowbreak{}先求总袋数，\allowbreak{}再扣已装袋数\par};
\end{scope}\end{tikzpicture}\newpage
\definecolor{answer}{rgb}{0.850980,0.047059,0.070588}
\null\begin{tikzpicture}[remember picture,overlay]\begin{scope}[shift={(current page.north west)},x=1bp,y=-1bp]
\node[anchor=north west,inner sep=0pt,text=answer,align=left,text width=229.0bp] at (15,15) {\fontsize{8.5bp}{10.62bp}\selectfont 答：\allowbreak{}还需要准备\ensuremath{23}个小袋子。\allowbreak{}\par};
\end{scope}\end{tikzpicture}\newpage
\definecolor{answer}{rgb}{0.850980,0.047059,0.070588}
\null\begin{tikzpicture}[remember picture,overlay]\begin{scope}[shift={(current page.north west)},x=1bp,y=-1bp]
\node[anchor=north west,inner sep=0pt,text=answer,align=left,text width=312.447509765625bp] at (15,15) {\fontsize{7.5bp}{9.5bp}\selectfont 析：\allowbreak{}先求和，\allowbreak{}再求商：\allowbreak{}(\ensuremath{56} + \ensuremath{16}) ÷ \ensuremath{8} = \ensuremath{9}。\allowbreak{}\par};
\end{scope}\end{tikzpicture}\newpage
\definecolor{answer}{rgb}{0.850980,0.047059,0.070588}
\null\begin{tikzpicture}[remember picture,overlay]\begin{scope}[shift={(current page.north west)},x=1bp,y=-1bp]
\node[anchor=north west,inner sep=0pt,text=answer,align=left,text width=324.74700927734375bp] at (15,15) {\fontsize{7.5bp}{9.5bp}\selectfont 析：\allowbreak{}原式先乘后加再除；\allowbreak{}要先加，\allowbreak{}用\ensuremath{208} ÷ [(\ensuremath{5} + \ensuremath{7}) × \ensuremath{3}]。\allowbreak{}\par};
\end{scope}\end{tikzpicture}\newpage
\definecolor{answer}{rgb}{0.850980,0.047059,0.070588}
\null\begin{tikzpicture}[remember picture,overlay]\begin{scope}[shift={(current page.north west)},x=1bp,y=-1bp]
\node[anchor=north west,inner sep=0pt,text=answer,align=left,text width=303.7879943847656bp] at (15,15) {\fontsize{7.5bp}{9.5bp}\selectfont 析：\allowbreak{}先\ensuremath{800} \ensuremath{-}\ensuremath{300},\allowbreak{}再÷\ensuremath{50},\allowbreak{}最后×\ensuremath{150}。\allowbreak{}\par};
\end{scope}\end{tikzpicture}\newpage
\definecolor{answer}{rgb}{0.850980,0.047059,0.070588}
\null\begin{tikzpicture}[remember picture,overlay]\begin{scope}[shift={(current page.north west)},x=1bp,y=-1bp]
\node[anchor=north west,inner sep=0pt,text=answer,align=left,text width=329.02020263671875bp] at (15,15) {\fontsize{7.5bp}{9.5bp}\selectfont 析：\allowbreak{}分别算两边：\allowbreak{}\ensuremath{8} = \ensuremath{8};\allowbreak{}\ensuremath{6} < \ensuremath{8};\allowbreak{} \ensuremath{13} < \ensuremath{75};\allowbreak{} \ensuremath{16} > \ensuremath{1}。\allowbreak{}\par};
\end{scope}\end{tikzpicture}\newpage
\definecolor{answer}{rgb}{0.850980,0.047059,0.070588}
\null\begin{tikzpicture}[remember picture,overlay]\begin{scope}[shift={(current page.north west)},x=1bp,y=-1bp]
\node[anchor=north west,inner sep=0pt,text=answer,align=left,text width=422.33099365234375bp] at (15,15) {\fontsize{7.5bp}{9.5bp}\selectfont 析：\allowbreak{}\ensuremath{116} ÷ \ensuremath{2} = \ensuremath{58}(千米/时)；\allowbreak{}\ensuremath{464} ÷ \ensuremath{58} \ensuremath{-}\ensuremath{2} = \ensuremath{6}(时)，\allowbreak{}也可(\ensuremath{464} \ensuremath{-}\ensuremath{116}) ÷ \ensuremath{58} = \ensuremath{6}。\allowbreak{}\par};
\end{scope}\end{tikzpicture}\newpage
\definecolor{answer}{rgb}{0.850980,0.047059,0.070588}
\null\begin{tikzpicture}[remember picture,overlay]\begin{scope}[shift={(current page.north west)},x=1bp,y=-1bp]
\node[anchor=north west,inner sep=0pt,text=answer,align=left,text width=373.301513671875bp] at (15,15) {\fontsize{7.5bp}{9.5bp}\selectfont 析：\allowbreak{}错式：\allowbreak{}(\ensuremath{800}\ensuremath{-}□)÷\ensuremath{40} = \ensuremath{15},\allowbreak{}故□=\ensuremath{200};\allowbreak{}正确算\ensuremath{600} ÷ \ensuremath{15} + \ensuremath{25} = \ensuremath{65}。\allowbreak{}\par};
\end{scope}\end{tikzpicture}\newpage
\definecolor{answer}{rgb}{0.850980,0.047059,0.070588}
\null\begin{tikzpicture}[remember picture,overlay]\begin{scope}[shift={(current page.north west)},x=1bp,y=-1bp]
\node[anchor=north west,inner sep=0pt,text=answer,align=left,text width=351.2340087890625bp] at (15,15) {\fontsize{7.5bp}{9.5bp}\selectfont 析：\allowbreak{}\ensuremath{0}不能作除数，\allowbreak{}所以“\ensuremath{0}除以任何数”缺少限制。\allowbreak{}\par};
\end{scope}\end{tikzpicture}\newpage
\definecolor{answer}{rgb}{0.850980,0.047059,0.070588}
\null\begin{tikzpicture}[remember picture,overlay]\begin{scope}[shift={(current page.north west)},x=1bp,y=-1bp]
\node[anchor=north west,inner sep=0pt,text=answer,align=left,text width=306.7869873046875bp] at (15,15) {\fontsize{7.5bp}{9.5bp}\selectfont 析：\allowbreak{}括号内\ensuremath{81} \ensuremath{-}\ensuremath{41}先算，\allowbreak{}不能先算\ensuremath{41} + \ensuremath{32}。\allowbreak{}\par};
\end{scope}\end{tikzpicture}\newpage
\definecolor{answer}{rgb}{0.850980,0.047059,0.070588}
\null\begin{tikzpicture}[remember picture,overlay]\begin{scope}[shift={(current page.north west)},x=1bp,y=-1bp]
\node[anchor=north west,inner sep=0pt,text=answer,align=left,text width=279.5bp] at (15,15) {\fontsize{7.5bp}{9.5bp}\selectfont 析：\allowbreak{}加减互逆，\allowbreak{}乘除互逆。\allowbreak{}\par};
\end{scope}\end{tikzpicture}\newpage
\definecolor{answer}{rgb}{0.850980,0.047059,0.070588}
\null\begin{tikzpicture}[remember picture,overlay]\begin{scope}[shift={(current page.north west)},x=1bp,y=-1bp]
\node[anchor=north west,inner sep=0pt,text=answer,align=left,text width=358.78900146484375bp] at (15,15) {\fontsize{7.5bp}{9.5bp}\selectfont 析：\allowbreak{}去掉中括号仍先算小括号，\allowbreak{}再从左到右除\ensuremath{15}、\allowbreak{}乘\ensuremath{35},\allowbreak{}顺序不变。\allowbreak{}\par};
\end{scope}\end{tikzpicture}\newpage
\definecolor{answer}{rgb}{0.850980,0.047059,0.070588}
\null\begin{tikzpicture}[remember picture,overlay]\begin{scope}[shift={(current page.north west)},x=1bp,y=-1bp]
\node[anchor=north west,inner sep=0pt,text=answer,align=left,text width=494.0bp] at (15,15) {\fontsize{7.5bp}{9.5bp}\selectfont 析：\allowbreak{}先算小括号\ensuremath{9} + \ensuremath{7},\allowbreak{}是加法，\allowbreak{}选A。\allowbreak{}\par};
\end{scope}\end{tikzpicture}\newpage
\definecolor{answer}{rgb}{0.850980,0.047059,0.070588}
\null\begin{tikzpicture}[remember picture,overlay]\begin{scope}[shift={(current page.north west)},x=1bp,y=-1bp]
\node[anchor=north west,inner sep=0pt,text=answer,align=left,text width=494.0bp] at (15,15) {\fontsize{7.5bp}{9.5bp}\selectfont 析：\allowbreak{}和作被除数、\allowbreak{}差作除数，\allowbreak{}两者都括起：\allowbreak{}(\ensuremath{352} + \ensuremath{146}) ÷ (\ensuremath{50} \ensuremath{-}\ensuremath{35}),\allowbreak{}选B。\allowbreak{}\par};
\end{scope}\end{tikzpicture}\newpage
\definecolor{answer}{rgb}{0.850980,0.047059,0.070588}
\null\begin{tikzpicture}[remember picture,overlay]\begin{scope}[shift={(current page.north west)},x=1bp,y=-1bp]
\node[anchor=north west,inner sep=0pt,text=answer,align=left,text width=494.0bp] at (15,15) {\fontsize{7.5bp}{9.5bp}\selectfont 析：\allowbreak{}A去中括号后仍先减，\allowbreak{}再从左到右乘、\allowbreak{}除，\allowbreak{}选A;\allowbreak{}其余会改变顺序。\allowbreak{}\par};
\end{scope}\end{tikzpicture}\newpage
\definecolor{answer}{rgb}{0.850980,0.047059,0.070588}
\null\begin{tikzpicture}[remember picture,overlay]\begin{scope}[shift={(current page.north west)},x=1bp,y=-1bp]
\node[anchor=north west,inner sep=0pt,text=answer,align=left,text width=494.0bp] at (15,15) {\fontsize{7.5bp}{9.5bp}\selectfont 析：\allowbreak{}五年级\ensuremath{28} × \ensuremath{3},\allowbreak{}四年级(\ensuremath{28} \ensuremath{-}\ensuremath{5}) × \ensuremath{4},\allowbreak{}相加得\ensuremath{176}棵，\allowbreak{}选A。\allowbreak{}\par};
\end{scope}\end{tikzpicture}\newpage
\definecolor{answer}{rgb}{0.850980,0.047059,0.070588}
\null\begin{tikzpicture}[remember picture,overlay]\begin{scope}[shift={(current page.north west)},x=1bp,y=-1bp]
\node[anchor=north west,inner sep=0pt,text=answer,align=left,text width=494.0bp] at (15,15) {\fontsize{7.5bp}{9.5bp}\selectfont 析：\allowbreak{}▲×▲=●且▲≠\ensuremath{0},\allowbreak{}可得●÷▲=▲，\allowbreak{}选C。\allowbreak{}\par};
\end{scope}\end{tikzpicture}\newpage
\definecolor{answer}{rgb}{0.850980,0.047059,0.070588}
\null\begin{tikzpicture}[remember picture,overlay]\begin{scope}[shift={(current page.north west)},x=1bp,y=-1bp]
\node[anchor=north west,inner sep=0pt,text=answer,align=left,text width=494.0bp] at (15,15) {\fontsize{7.5bp}{9.5bp}\selectfont 析：\allowbreak{}第二组减第一组：\allowbreak{}\ensuremath{180} ÷ \ensuremath{6} \ensuremath{-}\ensuremath{120} ÷ \ensuremath{6} = \ensuremath{10};\allowbreak{}也可(\ensuremath{180} \ensuremath{-}\ensuremath{120}) ÷ \ensuremath{6} = \ensuremath{10},\allowbreak{}选B。\allowbreak{}\par};
\end{scope}\end{tikzpicture}\newpage
\definecolor{answer}{rgb}{0.850980,0.047059,0.070588}
\null\begin{tikzpicture}[remember picture,overlay]\begin{scope}[shift={(current page.north west)},x=1bp,y=-1bp]
\node[anchor=north west,inner sep=0pt,text=answer,align=left,text width=311.4809875488281bp] at (15,15) {\fontsize{8.5bp}{10.62bp}\selectfont 总件数不变，\allowbreak{}实际用\ensuremath{10} \ensuremath{-}\ensuremath{2} = \ensuremath{8}天。\allowbreak{}\par};
\end{scope}\end{tikzpicture}\newpage
\definecolor{answer}{rgb}{0.850980,0.047059,0.070588}
\null\begin{tikzpicture}[remember picture,overlay]\begin{scope}[shift={(current page.north west)},x=1bp,y=-1bp]
\node[anchor=north west,inner sep=0pt,text=answer,align=left,text width=467.0bp] at (15,15) {\fontsize{8.5bp}{10.62bp}\selectfont 已生产的数量等于总任务减剩余数量，\allowbreak{}再除以\ensuremath{25}小时。\allowbreak{}\par};
\end{scope}\end{tikzpicture}\newpage
\definecolor{answer}{rgb}{0.850980,0.047059,0.070588}
\null\begin{tikzpicture}[remember picture,overlay]\begin{scope}[shift={(current page.north west)},x=1bp,y=-1bp]
\node[anchor=north west,inner sep=0pt,text=answer,align=left,text width=467.0bp] at (15,15) {\fontsize{8.5bp}{10.62bp}\selectfont 周五总销量为上午与下午销量之和，\allowbreak{}周六为周五的\ensuremath{2}倍少\ensuremath{15}支。\allowbreak{}\par};
\end{scope}\end{tikzpicture}\newpage
\definecolor{answer}{rgb}{0.850980,0.047059,0.070588}
\null\begin{tikzpicture}[remember picture,overlay]\begin{scope}[shift={(current page.north west)},x=1bp,y=-1bp]
\node[anchor=north west,inner sep=0pt,text=answer,align=left,text width=254.0bp] at (15,15) {\fontsize{7.5bp}{9.5bp}\selectfont 析：\allowbreak{}定义有两个条件：\allowbreak{}同一平面内、\allowbreak{}不相交。\allowbreak{}\par};
\end{scope}\end{tikzpicture}\newpage
\definecolor{answer}{rgb}{0.850980,0.047059,0.070588}
\null\begin{tikzpicture}[remember picture,overlay]\begin{scope}[shift={(current page.north west)},x=1bp,y=-1bp]
\node[anchor=north west,inner sep=0pt,text=answer,align=left,text width=254.0bp] at (15,15) {\fontsize{7.5bp}{9.5bp}\selectfont 析：\allowbreak{}相交成直角为互相垂直，\allowbreak{}交点为垂足。\allowbreak{}\par};
\end{scope}\end{tikzpicture}\newpage
\definecolor{answer}{rgb}{0.850980,0.047059,0.070588}
\null\begin{tikzpicture}[remember picture,overlay]\begin{scope}[shift={(current page.north west)},x=1bp,y=-1bp]
\node[anchor=north west,inner sep=0pt,text=answer,align=left,text width=231.0bp] at (15,15) {\fontsize{7.5bp}{9.5bp}\selectfont 析：\allowbreak{}②④⑧形成直角；\allowbreak{}⑤⑦延长也不相交。\allowbreak{}其余不能仅凭线段没碰到判平行。\allowbreak{}\par};
\end{scope}\end{tikzpicture}\newpage
\definecolor{answer}{rgb}{0.850980,0.047059,0.070588}
\null\begin{tikzpicture}[remember picture,overlay]\begin{scope}[shift={(current page.north west)},x=1bp,y=-1bp]
\node[anchor=north west,inner sep=0pt,text=answer,align=left,text width=254.0bp] at (15,15) {\fontsize{7.5bp}{9.5bp}\selectfont 析：\allowbreak{}\ensuremath{c}与d垂直，\allowbreak{}记作\ensuremath{c} ⊥d。\allowbreak{}\par};
\end{scope}\end{tikzpicture}\newpage
\definecolor{answer}{rgb}{0.850980,0.047059,0.070588}
\null\begin{tikzpicture}[remember picture,overlay]\begin{scope}[shift={(current page.north west)},x=1bp,y=-1bp]
\node[anchor=north west,inner sep=0pt,text=answer,align=left,text width=254.0bp] at (15,15) {\fontsize{7.5bp}{9.5bp}\selectfont 析：\allowbreak{}整时分针指\ensuremath{12};\allowbreak{}时针指\ensuremath{3}或\ensuremath{9}时夹角\ensuremath{90}\ensuremath{^\circ}。\allowbreak{}\par};
\end{scope}\end{tikzpicture}\newpage
\definecolor{answer}{rgb}{0.850980,0.047059,0.070588}
\null\begin{tikzpicture}[remember picture,overlay]\begin{scope}[shift={(current page.north west)},x=1bp,y=-1bp]
\node[anchor=north west,inner sep=0pt,text=answer,align=left,text width=254.0bp] at (15,15) {\fontsize{7.5bp}{9.5bp}\selectfont 析：\allowbreak{}推拉门轨道同向且不相交，\allowbreak{}所以平行。\allowbreak{}\par};
\end{scope}\end{tikzpicture}\newpage
\definecolor{answer}{rgb}{0.850980,0.047059,0.070588}
\null\begin{tikzpicture}[remember picture,overlay]\begin{scope}[shift={(current page.north west)},x=1bp,y=-1bp]
\node[anchor=north west,inner sep=0pt,text=answer,align=left,text width=222.3333740234375bp] at (15,15) {\fontsize{8.0bp}{10.0bp}\selectfont 红线:\allowbreak{} 平行;\allowbreak{}直角标记处:\allowbreak{} 垂直。\allowbreak{}\par};
\end{scope}\end{tikzpicture}\newpage
\definecolor{answer}{rgb}{0.850980,0.047059,0.070588}
\null\begin{tikzpicture}[remember picture,overlay]\begin{scope}[shift={(current page.north west)},x=1bp,y=-1bp]
\node[anchor=north west,inner sep=0pt,text=answer,align=left,text width=235.0bp] at (15,15) {\fontsize{7.5bp}{9.5bp}\selectfont 析：\allowbreak{}缺“同一平面内”；\allowbreak{}空间中不相交也可能不平行。\allowbreak{}\par};
\end{scope}\end{tikzpicture}\newpage
\definecolor{answer}{rgb}{0.850980,0.047059,0.070588}
\null\begin{tikzpicture}[remember picture,overlay]\begin{scope}[shift={(current page.north west)},x=1bp,y=-1bp]
\node[anchor=north west,inner sep=0pt,text=answer,align=left,text width=235.0bp] at (15,15) {\fontsize{7.5bp}{9.5bp}\selectfont 析：\allowbreak{}正方形是特殊长方形，\allowbreak{}两组对边均平行。\allowbreak{}\par};
\end{scope}\end{tikzpicture}\newpage
\definecolor{answer}{rgb}{0.850980,0.047059,0.070588}
\null\begin{tikzpicture}[remember picture,overlay]\begin{scope}[shift={(current page.north west)},x=1bp,y=-1bp]
\node[anchor=north west,inner sep=0pt,text=answer,align=left,text width=235.0bp] at (15,15) {\fontsize{7.5bp}{9.5bp}\selectfont 析：\allowbreak{}不平行可成锐角或钝角，\allowbreak{}不一定垂直。\allowbreak{}\par};
\end{scope}\end{tikzpicture}\newpage
\definecolor{answer}{rgb}{0.850980,0.047059,0.070588}
\null\begin{tikzpicture}[remember picture,overlay]\begin{scope}[shift={(current page.north west)},x=1bp,y=-1bp]
\node[anchor=north west,inner sep=0pt,text=answer,align=left,text width=235.0bp] at (15,15) {\fontsize{7.5bp}{9.5bp}\selectfont 析：\allowbreak{}相交成直角正是互相垂直的定义。\allowbreak{}\par};
\end{scope}\end{tikzpicture}\newpage
\definecolor{answer}{rgb}{0.850980,0.047059,0.070588}
\null\begin{tikzpicture}[remember picture,overlay]\begin{scope}[shift={(current page.north west)},x=1bp,y=-1bp]
\node[anchor=north west,inner sep=0pt,text=answer,align=left,text width=233.0bp] at (15,15) {\fontsize{7.5bp}{9.5bp}\selectfont 析：\allowbreak{}F 、\allowbreak{}H、\allowbreak{}N、\allowbreak{}Z含平行线段：\allowbreak{}①②④⑤。\allowbreak{}\par};
\end{scope}\end{tikzpicture}\newpage
\definecolor{answer}{rgb}{0.850980,0.047059,0.070588}
\null\begin{tikzpicture}[remember picture,overlay]\begin{scope}[shift={(current page.north west)},x=1bp,y=-1bp]
\node[anchor=north west,inner sep=0pt,text=answer,align=left,text width=233.0bp] at (15,15) {\fontsize{7.5bp}{9.5bp}\selectfont 析：\allowbreak{}F 、\allowbreak{}H、\allowbreak{}L、\allowbreak{}T 有直角：\allowbreak{}①②③⑥。\allowbreak{}\par};
\end{scope}\end{tikzpicture}\newpage
\definecolor{answer}{rgb}{0.850980,0.047059,0.070588}
\null\begin{tikzpicture}[remember picture,overlay]\begin{scope}[shift={(current page.north west)},x=1bp,y=-1bp]
\node[anchor=north west,inner sep=0pt,text=answer,align=left,text width=235.0bp] at (15,15) {\fontsize{7.5bp}{9.5bp}\selectfont 析：\allowbreak{}两侧履带印迹同向不相交，\allowbreak{}选A。\allowbreak{}\par};
\end{scope}\end{tikzpicture}\newpage
\definecolor{answer}{rgb}{0.850980,0.047059,0.070588}
\null\begin{tikzpicture}[remember picture,overlay]\begin{scope}[shift={(current page.north west)},x=1bp,y=-1bp]
\node[anchor=north west,inner sep=0pt,text=answer,align=left,text width=233.0bp] at (15,15) {\fontsize{7.5bp}{9.5bp}\selectfont 析：\allowbreak{}同时满足两种条件的是F 、\allowbreak{}H,\allowbreak{}即①②。\allowbreak{}\par};
\end{scope}\end{tikzpicture}\newpage
\definecolor{answer}{rgb}{0.850980,0.039216,0.070588}
\null\begin{tikzpicture}[remember picture,overlay]\begin{scope}[shift={(current page.north west)},x=1bp,y=-1bp]
\node[anchor=north west,inner sep=0pt,text=answer,align=left,text width=219.0bp] at (15,15) {\fontsize{10.0bp}{12.5bp}\selectfont 互相垂直。\allowbreak{}\ensuremath{b}与\ensuremath{c}所成的角为\par};
\end{scope}\end{tikzpicture}\newpage
\definecolor{answer}{rgb}{0.850980,0.047059,0.070588}
\null\begin{tikzpicture}[remember picture,overlay]\begin{scope}[shift={(current page.north west)},x=1bp,y=-1bp]
\node[anchor=north west,inner sep=0pt,text=answer,align=left,text width=235.0bp] at (15,15) {\fontsize{7.5bp}{9.5bp}\selectfont 析：\allowbreak{}A漏同一平面；\allowbreak{}B一般交点不叫垂足；\allowbreak{}C正确，\allowbreak{}选C。\allowbreak{}\par};
\end{scope}\end{tikzpicture}\newpage
\definecolor{answer}{rgb}{0.850980,0.047059,0.070588}
\null\begin{tikzpicture}[remember picture,overlay]\begin{scope}[shift={(current page.north west)},x=1bp,y=-1bp]
\node[anchor=north west,inner sep=0pt,text=answer,align=left,text width=235.0bp] at (15,15) {\fontsize{7.5bp}{9.5bp}\selectfont 析：\allowbreak{}两次折叠可同向，\allowbreak{}也可互相垂直，\allowbreak{}选C。\allowbreak{}\par};
\end{scope}\end{tikzpicture}\newpage
\definecolor{answer}{rgb}{0.850980,0.047059,0.070588}
\null\begin{tikzpicture}[remember picture,overlay]\begin{scope}[shift={(current page.north west)},x=1bp,y=-1bp]
\node[anchor=north west,inner sep=0pt,text=answer,align=left,text width=219.0bp] at (15,15) {\fontsize{10.0bp}{12.5bp}\selectfont 例如:\allowbreak{} 铁轨的两根钢轨互相平行;\allowbreak{}方格纸上横线与竖线互相垂直。\allowbreak{}\par};
\end{scope}\end{tikzpicture}\newpage
\definecolor{answer}{rgb}{0.850980,0.047059,0.070588}
\null\begin{tikzpicture}[remember picture,overlay]\begin{scope}[shift={(current page.north west)},x=1bp,y=-1bp]
\node[anchor=north west,inner sep=0pt,text=answer,align=left,text width=235.0bp] at (15,15) {\fontsize{7.5bp}{9.5bp}\selectfont 析：\allowbreak{}按道路中心线判断：\allowbreak{}①⊥③，\allowbreak{}②∥④。\allowbreak{}\par};
\end{scope}\end{tikzpicture}\newpage
\definecolor{answer}{rgb}{0.850980,0.047059,0.070588}
\null\begin{tikzpicture}[remember picture,overlay]\begin{scope}[shift={(current page.north west)},x=1bp,y=-1bp]
\node[anchor=north west,inner sep=0pt,text=answer,align=left,text width=231.0bp] at (15,15) {\fontsize{7.5bp}{9.5bp}\selectfont 析：\allowbreak{}横、\allowbreak{}竖、\allowbreak{}斜各\ensuremath{1}组平行，\allowbreak{}共\ensuremath{3}组；\allowbreak{}\ensuremath{2}条横与\ensuremath{2}条竖组成\ensuremath{2} × \ensuremath{2} = \ensuremath{4}组垂直。\allowbreak{}\par};
\end{scope}\end{tikzpicture}\newpage
\definecolor{answer}{rgb}{0.850980,0.047059,0.070588}
\null\begin{tikzpicture}[remember picture,overlay]\begin{scope}[shift={(current page.north west)},x=1bp,y=-1bp]
\node[anchor=north west,inner sep=0pt,text=answer,align=left,text width=243.0bp] at (15,15) {\fontsize{7.5bp}{9.5bp}\selectfont 析：\allowbreak{}垂直线段最短，\allowbreak{}其长度是点到直线的距离。\allowbreak{}\par};
\end{scope}\end{tikzpicture}\newpage
\definecolor{answer}{rgb}{0.850980,0.047059,0.070588}
\null\begin{tikzpicture}[remember picture,overlay]\begin{scope}[shift={(current page.north west)},x=1bp,y=-1bp]
\node[anchor=north west,inner sep=0pt,text=answer,align=left,text width=243.0bp] at (15,15) {\fontsize{7.5bp}{9.5bp}\selectfont 析：\allowbreak{}平行线间处处等距，\allowbreak{}可作无数条等长垂线段。\allowbreak{}\par};
\end{scope}\end{tikzpicture}\newpage
\definecolor{answer}{rgb}{0.850980,0.047059,0.070588}
\null\begin{tikzpicture}[remember picture,overlay]\begin{scope}[shift={(current page.north west)},x=1bp,y=-1bp]
\node[anchor=north west,inner sep=0pt,text=answer,align=left,text width=243.0bp] at (15,15) {\fontsize{7.5bp}{9.5bp}\selectfont 析：\allowbreak{}无论点在线上或线外，\allowbreak{}过一点的垂线都只有\ensuremath{1}条。\allowbreak{}\par};
\end{scope}\end{tikzpicture}\newpage
\definecolor{answer}{rgb}{0.850980,0.047059,0.070588}
\null\begin{tikzpicture}[remember picture,overlay]\begin{scope}[shift={(current page.north west)},x=1bp,y=-1bp]
\node[anchor=north west,inner sep=0pt,text=answer,align=left,text width=243.0bp] at (15,15) {\fontsize{7.5bp}{9.5bp}\selectfont 析：\allowbreak{}距离取垂直线段长度，\allowbreak{}即三条中的最短\ensuremath{110}米。\allowbreak{}\par};
\end{scope}\end{tikzpicture}\newpage
\definecolor{answer}{rgb}{0.850980,0.047059,0.070588}
\null\begin{tikzpicture}[remember picture,overlay]\begin{scope}[shift={(current page.north west)},x=1bp,y=-1bp]
\node[anchor=north west,inner sep=0pt,text=answer,align=left,text width=243.0bp] at (15,15) {\fontsize{7.5bp}{9.5bp}\selectfont 析：\allowbreak{}直角边贴直线，\allowbreak{}另一直角边过指定点；\allowbreak{}A符合。\allowbreak{}\par};
\end{scope}\end{tikzpicture}\newpage
\definecolor{answer}{rgb}{0.850980,0.047059,0.070588}
\null\begin{tikzpicture}[remember picture,overlay]\begin{scope}[shift={(current page.north west)},x=1bp,y=-1bp]
\node[anchor=north west,inner sep=0pt,text=answer,align=left,text width=243.0bp] at (15,15) {\fontsize{7.5bp}{9.5bp}\selectfont 析：\allowbreak{}BA ⊥AC,\allowbreak{} B到AC为\ensuremath{6}厘米；\allowbreak{}CA ⊥AB,\allowbreak{} C到AB为\ensuremath{8}厘米。\allowbreak{}\par};
\end{scope}\end{tikzpicture}\newpage
\definecolor{answer}{rgb}{0.850980,0.047059,0.070588}
\null\begin{tikzpicture}[remember picture,overlay]\begin{scope}[shift={(current page.north west)},x=1bp,y=-1bp]
\node[anchor=north west,inner sep=0pt,text=answer,align=left,text width=229.0bp] at (15,15) {\fontsize{7.5bp}{9.5bp}\selectfont 析：\allowbreak{}平行：\allowbreak{}\ensuremath{a} ∥\ensuremath{c},\allowbreak{} \ensuremath{b} ∥d ∥e;\allowbreak{}垂直：\allowbreak{}\ensuremath{a} ⊥l、\allowbreak{}\ensuremath{c} ⊥l。\allowbreak{}同一平面内，\allowbreak{}平行于同一直线或垂直于同一直线的两线均平行。\allowbreak{}\par};
\end{scope}\end{tikzpicture}\newpage
\definecolor{answer}{rgb}{0.850980,0.047059,0.070588}
\null\begin{tikzpicture}[remember picture,overlay]\begin{scope}[shift={(current page.north west)},x=1bp,y=-1bp]
\node[anchor=north west,inner sep=0pt,text=answer,align=left,text width=226.33334350585938bp] at (15,15) {\fontsize{10.0bp}{12.5bp}\selectfont 不公平。\allowbreak{}佳佳离凳子最近;\allowbreak{}点到直线的垂线段最短。\allowbreak{}可让大家站在以凳子为圆心的同一圆上。\allowbreak{}\par};
\end{scope}\end{tikzpicture}\newpage
\definecolor{answer}{rgb}{0.850980,0.047059,0.070588}
\null\begin{tikzpicture}[remember picture,overlay]\begin{scope}[shift={(current page.north west)},x=1bp,y=-1bp]
\node[anchor=north west,inner sep=0pt,text=answer,align=left,text width=234.0bp] at (15,15) {\fontsize{7.5bp}{9.5bp}\selectfont 析：\allowbreak{}佳佳沿垂线距离最短，\allowbreak{}原站法不公平；\allowbreak{}公平站法应等距，\allowbreak{}可都站在以凳子为圆心的同一圆周上。\allowbreak{}\par};
\end{scope}\end{tikzpicture}\newpage
\definecolor{answer}{rgb}{0.850980,0.047059,0.070588}
\null\begin{tikzpicture}[remember picture,overlay]\begin{scope}[shift={(current page.north west)},x=1bp,y=-1bp]
\node[anchor=north west,inner sep=0pt,text=answer,align=left,text width=225.6666717529297bp] at (15,15) {\fontsize{10.0bp}{12.5bp}\selectfont 这些垂线段的长度相等。\allowbreak{}两条平行线间的距离处处相等。\allowbreak{}\par};
\end{scope}\end{tikzpicture}\newpage
\definecolor{answer}{rgb}{0.850980,0.047059,0.070588}
\null\begin{tikzpicture}[remember picture,overlay]\begin{scope}[shift={(current page.north west)},x=1bp,y=-1bp]
\node[anchor=north west,inner sep=0pt,text=answer,align=left,text width=234.0bp] at (15,15) {\fontsize{7.5bp}{9.5bp}\selectfont 析：\allowbreak{}平行线间处处等距，\allowbreak{}所以各条垂直线段等长。\allowbreak{}\par};
\end{scope}\end{tikzpicture}\newpage
\definecolor{answer}{rgb}{0.850980,0.047059,0.070588}
\null\begin{tikzpicture}[remember picture,overlay]\begin{scope}[shift={(current page.north west)},x=1bp,y=-1bp]
\node[anchor=north west,inner sep=0pt,text=answer,align=left,text width=229.0bp] at (15,15) {\fontsize{8.5bp}{10.62bp}\selectfont 求单价，\allowbreak{}应用总价除以数量。\allowbreak{}单价=总价÷数量。\allowbreak{}\par};
\end{scope}\end{tikzpicture}\newpage
\definecolor{answer}{rgb}{0.850980,0.047059,0.070588}
\null\begin{tikzpicture}[remember picture,overlay]\begin{scope}[shift={(current page.north west)},x=1bp,y=-1bp]
\node[anchor=north west,inner sep=0pt,text=answer,align=left,text width=229.0bp] at (15,15) {\fontsize{8.5bp}{10.62bp}\selectfont 求数量，\allowbreak{}应用总价除以单价。\allowbreak{}数量=总价÷单价。\allowbreak{}\par};
\end{scope}\end{tikzpicture}\newpage
\definecolor{answer}{rgb}{0.850980,0.047059,0.070588}
\null\begin{tikzpicture}[remember picture,overlay]\begin{scope}[shift={(current page.north west)},x=1bp,y=-1bp]
\node[anchor=north west,inner sep=0pt,text=answer,align=left,text width=239.0bp] at (15,15) {\fontsize{8.5bp}{10.62bp}\selectfont 已知每罐\ensuremath{63}元、\allowbreak{}买\ensuremath{5}罐，\allowbreak{}求总价，\allowbreak{}用单价乘数量。\allowbreak{}数量关系：\allowbreak{}总价=单价×数量。\allowbreak{}\par};
\end{scope}\end{tikzpicture}\newpage
\definecolor{answer}{rgb}{0.850980,0.047059,0.070588}
\null\begin{tikzpicture}[remember picture,overlay]\begin{scope}[shift={(current page.north west)},x=1bp,y=-1bp]
\node[anchor=north west,inner sep=0pt,text=answer,align=left,text width=239.0bp] at (15,15) {\fontsize{8.5bp}{10.62bp}\selectfont 检验：\allowbreak{}\ensuremath{315} ÷ \ensuremath{5} = \ensuremath{63}(元)，\allowbreak{}与已知单价相同。\allowbreak{}\par};
\end{scope}\end{tikzpicture}\newpage
\definecolor{answer}{rgb}{0.850980,0.047059,0.070588}
\null\begin{tikzpicture}[remember picture,overlay]\begin{scope}[shift={(current page.north west)},x=1bp,y=-1bp]
\node[anchor=north west,inner sep=0pt,text=answer,align=left,text width=226.3333282470703bp] at (15,15) {\fontsize{9.0bp}{11.25bp}\selectfont \ensuremath{315} ÷ \ensuremath{5} = \ensuremath{63}(元),\allowbreak{}与已知单价相同。\allowbreak{}\par};
\end{scope}\end{tikzpicture}\newpage
\definecolor{answer}{rgb}{0.850980,0.047059,0.070588}
\null\begin{tikzpicture}[remember picture,overlay]\begin{scope}[shift={(current page.north west)},x=1bp,y=-1bp]
\node[anchor=north west,inner sep=0pt,text=answer,align=left,text width=229.0bp] at (15,15) {\fontsize{8.5bp}{10.62bp}\selectfont 每副\ensuremath{36}元，\allowbreak{}共\ensuremath{15}副，\allowbreak{}用单价乘数量。\allowbreak{}\par};
\end{scope}\end{tikzpicture}\newpage
\definecolor{answer}{rgb}{0.850980,0.047059,0.070588}
\null\begin{tikzpicture}[remember picture,overlay]\begin{scope}[shift={(current page.north west)},x=1bp,y=-1bp]
\node[anchor=north west,inner sep=0pt,text=answer,align=left,text width=235.0bp] at (15,15) {\fontsize{7.5bp}{9.5bp}\selectfont 析：\allowbreak{}按总价=单价×数量：\allowbreak{}\ensuremath{268} × \ensuremath{17} = \ensuremath{4556};\allowbreak{}\ensuremath{882} ÷ \ensuremath{98} = \ensuremath{9};\allowbreak{}\ensuremath{900} ÷\ensuremath{12} = \ensuremath{75}。\allowbreak{}\par};
\end{scope}\end{tikzpicture}\newpage
\definecolor{answer}{rgb}{0.850980,0.047059,0.070588}
\null\begin{tikzpicture}[remember picture,overlay]\begin{scope}[shift={(current page.north west)},x=1bp,y=-1bp]
\node[anchor=north west,inner sep=0pt,text=answer,align=left,text width=229.0bp] at (15,15) {\fontsize{8.5bp}{10.62bp}\selectfont 答：\allowbreak{}买\ensuremath{15}副中国象棋需要\ensuremath{540}元。\allowbreak{}\par};
\end{scope}\end{tikzpicture}\newpage
\definecolor{answer}{rgb}{0.850980,0.047059,0.070588}
\null\begin{tikzpicture}[remember picture,overlay]\begin{scope}[shift={(current page.north west)},x=1bp,y=-1bp]
\node[anchor=north west,inner sep=0pt,text=answer,align=left,text width=234.0bp] at (15,15) {\fontsize{8.5bp}{10.62bp}\selectfont 每盒\ensuremath{20}支，\allowbreak{}共\ensuremath{12}盒，\allowbreak{}用每盒支数乘盒数。\allowbreak{}\par};
\end{scope}\end{tikzpicture}\newpage
\definecolor{answer}{rgb}{0.850980,0.047059,0.070588}
\null\begin{tikzpicture}[remember picture,overlay]\begin{scope}[shift={(current page.north west)},x=1bp,y=-1bp]
\node[anchor=north west,inner sep=0pt,text=answer,align=left,text width=234.0bp] at (15,15) {\fontsize{8.5bp}{10.62bp}\selectfont 答：\allowbreak{}刘老师一共买了\ensuremath{240}支钢笔。\allowbreak{}\par};
\end{scope}\end{tikzpicture}\newpage
\definecolor{answer}{rgb}{0.850980,0.047059,0.070588}
\null\begin{tikzpicture}[remember picture,overlay]\begin{scope}[shift={(current page.north west)},x=1bp,y=-1bp]
\node[anchor=north west,inner sep=0pt,text=answer,align=left,text width=234.0bp] at (15,15) {\fontsize{8.5bp}{10.62bp}\selectfont 每盒\ensuremath{80}元，\allowbreak{}共\ensuremath{12}盒；\allowbreak{}本题按整盒单价计算，\allowbreak{}不用再乘每盒支数。\allowbreak{}\par};
\end{scope}\end{tikzpicture}\newpage
\definecolor{answer}{rgb}{0.850980,0.047059,0.070588}
\null\begin{tikzpicture}[remember picture,overlay]\begin{scope}[shift={(current page.north west)},x=1bp,y=-1bp]
\node[anchor=north west,inner sep=0pt,text=answer,align=left,text width=229.0bp] at (15,15) {\fontsize{8.5bp}{10.62bp}\selectfont 每个模型\ensuremath{80}元，\allowbreak{}总价随数量成倍增加；\allowbreak{}分别用\ensuremath{80}乘\ensuremath{2}、\allowbreak{}\ensuremath{3}、\allowbreak{}\ensuremath{4}、\allowbreak{}\ensuremath{5}、\allowbreak{}\ensuremath{6}。\allowbreak{}\par};
\end{scope}\end{tikzpicture}\newpage
\definecolor{answer}{rgb}{0.850980,0.047059,0.070588}
\null\begin{tikzpicture}[remember picture,overlay]\begin{scope}[shift={(current page.north west)},x=1bp,y=-1bp]
\node[anchor=north west,inner sep=0pt,text=answer,align=left,text width=229.0bp] at (15,15) {\fontsize{8.5bp}{10.62bp}\selectfont 答：\allowbreak{}买\ensuremath{2}、\allowbreak{}\ensuremath{3}、\allowbreak{}\ensuremath{4}、\allowbreak{}\ensuremath{5}、\allowbreak{}\ensuremath{6}个依次需\ensuremath{160}、\allowbreak{}\ensuremath{240}、\allowbreak{}\ensuremath{320}、\allowbreak{}\ensuremath{400}、\allowbreak{}\ensuremath{480}元。\allowbreak{}\par};
\end{scope}\end{tikzpicture}\newpage
\definecolor{answer}{rgb}{0.850980,0.047059,0.070588}
\null\begin{tikzpicture}[remember picture,overlay]\begin{scope}[shift={(current page.north west)},x=1bp,y=-1bp]
\node[anchor=north west,inner sep=0pt,text=answer,align=left,text width=234.0bp] at (15,15) {\fontsize{8.5bp}{10.62bp}\selectfont 每筐\ensuremath{30}千克，\allowbreak{}共\ensuremath{18}筐，\allowbreak{}用每筐质量乘筐数。\allowbreak{}\par};
\end{scope}\end{tikzpicture}\newpage
\definecolor{answer}{rgb}{0.850980,0.047059,0.070588}
\null\begin{tikzpicture}[remember picture,overlay]\begin{scope}[shift={(current page.north west)},x=1bp,y=-1bp]
\node[anchor=north west,inner sep=0pt,text=answer,align=left,text width=215.73748779296875bp] at (15,15) {\fontsize{7.5bp}{9.5bp}\selectfont 析：\allowbreak{}三题均为“每份数×份数=总数”；\allowbreak{}每份数对应单价，\allowbreak{}份数对应数量，\allowbreak{}总数对应总价。\allowbreak{}\par};
\end{scope}\end{tikzpicture}\newpage
\definecolor{answer}{rgb}{0.850980,0.047059,0.070588}
\null\begin{tikzpicture}[remember picture,overlay]\begin{scope}[shift={(current page.north west)},x=1bp,y=-1bp]
\node[anchor=north west,inner sep=0pt,text=answer,align=left,text width=234.0bp] at (15,15) {\fontsize{8.5bp}{10.62bp}\selectfont 答：\allowbreak{}一共采购了\ensuremath{540}千克苹果。\allowbreak{}\par};
\end{scope}\end{tikzpicture}\newpage
\definecolor{answer}{rgb}{0.850980,0.047059,0.070588}
\null\begin{tikzpicture}[remember picture,overlay]\begin{scope}[shift={(current page.north west)},x=1bp,y=-1bp]
\node[anchor=north west,inner sep=0pt,text=answer,align=left,text width=234.0bp] at (15,15) {\fontsize{8.5bp}{10.62bp}\selectfont 每分钟洒水\ensuremath{360}米，\allowbreak{}\ensuremath{5}分钟洒水长度用\ensuremath{360}乘\ensuremath{5}。\allowbreak{}\par};
\end{scope}\end{tikzpicture}\newpage
\definecolor{answer}{rgb}{0.850980,0.047059,0.070588}
\null\begin{tikzpicture}[remember picture,overlay]\begin{scope}[shift={(current page.north west)},x=1bp,y=-1bp]
\node[anchor=north west,inner sep=0pt,text=answer,align=left,text width=229.0bp] at (15,15) {\fontsize{8.5bp}{10.62bp}\selectfont 先求\ensuremath{200}块地砖总面积，\allowbreak{}再把平方分米换成平方米。\allowbreak{}\par};
\end{scope}\end{tikzpicture}\newpage
\definecolor{answer}{rgb}{0.850980,0.047059,0.070588}
\null\begin{tikzpicture}[remember picture,overlay]\begin{scope}[shift={(current page.north west)},x=1bp,y=-1bp]
\node[anchor=north west,inner sep=0pt,text=answer,align=left,text width=234.00000762939453bp] at (15,15) {\fontsize{8.5bp}{10.62bp}\selectfont 每公顷一年吸收\ensuremath{15}吨，\allowbreak{}总吸收量\ensuremath{945}吨，\allowbreak{}面积用总吸收量除以每公顷吸收量。\allowbreak{}\par};
\end{scope}\end{tikzpicture}\newpage
\definecolor{answer}{rgb}{0.850980,0.047059,0.070588}
\null\begin{tikzpicture}[remember picture,overlay]\begin{scope}[shift={(current page.north west)},x=1bp,y=-1bp]
\node[anchor=north west,inner sep=0pt,text=answer,align=left,text width=234.0bp] at (15,15) {\fontsize{8.5bp}{10.62bp}\selectfont 答：\allowbreak{}这个湿地公园面积约\ensuremath{63}公顷。\allowbreak{}\par};
\end{scope}\end{tikzpicture}\newpage
\definecolor{answer}{rgb}{0.850980,0.047059,0.070588}
\null\begin{tikzpicture}[remember picture,overlay]\begin{scope}[shift={(current page.north west)},x=1bp,y=-1bp]
\node[anchor=north west,inner sep=0pt,text=answer,align=left,text width=229.0bp] at (15,15) {\fontsize{8.5bp}{10.62bp}\selectfont 求速度，\allowbreak{}应用路程除以时间。\allowbreak{}\par};
\end{scope}\end{tikzpicture}\newpage
\definecolor{answer}{rgb}{0.850980,0.047059,0.070588}
\null\begin{tikzpicture}[remember picture,overlay]\begin{scope}[shift={(current page.north west)},x=1bp,y=-1bp]
\node[anchor=north west,inner sep=0pt,text=answer,align=left,text width=229.0bp] at (15,15) {\fontsize{8.5bp}{10.62bp}\selectfont 答：\allowbreak{}货车的行驶速度是\ensuremath{72}千米/时。\allowbreak{}\par};
\end{scope}\end{tikzpicture}\newpage
\definecolor{answer}{rgb}{0.850980,0.047059,0.070588}
\null\begin{tikzpicture}[remember picture,overlay]\begin{scope}[shift={(current page.north west)},x=1bp,y=-1bp]
\node[anchor=north west,inner sep=0pt,text=answer,align=left,text width=224.0bp] at (15,15) {\fontsize{7.5bp}{9.5bp}\selectfont 析：\allowbreak{}路程=速度×时间：\allowbreak{}\ensuremath{140} ÷ \ensuremath{2} = \ensuremath{70}(千米/时)；\allowbreak{}\ensuremath{320} ÷ \ensuremath{80} =\ensuremath{4}(时)；\allowbreak{}\ensuremath{75} × \ensuremath{5} = \ensuremath{375}(千米)。\allowbreak{}\par};
\end{scope}\end{tikzpicture}\newpage
\definecolor{answer}{rgb}{0.850980,0.047059,0.070588}
\null\begin{tikzpicture}[remember picture,overlay]\begin{scope}[shift={(current page.north west)},x=1bp,y=-1bp]
\node[anchor=north west,inner sep=0pt,text=answer,align=left,text width=234.0bp] at (15,15) {\fontsize{8.5bp}{10.62bp}\selectfont 甲船速度\ensuremath{20}千米/时，\allowbreak{}乙船\ensuremath{18}千米/时，\allowbreak{}甲船速度较大。\allowbreak{}\par};
\end{scope}\end{tikzpicture}\newpage
\definecolor{answer}{rgb}{0.850980,0.047059,0.070588}
\null\begin{tikzpicture}[remember picture,overlay]\begin{scope}[shift={(current page.north west)},x=1bp,y=-1bp]
\node[anchor=north west,inner sep=0pt,text=answer,align=left,text width=234.0bp] at (15,15) {\fontsize{8.5bp}{10.62bp}\selectfont \ensuremath{15}小时内甲船走\ensuremath{300}千米、\allowbreak{}乙船走\ensuremath{270}千米，\allowbreak{}甲船更远。\allowbreak{}\par};
\end{scope}\end{tikzpicture}\newpage
\definecolor{answer}{rgb}{0.850980,0.047059,0.070588}
\null\begin{tikzpicture}[remember picture,overlay]\begin{scope}[shift={(current page.north west)},x=1bp,y=-1bp]
\node[anchor=north west,inner sep=0pt,text=answer,align=left,text width=233.99999237060547bp] at (15,15) {\fontsize{8.5bp}{10.62bp}\selectfont 法\ensuremath{3} :\allowbreak{}比较航行同样\ensuremath{180}千米所用时间\par};
\end{scope}\end{tikzpicture}\newpage
\definecolor{answer}{rgb}{0.850980,0.047059,0.070588}
\null\begin{tikzpicture}[remember picture,overlay]\begin{scope}[shift={(current page.north west)},x=1bp,y=-1bp]
\node[anchor=north west,inner sep=0pt,text=answer,align=left,text width=229.0bp] at (15,15) {\fontsize{8.5bp}{10.62bp}\selectfont 答：\allowbreak{}宇宙飞船\ensuremath{1}分钟行驶\ensuremath{720}千米。\allowbreak{}\par};
\end{scope}\end{tikzpicture}\newpage
\definecolor{answer}{rgb}{0.850980,0.047059,0.070588}
\null\begin{tikzpicture}[remember picture,overlay]\begin{scope}[shift={(current page.north west)},x=1bp,y=-1bp]
\node[anchor=north west,inner sep=0pt,text=answer,align=left,text width=234.0bp] at (15,15) {\fontsize{8.5bp}{10.62bp}\selectfont 相同\ensuremath{180}千米，\allowbreak{}甲船用\ensuremath{9}小时，\allowbreak{}乙船用\ensuremath{10}小时，\allowbreak{}甲船少用\ensuremath{1}小时。\allowbreak{}\par};
\end{scope}\end{tikzpicture}\newpage
\definecolor{answer}{rgb}{0.850980,0.047059,0.070588}
\null\begin{tikzpicture}[remember picture,overlay]\begin{scope}[shift={(current page.north west)},x=1bp,y=-1bp]
\node[anchor=north west,inner sep=0pt,text=answer,align=left,text width=229.0bp] at (15,15) {\fontsize{10.0bp}{12.5bp}\selectfont 比较快慢要比较相同时间内的路程,\allowbreak{}也就是比较速度。\allowbreak{}速度= 路程÷ 时间。\allowbreak{}\par};
\end{scope}\end{tikzpicture}\newpage
\definecolor{answer}{rgb}{0.850980,0.047059,0.070588}
\null\begin{tikzpicture}[remember picture,overlay]\begin{scope}[shift={(current page.north west)},x=1bp,y=-1bp]
\node[anchor=north west,inner sep=0pt,text=answer,align=left,text width=229.0bp] at (15,15) {\fontsize{8.5bp}{10.62bp}\selectfont 每小时生产\ensuremath{153}箱，\allowbreak{}\ensuremath{5}小时总产量用工作效率乘工作时间。\allowbreak{}\par};
\end{scope}\end{tikzpicture}\newpage
\definecolor{answer}{rgb}{0.850980,0.047059,0.070588}
\null\begin{tikzpicture}[remember picture,overlay]\begin{scope}[shift={(current page.north west)},x=1bp,y=-1bp]
\node[anchor=north west,inner sep=0pt,text=answer,align=left,text width=234.0bp] at (15,15) {\fontsize{8.5bp}{10.62bp}\selectfont 声音从船到暗礁再返回，\allowbreak{}\ensuremath{2}秒为往返时间；\allowbreak{}到暗礁的单程路程是总路程的一半。\allowbreak{}\par};
\end{scope}\end{tikzpicture}\newpage
\definecolor{answer}{rgb}{0.850980,0.047059,0.070588}
\null\begin{tikzpicture}[remember picture,overlay]\begin{scope}[shift={(current page.north west)},x=1bp,y=-1bp]
\node[anchor=north west,inner sep=0pt,text=answer,align=left,text width=233.99999237060547bp] at (15,15) {\fontsize{8.5bp}{10.62bp}\selectfont 答：\allowbreak{}暗礁到探测船的距离是\ensuremath{1530}米。\allowbreak{}\par};
\end{scope}\end{tikzpicture}\newpage
\definecolor{answer}{rgb}{0.850980,0.047059,0.070588}
\null\begin{tikzpicture}[remember picture,overlay]\begin{scope}[shift={(current page.north west)},x=1bp,y=-1bp]
\node[anchor=north west,inner sep=0pt,text=answer,align=left,text width=229.0bp] at (15,15) {\fontsize{8.5bp}{10.62bp}\selectfont 按每月\ensuremath{30}天计算，\allowbreak{}总页数是每天页数乘天数。\allowbreak{}\par};
\end{scope}\end{tikzpicture}\newpage
\definecolor{answer}{rgb}{0.850980,0.047059,0.070588}
\null\begin{tikzpicture}[remember picture,overlay]\begin{scope}[shift={(current page.north west)},x=1bp,y=-1bp]
\node[anchor=north west,inner sep=0pt,text=answer,align=left,text width=234.0bp] at (15,15) {\fontsize{8.5bp}{10.62bp}\selectfont 已知路程和速度，\allowbreak{}时间用路程除以速度。\allowbreak{}\par};
\end{scope}\end{tikzpicture}\newpage
\definecolor{answer}{rgb}{0.850980,0.047059,0.070588}
\null\begin{tikzpicture}[remember picture,overlay]\begin{scope}[shift={(current page.north west)},x=1bp,y=-1bp]
\node[anchor=north west,inner sep=0pt,text=answer,align=left,text width=229.0bp] at (15,15) {\fontsize{8.5bp}{10.62bp}\selectfont 每分钟\ensuremath{105}个，\allowbreak{}\ensuremath{4}分钟打字总数用工作效率乘时间。\allowbreak{}\par};
\end{scope}\end{tikzpicture}\newpage
\definecolor{answer}{rgb}{0.850980,0.047059,0.070588}
\null\begin{tikzpicture}[remember picture,overlay]\begin{scope}[shift={(current page.north west)},x=1bp,y=-1bp]
\node[anchor=north west,inner sep=0pt,text=answer,align=left,text width=229.0bp] at (15,15) {\fontsize{8.5bp}{10.62bp}\selectfont 图书馆到学校\ensuremath{840}米，\allowbreak{}天天速度\ensuremath{70}米/分，\allowbreak{}用路程除以速度求时间。\allowbreak{}\par};
\end{scope}\end{tikzpicture}\newpage
\definecolor{answer}{rgb}{0.850980,0.047059,0.070588}
\null\begin{tikzpicture}[remember picture,overlay]\begin{scope}[shift={(current page.north west)},x=1bp,y=-1bp]
\node[anchor=north west,inner sep=0pt,text=answer,align=left,text width=227.0bp] at (15,15) {\fontsize{7.5bp}{9.5bp}\selectfont 析：\allowbreak{}工作效率×工作时间=工作总量；\allowbreak{}对应“速度×时间=路程”，\allowbreak{}数量关系结构相同。\allowbreak{}\par};
\end{scope}\end{tikzpicture}\newpage
\definecolor{answer}{rgb}{0.850980,0.047059,0.070588}
\null\begin{tikzpicture}[remember picture,overlay]\begin{scope}[shift={(current page.north west)},x=1bp,y=-1bp]
\node[anchor=north west,inner sep=0pt,text=answer,align=left,text width=229.0bp] at (15,15) {\fontsize{8.5bp}{10.62bp}\selectfont 开放提问，\allowbreak{}只要利用图中信息、\allowbreak{}条件充分且数量关系正确即可；\allowbreak{}下面列两个不同问题示例。\allowbreak{}\par};
\end{scope}\end{tikzpicture}\newpage
\definecolor{answer}{rgb}{0.850980,0.047059,0.070588}
\null\begin{tikzpicture}[remember picture,overlay]\begin{scope}[shift={(current page.north west)},x=1bp,y=-1bp]
\node[anchor=north west,inner sep=0pt,text=answer,align=left,text width=234.0bp] at (15,15) {\fontsize{8.5bp}{10.62bp}\selectfont 文文步行速度\ensuremath{65}米/分，\allowbreak{}走\ensuremath{12}分钟，\allowbreak{}用速度乘时间求家到学校的距离。\allowbreak{}\par};
\end{scope}\end{tikzpicture}\newpage
\definecolor{answer}{rgb}{0.850980,0.047059,0.070588}
\null\begin{tikzpicture}[remember picture,overlay]\begin{scope}[shift={(current page.north west)},x=1bp,y=-1bp]
\node[anchor=north west,inner sep=0pt,text=answer,align=left,text width=229.0bp] at (15,15) {\fontsize{8.5bp}{10.62bp}\selectfont 从第一天\ensuremath{8} :\allowbreak{} \ensuremath{30}到次日\ensuremath{8} :\allowbreak{} \ensuremath{30}是\ensuremath{24}小时，\allowbreak{}再到次日\ensuremath{9} :\allowbreak{} \ensuremath{30}多\ensuremath{1}小时，\allowbreak{}共\ensuremath{25}小时。\allowbreak{}\par};
\end{scope}\end{tikzpicture}\newpage
\definecolor{answer}{rgb}{0.850980,0.047059,0.070588}
\null\begin{tikzpicture}[remember picture,overlay]\begin{scope}[shift={(current page.north west)},x=1bp,y=-1bp]
\node[anchor=north west,inner sep=0pt,text=answer,align=left,text width=234.0bp] at (15,15) {\fontsize{8.5bp}{10.62bp}\selectfont 答：\allowbreak{}文文家与学校相距\ensuremath{780}米。\allowbreak{}\par};
\end{scope}\end{tikzpicture}\newpage
\definecolor{answer}{rgb}{0.850980,0.047059,0.070588}
\null\begin{tikzpicture}[remember picture,overlay]\begin{scope}[shift={(current page.north west)},x=1bp,y=-1bp]
\node[anchor=north west,inner sep=0pt,text=answer,align=left,text width=229.0bp] at (15,15) {\fontsize{8.5bp}{10.62bp}\selectfont 答：\allowbreak{}甲、\allowbreak{}乙两港相距\ensuremath{1000}千米。\allowbreak{}\par};
\end{scope}\end{tikzpicture}\newpage
\definecolor{answer}{rgb}{0.850980,0.047059,0.070588}
\null\begin{tikzpicture}[remember picture,overlay]\begin{scope}[shift={(current page.north west)},x=1bp,y=-1bp]
\node[anchor=north west,inner sep=0pt,text=answer,align=left,text width=229.0bp] at (15,15) {\fontsize{8.5bp}{10.62bp}\selectfont 两家在栈桥两侧，\allowbreak{}同走\ensuremath{6}分钟，\allowbreak{}可用速度和乘时间。\allowbreak{}\par};
\end{scope}\end{tikzpicture}\newpage
\definecolor{answer}{rgb}{0.850980,0.047059,0.070588}
\null\begin{tikzpicture}[remember picture,overlay]\begin{scope}[shift={(current page.north west)},x=1bp,y=-1bp]
\node[anchor=north west,inner sep=0pt,text=answer,align=left,text width=223.6666717529297bp] at (15,15) {\fontsize{10.0bp}{12.5bp}\selectfont 学校在两家之间;\allowbreak{} 刘华\ensuremath{60}米/分,\allowbreak{} 走\ensuremath{4}分;\allowbreak{}李丽\ensuremath{75}米/分,\allowbreak{} 走\ensuremath{5}分。\allowbreak{}求两家总距离。\allowbreak{}\par};
\end{scope}\end{tikzpicture}\newpage
\definecolor{answer}{rgb}{0.850980,0.047059,0.070588}
\null\begin{tikzpicture}[remember picture,overlay]\begin{scope}[shift={(current page.north west)},x=1bp,y=-1bp]
\node[anchor=north west,inner sep=0pt,text=answer,align=left,text width=234.0bp] at (15,15) {\fontsize{8.5bp}{10.62bp}\selectfont 学校在两家之间，\allowbreak{}总距离为两家各自到学校的距离之和；\allowbreak{}两人所走时间不同，\allowbreak{}应分别计算。\allowbreak{}\par};
\end{scope}\end{tikzpicture}\newpage
\definecolor{answer}{rgb}{0.850980,0.047059,0.070588}
\null\begin{tikzpicture}[remember picture,overlay]\begin{scope}[shift={(current page.north west)},x=1bp,y=-1bp]
\node[anchor=north west,inner sep=0pt,text=answer,align=left,text width=229.0bp] at (15,15) {\fontsize{8.5bp}{10.62bp}\selectfont 先分别求两个教室能展出的幅数，\allowbreak{}再相加。\allowbreak{}\par};
\end{scope}\end{tikzpicture}\newpage
\definecolor{answer}{rgb}{0.850980,0.047059,0.070588}
\null\begin{tikzpicture}[remember picture,overlay]\begin{scope}[shift={(current page.north west)},x=1bp,y=-1bp]
\node[anchor=north west,inner sep=0pt,text=answer,align=left,text width=229.0bp] at (15,15) {\fontsize{8.5bp}{10.62bp}\selectfont 法\ensuremath{1} :\allowbreak{}分别求两个教室展出的幅数\par};
\end{scope}\end{tikzpicture}\newpage
\definecolor{answer}{rgb}{0.850980,0.047059,0.070588}
\null\begin{tikzpicture}[remember picture,overlay]\begin{scope}[shift={(current page.north west)},x=1bp,y=-1bp]
\node[anchor=north west,inner sep=0pt,text=answer,align=left,text width=234.0bp] at (15,15) {\fontsize{8.5bp}{10.62bp}\selectfont 法\ensuremath{2} :\allowbreak{}统一时间后补足或扣除多算路程\par};
\end{scope}\end{tikzpicture}\newpage
\definecolor{answer}{rgb}{0.850980,0.047059,0.070588}
\null\begin{tikzpicture}[remember picture,overlay]\begin{scope}[shift={(current page.north west)},x=1bp,y=-1bp]
\node[anchor=north west,inner sep=0pt,text=answer,align=left,text width=229.0bp] at (15,15) {\fontsize{8.5bp}{10.62bp}\selectfont \ensuremath{7}块展板先都按\ensuremath{32}幅算，\allowbreak{}第二个教室\ensuremath{3}块每块再补\ensuremath{8}幅。\allowbreak{}\par};
\end{scope}\end{tikzpicture}\newpage
\definecolor{answer}{rgb}{0.850980,0.047059,0.070588}
\null\begin{tikzpicture}[remember picture,overlay]\begin{scope}[shift={(current page.north west)},x=1bp,y=-1bp]
\node[anchor=north west,inner sep=0pt,text=answer,align=left,text width=234.0bp] at (15,15) {\fontsize{8.5bp}{10.62bp}\selectfont 先按两人都走\ensuremath{4}分钟算，\allowbreak{}再补李丽多走的\ensuremath{1}分钟。\allowbreak{}也可先按两人都走\ensuremath{5}分钟算，\allowbreak{}再扣刘华多算的\ensuremath{1}分钟；\allowbreak{}属于同一补差思路。\allowbreak{}\par};
\end{scope}\end{tikzpicture}\newpage
\definecolor{answer}{rgb}{0.850980,0.047059,0.070588}
\null\begin{tikzpicture}[remember picture,overlay]\begin{scope}[shift={(current page.north west)},x=1bp,y=-1bp]
\node[anchor=north west,inner sep=0pt,text=answer,align=left,text width=229.0bp] at (15,15) {\fontsize{8.5bp}{10.62bp}\selectfont 每人一件上衣、\allowbreak{}一条裙子，\allowbreak{}先求每套价格再乘\ensuremath{40}人。\allowbreak{}\par};
\end{scope}\end{tikzpicture}\newpage
\definecolor{answer}{rgb}{0.850980,0.047059,0.070588}
\null\begin{tikzpicture}[remember picture,overlay]\begin{scope}[shift={(current page.north west)},x=1bp,y=-1bp]
\node[anchor=north west,inner sep=0pt,text=answer,align=left,text width=223.6666717529297bp] at (15,15) {\fontsize{10.0bp}{12.5bp}\selectfont 先分别求两人各自走的路程,\allowbreak{} 再相加。\allowbreak{}两人所走时间不同,\allowbreak{} 要分别计算。\allowbreak{}\par};
\end{scope}\end{tikzpicture}\newpage
\definecolor{answer}{rgb}{0.850980,0.047059,0.070588}
\null\begin{tikzpicture}[remember picture,overlay]\begin{scope}[shift={(current page.north west)},x=1bp,y=-1bp]
\node[anchor=north west,inner sep=0pt,text=answer,align=left,text width=234.0bp] at (15,15) {\fontsize{8.5bp}{10.62bp}\selectfont 分别求\ensuremath{3}辆大客车与\ensuremath{2}辆小客车最多载的人数再相加。\allowbreak{}\par};
\end{scope}\end{tikzpicture}\newpage
\definecolor{answer}{rgb}{0.850980,0.047059,0.070588}
\null\begin{tikzpicture}[remember picture,overlay]\begin{scope}[shift={(current page.north west)},x=1bp,y=-1bp]
\node[anchor=north west,inner sep=0pt,text=answer,align=left,text width=235.0009765625bp] at (15,15) {\fontsize{8.5bp}{10.62bp}\selectfont 从同一地点相背而行，\allowbreak{}两人间距为各自离出发点路程之和。\allowbreak{}\par};
\end{scope}\end{tikzpicture}\newpage
\definecolor{answer}{rgb}{0.850980,0.047059,0.070588}
\null\begin{tikzpicture}[remember picture,overlay]\begin{scope}[shift={(current page.north west)},x=1bp,y=-1bp]
\node[anchor=north west,inner sep=0pt,text=answer,align=left,text width=234.0bp] at (15,15) {\fontsize{8.5bp}{10.62bp}\selectfont 两辆小客车每辆比大客车少\ensuremath{9}人，\allowbreak{}要扣去\ensuremath{18}人。\allowbreak{}\par};
\end{scope}\end{tikzpicture}\newpage
\definecolor{answer}{rgb}{0.850980,0.047059,0.070588}
\null\begin{tikzpicture}[remember picture,overlay]\begin{scope}[shift={(current page.north west)},x=1bp,y=-1bp]
\node[anchor=north west,inner sep=0pt,text=answer,align=left,text width=235.0009765625bp] at (15,15) {\fontsize{8.5bp}{10.62bp}\selectfont 两台收割机共同工作\ensuremath{4}小时，\allowbreak{}总面积为工作效率和乘时间。\allowbreak{}\par};
\end{scope}\end{tikzpicture}\newpage
\definecolor{answer}{rgb}{0.850980,0.047059,0.070588}
\null\begin{tikzpicture}[remember picture,overlay]\begin{scope}[shift={(current page.north west)},x=1bp,y=-1bp]
\node[anchor=north west,inner sep=0pt,text=answer,align=left,text width=229.0bp] at (15,15) {\fontsize{8.5bp}{10.62bp}\selectfont 答：\allowbreak{}一共能收割\ensuremath{47200}平方米。\allowbreak{}\par};
\end{scope}\end{tikzpicture}\newpage
\definecolor{answer}{rgb}{0.850980,0.047059,0.070588}
\null\begin{tikzpicture}[remember picture,overlay]\begin{scope}[shift={(current page.north west)},x=1bp,y=-1bp]
\node[anchor=north west,inner sep=0pt,text=answer,align=left,text width=234.0bp] at (15,15) {\fontsize{8.5bp}{10.62bp}\selectfont 前三天与后两天的每天数量不同，\allowbreak{}先分别求两个时段总量再相加。\allowbreak{}\par};
\end{scope}\end{tikzpicture}\newpage
\definecolor{answer}{rgb}{0.850980,0.047059,0.070588}
\null\begin{tikzpicture}[remember picture,overlay]\begin{scope}[shift={(current page.north west)},x=1bp,y=-1bp]
\node[anchor=north west,inner sep=0pt,text=answer,align=left,text width=234.0bp] at (15,15) {\fontsize{8.5bp}{10.62bp}\selectfont 答：\allowbreak{}五天一共回收\ensuremath{1068}个塑料瓶。\allowbreak{}\par};
\end{scope}\end{tikzpicture}\newpage
\definecolor{answer}{rgb}{0.850980,0.047059,0.070588}
\null\begin{tikzpicture}[remember picture,overlay]\begin{scope}[shift={(current page.north west)},x=1bp,y=-1bp]
\node[anchor=north west,inner sep=0pt,text=answer,align=left,text width=229.0bp] at (15,15) {\fontsize{8.5bp}{10.62bp}\selectfont 货车先行\ensuremath{2}小时又行\ensuremath{5}小时，\allowbreak{}共\ensuremath{7}小时；\allowbreak{}客车只行\ensuremath{5}小时，\allowbreak{}两段路程之和为两地距离。\allowbreak{}\par};
\end{scope}\end{tikzpicture}\newpage
\definecolor{answer}{rgb}{0.850980,0.047059,0.070588}
\null\begin{tikzpicture}[remember picture,overlay]\begin{scope}[shift={(current page.north west)},x=1bp,y=-1bp]
\node[anchor=north west,inner sep=0pt,text=answer,align=left,text width=234.0bp] at (15,15) {\fontsize{8.5bp}{10.62bp}\selectfont 每箱盒数不同，\allowbreak{}应分别用每箱盒数乘箱数，\allowbreak{}再相加。\allowbreak{}\par};
\end{scope}\end{tikzpicture}\newpage
\definecolor{answer}{rgb}{0.850980,0.047059,0.070588}
\null\begin{tikzpicture}[remember picture,overlay]\begin{scope}[shift={(current page.north west)},x=1bp,y=-1bp]
\node[anchor=north west,inner sep=0pt,text=answer,align=left,text width=229.0bp] at (15,15) {\fontsize{8.5bp}{10.62bp}\selectfont 答：\allowbreak{}甲、\allowbreak{}乙两地相距\ensuremath{703}千米。\allowbreak{}\par};
\end{scope}\end{tikzpicture}\newpage
\definecolor{answer}{rgb}{0.850980,0.047059,0.070588}
\null\begin{tikzpicture}[remember picture,overlay]\begin{scope}[shift={(current page.north west)},x=1bp,y=-1bp]
\node[anchor=north west,inner sep=0pt,text=answer,align=left,text width=234.0bp] at (15,15) {\fontsize{8.5bp}{10.62bp}\selectfont 答：\allowbreak{}希望小学每月收到\ensuremath{1800}盒牛奶。\allowbreak{}\par};
\end{scope}\end{tikzpicture}\newpage
\definecolor{answer}{rgb}{0.850980,0.047059,0.070588}
\null\begin{tikzpicture}[remember picture,overlay]\begin{scope}[shift={(current page.north west)},x=1bp,y=-1bp]
\node[anchor=north west,inner sep=0pt,text=answer,align=left,text width=229.0bp] at (15,15) {\fontsize{8.5bp}{10.62bp}\selectfont 同地相背而行，\allowbreak{}两人间距每分钟增加\ensuremath{70} + \ensuremath{60}米。\allowbreak{}\par};
\end{scope}\end{tikzpicture}\newpage
\definecolor{answer}{rgb}{0.850980,0.047059,0.070588}
\null\begin{tikzpicture}[remember picture,overlay]\begin{scope}[shift={(current page.north west)},x=1bp,y=-1bp]
\node[anchor=north west,inner sep=0pt,text=answer,align=left,text width=234.0bp] at (15,15) {\fontsize{8.5bp}{10.62bp}\selectfont 两车同时相向行驶，\allowbreak{}相遇前合走的路程等于两城间路程；\allowbreak{}用速度和乘相遇时间。\allowbreak{}\par};
\end{scope}\end{tikzpicture}\newpage
\definecolor{answer}{rgb}{0.850980,0.047059,0.070588}
\null\begin{tikzpicture}[remember picture,overlay]\begin{scope}[shift={(current page.north west)},x=1bp,y=-1bp]
\node[anchor=north west,inner sep=0pt,text=answer,align=left,text width=229.0bp] at (15,15) {\fontsize{8.5bp}{10.62bp}\selectfont 甲、\allowbreak{}乙从同一地点向相反方向修路，\allowbreak{}\ensuremath{4}天总长度是两队工作量之和。\allowbreak{}\par};
\end{scope}\end{tikzpicture}\newpage
\definecolor{answer}{rgb}{0.850980,0.047059,0.070588}
\null\begin{tikzpicture}[remember picture,overlay]\begin{scope}[shift={(current page.north west)},x=1bp,y=-1bp]
\node[anchor=north west,inner sep=0pt,text=answer,align=left,text width=234.0bp] at (15,15) {\fontsize{8.5bp}{10.62bp}\selectfont 答：\allowbreak{}甲城到乙城的公路长\ensuremath{520}千米。\allowbreak{}\par};
\end{scope}\end{tikzpicture}\newpage
\definecolor{answer}{rgb}{0.850980,0.047059,0.070588}
\null\begin{tikzpicture}[remember picture,overlay]\begin{scope}[shift={(current page.north west)},x=1bp,y=-1bp]
\node[anchor=north west,inner sep=0pt,text=answer,align=left,text width=221.6666717529297bp] at (15,15) {\fontsize{11.0bp}{13.75bp}\selectfont 两车同时相向而行,\allowbreak{}相遇前合走的路程就是两城间的路程。\allowbreak{}总路程= 速度和× 相遇时间。\allowbreak{}\par};
\end{scope}\end{tikzpicture}\newpage
\definecolor{answer}{rgb}{0.850980,0.047059,0.070588}
\null\begin{tikzpicture}[remember picture,overlay]\begin{scope}[shift={(current page.north west)},x=1bp,y=-1bp]
\node[anchor=north west,inner sep=0pt,text=answer,align=left,text width=234.0bp] at (15,15) {\fontsize{8.5bp}{10.62bp}\selectfont 两队从隧道两头同时开凿，\allowbreak{}\ensuremath{6}个月合计工作量等于隧道长度。\allowbreak{}\par};
\end{scope}\end{tikzpicture}\newpage
\definecolor{answer}{rgb}{0.850980,0.047059,0.070588}
\null\begin{tikzpicture}[remember picture,overlay]\begin{scope}[shift={(current page.north west)},x=1bp,y=-1bp]
\node[anchor=north west,inner sep=0pt,text=answer,align=left,text width=229.0bp] at (15,15) {\fontsize{8.5bp}{10.62bp}\selectfont 飞飞与乐乐从两家相向走\ensuremath{6}分钟到学校，\allowbreak{}两家距离是两人路程之和。\allowbreak{}\par};
\end{scope}\end{tikzpicture}\newpage
\definecolor{answer}{rgb}{0.850980,0.047059,0.070588}
\null\begin{tikzpicture}[remember picture,overlay]\begin{scope}[shift={(current page.north west)},x=1bp,y=-1bp]
\node[anchor=north west,inner sep=0pt,text=answer,align=left,text width=235.3585205078125bp] at (15,15) {\fontsize{8.5bp}{10.62bp}\selectfont 两人录入时间相同，\allowbreak{}字数之差等于每分钟字数差乘\ensuremath{20}分钟。\allowbreak{}\par};
\end{scope}\end{tikzpicture}\newpage
\definecolor{answer}{rgb}{0.850980,0.047059,0.070588}
\null\begin{tikzpicture}[remember picture,overlay]\begin{scope}[shift={(current page.north west)},x=1bp,y=-1bp]
\node[anchor=north west,inner sep=0pt,text=answer,align=left,text width=229.0bp] at (15,15) {\fontsize{8.5bp}{10.62bp}\selectfont 两人同时同路从学校出发，\allowbreak{}乐乐速度快，\allowbreak{}两人间距为同时间内路程差。\allowbreak{}\par};
\end{scope}\end{tikzpicture}\newpage
\definecolor{answer}{rgb}{0.850980,0.047059,0.070588}
\null\begin{tikzpicture}[remember picture,overlay]\begin{scope}[shift={(current page.north west)},x=1bp,y=-1bp]
\node[anchor=north west,inner sep=0pt,text=answer,align=left,text width=234.0bp] at (15,15) {\fontsize{8.5bp}{10.62bp}\selectfont 答：\allowbreak{}小青比小亮少录入\ensuremath{200}字。\allowbreak{}\par};
\end{scope}\end{tikzpicture}\newpage
\definecolor{answer}{rgb}{0.850980,0.047059,0.070588}
\null\begin{tikzpicture}[remember picture,overlay]\begin{scope}[shift={(current page.north west)},x=1bp,y=-1bp]
\node[anchor=north west,inner sep=0pt,text=answer,align=left,text width=209.0bp] at (15,15) {\fontsize{7.5bp}{9.5bp}\selectfont 析：\allowbreak{}工作效率、\allowbreak{}时间、\allowbreak{}工作总量分别对应速度、\allowbreak{}时间、\allowbreak{}路程；\allowbreak{}求合计用和，\allowbreak{}求相差用差。\allowbreak{}\par};
\end{scope}\end{tikzpicture}\newpage
\definecolor{answer}{rgb}{0.850980,0.047059,0.070588}
\null\begin{tikzpicture}[remember picture,overlay]\begin{scope}[shift={(current page.north west)},x=1bp,y=-1bp]
\node[anchor=north west,inner sep=0pt,text=answer,align=left,text width=235.0009765625bp] at (15,15) {\fontsize{8.5bp}{10.62bp}\selectfont 两人同时同路骑行，\allowbreak{}路程差除以速度差就是共同骑行时间。\allowbreak{}\par};
\end{scope}\end{tikzpicture}\newpage
\definecolor{answer}{rgb}{0.850980,0.047059,0.070588}
\null\begin{tikzpicture}[remember picture,overlay]\begin{scope}[shift={(current page.north west)},x=1bp,y=-1bp]
\node[anchor=north west,inner sep=0pt,text=answer,align=left,text width=229.0bp] at (15,15) {\fontsize{8.5bp}{10.62bp}\selectfont 答：\allowbreak{}此时两人已经骑行\ensuremath{3}小时。\allowbreak{}\par};
\end{scope}\end{tikzpicture}\newpage
\definecolor{answer}{rgb}{0.850980,0.047059,0.070588}
\null\begin{tikzpicture}[remember picture,overlay]\begin{scope}[shift={(current page.north west)},x=1bp,y=-1bp]
\node[anchor=north west,inner sep=0pt,text=answer,align=left,text width=234.0bp] at (15,15) {\fontsize{8.5bp}{10.62bp}\selectfont 环形跑道相背出发，\allowbreak{}第一次相遇时，\allowbreak{}两人路程之和恰为跑道一圈\ensuremath{900}米。\allowbreak{}\par};
\end{scope}\end{tikzpicture}\newpage
\definecolor{answer}{rgb}{0.850980,0.047059,0.070588}
\null\begin{tikzpicture}[remember picture,overlay]\begin{scope}[shift={(current page.north west)},x=1bp,y=-1bp]
\node[anchor=north west,inner sep=0pt,text=answer,align=left,text width=234.00000762939453bp] at (15,15) {\fontsize{8.5bp}{10.62bp}\selectfont 答：\allowbreak{}经过\ensuremath{10}分钟两人第一次相遇。\allowbreak{}\par};
\end{scope}\end{tikzpicture}\newpage
\definecolor{answer}{rgb}{0.850980,0.047059,0.070588}
\null\begin{tikzpicture}[remember picture,overlay]\begin{scope}[shift={(current page.north west)},x=1bp,y=-1bp]
\node[anchor=north west,inner sep=0pt,text=answer,align=left,text width=229.0bp] at (15,15) {\fontsize{7.5bp}{9.5bp}\selectfont 法\ensuremath{1}：\allowbreak{}先求每天合飞路程，\allowbreak{}再作有余数除法\par};
\end{scope}\end{tikzpicture}\newpage
\definecolor{answer}{rgb}{0.850980,0.047059,0.070588}
\null\begin{tikzpicture}[remember picture,overlay]\begin{scope}[shift={(current page.north west)},x=1bp,y=-1bp]
\node[anchor=north west,inner sep=0pt,text=answer,align=left,text width=229.0bp] at (15,15) {\fontsize{7.5bp}{9.5bp}\selectfont \ensuremath{22}天后还剩\ensuremath{20}千米；\allowbreak{}两鸟半天合飞\ensuremath{45}千米，\allowbreak{}剩下的路程不到半天，\allowbreak{}所以相遇时间更接近\ensuremath{22}天。\allowbreak{}\par};
\end{scope}\end{tikzpicture}\newpage
\definecolor{answer}{rgb}{0.850980,0.047059,0.070588}
\null\begin{tikzpicture}[remember picture,overlay]\begin{scope}[shift={(current page.north west)},x=1bp,y=-1bp]
\node[anchor=north west,inner sep=0pt,text=answer,align=left,text width=229.0bp] at (15,15) {\fontsize{7.5bp}{9.5bp}\selectfont \ensuremath{23}天合飞\ensuremath{2070}千米，\allowbreak{}已经超过\ensuremath{2000}千米。\allowbreak{}\ensuremath{22}天后剩\ensuremath{20}千米，\allowbreak{}不足半天合飞的\ensuremath{45}千米，\allowbreak{}所以约\ensuremath{22}天。\allowbreak{}\par};
\end{scope}\end{tikzpicture}\newpage
\definecolor{answer}{rgb}{0.850980,0.047059,0.070588}
\null\begin{tikzpicture}[remember picture,overlay]\begin{scope}[shift={(current page.north west)},x=1bp,y=-1bp]
\node[anchor=north west,inner sep=0pt,text=answer,align=left,text width=229.0bp] at (15,15) {\fontsize{7.5bp}{9.5bp}\selectfont 第二天大鼠挖\ensuremath{2}尺，\allowbreak{}小鼠挖半尺，\allowbreak{}即\ensuremath{0.5}尺。\allowbreak{}\par};
\end{scope}\end{tikzpicture}\newpage
\definecolor{answer}{rgb}{0.850980,0.047059,0.070588}
\null\begin{tikzpicture}[remember picture,overlay]\begin{scope}[shift={(current page.north west)},x=1bp,y=-1bp]
\node[anchor=north west,inner sep=0pt,text=answer,align=left,text width=229.0bp] at (15,15) {\fontsize{7.5bp}{9.5bp}\selectfont 第三天大鼠一天可挖\ensuremath{4}尺，\allowbreak{}小鼠一天可挖\ensuremath{0.25}尺。\allowbreak{}\par};
\end{scope}\end{tikzpicture}\newpage
\definecolor{answer}{rgb}{0.850980,0.047059,0.070588}
\null\begin{tikzpicture}[remember picture,overlay]\begin{scope}[shift={(current page.north west)},x=1bp,y=-1bp]
\node[anchor=north west,inner sep=0pt,text=answer,align=left,text width=229.0bp] at (15,15) {\fontsize{7.5bp}{9.5bp}\selectfont \ensuremath{4.25}尺大于剩下的\ensuremath{0.5}尺，\allowbreak{}因此在第三天结束以前就会相逢。\allowbreak{}\par};
\end{scope}\end{tikzpicture}\newpage
\definecolor{answer}{rgb}{0.850980,0.047059,0.070588}
\null\begin{tikzpicture}[remember picture,overlay]\begin{scope}[shift={(current page.north west)},x=1bp,y=-1bp]
\node[anchor=north west,inner sep=0pt,text=answer,align=left,text width=229.0bp] at (15,15) {\fontsize{7.5bp}{9.5bp}\selectfont 答：\allowbreak{}两鼠在第\ensuremath{3}天相逢，\allowbreak{}也就是\ensuremath{2}天多一点。\allowbreak{}\par};
\end{scope}\end{tikzpicture}\newpage
\definecolor{answer}{rgb}{0.850980,0.047059,0.070588}
\null\begin{tikzpicture}[remember picture,overlay]\begin{scope}[shift={(current page.north west)},x=1bp,y=-1bp]
\node[anchor=north west,inner sep=0pt,text=answer,align=left,text width=235.0bp] at (15,15) {\fontsize{7.5bp}{9.5bp}\selectfont 析：\allowbreak{}交换律只换位置：\allowbreak{}B、\allowbreak{}D;\allowbreak{}结合律只变分组：\allowbreak{}A、\allowbreak{}E;\allowbreak{}两者都用：\allowbreak{}C、\allowbreak{}F。\allowbreak{}\par};
\end{scope}\end{tikzpicture}\newpage
\definecolor{answer}{rgb}{0.850980,0.047059,0.070588}
\null\begin{tikzpicture}[remember picture,overlay]\begin{scope}[shift={(current page.north west)},x=1bp,y=-1bp]
\node[anchor=north west,inner sep=0pt,text=answer,align=left,text width=235.0bp] at (15,15) {\fontsize{7.5bp}{9.5bp}\selectfont 析：\allowbreak{}加数不变：\allowbreak{}交换律换位置，\allowbreak{}结合律换分组。\allowbreak{}各空依等式对应加数填写。\allowbreak{}\par};
\end{scope}\end{tikzpicture}\newpage
\definecolor{answer}{rgb}{0.850980,0.047059,0.070588}
\null\begin{tikzpicture}[remember picture,overlay]\begin{scope}[shift={(current page.north west)},x=1bp,y=-1bp]
\node[anchor=north west,inner sep=0pt,text=answer,align=left,text width=229.0bp] at (15,15) {\fontsize{8.5bp}{10.62bp}\selectfont 法\ensuremath{1} :\allowbreak{} \ensuremath{480} + \ensuremath{173} + \ensuremath{227} + \ensuremath{520} :\allowbreak{}凑\ensuremath{1000}与\ensuremath{400}\par};
\end{scope}\end{tikzpicture}\newpage
\definecolor{answer}{rgb}{0.850980,0.047059,0.070588}
\null\begin{tikzpicture}[remember picture,overlay]\begin{scope}[shift={(current page.north west)},x=1bp,y=-1bp]
\node[anchor=north west,inner sep=0pt,text=answer,align=left,text width=229.0bp] at (15,15) {\fontsize{8.5bp}{10.62bp}\selectfont 法\ensuremath{2} :\allowbreak{} \ensuremath{480} + \ensuremath{173} + \ensuremath{227} + \ensuremath{520} :\allowbreak{}按整百数与剩余部分分组\par};
\end{scope}\end{tikzpicture}\newpage
\definecolor{answer}{rgb}{0.850980,0.047059,0.070588}
\null\begin{tikzpicture}[remember picture,overlay]\begin{scope}[shift={(current page.north west)},x=1bp,y=-1bp]
\node[anchor=north west,inner sep=0pt,text=answer,align=left,text width=234.25bp] at (15,15) {\fontsize{8.5bp}{10.62bp}\selectfont 按原图第三段\ensuremath{153}千米计算，\allowbreak{}先把\ensuremath{238}与\ensuremath{262}凑成\ensuremath{500}\par};
\end{scope}\end{tikzpicture}\newpage
\definecolor{answer}{rgb}{0.850980,0.047059,0.070588}
\null\begin{tikzpicture}[remember picture,overlay]\begin{scope}[shift={(current page.north west)},x=1bp,y=-1bp]
\node[anchor=north west,inner sep=0pt,text=answer,align=left,text width=234.0bp] at (15,15) {\fontsize{8.5bp}{10.62bp}\selectfont 先求灯笼和春联总价，\allowbreak{}再与糖果价钱比较。\allowbreak{}\par};
\end{scope}\end{tikzpicture}\newpage
\definecolor{answer}{rgb}{0.850980,0.047059,0.070588}
\null\begin{tikzpicture}[remember picture,overlay]\begin{scope}[shift={(current page.north west)},x=1bp,y=-1bp]
\node[anchor=north west,inner sep=0pt,text=answer,align=left,text width=234.0bp] at (15,15) {\fontsize{8.5bp}{10.62bp}\selectfont 答：\allowbreak{}买糖果比买灯笼和春联的总钱数多\ensuremath{12}元。\allowbreak{}\par};
\end{scope}\end{tikzpicture}\newpage
\definecolor{answer}{rgb}{0.850980,0.047059,0.070588}
\null\begin{tikzpicture}[remember picture,overlay]\begin{scope}[shift={(current page.north west)},x=1bp,y=-1bp]
\node[anchor=north west,inner sep=0pt,text=answer,align=left,text width=232.3333282470703bp] at (15,15) {\fontsize{11.0bp}{13.75bp}\selectfont \ensuremath{a}不变,\allowbreak{} \ensuremath{b}每增加\ensuremath{10},\allowbreak{}和也增加\ensuremath{10},\allowbreak{}差就减少\ensuremath{10}。\allowbreak{}\par};
\end{scope}\end{tikzpicture}\newpage
\definecolor{answer}{rgb}{0.850980,0.047059,0.070588}
\null\begin{tikzpicture}[remember picture,overlay]\begin{scope}[shift={(current page.north west)},x=1bp,y=-1bp]
\node[anchor=north west,inner sep=0pt,text=answer,align=left,text width=235.0bp] at (15,15) {\fontsize{7.5bp}{9.5bp}\selectfont 析：\allowbreak{}\ensuremath{a} = \ensuremath{500}不变，\allowbreak{}\ensuremath{b}每次增加\ensuremath{10};\allowbreak{}和就增加\ensuremath{10},\allowbreak{}差就减少\ensuremath{10}。\allowbreak{}表内逐列用\ensuremath{500}加减\ensuremath{b}。\allowbreak{}\par};
\end{scope}\end{tikzpicture}\newpage
\definecolor{answer}{rgb}{0.850980,0.047059,0.070588}
\null\begin{tikzpicture}[remember picture,overlay]\begin{scope}[shift={(current page.north west)},x=1bp,y=-1bp]
\node[anchor=north west,inner sep=0pt,text=answer,align=left,text width=234.0bp] at (15,15) {\fontsize{7.5bp}{9.5bp}\selectfont 析：\allowbreak{}交换两因数位置，\allowbreak{}积不变：\allowbreak{}\ensuremath{a} × \ensuremath{b} = \ensuremath{b} × \ensuremath{a}。\allowbreak{}\par};
\end{scope}\end{tikzpicture}\newpage
\definecolor{answer}{rgb}{0.850980,0.047059,0.070588}
\null\begin{tikzpicture}[remember picture,overlay]\begin{scope}[shift={(current page.north west)},x=1bp,y=-1bp]
\node[anchor=north west,inner sep=0pt,text=answer,align=left,text width=234.0bp] at (15,15) {\fontsize{7.5bp}{9.5bp}\selectfont 析：\allowbreak{}只改变分组：\allowbreak{}(\ensuremath{a} × \ensuremath{b}) × \ensuremath{c} = \ensuremath{a} × (\ensuremath{b} × \ensuremath{c})。\allowbreak{}\par};
\end{scope}\end{tikzpicture}\newpage
\definecolor{answer}{rgb}{0.850980,0.047059,0.070588}
\null\begin{tikzpicture}[remember picture,overlay]\begin{scope}[shift={(current page.north west)},x=1bp,y=-1bp]
\node[anchor=north west,inner sep=0pt,text=answer,align=left,text width=234.0bp] at (15,15) {\fontsize{7.5bp}{9.5bp}\selectfont 析：\allowbreak{}因数次序不变，\allowbreak{}只把\ensuremath{5}与\ensuremath{20}先乘，\allowbreak{}是乘法结合律。\allowbreak{}\par};
\end{scope}\end{tikzpicture}\newpage
\definecolor{answer}{rgb}{0.850980,0.047059,0.070588}
\null\begin{tikzpicture}[remember picture,overlay]\begin{scope}[shift={(current page.north west)},x=1bp,y=-1bp]
\node[anchor=north west,inner sep=0pt,text=answer,align=left,text width=224.95147705078125bp] at (15,15) {\fontsize{7.5bp}{9.5bp}\selectfont 析：\allowbreak{}依次填\ensuremath{25}、\allowbreak{}\ensuremath{4};\allowbreak{} \ensuremath{28};\allowbreak{} \ensuremath{125}、\allowbreak{}\ensuremath{7};\allowbreak{}\ensuremath{19}、\allowbreak{}\ensuremath{26}。\allowbreak{}交换律变次序，\allowbreak{}结合律变分组。\allowbreak{}\par};
\end{scope}\end{tikzpicture}\newpage
\definecolor{answer}{rgb}{0.850980,0.047059,0.070588}
\null\begin{tikzpicture}[remember picture,overlay]\begin{scope}[shift={(current page.north west)},x=1bp,y=-1bp]
\node[anchor=north west,inner sep=0pt,text=answer,align=left,text width=234.0bp] at (15,15) {\fontsize{7.5bp}{9.5bp}\selectfont 析：\allowbreak{}交换\ensuremath{13}和\ensuremath{4}的位置：\allowbreak{}\ensuremath{25} × \ensuremath{4} × \ensuremath{13},\allowbreak{}乘法交换律。\allowbreak{}\par};
\end{scope}\end{tikzpicture}\newpage
\definecolor{answer}{rgb}{0.850980,0.047059,0.070588}
\null\begin{tikzpicture}[remember picture,overlay]\begin{scope}[shift={(current page.north west)},x=1bp,y=-1bp]
\node[anchor=north west,inner sep=0pt,text=answer,align=left,text width=234.0bp] at (15,15) {\fontsize{7.5bp}{9.5bp}\selectfont 析：\allowbreak{}\ensuremath{125} × \ensuremath{8} = \ensuremath{1000},\allowbreak{}先凑整，\allowbreak{}运用乘法结合律。\allowbreak{}\par};
\end{scope}\end{tikzpicture}\newpage
\definecolor{answer}{rgb}{0.850980,0.047059,0.070588}
\null\begin{tikzpicture}[remember picture,overlay]\begin{scope}[shift={(current page.north west)},x=1bp,y=-1bp]
\node[anchor=north west,inner sep=0pt,text=answer,align=left,text width=234.00000762939453bp] at (15,15) {\fontsize{7.5bp}{9.5bp}\selectfont 析：\allowbreak{}交换律连接\ensuremath{12} × \ensuremath{25}与\ensuremath{25} × \ensuremath{12}、\allowbreak{}\ensuremath{a} × \ensuremath{b}与\ensuremath{b} × \ensuremath{a};\allowbreak{}结合律连接改变三因数分组的各式。\allowbreak{}\par};
\end{scope}\end{tikzpicture}\newpage
\definecolor{answer}{rgb}{0.850980,0.047059,0.070588}
\null\begin{tikzpicture}[remember picture,overlay]\begin{scope}[shift={(current page.north west)},x=1bp,y=-1bp]
\node[anchor=north west,inner sep=0pt,text=answer,align=left,text width=225.75bp] at (15,15) {\fontsize{8.5bp}{10.62bp}\selectfont 每场收入乘每天演出场数；\allowbreak{}也可先合并一天售出的票数。\allowbreak{}\par};
\end{scope}\end{tikzpicture}\newpage
\definecolor{answer}{rgb}{0.850980,0.047059,0.070588}
\null\begin{tikzpicture}[remember picture,overlay]\begin{scope}[shift={(current page.north west)},x=1bp,y=-1bp]
\node[anchor=north west,inner sep=0pt,text=answer,align=left,text width=223.0bp] at (15,15) {\fontsize{8.5bp}{10.62bp}\selectfont 答：\allowbreak{}一天的门票收入共\ensuremath{3000}元。\allowbreak{}\par};
\end{scope}\end{tikzpicture}\newpage
\definecolor{answer}{rgb}{0.850980,0.047059,0.070588}
\null\begin{tikzpicture}[remember picture,overlay]\begin{scope}[shift={(current page.north west)},x=1bp,y=-1bp]
\node[anchor=north west,inner sep=0pt,text=answer,align=left,text width=217.0bp] at (15,15) {\fontsize{8.5bp}{10.62bp}\selectfont 平均每架每层的本数，\allowbreak{}可逐级平均分，\allowbreak{}也可除以书架总层数。\allowbreak{}\par};
\end{scope}\end{tikzpicture}\newpage
\definecolor{answer}{rgb}{0.850980,0.047059,0.070588}
\null\begin{tikzpicture}[remember picture,overlay]\begin{scope}[shift={(current page.north west)},x=1bp,y=-1bp]
\node[anchor=north west,inner sep=0pt,text=answer,align=left,text width=225.2650146484375bp] at (15,15) {\fontsize{7.5bp}{9.5bp}\selectfont 每槽\ensuremath{4}盏，\allowbreak{}每区\ensuremath{7}槽，\allowbreak{}共\ensuremath{25}区。\allowbreak{}既可按每区灯数或全部槽数计算，\allowbreak{}也可把\ensuremath{25}个区中位置对应的水培槽合成一组，\allowbreak{}先算\ensuremath{4} × \ensuremath{25}。\allowbreak{}\par};
\end{scope}\end{tikzpicture}\newpage
\definecolor{answer}{rgb}{0.850980,0.047059,0.070588}
\null\begin{tikzpicture}[remember picture,overlay]\begin{scope}[shift={(current page.north west)},x=1bp,y=-1bp]
\node[anchor=north west,inner sep=0pt,text=answer,align=left,text width=217.0bp] at (15,15) {\fontsize{8.5bp}{10.62bp}\selectfont 答：\allowbreak{}平均每个书架每层放\ensuremath{36}本书。\allowbreak{}\par};
\end{scope}\end{tikzpicture}\newpage
\definecolor{answer}{rgb}{0.850980,0.047059,0.070588}
\null\begin{tikzpicture}[remember picture,overlay]\begin{scope}[shift={(current page.north west)},x=1bp,y=-1bp]
\node[anchor=north west,inner sep=0pt,text=answer,align=left,text width=223.0bp] at (15,15) {\fontsize{8.5bp}{10.62bp}\selectfont 一至五年级是\ensuremath{5}个年级，\allowbreak{}每班\ensuremath{32}人，\allowbreak{}每个年级\ensuremath{6}个班。\allowbreak{}\par};
\end{scope}\end{tikzpicture}\newpage
\definecolor{answer}{rgb}{0.850980,0.047059,0.070588}
\null\begin{tikzpicture}[remember picture,overlay]\begin{scope}[shift={(current page.north west)},x=1bp,y=-1bp]
\node[anchor=north west,inner sep=0pt,text=answer,align=left,text width=223.0bp] at (15,15) {\fontsize{8.5bp}{10.62bp}\selectfont 答：\allowbreak{}一共有\ensuremath{960}人参加答题活动。\allowbreak{}\par};
\end{scope}\end{tikzpicture}\newpage
\definecolor{answer}{rgb}{0.850980,0.047059,0.070588}
\null\begin{tikzpicture}[remember picture,overlay]\begin{scope}[shift={(current page.north west)},x=1bp,y=-1bp]
\node[anchor=north west,inner sep=0pt,text=answer,align=left,text width=227.0bp] at (15,15) {\fontsize{7.5bp}{9.5bp}\selectfont 析：\allowbreak{}选B。\allowbreak{}两行\ensuremath{15}个点为一组，\allowbreak{}共\ensuremath{7}组，\allowbreak{}表示\ensuremath{7} × (\ensuremath{2} × \ensuremath{15})。\allowbreak{}\par};
\end{scope}\end{tikzpicture}\newpage
\definecolor{answer}{rgb}{0.850980,0.047059,0.070588}
\null\begin{tikzpicture}[remember picture,overlay]\begin{scope}[shift={(current page.north west)},x=1bp,y=-1bp]
\node[anchor=north west,inner sep=0pt,text=answer,align=left,text width=229.0bp] at (15,15) {\fontsize{7.5bp}{9.5bp}\selectfont 析：\allowbreak{}分配律：\allowbreak{}括号内每项都乘外面的因数。\allowbreak{}依次填\ensuremath{2}、\allowbreak{}\ensuremath{2};\allowbreak{} \ensuremath{6}、\allowbreak{}\ensuremath{52}、\allowbreak{}\ensuremath{48};\allowbreak{} x、\allowbreak{}y;\allowbreak{}m、\allowbreak{}\ensuremath{12}、\allowbreak{}\ensuremath{18}。\allowbreak{}\par};
\end{scope}\end{tikzpicture}\newpage
\definecolor{answer}{rgb}{0.850980,0.047059,0.070588}
\null\begin{tikzpicture}[remember picture,overlay]\begin{scope}[shift={(current page.north west)},x=1bp,y=-1bp]
\node[anchor=north west,inner sep=0pt,text=answer,align=left,text width=230.60650634765625bp] at (15,15) {\fontsize{7.5bp}{9.5bp}\selectfont 析：\allowbreak{}选D。\allowbreak{}竖式表示(\ensuremath{60} + \ensuremath{4}) × \ensuremath{3};\allowbreak{}A、\allowbreak{}B、\allowbreak{}C均表示(\ensuremath{6} + \ensuremath{4}) × \ensuremath{3}。\allowbreak{}\par};
\end{scope}\end{tikzpicture}\newpage
\definecolor{answer}{rgb}{0.850980,0.047059,0.070588}
\null\begin{tikzpicture}[remember picture,overlay]\begin{scope}[shift={(current page.north west)},x=1bp,y=-1bp]
\node[anchor=north west,inner sep=0pt,text=answer,align=left,text width=229.0bp] at (15,15) {\fontsize{7.5bp}{9.5bp}\selectfont 析：\allowbreak{}选C。\allowbreak{}提取共同因数\ensuremath{38},\allowbreak{}得\ensuremath{38} × (\ensuremath{26} + \ensuremath{74});\allowbreak{} A用结合律。\allowbreak{}\par};
\end{scope}\end{tikzpicture}\newpage
\definecolor{answer}{rgb}{0.850980,0.047059,0.070588}
\null\begin{tikzpicture}[remember picture,overlay]\begin{scope}[shift={(current page.north west)},x=1bp,y=-1bp]
\node[anchor=north west,inner sep=0pt,text=answer,align=left,text width=229.0bp] at (15,15) {\fontsize{7.5bp}{9.5bp}\selectfont 析：\allowbreak{}选A。\allowbreak{}两个加数都须乘\ensuremath{3};\allowbreak{}B减数符号错，\allowbreak{}C漏乘\ensuremath{8}。\allowbreak{}\par};
\end{scope}\end{tikzpicture}\newpage
\definecolor{answer}{rgb}{0.850980,0.047059,0.070588}
\null\begin{tikzpicture}[remember picture,overlay]\begin{scope}[shift={(current page.north west)},x=1bp,y=-1bp]
\node[anchor=north west,inner sep=0pt,text=answer,align=left,text width=244.8079833984375bp] at (15,15) {\fontsize{8.5bp}{10.62bp}\selectfont 先合并两人每小时的工作量，\allowbreak{}再乘\ensuremath{3}小时；\allowbreak{}也可分别求总数。\allowbreak{}\par};
\end{scope}\end{tikzpicture}\newpage
\definecolor{answer}{rgb}{0.850980,0.047059,0.070588}
\null\begin{tikzpicture}[remember picture,overlay]\begin{scope}[shift={(current page.north west)},x=1bp,y=-1bp]
\node[anchor=north west,inner sep=0pt,text=answer,align=left,text width=229.0bp] at (15,15) {\fontsize{7.5bp}{9.5bp}\selectfont 析：\allowbreak{}选C。\allowbreak{}\ensuremath{56} = \ensuremath{56} × \ensuremath{1},\allowbreak{}合并成(\ensuremath{1} + \ensuremath{99})个\ensuremath{56}。\allowbreak{}\par};
\end{scope}\end{tikzpicture}\newpage
\definecolor{answer}{rgb}{0.850980,0.047059,0.070588}
\null\begin{tikzpicture}[remember picture,overlay]\begin{scope}[shift={(current page.north west)},x=1bp,y=-1bp]
\node[anchor=north west,inner sep=0pt,text=answer,align=left,text width=229.0bp] at (15,15) {\fontsize{7.5bp}{9.5bp}\selectfont 析：\allowbreak{}选B。\allowbreak{}\ensuremath{45} × \ensuremath{101}比\ensuremath{45} × \ensuremath{100} = \ensuremath{4500}多\ensuremath{45}。\allowbreak{}\par};
\end{scope}\end{tikzpicture}\newpage
\definecolor{answer}{rgb}{0.850980,0.047059,0.070588}
\null\begin{tikzpicture}[remember picture,overlay]\begin{scope}[shift={(current page.north west)},x=1bp,y=-1bp]
\node[anchor=north west,inner sep=0pt,text=answer,align=left,text width=234.35748291015625bp] at (15,15) {\fontsize{8.5bp}{10.62bp}\selectfont 一套包括一张桌子和一把椅子，\allowbreak{}先求一套价格，\allowbreak{}再乘\ensuremath{53}套。\allowbreak{}\par};
\end{scope}\end{tikzpicture}\newpage
\definecolor{answer}{rgb}{0.850980,0.047059,0.070588}
\null\begin{tikzpicture}[remember picture,overlay]\begin{scope}[shift={(current page.north west)},x=1bp,y=-1bp]
\node[anchor=north west,inner sep=0pt,text=answer,align=left,text width=214.08599853515625bp] at (15,15) {\fontsize{8.5bp}{10.62bp}\selectfont 正确式中括号内的\ensuremath{4}也应乘\ensuremath{45},\allowbreak{}错误式只加了一个\ensuremath{4},\allowbreak{}因此少了\ensuremath{44}个\ensuremath{4}。\allowbreak{}\par};
\end{scope}\end{tikzpicture}\newpage
\definecolor{answer}{rgb}{0.850980,0.047059,0.070588}
\null\begin{tikzpicture}[remember picture,overlay]\begin{scope}[shift={(current page.north west)},x=1bp,y=-1bp]
\node[anchor=north west,inner sep=0pt,text=answer,align=left,text width=214.0bp] at (15,15) {\fontsize{8.5bp}{10.62bp}\selectfont 答：\allowbreak{}错误结果比正确结果少\ensuremath{176}。\allowbreak{}\par};
\end{scope}\end{tikzpicture}\newpage
\definecolor{answer}{rgb}{0.850980,0.047059,0.070588}
\null\begin{tikzpicture}[remember picture,overlay]\begin{scope}[shift={(current page.north west)},x=1bp,y=-1bp]
\node[anchor=north west,inner sep=0pt,text=answer,align=left,text width=229.0bp] at (15,15) {\fontsize{7.5bp}{9.5bp}\selectfont 析：\allowbreak{}三角形有稳定性，\allowbreak{}三边分别作底，\allowbreak{}共有\ensuremath{3}条高。\allowbreak{}\par};
\end{scope}\end{tikzpicture}\newpage
\definecolor{answer}{rgb}{0.850980,0.047059,0.070588}
\null\begin{tikzpicture}[remember picture,overlay]\begin{scope}[shift={(current page.north west)},x=1bp,y=-1bp]
\node[anchor=north west,inner sep=0pt,text=answer,align=left,text width=229.0bp] at (15,15) {\fontsize{7.5bp}{9.5bp}\selectfont 析：\allowbreak{}每条底边各对应一条高，\allowbreak{}共\ensuremath{3}组底和高。\allowbreak{}\par};
\end{scope}\end{tikzpicture}\newpage
\definecolor{answer}{rgb}{0.850980,0.047059,0.070588}
\null\begin{tikzpicture}[remember picture,overlay]\begin{scope}[shift={(current page.north west)},x=1bp,y=-1bp]
\node[anchor=north west,inner sep=0pt,text=answer,align=left,text width=224.0bp] at (15,15) {\fontsize{7.5bp}{9.5bp}\selectfont 析：\allowbreak{}三线段须首尾顺次相接，\allowbreak{}围成封闭图形。\allowbreak{}\par};
\end{scope}\end{tikzpicture}\newpage
\definecolor{answer}{rgb}{0.850980,0.047059,0.070588}
\null\begin{tikzpicture}[remember picture,overlay]\begin{scope}[shift={(current page.north west)},x=1bp,y=-1bp]
\node[anchor=north west,inner sep=0pt,text=answer,align=left,text width=229.0bp] at (15,15) {\fontsize{7.5bp}{9.5bp}\selectfont 析：\allowbreak{}CD ⊥AB,\allowbreak{} BE ⊥AC,\allowbreak{}故相应高是CD、\allowbreak{}BE。\allowbreak{}\par};
\end{scope}\end{tikzpicture}\newpage
\definecolor{answer}{rgb}{0.850980,0.047059,0.070588}
\null\begin{tikzpicture}[remember picture,overlay]\begin{scope}[shift={(current page.north west)},x=1bp,y=-1bp]
\node[anchor=north west,inner sep=0pt,text=answer,align=left,text width=224.0bp] at (15,15) {\fontsize{7.5bp}{9.5bp}\selectfont 析：\allowbreak{}依图填角、\allowbreak{}高、\allowbreak{}底、\allowbreak{}顶点；\allowbreak{}高与底垂直。\allowbreak{}\par};
\end{scope}\end{tikzpicture}\newpage
\definecolor{answer}{rgb}{0.850980,0.047059,0.070588}
\null\begin{tikzpicture}[remember picture,overlay]\begin{scope}[shift={(current page.north west)},x=1bp,y=-1bp]
\node[anchor=north west,inner sep=0pt,text=answer,align=left,text width=224.0bp] at (15,15) {\fontsize{7.5bp}{9.5bp}\selectfont 析：\allowbreak{}过顶点作底边的垂线；\allowbreak{}垂线段叫高，\allowbreak{}相应边叫底。\allowbreak{}\par};
\end{scope}\end{tikzpicture}\newpage
\definecolor{answer}{rgb}{0.850980,0.047059,0.070588}
\null\begin{tikzpicture}[remember picture,overlay]\begin{scope}[shift={(current page.north west)},x=1bp,y=-1bp]
\node[anchor=north west,inner sep=0pt,text=answer,align=left,text width=224.0bp] at (15,15) {\fontsize{7.5bp}{9.5bp}\selectfont 析：\allowbreak{}三角形有\ensuremath{3}条边、\allowbreak{}\ensuremath{3}个顶点和\ensuremath{3}个角。\allowbreak{}\par};
\end{scope}\end{tikzpicture}\newpage
\definecolor{answer}{rgb}{0.850980,0.047059,0.070588}
\null\begin{tikzpicture}[remember picture,overlay]\begin{scope}[shift={(current page.north west)},x=1bp,y=-1bp]
\node[anchor=north west,inner sep=0pt,text=answer,align=left,text width=235.5bp] at (15,15) {\fontsize{7.5bp}{9.5bp}\selectfont 指定底边上的高从对面顶点出发，\allowbreak{}垂直于底边所在直线。\allowbreak{}必要时延长底边，\allowbreak{}高用虚线，\allowbreak{}旁边标高和直角。\allowbreak{}\par};
\end{scope}\end{tikzpicture}\newpage
\definecolor{answer}{rgb}{0.850980,0.047059,0.070588}
\null\begin{tikzpicture}[remember picture,overlay]\begin{scope}[shift={(current page.north west)},x=1bp,y=-1bp]
\node[anchor=north west,inner sep=0pt,text=answer,align=left,text width=227.0bp] at (15,15) {\fontsize{7.5bp}{9.5bp}\selectfont 析：\allowbreak{}选A。\allowbreak{}三角形具有稳定性，\allowbreak{}使桥梁支架稳固。\allowbreak{}\par};
\end{scope}\end{tikzpicture}\newpage
\definecolor{answer}{rgb}{0.850980,0.047059,0.070588}
\null\begin{tikzpicture}[remember picture,overlay]\begin{scope}[shift={(current page.north west)},x=1bp,y=-1bp]
\node[anchor=north west,inner sep=0pt,text=answer,align=left,text width=227.0bp] at (15,15) {\fontsize{7.5bp}{9.5bp}\selectfont 析：\allowbreak{}选C。\allowbreak{}斜撑把四边形分成三角形，\allowbreak{}结构最稳固。\allowbreak{}\par};
\end{scope}\end{tikzpicture}\newpage
\definecolor{answer}{rgb}{0.850980,0.047059,0.070588}
\null\begin{tikzpicture}[remember picture,overlay]\begin{scope}[shift={(current page.north west)},x=1bp,y=-1bp]
\node[anchor=north west,inner sep=0pt,text=answer,align=left,text width=227.0bp] at (15,15) {\fontsize{7.5bp}{9.5bp}\selectfont 析：\allowbreak{}选C。\allowbreak{}高既须垂直于底，\allowbreak{}也须从对面的顶点出发。\allowbreak{}\par};
\end{scope}\end{tikzpicture}\newpage
\definecolor{answer}{rgb}{0.850980,0.047059,0.070588}
\null\begin{tikzpicture}[remember picture,overlay]\begin{scope}[shift={(current page.north west)},x=1bp,y=-1bp]
\node[anchor=north west,inner sep=0pt,text=answer,align=left,text width=229.0bp] at (15,15) {\fontsize{7.5bp}{9.5bp}\selectfont 析：\allowbreak{}每多\ensuremath{1}个三角形只添\ensuremath{2}根；\allowbreak{}\ensuremath{5}个用\ensuremath{11}根，\allowbreak{}\ensuremath{6}个用\ensuremath{13}根。\allowbreak{}\par};
\end{scope}\end{tikzpicture}\newpage
\definecolor{answer}{rgb}{0.850980,0.047059,0.070588}
\null\begin{tikzpicture}[remember picture,overlay]\begin{scope}[shift={(current page.north west)},x=1bp,y=-1bp]
\node[anchor=north west,inner sep=0pt,text=answer,align=left,text width=229.0bp] at (15,15) {\fontsize{7.5bp}{9.5bp}\selectfont 析：\allowbreak{}按最大角分类：\allowbreak{}锐②④⑥⑨；\allowbreak{}直①⑤⑦；\allowbreak{}钝③⑧。\allowbreak{}\par};
\end{scope}\end{tikzpicture}\newpage
\definecolor{answer}{rgb}{0.850980,0.047059,0.070588}
\null\begin{tikzpicture}[remember picture,overlay]\begin{scope}[shift={(current page.north west)},x=1bp,y=-1bp]
\node[anchor=north west,inner sep=0pt,text=answer,align=left,text width=227.0bp] at (15,15) {\fontsize{7.5bp}{9.5bp}\selectfont 析：\allowbreak{}按角分锐角、\allowbreak{}直角、\allowbreak{}钝角三角形，\allowbreak{}三类不重不漏。\allowbreak{}\par};
\end{scope}\end{tikzpicture}\newpage
\definecolor{answer}{rgb}{0.850980,0.047059,0.070588}
\null\begin{tikzpicture}[remember picture,overlay]\begin{scope}[shift={(current page.north west)},x=1bp,y=-1bp]
\node[anchor=north west,inner sep=0pt,text=answer,align=left,text width=227.0bp] at (15,15) {\fontsize{7.5bp}{9.5bp}\selectfont 析：\allowbreak{}填\ensuremath{1};\allowbreak{}两个钝角之和已超过\ensuremath{180}\ensuremath{^\circ},\allowbreak{}不可能。\allowbreak{}\par};
\end{scope}\end{tikzpicture}\newpage
\definecolor{answer}{rgb}{0.850980,0.047059,0.070588}
\null\begin{tikzpicture}[remember picture,overlay]\begin{scope}[shift={(current page.north west)},x=1bp,y=-1bp]
\node[anchor=north west,inner sep=0pt,text=answer,align=left,text width=237.0bp] at (15,15) {\fontsize{5.8bp}{7.8bp}\selectfont 析：\allowbreak{}两底角相等且和小于\ensuremath{180}°，\allowbreak{}所以都小于\ensuremath{90}°，\allowbreak{}判√。\allowbreak{}\par};
\end{scope}\end{tikzpicture}\newpage
\definecolor{answer}{rgb}{0.850980,0.047059,0.070588}
\null\begin{tikzpicture}[remember picture,overlay]\begin{scope}[shift={(current page.north west)},x=1bp,y=-1bp]
\node[anchor=north west,inner sep=0pt,text=answer,align=left,text width=227.0bp] at (15,15) {\fontsize{7.5bp}{9.5bp}\selectfont 析：\allowbreak{}等腰三角形有两条相等的边，\allowbreak{}都叫腰。\allowbreak{}\par};
\end{scope}\end{tikzpicture}\newpage
\definecolor{answer}{rgb}{0.850980,0.047059,0.070588}
\null\begin{tikzpicture}[remember picture,overlay]\begin{scope}[shift={(current page.north west)},x=1bp,y=-1bp]
\node[anchor=north west,inner sep=0pt,text=answer,align=left,text width=227.0bp] at (15,15) {\fontsize{7.5bp}{9.5bp}\selectfont 析：\allowbreak{}锐角三角形三角都锐；\allowbreak{}直角三角形有一直角。\allowbreak{}\par};
\end{scope}\end{tikzpicture}\newpage
\definecolor{answer}{rgb}{0.850980,0.047059,0.070588}
\null\begin{tikzpicture}[remember picture,overlay]\begin{scope}[shift={(current page.north west)},x=1bp,y=-1bp]
\node[anchor=north west,inner sep=0pt,text=answer,align=left,text width=227.0bp] at (15,15) {\fontsize{7.5bp}{9.5bp}\selectfont 析：\allowbreak{}锐角最多\ensuremath{3}个、\allowbreak{}最少\ensuremath{2}个；\allowbreak{}钝角最多\ensuremath{1}个、\allowbreak{}最少\ensuremath{0}个。\allowbreak{}\par};
\end{scope}\end{tikzpicture}\newpage
\definecolor{answer}{rgb}{0.850980,0.047059,0.070588}
\null\begin{tikzpicture}[remember picture,overlay]\begin{scope}[shift={(current page.north west)},x=1bp,y=-1bp]
\node[anchor=north west,inner sep=0pt,text=answer,align=left,text width=227.0bp] at (15,15) {\fontsize{7.5bp}{9.5bp}\selectfont 析：\allowbreak{}选B。\allowbreak{}三条边都相等，\allowbreak{}叫等边三角形。\allowbreak{}\par};
\end{scope}\end{tikzpicture}\newpage
\definecolor{answer}{rgb}{0.850980,0.047059,0.070588}
\null\begin{tikzpicture}[remember picture,overlay]\begin{scope}[shift={(current page.north west)},x=1bp,y=-1bp]
\node[anchor=north west,inner sep=0pt,text=answer,align=left,text width=234.75bp] at (15,15) {\fontsize{7.5bp}{9.5bp}\selectfont 析：\allowbreak{}每个三角形都保留了长方形的一个直角，\allowbreak{}故均为直角三角形。\allowbreak{}\par};
\end{scope}\end{tikzpicture}\newpage
\definecolor{answer}{rgb}{0.850980,0.047059,0.070588}
\null\begin{tikzpicture}[remember picture,overlay]\begin{scope}[shift={(current page.north west)},x=1bp,y=-1bp]
\node[anchor=north west,inner sep=0pt,text=answer,align=left,text width=232.84149169921875bp] at (15,15) {\fontsize{7.5bp}{9.5bp}\selectfont 析：\allowbreak{}选D。\allowbreak{}只露一个锐角，\allowbreak{}另两角仍可组成锐、\allowbreak{}直、\allowbreak{}钝三种三角形。\allowbreak{}\par};
\end{scope}\end{tikzpicture}\newpage
\definecolor{answer}{rgb}{0.850980,0.047059,0.070588}
\null\begin{tikzpicture}[remember picture,overlay]\begin{scope}[shift={(current page.north west)},x=1bp,y=-1bp]
\node[anchor=north west,inner sep=0pt,text=answer,align=left,text width=227.0bp] at (15,15) {\fontsize{7.5bp}{9.5bp}\selectfont 析：\allowbreak{}两直角边是正方形的两邻边，\allowbreak{}等长且垂直，\allowbreak{}故均为等腰直角三角形。\allowbreak{}\par};
\end{scope}\end{tikzpicture}\newpage
\definecolor{answer}{rgb}{0.850980,0.047059,0.070588}
\null\begin{tikzpicture}[remember picture,overlay]\begin{scope}[shift={(current page.north west)},x=1bp,y=-1bp]
\node[anchor=north west,inner sep=0pt,text=answer,align=left,text width=219.0bp] at (15,15) {\fontsize{7.5bp}{9.5bp}\selectfont 析：\allowbreak{}选A。\allowbreak{}两段长腰加一段短底，\allowbreak{}首尾相接，\allowbreak{}总长等于周长。\allowbreak{}\par};
\end{scope}\end{tikzpicture}\newpage
\definecolor{answer}{rgb}{0.850980,0.047059,0.070588}
\null\begin{tikzpicture}[remember picture,overlay]\begin{scope}[shift={(current page.north west)},x=1bp,y=-1bp]
\node[anchor=north west,inner sep=0pt,text=answer,align=left,text width=232.49899291992188bp] at (15,15) {\fontsize{7.5bp}{9.5bp}\selectfont 析：\allowbreak{}先画线段；\allowbreak{}两端作圆心、\allowbreak{}\ensuremath{a}作半径，\allowbreak{}交点为第三顶点，\allowbreak{}再连接。\allowbreak{}\par};
\end{scope}\end{tikzpicture}\newpage
\definecolor{answer}{rgb}{0.850980,0.047059,0.070588}
\null\begin{tikzpicture}[remember picture,overlay]\begin{scope}[shift={(current page.north west)},x=1bp,y=-1bp]
\node[anchor=north west,inner sep=0pt,text=answer,align=left,text width=217.375bp] at (15,15) {\fontsize{8.5bp}{10.62bp}\selectfont 共用边在内部，\allowbreak{}外框有\ensuremath{4}条等长的边。\allowbreak{}题目问米，\allowbreak{}要把厘米换成米。\allowbreak{}\par};
\end{scope}\end{tikzpicture}\newpage
\definecolor{answer}{rgb}{0.850980,0.047059,0.070588}
\null\begin{tikzpicture}[remember picture,overlay]\begin{scope}[shift={(current page.north west)},x=1bp,y=-1bp]
\node[anchor=north west,inner sep=0pt,text=answer,align=left,text width=227.0bp] at (15,15) {\fontsize{7.5bp}{9.5bp}\selectfont 析：\allowbreak{}选B。\allowbreak{}三边等长，\allowbreak{}周长是边长的\ensuremath{3}倍；\allowbreak{}底与高一般不等。\allowbreak{}\par};
\end{scope}\end{tikzpicture}\newpage
\definecolor{answer}{rgb}{0.850980,0.047059,0.070588}
\null\begin{tikzpicture}[remember picture,overlay]\begin{scope}[shift={(current page.north west)},x=1bp,y=-1bp]
\node[anchor=north west,inner sep=0pt,text=answer,align=left,text width=221.0bp] at (15,15) {\fontsize{7.5bp}{9.5bp}\selectfont 析：\allowbreak{}第三边须大于\ensuremath{12} \ensuremath{-}\ensuremath{8} = \ensuremath{4}、\allowbreak{}小于\ensuremath{12} + \ensuremath{8} = \ensuremath{20};\allowbreak{}取整数最短\ensuremath{5}、\allowbreak{}最长\ensuremath{19}。\allowbreak{}\par};
\end{scope}\end{tikzpicture}\newpage
\definecolor{answer}{rgb}{0.850980,0.047059,0.070588}
\null\begin{tikzpicture}[remember picture,overlay]\begin{scope}[shift={(current page.north west)},x=1bp,y=-1bp]
\node[anchor=north west,inner sep=0pt,text=answer,align=left,text width=229.8900146484375bp] at (15,15) {\fontsize{7.5bp}{9.5bp}\selectfont 析：\allowbreak{}选B。\allowbreak{}比较三条边能否首尾相接，\allowbreak{}研究两边之和大于第三边。\allowbreak{}\par};
\end{scope}\end{tikzpicture}\newpage
\definecolor{answer}{rgb}{0.850980,0.047059,0.070588}
\null\begin{tikzpicture}[remember picture,overlay]\begin{scope}[shift={(current page.north west)},x=1bp,y=-1bp]
\node[anchor=north west,inner sep=0pt,text=answer,align=left,text width=219.0bp] at (15,15) {\fontsize{8.5bp}{10.62bp}\selectfont \ensuremath{60}厘米竹条减去\ensuremath{18}厘米底边，\allowbreak{}剩余\ensuremath{42}厘米平分给两条相等的腰。\allowbreak{}\par};
\end{scope}\end{tikzpicture}\newpage
\definecolor{answer}{rgb}{0.850980,0.047059,0.070588}
\null\begin{tikzpicture}[remember picture,overlay]\begin{scope}[shift={(current page.north west)},x=1bp,y=-1bp]
\node[anchor=north west,inner sep=0pt,text=answer,align=left,text width=227.0bp] at (15,15) {\fontsize{7.5bp}{9.5bp}\selectfont 析：\allowbreak{}选B。\allowbreak{}\ensuremath{20} <第三边<\ensuremath{80},\allowbreak{}选项中\ensuremath{60}符合，\allowbreak{}\ensuremath{10}与\ensuremath{90}不符合。\allowbreak{}\par};
\end{scope}\end{tikzpicture}\newpage
\definecolor{answer}{rgb}{0.850980,0.047059,0.070588}
\null\begin{tikzpicture}[remember picture,overlay]\begin{scope}[shift={(current page.north west)},x=1bp,y=-1bp]
\node[anchor=north west,inner sep=0pt,text=answer,align=left,text width=221.0bp] at (15,15) {\fontsize{7.5bp}{9.5bp}\selectfont 析：\allowbreak{}选D。\allowbreak{}\ensuremath{5} <第三边<\ensuremath{11},\allowbreak{}整数\ensuremath{6}、\allowbreak{}\ensuremath{7}、\allowbreak{}\ensuremath{8}、\allowbreak{}\ensuremath{9}、\allowbreak{}\ensuremath{10},\allowbreak{}共\ensuremath{5}种。\allowbreak{}\par};
\end{scope}\end{tikzpicture}\newpage
\definecolor{answer}{rgb}{0.850980,0.047059,0.070588}
\null\begin{tikzpicture}[remember picture,overlay]\begin{scope}[shift={(current page.north west)},x=1bp,y=-1bp]
\node[anchor=north west,inner sep=0pt,text=answer,align=left,text width=213.0bp] at (15,15) {\fontsize{8.5bp}{10.62bp}\selectfont 去程为\ensuremath{168} + \ensuremath{132}千米，\allowbreak{}先用总路程除以\ensuremath{5}小时求速度；\allowbreak{}返程慢\ensuremath{10}千米/时，\allowbreak{}取\ensuremath{250}千米近路。\allowbreak{}\par};
\end{scope}\end{tikzpicture}\newpage
\definecolor{answer}{rgb}{0.850980,0.047059,0.070588}
\null\begin{tikzpicture}[remember picture,overlay]\begin{scope}[shift={(current page.north west)},x=1bp,y=-1bp]
\node[anchor=north west,inner sep=0pt,text=answer,align=left,text width=213.0bp] at (15,15) {\fontsize{8.5bp}{10.62bp}\selectfont 法\ensuremath{2} :\allowbreak{}假设返程也行\ensuremath{5}小时作比较\par};
\end{scope}\end{tikzpicture}\newpage
\definecolor{answer}{rgb}{0.850980,0.047059,0.070588}
\null\begin{tikzpicture}[remember picture,overlay]\begin{scope}[shift={(current page.north west)},x=1bp,y=-1bp]
\node[anchor=north west,inner sep=0pt,text=answer,align=left,text width=218.47149658203125bp] at (15,15) {\fontsize{8.5bp}{10.62bp}\selectfont 近路比去程少\ensuremath{50}千米，\allowbreak{}速度每小时少\ensuremath{10}千米，\allowbreak{}\ensuremath{5}小时恰好少走\ensuremath{50}千米，\allowbreak{}所以走近路也恰好需要\ensuremath{5}小时。\allowbreak{}\par};
\end{scope}\end{tikzpicture}\newpage
\definecolor{answer}{rgb}{0.850980,0.047059,0.070588}
\null\begin{tikzpicture}[remember picture,overlay]\begin{scope}[shift={(current page.north west)},x=1bp,y=-1bp]
\node[anchor=north west,inner sep=0pt,text=answer,align=left,text width=213.0bp] at (15,15) {\fontsize{8.5bp}{10.62bp}\selectfont 答：\allowbreak{}从C城返回A城至少需要\ensuremath{5}小时。\allowbreak{}\par};
\end{scope}\end{tikzpicture}\newpage
\definecolor{answer}{rgb}{0.850980,0.047059,0.070588}
\null\begin{tikzpicture}[remember picture,overlay]\begin{scope}[shift={(current page.north west)},x=1bp,y=-1bp]
\node[anchor=north west,inner sep=0pt,text=answer,align=left,text width=212.5714111328125bp] at (15,15) {\fontsize{8.0bp}{10.0bp}\selectfont 如图,\allowbreak{} 从B城向AC作垂线,\allowbreak{}垂线段最短。\allowbreak{}\par};
\end{scope}\end{tikzpicture}\newpage
\definecolor{answer}{rgb}{0.850980,0.047059,0.070588}
\null\begin{tikzpicture}[remember picture,overlay]\begin{scope}[shift={(current page.north west)},x=1bp,y=-1bp]
\node[anchor=north west,inner sep=0pt,text=answer,align=left,text width=219.0bp] at (15,15) {\fontsize{8.5bp}{10.62bp}\selectfont 答：\allowbreak{}把\ensuremath{7}厘米竹签剪成\ensuremath{2}厘米和\ensuremath{5}厘米两段。\allowbreak{}\par};
\end{scope}\end{tikzpicture}\newpage
\definecolor{answer}{rgb}{0.850980,0.047059,0.070588}
\null\begin{tikzpicture}[remember picture,overlay]\begin{scope}[shift={(current page.north west)},x=1bp,y=-1bp]
\node[anchor=north west,inner sep=0pt,text=answer,align=left,text width=218.99999237060547bp] at (15,15) {\fontsize{8.5bp}{10.62bp}\selectfont 答：\allowbreak{}把\ensuremath{7}厘米竹签剪成\ensuremath{3}厘米和\ensuremath{4}厘米两段。\allowbreak{}\par};
\end{scope}\end{tikzpicture}\newpage
\definecolor{answer}{rgb}{0.850980,0.047059,0.070588}
\null\begin{tikzpicture}[remember picture,overlay]\begin{scope}[shift={(current page.north west)},x=1bp,y=-1bp]
\node[anchor=north west,inner sep=0pt,text=answer,align=left,text width=223.75bp] at (15,15) {\fontsize{8.5bp}{10.62bp}\selectfont 第三边必须大于两边之差、\allowbreak{}小于两边之和，\allowbreak{}且取整米数。\allowbreak{}\par};
\end{scope}\end{tikzpicture}\newpage
\definecolor{answer}{rgb}{0.850980,0.047059,0.070588}
\null\begin{tikzpicture}[remember picture,overlay]\begin{scope}[shift={(current page.north west)},x=1bp,y=-1bp]
\node[anchor=north west,inner sep=0pt,text=answer,align=left,text width=222.4630126953125bp] at (15,15) {\fontsize{7.5bp}{9.5bp}\selectfont 析：\allowbreak{}选B。\allowbreak{}小段当\ensuremath{1}份，\allowbreak{}B得\ensuremath{3}、\allowbreak{}\ensuremath{4}、\allowbreak{}\ensuremath{4}可围成；\allowbreak{}A为\ensuremath{3}、\allowbreak{}\ensuremath{2}、\allowbreak{}\ensuremath{6},\allowbreak{} C为\ensuremath{3}、\allowbreak{}\ensuremath{7}、\allowbreak{}\ensuremath{1},\allowbreak{}均不可。\allowbreak{}\par};
\end{scope}\end{tikzpicture}\newpage
\definecolor{answer}{rgb}{0.850980,0.047059,0.070588}
\null\begin{tikzpicture}[remember picture,overlay]\begin{scope}[shift={(current page.north west)},x=1bp,y=-1bp]
\node[anchor=north west,inner sep=0pt,text=answer,align=left,text width=226.76901245117188bp] at (15,15) {\fontsize{8.5bp}{10.62bp}\selectfont 答：\allowbreak{}可用\ensuremath{4}至\ensuremath{16}米之间任一整数米长的竹竿，\allowbreak{}共\ensuremath{13}种长度。\allowbreak{}\par};
\end{scope}\end{tikzpicture}\newpage
\definecolor{answer}{rgb}{0.850980,0.047059,0.070588}
\null\begin{tikzpicture}[remember picture,overlay]\begin{scope}[shift={(current page.north west)},x=1bp,y=-1bp]
\node[anchor=north west,inner sep=0pt,text=answer,align=left,text width=229.0bp] at (15,15) {\fontsize{7.5bp}{9.5bp}\selectfont 析：\allowbreak{}三个内角拼成平角，\allowbreak{}故填平、\allowbreak{}\ensuremath{180}。\allowbreak{}\par};
\end{scope}\end{tikzpicture}\newpage
\definecolor{answer}{rgb}{0.850980,0.047059,0.070588}
\null\begin{tikzpicture}[remember picture,overlay]\begin{scope}[shift={(current page.north west)},x=1bp,y=-1bp]
\node[anchor=north west,inner sep=0pt,text=answer,align=left,text width=227.0bp] at (15,15) {\fontsize{7.5bp}{9.5bp}\selectfont 析：\allowbreak{}\ensuremath{180}\ensuremath{^\circ}\ensuremath{-}\ensuremath{65}\ensuremath{^\circ}\ensuremath{-}\ensuremath{42}\ensuremath{^\circ}= \ensuremath{73}\ensuremath{^\circ},\allowbreak{}三个角均锐，\allowbreak{}填\ensuremath{73}、\allowbreak{}锐角。\allowbreak{}\par};
\end{scope}\end{tikzpicture}\newpage
\definecolor{answer}{rgb}{0.850980,0.047059,0.070588}
\null\begin{tikzpicture}[remember picture,overlay]\begin{scope}[shift={(current page.north west)},x=1bp,y=-1bp]
\node[anchor=north west,inner sep=0pt,text=answer,align=left,text width=229.0bp] at (15,15) {\fontsize{7.5bp}{9.5bp}\selectfont 析：\allowbreak{}每个小三角形的内角和仍为\ensuremath{180}\ensuremath{^\circ},\allowbreak{}不是原角和的一半。\allowbreak{}\par};
\end{scope}\end{tikzpicture}\newpage
\definecolor{answer}{rgb}{0.850980,0.047059,0.070588}
\null\begin{tikzpicture}[remember picture,overlay]\begin{scope}[shift={(current page.north west)},x=1bp,y=-1bp]
\node[anchor=north west,inner sep=0pt,text=answer,align=left,text width=229.0bp] at (15,15) {\fontsize{7.5bp}{9.5bp}\selectfont 析：\allowbreak{}\ensuremath{180}\ensuremath{^\circ}\ensuremath{-}\ensuremath{70}\ensuremath{^\circ}\ensuremath{-}\ensuremath{45}\ensuremath{^\circ}= \ensuremath{65}\ensuremath{^\circ},\allowbreak{}填\ensuremath{65}。\allowbreak{}\par};
\end{scope}\end{tikzpicture}\newpage
\definecolor{answer}{rgb}{0.850980,0.047059,0.070588}
\null\begin{tikzpicture}[remember picture,overlay]\begin{scope}[shift={(current page.north west)},x=1bp,y=-1bp]
\node[anchor=north west,inner sep=0pt,text=answer,align=left,text width=227.0bp] at (15,15) {\fontsize{7.5bp}{9.5bp}\selectfont 析：\allowbreak{}\ensuremath{90}\ensuremath{^\circ}\ensuremath{-}\ensuremath{55}\ensuremath{^\circ}= \ensuremath{35}\ensuremath{^\circ},\allowbreak{}有一直角，\allowbreak{}填\ensuremath{35}、\allowbreak{}直角。\allowbreak{}\par};
\end{scope}\end{tikzpicture}\newpage
\definecolor{answer}{rgb}{0.850980,0.047059,0.070588}
\null\begin{tikzpicture}[remember picture,overlay]\begin{scope}[shift={(current page.north west)},x=1bp,y=-1bp]
\node[anchor=north west,inner sep=0pt,text=answer,align=left,text width=229.0bp] at (15,15) {\fontsize{7.5bp}{9.5bp}\selectfont 析：\allowbreak{}两锐角和\ensuremath{90}\ensuremath{^\circ},\allowbreak{} \ensuremath{90}\ensuremath{^\circ}\ensuremath{-}\ensuremath{42}\ensuremath{^\circ}= \ensuremath{48}\ensuremath{^\circ},\allowbreak{}填\ensuremath{48}。\allowbreak{}\par};
\end{scope}\end{tikzpicture}\newpage
\definecolor{answer}{rgb}{0.850980,0.047059,0.070588}
\null\begin{tikzpicture}[remember picture,overlay]\begin{scope}[shift={(current page.north west)},x=1bp,y=-1bp]
\node[anchor=north west,inner sep=0pt,text=answer,align=left,text width=227.0bp] at (15,15) {\fontsize{7.5bp}{9.5bp}\selectfont 析：\allowbreak{}\ensuremath{180}\ensuremath{^\circ}\ensuremath{-}\ensuremath{40}\ensuremath{^\circ}\ensuremath{-}\ensuremath{112}\ensuremath{^\circ}= \ensuremath{28}\ensuremath{^\circ};\allowbreak{}有\ensuremath{112}\ensuremath{^\circ}钝角，\allowbreak{}填\ensuremath{28}、\allowbreak{}钝角。\allowbreak{}\par};
\end{scope}\end{tikzpicture}\newpage
\definecolor{answer}{rgb}{0.850980,0.047059,0.070588}
\null\begin{tikzpicture}[remember picture,overlay]\begin{scope}[shift={(current page.north west)},x=1bp,y=-1bp]
\node[anchor=north west,inner sep=0pt,text=answer,align=left,text width=221.0bp] at (15,15) {\fontsize{8.5bp}{10.62bp}\selectfont 等腰三角形两个底角相等，\allowbreak{}均为\ensuremath{42}\ensuremath{^\circ}。\allowbreak{}\par};
\end{scope}\end{tikzpicture}\newpage
\definecolor{answer}{rgb}{0.850980,0.047059,0.070588}
\null\begin{tikzpicture}[remember picture,overlay]\begin{scope}[shift={(current page.north west)},x=1bp,y=-1bp]
\node[anchor=north west,inner sep=0pt,text=answer,align=left,text width=227.0bp] at (15,15) {\fontsize{7.5bp}{9.5bp}\selectfont 析：\allowbreak{}第三角依次\ensuremath{101}\ensuremath{^\circ}、\allowbreak{}\ensuremath{70}\ensuremath{^\circ}、\allowbreak{}\ensuremath{90}\ensuremath{^\circ},\allowbreak{}故填钝角、\allowbreak{}锐角、\allowbreak{}直角。\allowbreak{}\par};
\end{scope}\end{tikzpicture}\newpage
\definecolor{answer}{rgb}{0.850980,0.047059,0.070588}
\null\begin{tikzpicture}[remember picture,overlay]\begin{scope}[shift={(current page.north west)},x=1bp,y=-1bp]
\node[anchor=north west,inner sep=0pt,text=answer,align=left,text width=220.0bp] at (15,15) {\fontsize{8.5bp}{10.62bp}\selectfont 先从\ensuremath{180}\ensuremath{^\circ}中扣除\ensuremath{120}\ensuremath{^\circ}顶角，\allowbreak{}剩余两个相等的底角平分。\allowbreak{}\par};
\end{scope}\end{tikzpicture}\newpage
\definecolor{answer}{rgb}{0.850980,0.047059,0.070588}
\null\begin{tikzpicture}[remember picture,overlay]\begin{scope}[shift={(current page.north west)},x=1bp,y=-1bp]
\node[anchor=north west,inner sep=0pt,text=answer,align=left,text width=219.0bp] at (15,15) {\fontsize{8.5bp}{10.62bp}\selectfont 答：\allowbreak{}两个底角分别是\ensuremath{30}\ensuremath{^\circ}、\allowbreak{}\ensuremath{30}\ensuremath{^\circ}。\allowbreak{}\par};
\end{scope}\end{tikzpicture}\newpage
\definecolor{answer}{rgb}{0.850980,0.047059,0.070588}
\null\begin{tikzpicture}[remember picture,overlay]\begin{scope}[shift={(current page.north west)},x=1bp,y=-1bp]
\node[anchor=north west,inner sep=0pt,text=answer,align=left,text width=224.0bp] at (15,15) {\fontsize{7.5bp}{9.5bp}\selectfont 析：\allowbreak{}选B。\allowbreak{}第三角等于前两角和，\allowbreak{}占\ensuremath{180}\ensuremath{^\circ}一半，\allowbreak{}即\ensuremath{90}\ensuremath{^\circ}。\allowbreak{}\par};
\end{scope}\end{tikzpicture}\newpage
\definecolor{answer}{rgb}{0.850980,0.047059,0.070588}
\null\begin{tikzpicture}[remember picture,overlay]\begin{scope}[shift={(current page.north west)},x=1bp,y=-1bp]
\node[anchor=north west,inner sep=0pt,text=answer,align=left,text width=224.0bp] at (15,15) {\fontsize{7.5bp}{9.5bp}\selectfont 析：\allowbreak{}选B。\allowbreak{}另两角须合\ensuremath{115}\ensuremath{^\circ};\allowbreak{}\ensuremath{70}\ensuremath{^\circ}+ \ensuremath{55}\ensuremath{^\circ}= \ensuremath{125}\ensuremath{^\circ},\allowbreak{}不符合。\allowbreak{}\par};
\end{scope}\end{tikzpicture}\newpage
\definecolor{answer}{rgb}{0.850980,0.047059,0.070588}
\null\begin{tikzpicture}[remember picture,overlay]\begin{scope}[shift={(current page.north west)},x=1bp,y=-1bp]
\node[anchor=north west,inner sep=0pt,text=answer,align=left,text width=225.53048706054688bp] at (15,15) {\fontsize{7.5bp}{9.5bp}\selectfont 析：\allowbreak{}选C。\allowbreak{}等腰钝角三角形的两底角和小于\ensuremath{90}\ensuremath{^\circ},\allowbreak{}每角须小于\ensuremath{45}\ensuremath{^\circ}。\allowbreak{}\par};
\end{scope}\end{tikzpicture}\newpage
\definecolor{answer}{rgb}{0.850980,0.047059,0.070588}
\null\begin{tikzpicture}[remember picture,overlay]\begin{scope}[shift={(current page.north west)},x=1bp,y=-1bp]
\node[anchor=north west,inner sep=0pt,text=answer,align=left,text width=167.0bp] at (15,15) {\fontsize{7.5bp}{9.5bp}\selectfont 析：\allowbreak{}选C。\allowbreak{}只剪开，\allowbreak{}未拼成平角，\allowbreak{}不能证明。\allowbreak{}\par};
\end{scope}\end{tikzpicture}\newpage
\definecolor{answer}{rgb}{0.850980,0.047059,0.070588}
\null\begin{tikzpicture}[remember picture,overlay]\begin{scope}[shift={(current page.north west)},x=1bp,y=-1bp]
\node[anchor=north west,inner sep=0pt,text=answer,align=left,text width=227.0bp] at (15,15) {\fontsize{7.5bp}{9.5bp}\selectfont 析：\allowbreak{}两组对边平行是平行四边形；\allowbreak{}只有一组平行是梯形。\allowbreak{}\par};
\end{scope}\end{tikzpicture}\newpage
\definecolor{answer}{rgb}{0.850980,0.047059,0.070588}
\null\begin{tikzpicture}[remember picture,overlay]\begin{scope}[shift={(current page.north west)},x=1bp,y=-1bp]
\node[anchor=north west,inner sep=0pt,text=answer,align=left,text width=227.0bp] at (15,15) {\fontsize{7.5bp}{9.5bp}\selectfont 析：\allowbreak{}梯形平行的两边叫底，\allowbreak{}不平行的两边叫腰。\allowbreak{}\par};
\end{scope}\end{tikzpicture}\newpage
\definecolor{answer}{rgb}{0.850980,0.047059,0.070588}
\null\begin{tikzpicture}[remember picture,overlay]\begin{scope}[shift={(current page.north west)},x=1bp,y=-1bp]
\node[anchor=north west,inner sep=0pt,text=answer,align=left,text width=227.0bp] at (15,15) {\fontsize{7.5bp}{9.5bp}\selectfont 析：\allowbreak{}填易变形。\allowbreak{}伸缩门利用平行四边形角度、\allowbreak{}高可改变的特点。\allowbreak{}\par};
\end{scope}\end{tikzpicture}\newpage
\definecolor{answer}{rgb}{0.850980,0.047059,0.070588}
\null\begin{tikzpicture}[remember picture,overlay]\begin{scope}[shift={(current page.north west)},x=1bp,y=-1bp]
\node[anchor=north west,inner sep=0pt,text=answer,align=left,text width=229.0bp] at (15,15) {\fontsize{7.5bp}{9.5bp}\selectfont 析：\allowbreak{}选C。\allowbreak{}平行线间能画无数条垂线段，\allowbreak{}都是对应的高。\allowbreak{}\par};
\end{scope}\end{tikzpicture}\newpage
\definecolor{answer}{rgb}{0.850980,0.047059,0.070588}
\null\begin{tikzpicture}[remember picture,overlay]\begin{scope}[shift={(current page.north west)},x=1bp,y=-1bp]
\node[anchor=north west,inner sep=0pt,text=answer,align=left,text width=227.00001525878906bp] at (15,15) {\fontsize{7.5bp}{9.5bp}\selectfont 析：\allowbreak{}底\ensuremath{10}对应垂直高\ensuremath{8},\allowbreak{}底\ensuremath{12}对应垂直高\ensuremath{7};\allowbreak{}对角相等，\allowbreak{}故填\ensuremath{8}、\allowbreak{}\ensuremath{7}、\allowbreak{}\ensuremath{45}。\allowbreak{}\par};
\end{scope}\end{tikzpicture}\newpage
\definecolor{answer}{rgb}{0.850980,0.047059,0.070588}
\null\begin{tikzpicture}[remember picture,overlay]\begin{scope}[shift={(current page.north west)},x=1bp,y=-1bp]
\node[anchor=north west,inner sep=0pt,text=answer,align=left,text width=229.0bp] at (15,15) {\fontsize{7.5bp}{9.5bp}\selectfont 析：\allowbreak{}选A。\allowbreak{}梯形唯一平行的一组对边是上底和下底。\allowbreak{}\par};
\end{scope}\end{tikzpicture}\newpage
\definecolor{answer}{rgb}{0.850980,0.047059,0.070588}
\null\begin{tikzpicture}[remember picture,overlay]\begin{scope}[shift={(current page.north west)},x=1bp,y=-1bp]
\node[anchor=north west,inner sep=0pt,text=answer,align=left,text width=227.0bp] at (15,15) {\fontsize{8.5bp}{10.62bp}\selectfont 利用三个等边三角形各边相等，\allowbreak{}分别比较图中三条A到E的道路。\allowbreak{}\par};
\end{scope}\end{tikzpicture}\newpage
\definecolor{answer}{rgb}{0.850980,0.047059,0.070588}
\null\begin{tikzpicture}[remember picture,overlay]\begin{scope}[shift={(current page.north west)},x=1bp,y=-1bp]
\node[anchor=north west,inner sep=0pt,text=answer,align=left,text width=220.0bp] at (15,15) {\fontsize{8.5bp}{10.62bp}\selectfont 平行四边形对边分别相等，\allowbreak{}周长是两个相邻边长之和的\ensuremath{2}倍。\allowbreak{}\par};
\end{scope}\end{tikzpicture}\newpage
\definecolor{answer}{rgb}{0.850980,0.047059,0.070588}
\null\begin{tikzpicture}[remember picture,overlay]\begin{scope}[shift={(current page.north west)},x=1bp,y=-1bp]
\node[anchor=north west,inner sep=0pt,text=answer,align=left,text width=228.1500244140625bp] at (15,15) {\fontsize{8.5bp}{10.62bp}\selectfont 答：\allowbreak{}ACE最近；\allowbreak{}ABCDE与AFE一样长，\allowbreak{}都是\ensuremath{280}米。\allowbreak{}\par};
\end{scope}\end{tikzpicture}\newpage
\definecolor{answer}{rgb}{0.850980,0.047059,0.070588}
\null\begin{tikzpicture}[remember picture,overlay]\begin{scope}[shift={(current page.north west)},x=1bp,y=-1bp]
\node[anchor=north west,inner sep=0pt,text=answer,align=left,text width=227.0bp] at (15,15) {\fontsize{8.5bp}{10.62bp}\selectfont 答：\allowbreak{}ACE最近，\allowbreak{}ABCDE与AFE一样长。\allowbreak{}\par};
\end{scope}\end{tikzpicture}\newpage
\definecolor{answer}{rgb}{0.850980,0.047059,0.070588}
\null\begin{tikzpicture}[remember picture,overlay]\begin{scope}[shift={(current page.north west)},x=1bp,y=-1bp]
\node[anchor=north west,inner sep=0pt,text=answer,align=left,text width=220.0bp] at (15,15) {\fontsize{8.5bp}{10.62bp}\selectfont 答：\allowbreak{}这个平行四边形的周长是\ensuremath{24}厘米。\allowbreak{}\par};
\end{scope}\end{tikzpicture}\newpage
\definecolor{answer}{rgb}{0.850980,0.047059,0.070588}
\null\begin{tikzpicture}[remember picture,overlay]\begin{scope}[shift={(current page.north west)},x=1bp,y=-1bp]
\node[anchor=north west,inner sep=0pt,text=answer,align=left,text width=219.0bp] at (15,15) {\fontsize{8.5bp}{10.62bp}\selectfont 等腰梯形的两腰相等，\allowbreak{}用周长扣除上、\allowbreak{}下底后平分。\allowbreak{}\par};
\end{scope}\end{tikzpicture}\newpage
\definecolor{answer}{rgb}{0.850980,0.047059,0.070588}
\null\begin{tikzpicture}[remember picture,overlay]\begin{scope}[shift={(current page.north west)},x=1bp,y=-1bp]
\node[anchor=north west,inner sep=0pt,text=answer,align=left,text width=227.0bp] at (15,15) {\fontsize{8.5bp}{10.62bp}\selectfont 只比较把原平行四边形分成\ensuremath{4}个形状相同小平行四边形的小路；\allowbreak{}有两种达到\ensuremath{90}米的方案。\allowbreak{}\par};
\end{scope}\end{tikzpicture}\newpage
\definecolor{answer}{rgb}{0.850980,0.047059,0.070588}
\null\begin{tikzpicture}[remember picture,overlay]\begin{scope}[shift={(current page.north west)},x=1bp,y=-1bp]
\node[anchor=north west,inner sep=0pt,text=answer,align=left,text width=489.0bp] at (15,15) {\fontsize{7.5bp}{9.5bp}\selectfont 析：\allowbreak{}正方形属于长方形，\allowbreak{}长方形属于平行四边形，\allowbreak{}平行四边形与梯形都属于四边形。\allowbreak{}\par};
\end{scope}\end{tikzpicture}\newpage
\definecolor{answer}{rgb}{0.850980,0.047059,0.070588}
\null\begin{tikzpicture}[remember picture,overlay]\begin{scope}[shift={(current page.north west)},x=1bp,y=-1bp]
\node[anchor=north west,inner sep=0pt,text=answer,align=left,text width=489.0bp] at (15,15) {\fontsize{7.5bp}{9.5bp}\selectfont 析：\allowbreak{}按最大角分类：\allowbreak{}锐⑨⑩；\allowbreak{}直①②④⑤⑥{\scfallback ⑪}；\allowbreak{}钝③⑦⑧{\scfallback ⑫}。\allowbreak{}\par};
\end{scope}\end{tikzpicture}\newpage
\definecolor{answer}{rgb}{0.850980,0.047059,0.070588}
\null\begin{tikzpicture}[remember picture,overlay]\begin{scope}[shift={(current page.north west)},x=1bp,y=-1bp]
\node[anchor=north west,inner sep=0pt,text=answer,align=left,text width=489.0bp] at (15,15) {\fontsize{7.5bp}{9.5bp}\selectfont 析：\allowbreak{}内角和\ensuremath{180}\ensuremath{^\circ},\allowbreak{}至少有\ensuremath{2}个锐角；\allowbreak{}钝角或直角都最多\ensuremath{1}个。\allowbreak{}\par};
\end{scope}\end{tikzpicture}\newpage
\definecolor{answer}{rgb}{0.850980,0.047059,0.070588}
\null\begin{tikzpicture}[remember picture,overlay]\begin{scope}[shift={(current page.north west)},x=1bp,y=-1bp]
\node[anchor=north west,inner sep=0pt,text=answer,align=left,text width=489.0bp] at (15,15) {\fontsize{7.5bp}{9.5bp}\selectfont 析：\allowbreak{}对边分别相等，\allowbreak{}周长是邻边和的\ensuremath{2}倍：\allowbreak{}\ensuremath{16} × \ensuremath{2} = \ensuremath{32}厘米。\allowbreak{}\par};
\end{scope}\end{tikzpicture}\newpage
\definecolor{answer}{rgb}{0.850980,0.047059,0.070588}
\null\begin{tikzpicture}[remember picture,overlay]\begin{scope}[shift={(current page.north west)},x=1bp,y=-1bp]
\node[anchor=north west,inner sep=0pt,text=answer,align=left,text width=489.0bp] at (15,15) {\fontsize{7.5bp}{9.5bp}\selectfont 析：\allowbreak{}等腰梯形两腰均长\ensuremath{5}厘米，\allowbreak{}\ensuremath{4} + \ensuremath{8} + \ensuremath{5} × \ensuremath{2} = \ensuremath{22}厘米。\allowbreak{}\par};
\end{scope}\end{tikzpicture}\newpage
\definecolor{answer}{rgb}{0.850980,0.047059,0.070588}
\null\begin{tikzpicture}[remember picture,overlay]\begin{scope}[shift={(current page.north west)},x=1bp,y=-1bp]
\node[anchor=north west,inner sep=0pt,text=answer,align=left,text width=489.0bp] at (15,15) {\fontsize{7.5bp}{9.5bp}\selectfont 析：\allowbreak{}三脚架、\allowbreak{}画架、\allowbreak{}马扎用三角形支撑，\allowbreak{}利用其稳定性。\allowbreak{}\par};
\end{scope}\end{tikzpicture}\newpage
\definecolor{answer}{rgb}{0.850980,0.047059,0.070588}
\null\begin{tikzpicture}[remember picture,overlay]\begin{scope}[shift={(current page.north west)},x=1bp,y=-1bp]
\node[anchor=north west,inner sep=0pt,text=answer,align=left,text width=489.0bp] at (15,15) {\fontsize{7.5bp}{9.5bp}\selectfont 析：\allowbreak{}\ensuremath{4}不能作腰，\allowbreak{}因为\ensuremath{4} + \ensuremath{4} < \ensuremath{9};\allowbreak{}腰长\ensuremath{9},\allowbreak{}周长\ensuremath{9} × \ensuremath{2} + \ensuremath{4} = \ensuremath{22}厘米。\allowbreak{}\par};
\end{scope}\end{tikzpicture}\newpage
\definecolor{answer}{rgb}{0.850980,0.047059,0.070588}
\null\begin{tikzpicture}[remember picture,overlay]\begin{scope}[shift={(current page.north west)},x=1bp,y=-1bp]
\node[anchor=north west,inner sep=0pt,text=answer,align=left,text width=489.125bp] at (15,15) {\fontsize{7.52bp}{9.52bp}\selectfont 析：\allowbreak{}\ensuremath{180}\ensuremath{^\circ}\ensuremath{-}(\ensuremath{45}\ensuremath{^\circ}+ \ensuremath{65}\ensuremath{^\circ}) = \ensuremath{180}\ensuremath{^\circ}\ensuremath{-}\ensuremath{110}\ensuremath{^\circ}= \ensuremath{70}\ensuremath{^\circ}，\allowbreak{}三个角都锐，\allowbreak{}填\ensuremath{70}、\allowbreak{}锐角\par};
\end{scope}\end{tikzpicture}\newpage
\definecolor{answer}{rgb}{0.850980,0.047059,0.070588}
\null\begin{tikzpicture}[remember picture,overlay]\begin{scope}[shift={(current page.north west)},x=1bp,y=-1bp]
\node[anchor=north west,inner sep=0pt,text=answer,align=left,text width=489.0bp] at (15,15) {\fontsize{7.5bp}{9.5bp}\selectfont 析：\allowbreak{}选A。\allowbreak{}同一梯形的高都是两平行底边间的垂线段，\allowbreak{}长度相等。\allowbreak{}\par};
\end{scope}\end{tikzpicture}\newpage
\definecolor{answer}{rgb}{0.850980,0.047059,0.070588}
\null\begin{tikzpicture}[remember picture,overlay]\begin{scope}[shift={(current page.north west)},x=1bp,y=-1bp]
\node[anchor=north west,inner sep=0pt,text=answer,align=left,text width=489.0bp] at (15,15) {\fontsize{7.5bp}{9.5bp}\selectfont 析：\allowbreak{}选B。\allowbreak{}平行四边形两组对角分别相等，\allowbreak{}邻角一般不相等。\allowbreak{}\par};
\end{scope}\end{tikzpicture}\newpage
\definecolor{answer}{rgb}{0.850980,0.047059,0.070588}
\null\begin{tikzpicture}[remember picture,overlay]\begin{scope}[shift={(current page.north west)},x=1bp,y=-1bp]
\node[anchor=north west,inner sep=0pt,text=answer,align=left,text width=489.0bp] at (15,15) {\fontsize{7.5bp}{9.5bp}\selectfont 析：\allowbreak{}选A。\allowbreak{}三块按正、\allowbreak{}倒、\allowbreak{}正排列，\allowbreak{}可围成上底\ensuremath{1}份、\allowbreak{}下底\ensuremath{2}份的梯形。\allowbreak{}\par};
\end{scope}\end{tikzpicture}\newpage
\definecolor{answer}{rgb}{0.850980,0.047059,0.070588}
\null\begin{tikzpicture}[remember picture,overlay]\begin{scope}[shift={(current page.north west)},x=1bp,y=-1bp]
\node[anchor=north west,inner sep=0pt,text=answer,align=left,text width=489.0bp] at (15,15) {\fontsize{7.5bp}{9.5bp}\selectfont 析：\allowbreak{}选C。\allowbreak{}两腰相等，\allowbreak{}(\ensuremath{40} \ensuremath{-}\ensuremath{28}) ÷ \ensuremath{2} = \ensuremath{6}厘米。\allowbreak{}\par};
\end{scope}\end{tikzpicture}\newpage
\definecolor{answer}{rgb}{0.850980,0.047059,0.070588}
\null\begin{tikzpicture}[remember picture,overlay]\begin{scope}[shift={(current page.north west)},x=1bp,y=-1bp]
\node[anchor=north west,inner sep=0pt,text=answer,align=left,text width=489.0bp] at (15,15) {\fontsize{7.5bp}{9.5bp}\selectfont 析：\allowbreak{}选C。\allowbreak{}拉动只改变角与高，\allowbreak{}四边长度不变，\allowbreak{}周长不变。\allowbreak{}\par};
\end{scope}\end{tikzpicture}\newpage
\definecolor{answer}{rgb}{0.850980,0.047059,0.070588}
\null\begin{tikzpicture}[remember picture,overlay]\begin{scope}[shift={(current page.north west)},x=1bp,y=-1bp]
\node[anchor=north west,inner sep=0pt,text=answer,align=left,text width=489.0bp] at (15,15) {\fontsize{7.5bp}{9.5bp}\selectfont 析：\allowbreak{}选D。\allowbreak{}C在E正上方可得直角，\allowbreak{}在EF 中点正上方可得等腰，\allowbreak{}移向远处可得钝角三角形。\allowbreak{}\par};
\end{scope}\end{tikzpicture}\newpage
\definecolor{answer}{rgb}{0.850980,0.047059,0.070588}
\null\begin{tikzpicture}[remember picture,overlay]\begin{scope}[shift={(current page.north west)},x=1bp,y=-1bp]
\node[anchor=north west,inner sep=0pt,text=answer,align=left,text width=489.0bp] at (15,15) {\fontsize{7.5bp}{9.5bp}\selectfont 析：\allowbreak{}选B。\allowbreak{}高应垂直于平行的底，\allowbreak{}B的虚线却垂直于腰。\allowbreak{}\par};
\end{scope}\end{tikzpicture}\newpage
\definecolor{answer}{rgb}{0.850980,0.047059,0.070588}
\null\begin{tikzpicture}[remember picture,overlay]\begin{scope}[shift={(current page.north west)},x=1bp,y=-1bp]
\node[anchor=north west,inner sep=0pt,text=answer,align=left,text width=489.0bp] at (15,15) {\fontsize{7.5bp}{9.5bp}\selectfont 析：\allowbreak{}选C。\allowbreak{}\ensuremath{8} <第三边<\ensuremath{20},\allowbreak{}只有\ensuremath{9}厘米符合，\allowbreak{}\ensuremath{20}不能取。\allowbreak{}\par};
\end{scope}\end{tikzpicture}\newpage
\definecolor{answer}{rgb}{0.850980,0.047059,0.070588}
\null\begin{tikzpicture}[remember picture,overlay]\begin{scope}[shift={(current page.north west)},x=1bp,y=-1bp]
\node[anchor=north west,inner sep=0pt,text=answer,align=left,text width=489.0bp] at (15,15) {\fontsize{7.5bp}{9.5bp}\selectfont 析：\allowbreak{}选A。\allowbreak{}第三角\ensuremath{180}\ensuremath{^\circ}\ensuremath{-}\ensuremath{50}\ensuremath{^\circ}\ensuremath{-}\ensuremath{60}\ensuremath{^\circ}= \ensuremath{70}\ensuremath{^\circ},\allowbreak{}三个角都锐。\allowbreak{}\par};
\end{scope}\end{tikzpicture}\newpage
\definecolor{answer}{rgb}{0.850980,0.047059,0.070588}
\null\begin{tikzpicture}[remember picture,overlay]\begin{scope}[shift={(current page.north west)},x=1bp,y=-1bp]
\node[anchor=north west,inner sep=0pt,text=answer,align=left,text width=489.0bp] at (15,15) {\fontsize{7.5bp}{9.5bp}\selectfont 析：\allowbreak{}选D。\allowbreak{}两底角相等，\allowbreak{}(\ensuremath{180}\ensuremath{^\circ}\ensuremath{-}\ensuremath{100}\ensuremath{^\circ}) ÷ \ensuremath{2} = \ensuremath{40}\ensuremath{^\circ}。\allowbreak{}\par};
\end{scope}\end{tikzpicture}\newpage
\definecolor{answer}{rgb}{0.850980,0.047059,0.070588}
\null\begin{tikzpicture}[remember picture,overlay]\begin{scope}[shift={(current page.north west)},x=1bp,y=-1bp]
\node[anchor=north west,inner sep=0pt,text=answer,align=left,text width=487.0bp] at (15,15) {\fontsize{6.2bp}{8.2bp}\selectfont 析：\allowbreak{}判×。\allowbreak{}\ensuremath{3}条线段须首尾相接围成封闭图形，\allowbreak{}仅有\ensuremath{3}条线段不一定能围成三角形。\allowbreak{}\par};
\end{scope}\end{tikzpicture}\newpage
\definecolor{answer}{rgb}{0.850980,0.047059,0.070588}
\null\begin{tikzpicture}[remember picture,overlay]\begin{scope}[shift={(current page.north west)},x=1bp,y=-1bp]
\node[anchor=north west,inner sep=0pt,text=answer,align=left,text width=487.0bp] at (15,15) {\fontsize{6.2bp}{8.2bp}\selectfont 析：\allowbreak{}判×。\allowbreak{}例如\ensuremath{30}°、\allowbreak{}\ensuremath{30}°、\allowbreak{}\ensuremath{120}°，\allowbreak{}有两个锐角，\allowbreak{}却是钝角三角形。\allowbreak{}\par};
\end{scope}\end{tikzpicture}\newpage
\definecolor{answer}{rgb}{0.850980,0.047059,0.070588}
\null\begin{tikzpicture}[remember picture,overlay]\begin{scope}[shift={(current page.north west)},x=1bp,y=-1bp]
\node[anchor=north west,inner sep=0pt,text=answer,align=left,text width=487.0bp] at (15,15) {\fontsize{6.2bp}{8.2bp}\selectfont 析：\allowbreak{}判×。\allowbreak{}例如\ensuremath{45}°、\allowbreak{}\ensuremath{45}°、\allowbreak{}\ensuremath{90}°，\allowbreak{}两底角相等，\allowbreak{}是等腰直角三角形。\allowbreak{}\par};
\end{scope}\end{tikzpicture}\newpage
\definecolor{answer}{rgb}{0.850980,0.047059,0.070588}
\null\begin{tikzpicture}[remember picture,overlay]\begin{scope}[shift={(current page.north west)},x=1bp,y=-1bp]
\node[anchor=north west,inner sep=0pt,text=answer,align=left,text width=487.0bp] at (15,15) {\fontsize{6.2bp}{8.2bp}\selectfont 析：\allowbreak{}判×。\allowbreak{}指定顶点到指定对边只能作\ensuremath{1}条垂线，\allowbreak{}所以相应的高只有\ensuremath{1}条。\allowbreak{}\par};
\end{scope}\end{tikzpicture}\newpage
\definecolor{answer}{rgb}{0.850980,0.047059,0.070588}
\null\begin{tikzpicture}[remember picture,overlay]\begin{scope}[shift={(current page.north west)},x=1bp,y=-1bp]
\node[anchor=north west,inner sep=0pt,text=answer,align=left,text width=487.0bp] at (15,15) {\fontsize{6.2bp}{8.2bp}\selectfont 析：\allowbreak{}判{\symbolfont ✓}。\allowbreak{}将一个梯形旋转\ensuremath{180}°，\allowbreak{}沿对应的一条腰与另一个拼接，\allowbreak{}可拼成平行四边形。\allowbreak{}\par};
\end{scope}\end{tikzpicture}\newpage
\definecolor{answer}{rgb}{0.850980,0.047059,0.070588}
\null\begin{tikzpicture}[remember picture,overlay]\begin{scope}[shift={(current page.north west)},x=1bp,y=-1bp]
\node[anchor=north west,inner sep=0pt,text=answer,align=left,text width=487.0bp] at (15,15) {\fontsize{6.2bp}{8.2bp}\selectfont 析：\allowbreak{}判×。\allowbreak{}三角形任意两边之和应大于第三边。\allowbreak{}\par};
\end{scope}\end{tikzpicture}\newpage
\definecolor{answer}{rgb}{0.850980,0.047059,0.070588}
\null\begin{tikzpicture}[remember picture,overlay]\begin{scope}[shift={(current page.north west)},x=1bp,y=-1bp]
\node[anchor=north west,inner sep=0pt,text=answer,align=left,text width=487.0bp] at (15,15) {\fontsize{6.2bp}{8.2bp}\selectfont 析：\allowbreak{}判×。\allowbreak{}任意三角形的内角和都是\ensuremath{180}°，\allowbreak{}钝角三角形也一样。\allowbreak{}\par};
\end{scope}\end{tikzpicture}\newpage
\definecolor{answer}{rgb}{0.850980,0.047059,0.070588}
\null\begin{tikzpicture}[remember picture,overlay]\begin{scope}[shift={(current page.north west)},x=1bp,y=-1bp]
\node[anchor=north west,inner sep=0pt,text=answer,align=left,text width=475.0bp] at (15,15) {\fontsize{8.5bp}{10.62bp}\selectfont 两根铁丝各长\ensuremath{60}厘米，\allowbreak{}分别围两个三角形，\allowbreak{}不能把两根铁丝合在一起计算。\allowbreak{}\par};
\end{scope}\end{tikzpicture}\newpage
\definecolor{answer}{rgb}{0.850980,0.047059,0.070588}
\null\begin{tikzpicture}[remember picture,overlay]\begin{scope}[shift={(current page.north west)},x=1bp,y=-1bp]
\node[anchor=north west,inner sep=0pt,text=answer,align=left,text width=484.708984375bp] at (15,15) {\fontsize{8.5bp}{10.62bp}\selectfont 三段总长\ensuremath{20}米。\allowbreak{}三角形较小两边之和必须大于最长边，\allowbreak{}所以最长边必须小于\ensuremath{10}米。\allowbreak{}原题只要求一种剪法，\allowbreak{}以下同时保留一种方便的方案和全部\ensuremath{8}种不重复整数方案。\allowbreak{}\par};
\end{scope}\end{tikzpicture}\newpage
\definecolor{answer}{rgb}{0.850980,0.047059,0.070588}
\null\begin{tikzpicture}[remember picture,overlay]\begin{scope}[shift={(current page.north west)},x=1bp,y=-1bp]
\node[anchor=north west,inner sep=0pt,text=answer,align=left,text width=475.0bp] at (15,15) {\fontsize{8.5bp}{10.62bp}\selectfont 答：\allowbreak{}以上\ensuremath{8}种方案都可行，\allowbreak{}任选一种；\allowbreak{}每组三边和\ensuremath{20}米，\allowbreak{}且最小两边和大于最长边。\allowbreak{}\par};
\end{scope}\end{tikzpicture}\newpage
\definecolor{answer}{rgb}{0.850980,0.047059,0.070588}
\null\begin{tikzpicture}[remember picture,overlay]\begin{scope}[shift={(current page.north west)},x=1bp,y=-1bp]
\node[anchor=north west,inner sep=0pt,text=answer,align=left,text width=470.25bp] at (15,15) {\fontsize{8.5bp}{10.62bp}\selectfont 展开括号后，\allowbreak{}先合并两个\ensuremath{7}，\allowbreak{}并把\ensuremath{4}与\ensuremath{6}凑成\ensuremath{10}\par};
\end{scope}\end{tikzpicture}\newpage
\definecolor{answer}{rgb}{0.850980,0.047059,0.070588}
\null\begin{tikzpicture}[remember picture,overlay]\begin{scope}[shift={(current page.north west)},x=1bp,y=-1bp]
\node[anchor=north west,inner sep=0pt,text=answer,align=left,text width=470.0bp] at (15,15) {\fontsize{8.5bp}{10.62bp}\selectfont 变成平行四边形后上下底相等，\allowbreak{}所加\ensuremath{9}厘米恰好是原上底的\ensuremath{4} \ensuremath{-}\ensuremath{1}份。\allowbreak{}\par};
\end{scope}\end{tikzpicture}\newpage
\definecolor{answer}{rgb}{0.850980,0.047059,0.070588}
\null\begin{tikzpicture}[remember picture,overlay]\begin{scope}[shift={(current page.north west)},x=1bp,y=-1bp]
\node[anchor=north west,inner sep=0pt,text=answer,align=left,text width=470.0bp] at (15,15) {\fontsize{8.5bp}{10.62bp}\selectfont 答：\allowbreak{}上底是\ensuremath{3}厘米，\allowbreak{}下底是\ensuremath{12}厘米。\allowbreak{}\par};
\end{scope}\end{tikzpicture}\newpage
\definecolor{answer}{rgb}{0.850980,0.047059,0.070588}
\null\begin{tikzpicture}[remember picture,overlay]\begin{scope}[shift={(current page.north west)},x=1bp,y=-1bp]
\node[anchor=north west,inner sep=0pt,text=answer,align=left,text width=470.0bp] at (15,15) {\fontsize{8.5bp}{10.62bp}\selectfont 两条竖边互相平行，\allowbreak{}它们之间的距离等于底部水平线段长；\allowbreak{}两个等腰直角三角形的水平直角边分别也长\ensuremath{6}厘米、\allowbreak{}\ensuremath{4}厘米。\allowbreak{}\par};
\end{scope}\end{tikzpicture}\newpage
\definecolor{answer}{rgb}{0.850980,0.047059,0.070588}
\null\begin{tikzpicture}[remember picture,overlay]\begin{scope}[shift={(current page.north west)},x=1bp,y=-1bp]
\node[anchor=north west,inner sep=0pt,text=answer,align=left,text width=489.0bp] at (15,15) {\fontsize{7.5bp}{9.5bp}\selectfont 先扣去第一天卖出的\ensuremath{40}千克，\allowbreak{}再把剩余的数量平均分给后\ensuremath{4}天。\allowbreak{}也可假设\ensuremath{5}天每天都卖\ensuremath{40}千克，\allowbreak{}把还多出的数量补给后\ensuremath{4}天。\allowbreak{}\par};
\end{scope}\end{tikzpicture}\newpage
\definecolor{answer}{rgb}{0.850980,0.047059,0.070588}
\null\begin{tikzpicture}[remember picture,overlay]\begin{scope}[shift={(current page.north west)},x=1bp,y=-1bp]
\node[anchor=north west,inner sep=0pt,text=answer,align=left,text width=489.0bp] at (15,15) {\fontsize{7.5bp}{9.5bp}\selectfont 析：\allowbreak{}\ensuremath{4}本笔记本共\ensuremath{9} × \ensuremath{4}元，\allowbreak{}(\ensuremath{9} × \ensuremath{4}) ÷ \ensuremath{12} = \ensuremath{3}支钢笔。\allowbreak{}\par};
\end{scope}\end{tikzpicture}\newpage
\definecolor{answer}{rgb}{0.850980,0.047059,0.070588}
\null\begin{tikzpicture}[remember picture,overlay]\begin{scope}[shift={(current page.north west)},x=1bp,y=-1bp]
\node[anchor=north west,inner sep=0pt,text=answer,align=left,text width=332.9679870605469bp] at (15,15) {\fontsize{7.5bp}{9.5bp}\selectfont 析：\allowbreak{}选B。\allowbreak{}先求多出的球数\ensuremath{240} \ensuremath{-}\ensuremath{180},\allowbreak{}再除每盒\ensuremath{12}个，\allowbreak{}求多装的盒数。\allowbreak{}\par};
\end{scope}\end{tikzpicture}\newpage
\definecolor{answer}{rgb}{0.850980,0.047059,0.070588}
\null\begin{tikzpicture}[remember picture,overlay]\begin{scope}[shift={(current page.north west)},x=1bp,y=-1bp]
\node[anchor=north west,inner sep=0pt,text=answer,align=left,text width=386.43701171875bp] at (15,15) {\fontsize{7.5bp}{9.5bp}\selectfont 析：\allowbreak{}选B。\allowbreak{}只需每袋\ensuremath{2}千克、\allowbreak{}每千克可长\ensuremath{10}千克豆芽；\allowbreak{}袋数\ensuremath{8}和售价\ensuremath{3}元无关。\allowbreak{}\par};
\end{scope}\end{tikzpicture}\newpage
\definecolor{answer}{rgb}{0.850980,0.047059,0.070588}
\null\begin{tikzpicture}[remember picture,overlay]\begin{scope}[shift={(current page.north west)},x=1bp,y=-1bp]
\node[anchor=north west,inner sep=0pt,text=answer,align=left,text width=467.0bp] at (15,15) {\fontsize{8.5bp}{10.62bp}\selectfont 正确式去括号后应减\ensuremath{70},\allowbreak{}错误式却加\ensuremath{70},\allowbreak{}从同一个基准看，\allowbreak{}两结果相差两个\ensuremath{70}。\allowbreak{}\par};
\end{scope}\end{tikzpicture}\newpage
\definecolor{answer}{rgb}{0.850980,0.047059,0.070588}
\null\begin{tikzpicture}[remember picture,overlay]\begin{scope}[shift={(current page.north west)},x=1bp,y=-1bp]
\node[anchor=north west,inner sep=0pt,text=answer,align=left,text width=468.0bp] at (15,15) {\fontsize{8.5bp}{10.62bp}\selectfont 每个会场都是\ensuremath{15}排，\allowbreak{}把两场座位总数与\ensuremath{48}校老师总数比较。\allowbreak{}\par};
\end{scope}\end{tikzpicture}\newpage
\definecolor{answer}{rgb}{0.850980,0.047059,0.070588}
\null\begin{tikzpicture}[remember picture,overlay]\begin{scope}[shift={(current page.north west)},x=1bp,y=-1bp]
\node[anchor=north west,inner sep=0pt,text=answer,align=left,text width=467.0bp] at (15,15) {\fontsize{8.5bp}{10.62bp}\selectfont B加入前已完成\ensuremath{200}辆，\allowbreak{}距\ensuremath{1000}辆还差\ensuremath{800}辆；\allowbreak{}加入后两车间每小时一共生产\ensuremath{100}辆。\allowbreak{}\par};
\end{scope}\end{tikzpicture}\newpage
\definecolor{answer}{rgb}{0.850980,0.047059,0.070588}
\null\begin{tikzpicture}[remember picture,overlay]\begin{scope}[shift={(current page.north west)},x=1bp,y=-1bp]
\node[anchor=north west,inner sep=0pt,text=answer,align=left,text width=467.0bp] at (15,15) {\fontsize{8.5bp}{10.62bp}\selectfont 按每人每天最低\ensuremath{60}克计算，\allowbreak{}三人每天至少吃\ensuremath{180}克；\allowbreak{}求最多整天数，\allowbreak{}余\ensuremath{100}克不足一天。\allowbreak{}\par};
\end{scope}\end{tikzpicture}\newpage
\definecolor{answer}{rgb}{0.850980,0.047059,0.070588}
\null\begin{tikzpicture}[remember picture,overlay]\begin{scope}[shift={(current page.north west)},x=1bp,y=-1bp]
\node[anchor=north west,inner sep=0pt,text=answer,align=left,text width=469.0bp] at (15,15) {\fontsize{8.5bp}{10.62bp}\selectfont 每买\ensuremath{4}盒获赠\ensuremath{1}盒，\allowbreak{}每组实际得到\ensuremath{5}盒；\allowbreak{}\ensuremath{15}盒正好\ensuremath{3}组，\allowbreak{}只需付\ensuremath{12}盒的钱。\allowbreak{}\par};
\end{scope}\end{tikzpicture}\newpage
\definecolor{answer}{rgb}{0.850980,0.047059,0.070588}
\null\begin{tikzpicture}[remember picture,overlay]\begin{scope}[shift={(current page.north west)},x=1bp,y=-1bp]
\node[anchor=north west,inner sep=0pt,text=answer,align=left,text width=467.0bp] at (15,15) {\fontsize{8.5bp}{10.62bp}\selectfont 先扣\ensuremath{10}辆不合格车，\allowbreak{}\ensuremath{840}辆合格车每\ensuremath{40}辆一车，\allowbreak{}恰好\ensuremath{21}车，\allowbreak{}没有不足一车的情况。\allowbreak{}\par};
\end{scope}\end{tikzpicture}\newpage
\definecolor{answer}{rgb}{0.850980,0.047059,0.070588}
\null\begin{tikzpicture}[remember picture,overlay]\begin{scope}[shift={(current page.north west)},x=1bp,y=-1bp]
\node[anchor=north west,inner sep=0pt,text=answer,align=left,text width=467.0bp] at (15,15) {\fontsize{8.5bp}{10.62bp}\selectfont 两城距离等于甲、\allowbreak{}乙\ensuremath{4}小时已行路程之和再加未行的\ensuremath{50}千米；\allowbreak{}按原图线段保留两车相向行驶示意。\allowbreak{}\par};
\end{scope}\end{tikzpicture}\newpage
\definecolor{answer}{rgb}{0.850980,0.047059,0.070588}
\null\begin{tikzpicture}[remember picture,overlay]\begin{scope}[shift={(current page.north west)},x=1bp,y=-1bp]
\node[anchor=north west,inner sep=0pt,text=answer,align=left,text width=474.0bp] at (15,15) {\fontsize{8.5bp}{10.62bp}\selectfont 先求\ensuremath{3}天完成量，\allowbreak{}再从计划总量中扣除，\allowbreak{}剩余任务平均分到\ensuremath{4}天。\allowbreak{}\par};
\end{scope}\end{tikzpicture}\newpage
\definecolor{answer}{rgb}{0.850980,0.047059,0.070588}
\null\begin{tikzpicture}[remember picture,overlay]\begin{scope}[shift={(current page.north west)},x=1bp,y=-1bp]
\node[anchor=north west,inner sep=0pt,text=answer,align=left,text width=474.0bp] at (15,15) {\fontsize{8.5bp}{10.62bp}\selectfont 每千克少收\ensuremath{16} \ensuremath{-}\ensuremath{13}即\ensuremath{3}元，\allowbreak{}少收\ensuremath{48}元可先求出水果质量，\allowbreak{}再求正确总价。\allowbreak{}\par};
\end{scope}\end{tikzpicture}\newpage
\definecolor{answer}{rgb}{0.850980,0.047059,0.070588}
\null\begin{tikzpicture}[remember picture,overlay]\begin{scope}[shift={(current page.north west)},x=1bp,y=-1bp]
\node[anchor=north west,inner sep=0pt,text=answer,align=left,text width=474.0bp] at (15,15) {\fontsize{8.5bp}{10.62bp}\selectfont 一年级人数是二年级的一半，\allowbreak{}把二年级\ensuremath{130}人平均分成\ensuremath{2}份。\allowbreak{}\par};
\end{scope}\end{tikzpicture}\newpage
\definecolor{answer}{rgb}{0.850980,0.047059,0.070588}
\null\begin{tikzpicture}[remember picture,overlay]\begin{scope}[shift={(current page.north west)},x=1bp,y=-1bp]
\node[anchor=north west,inner sep=0pt,text=answer,align=left,text width=474.375bp] at (15,15) {\fontsize{8.5bp}{10.62bp}\selectfont \ensuremath{260} = \ensuremath{130} × \ensuremath{2}，\allowbreak{}可用乘法分配律把相加再相减的两份\ensuremath{130}合并。\allowbreak{}\par};
\end{scope}\end{tikzpicture}\newpage
\definecolor{answer}{rgb}{0.850980,0.047059,0.070588}
\null\begin{tikzpicture}[remember picture,overlay]\begin{scope}[shift={(current page.north west)},x=1bp,y=-1bp]
\node[anchor=north west,inner sep=0pt,text=answer,align=left,text width=474.0bp] at (15,15) {\fontsize{8.5bp}{10.62bp}\selectfont 先合并二、\allowbreak{}三年级总人数，\allowbreak{}再依次平均分；\allowbreak{}也可直接除以总小组数。\allowbreak{}\par};
\end{scope}\end{tikzpicture}\newpage
\definecolor{answer}{rgb}{0.850980,0.047059,0.070588}
\null\begin{tikzpicture}[remember picture,overlay]\begin{scope}[shift={(current page.north west)},x=1bp,y=-1bp]
\node[anchor=north west,inner sep=0pt,text=answer,align=left,text width=489.0bp] at (15,15) {\fontsize{7.5bp}{9.5bp}\selectfont 析：\allowbreak{}选①②④。\allowbreak{}先由①求编程人数，\allowbreak{}再用②求手工人数，\allowbreak{}最后按④求篮球人数。\allowbreak{}\par};
\end{scope}\end{tikzpicture}\newpage
\definecolor{answer}{rgb}{0.850980,0.047059,0.070588}
\null\begin{tikzpicture}[remember picture,overlay]\begin{scope}[shift={(current page.north west)},x=1bp,y=-1bp]
\node[anchor=north west,inner sep=0pt,text=answer,align=left,text width=489.0bp] at (15,15) {\fontsize{7.5bp}{9.5bp}\selectfont 析：\allowbreak{}先\ensuremath{28} × \ensuremath{3}求编程人数，\allowbreak{}再减\ensuremath{18}得手工人数，\allowbreak{}最后乘\ensuremath{3}求篮球人数；\allowbreak{}依此逐层填框。\allowbreak{}\par};
\end{scope}\end{tikzpicture}\newpage
\definecolor{answer}{rgb}{0.850980,0.047059,0.070588}
\null\begin{tikzpicture}[remember picture,overlay]\begin{scope}[shift={(current page.north west)},x=1bp,y=-1bp]
\node[anchor=north west,inner sep=0pt,text=answer,align=left,text width=469.0bp] at (15,15) {\fontsize{8.5bp}{10.62bp}\selectfont 先依次求编程、\allowbreak{}手工、\allowbreak{}篮球社团人数；\allowbreak{}也可把“少\ensuremath{18}人”的影响先乘\ensuremath{3}后扣除。\allowbreak{}\par};
\end{scope}\end{tikzpicture}\newpage
\definecolor{answer}{rgb}{0.850980,0.047059,0.070588}
\null\begin{tikzpicture}[remember picture,overlay]\begin{scope}[shift={(current page.north west)},x=1bp,y=-1bp]
\node[anchor=north west,inner sep=0pt,text=answer,align=left,text width=494.0bp] at (15,15) {\fontsize{7.5bp}{9.5bp}\selectfont 析：\allowbreak{}速度写作\ensuremath{80}千米/时；\allowbreak{}路程\ensuremath{80} × \ensuremath{4} = \ensuremath{320}千米。\allowbreak{}\par};
\end{scope}\end{tikzpicture}\newpage
\definecolor{answer}{rgb}{0.850980,0.047059,0.070588}
\null\begin{tikzpicture}[remember picture,overlay]\begin{scope}[shift={(current page.north west)},x=1bp,y=-1bp]
\node[anchor=north west,inner sep=0pt,text=answer,align=left,text width=494.0bp] at (15,15) {\fontsize{7.5bp}{9.5bp}\selectfont 析：\allowbreak{}腰不能是\ensuremath{4}厘米（\ensuremath{4} + \ensuremath{4} = \ensuremath{8}),\allowbreak{}故两腰\ensuremath{8}厘米，\allowbreak{}周长\ensuremath{8} × \ensuremath{2} + \ensuremath{4} = \ensuremath{20}厘米。\allowbreak{}\par};
\end{scope}\end{tikzpicture}\newpage
\definecolor{answer}{rgb}{0.850980,0.047059,0.070588}
\null\begin{tikzpicture}[remember picture,overlay]\begin{scope}[shift={(current page.north west)},x=1bp,y=-1bp]
\node[anchor=north west,inner sep=0pt,text=answer,align=left,text width=494.0bp] at (15,15) {\fontsize{7.5bp}{9.5bp}\selectfont 析：\allowbreak{}每分钟\ensuremath{24} ÷ \ensuremath{3} = \ensuremath{8}道；\allowbreak{}\ensuremath{6}分钟\ensuremath{48}道；\allowbreak{}\ensuremath{120} ÷ \ensuremath{8} = \ensuremath{15}分钟。\allowbreak{}\par};
\end{scope}\end{tikzpicture}\newpage
\definecolor{answer}{rgb}{0.850980,0.047059,0.070588}
\null\begin{tikzpicture}[remember picture,overlay]\begin{scope}[shift={(current page.north west)},x=1bp,y=-1bp]
\node[anchor=north west,inner sep=0pt,text=answer,align=left,text width=494.0bp] at (15,15) {\fontsize{7.5bp}{9.5bp}\selectfont 析：\allowbreak{}两腰相等，\allowbreak{}腰长(\ensuremath{40} \ensuremath{-}\ensuremath{8} \ensuremath{-}\ensuremath{10}) ÷ \ensuremath{2} = \ensuremath{11}厘米。\allowbreak{}\par};
\end{scope}\end{tikzpicture}\newpage
\definecolor{answer}{rgb}{0.850980,0.047059,0.070588}
\null\begin{tikzpicture}[remember picture,overlay]\begin{scope}[shift={(current page.north west)},x=1bp,y=-1bp]
\node[anchor=north west,inner sep=0pt,text=answer,align=left,text width=494.0bp] at (15,15) {\fontsize{7.5bp}{9.5bp}\selectfont 析：\allowbreak{}只改变因数分组，\allowbreak{}用乘法结合律。\allowbreak{}\par};
\end{scope}\end{tikzpicture}\newpage
\definecolor{answer}{rgb}{0.850980,0.047059,0.070588}
\null\begin{tikzpicture}[remember picture,overlay]\begin{scope}[shift={(current page.north west)},x=1bp,y=-1bp]
\node[anchor=north west,inner sep=0pt,text=answer,align=left,text width=494.0bp] at (15,15) {\fontsize{7.5bp}{9.5bp}\selectfont 析：\allowbreak{}长方形对边平行，\allowbreak{}邻边垂直。\allowbreak{}\par};
\end{scope}\end{tikzpicture}\newpage
\definecolor{answer}{rgb}{0.850980,0.047059,0.070588}
\null\begin{tikzpicture}[remember picture,overlay]\begin{scope}[shift={(current page.north west)},x=1bp,y=-1bp]
\node[anchor=north west,inner sep=0pt,text=answer,align=left,text width=494.0bp] at (15,15) {\fontsize{7.5bp}{9.5bp}\selectfont 析：\allowbreak{}逆算：\allowbreak{}\ensuremath{240} ÷ \ensuremath{30} × \ensuremath{5} = \ensuremath{40},\allowbreak{}□=\ensuremath{100} \ensuremath{-}\ensuremath{40} = \ensuremath{60}。\allowbreak{}\par};
\end{scope}\end{tikzpicture}\newpage
\definecolor{answer}{rgb}{0.850980,0.047059,0.070588}
\null\begin{tikzpicture}[remember picture,overlay]\begin{scope}[shift={(current page.north west)},x=1bp,y=-1bp]
\node[anchor=north west,inner sep=0pt,text=answer,align=left,text width=494.0bp] at (15,15) {\fontsize{7.5bp}{9.5bp}\selectfont 析：\allowbreak{}每只(\ensuremath{100} \ensuremath{-}\ensuremath{52}) ÷ \ensuremath{12} = \ensuremath{4}元；\allowbreak{}\ensuremath{52} ÷ \ensuremath{4} = \ensuremath{13}只。\allowbreak{}\par};
\end{scope}\end{tikzpicture}\newpage
\definecolor{answer}{rgb}{0.850980,0.047059,0.070588}
\null\begin{tikzpicture}[remember picture,overlay]\begin{scope}[shift={(current page.north west)},x=1bp,y=-1bp]
\node[anchor=north west,inner sep=0pt,text=answer,align=left,text width=494.0bp] at (15,15) {\fontsize{7.5bp}{9.5bp}\selectfont 析：\allowbreak{}同路程用时少者快，\allowbreak{}\ensuremath{4} < \ensuremath{5},\allowbreak{}选A。\allowbreak{}\par};
\end{scope}\end{tikzpicture}\newpage
\definecolor{answer}{rgb}{0.850980,0.047059,0.070588}
\null\begin{tikzpicture}[remember picture,overlay]\begin{scope}[shift={(current page.north west)},x=1bp,y=-1bp]
\node[anchor=north west,inner sep=0pt,text=answer,align=left,text width=494.0bp] at (15,15) {\fontsize{7.5bp}{9.5bp}\selectfont 析：\allowbreak{}乘除本就优先于加法，\allowbreak{}C的括号可去；\allowbreak{}选C。\allowbreak{}\par};
\end{scope}\end{tikzpicture}\newpage
\definecolor{answer}{rgb}{0.850980,0.047059,0.070588}
\null\begin{tikzpicture}[remember picture,overlay]\begin{scope}[shift={(current page.north west)},x=1bp,y=-1bp]
\node[anchor=north west,inner sep=0pt,text=answer,align=left,text width=494.0bp] at (15,15) {\fontsize{7.5bp}{9.5bp}\selectfont 析：\allowbreak{}两次折痕分别水平、\allowbreak{}竖直，\allowbreak{}互相垂直；\allowbreak{}选A。\allowbreak{}\par};
\end{scope}\end{tikzpicture}\newpage
\definecolor{answer}{rgb}{0.850980,0.047059,0.070588}
\null\begin{tikzpicture}[remember picture,overlay]\begin{scope}[shift={(current page.north west)},x=1bp,y=-1bp]
\node[anchor=north west,inner sep=0pt,text=answer,align=left,text width=494.0bp] at (15,15) {\fontsize{7.5bp}{9.5bp}\selectfont 析：\allowbreak{}另两角之和\ensuremath{180}\ensuremath{^\circ}\ensuremath{-}\ensuremath{75}\ensuremath{^\circ}= \ensuremath{105}\ensuremath{^\circ},\allowbreak{} \ensuremath{30}\ensuremath{^\circ}+ \ensuremath{75}\ensuremath{^\circ}符合；\allowbreak{}选C。\allowbreak{}\par};
\end{scope}\end{tikzpicture}\newpage
\definecolor{answer}{rgb}{0.850980,0.047059,0.070588}
\null\begin{tikzpicture}[remember picture,overlay]\begin{scope}[shift={(current page.north west)},x=1bp,y=-1bp]
\node[anchor=north west,inner sep=0pt,text=answer,align=left,text width=494.0bp] at (15,15) {\fontsize{7.5bp}{9.5bp}\selectfont 析：\allowbreak{}路程需①速度及②⑤求时间：\allowbreak{}\ensuremath{7} :\allowbreak{} \ensuremath{55} \ensuremath{-}\ensuremath{7} :\allowbreak{} \ensuremath{40}为\ensuremath{15}分钟，\allowbreak{}选C。\allowbreak{}\par};
\end{scope}\end{tikzpicture}\newpage
\definecolor{answer}{rgb}{0.850980,0.047059,0.070588}
\null\begin{tikzpicture}[remember picture,overlay]\begin{scope}[shift={(current page.north west)},x=1bp,y=-1bp]
\node[anchor=north west,inner sep=0pt,text=answer,align=left,text width=494.0bp] at (15,15) {\fontsize{7.5bp}{9.5bp}\selectfont 析：\allowbreak{}直线两侧各有一条相距\ensuremath{3}厘米的平行线，\allowbreak{}共\ensuremath{2}条，\allowbreak{}判×。\allowbreak{}\par};
\end{scope}\end{tikzpicture}\newpage
\definecolor{answer}{rgb}{0.850980,0.047059,0.070588}
\null\begin{tikzpicture}[remember picture,overlay]\begin{scope}[shift={(current page.north west)},x=1bp,y=-1bp]
\node[anchor=north west,inner sep=0pt,text=answer,align=left,text width=494.0bp] at (15,15) {\fontsize{7.5bp}{9.5bp}\selectfont 析：\allowbreak{}加法两侧的除法与乘法可同时先算，\allowbreak{}判√。\allowbreak{}\par};
\end{scope}\end{tikzpicture}\newpage
\definecolor{answer}{rgb}{0.850980,0.047059,0.070588}
\null\begin{tikzpicture}[remember picture,overlay]\begin{scope}[shift={(current page.north west)},x=1bp,y=-1bp]
\node[anchor=north west,inner sep=0pt,text=answer,align=left,text width=334.3605041503906bp] at (15,15) {\fontsize{7.5bp}{9.5bp}\selectfont 析：\allowbreak{}若每千克\ensuremath{8}元，\allowbreak{}\ensuremath{5}千克应付\ensuremath{8} × \ensuremath{5} = \ensuremath{40}元；\allowbreak{}\ensuremath{36} < \ensuremath{40}，\allowbreak{}单价低于\ensuremath{8}元，\allowbreak{}判√。\allowbreak{}\par};
\end{scope}\end{tikzpicture}\newpage
\definecolor{answer}{rgb}{0.850980,0.047059,0.070588}
\null\begin{tikzpicture}[remember picture,overlay]\begin{scope}[shift={(current page.north west)},x=1bp,y=-1bp]
\node[anchor=north west,inner sep=0pt,text=answer,align=left,text width=302.8664855957031bp] at (15,15) {\fontsize{7.5bp}{9.5bp}\selectfont 析：\allowbreak{}射线无限长，\allowbreak{}不能量得\ensuremath{6}厘米，\allowbreak{}判×。\allowbreak{}\par};
\end{scope}\end{tikzpicture}\newpage
\definecolor{answer}{rgb}{0.850980,0.047059,0.070588}
\null\begin{tikzpicture}[remember picture,overlay]\begin{scope}[shift={(current page.north west)},x=1bp,y=-1bp]
\node[anchor=north west,inner sep=0pt,text=answer,align=left,text width=286.44500732421875bp] at (15,15) {\fontsize{7.5bp}{9.5bp}\selectfont 析：\allowbreak{}速度还须时间单位，\allowbreak{}如千米/时，\allowbreak{}判×。\allowbreak{}\par};
\end{scope}\end{tikzpicture}\newpage
\definecolor{answer}{rgb}{0.850980,0.047059,0.070588}
\null\begin{tikzpicture}[remember picture,overlay]\begin{scope}[shift={(current page.north west)},x=1bp,y=-1bp]
\node[anchor=north west,inner sep=0pt,text=answer,align=left,text width=219.0bp] at (15,15) {\fontsize{7.5bp}{9.5bp}\selectfont 析：\allowbreak{}依次\ensuremath{59}、\allowbreak{}\ensuremath{1062}、\allowbreak{}\ensuremath{592};\allowbreak{}先和后乘再减：\allowbreak{}\ensuremath{18} × (\ensuremath{34} + \ensuremath{25}) \ensuremath{-}\ensuremath{470}。\allowbreak{}\par};
\end{scope}\end{tikzpicture}\newpage
\definecolor{answer}{rgb}{0.850980,0.047059,0.070588}
\null\begin{tikzpicture}[remember picture,overlay]\begin{scope}[shift={(current page.north west)},x=1bp,y=-1bp]
\node[anchor=north west,inner sep=0pt,text=answer,align=left,text width=214.0bp] at (15,15) {\fontsize{7.5bp}{9.5bp}\selectfont 析：\allowbreak{}依次\ensuremath{465}、\allowbreak{}\ensuremath{15}、\allowbreak{}\ensuremath{450};\allowbreak{}先差后除再乘：\allowbreak{}\ensuremath{30} × [(\ensuremath{608} \ensuremath{-}\ensuremath{143}) ÷ \ensuremath{31}]。\allowbreak{}\par};
\end{scope}\end{tikzpicture}\newpage
\definecolor{answer}{rgb}{0.850980,0.047059,0.070588}
\null\begin{tikzpicture}[remember picture,overlay]\begin{scope}[shift={(current page.north west)},x=1bp,y=-1bp]
\node[anchor=north west,inner sep=0pt,text=answer,align=left,text width=484.0bp] at (15,15) {\fontsize{7.5bp}{9.5bp}\selectfont 析：\allowbreak{}B、\allowbreak{}C求出篮球\ensuremath{30} \ensuremath{-}\ensuremath{12} = \ensuremath{18}个；\allowbreak{}还需D的篮球单价及F 的排球总价。\allowbreak{}\par};
\end{scope}\end{tikzpicture}\newpage
\definecolor{answer}{rgb}{0.850980,0.047059,0.070588}
\null\begin{tikzpicture}[remember picture,overlay]\begin{scope}[shift={(current page.north west)},x=1bp,y=-1bp]
\node[anchor=north west,inner sep=0pt,text=answer,align=left,text width=484.0bp] at (15,15) {\fontsize{7.5bp}{9.5bp}\selectfont 析：\allowbreak{}路程=速度×时间。\allowbreak{}甲\ensuremath{250} × \ensuremath{3},\allowbreak{}乙\ensuremath{200} × \ensuremath{3};\allowbreak{}两站距离取和，\allowbreak{}少行路程取差。\allowbreak{}\par};
\end{scope}\end{tikzpicture}\newpage
\definecolor{answer}{rgb}{0.850980,0.047059,0.070588}
\null\begin{tikzpicture}[remember picture,overlay]\begin{scope}[shift={(current page.north west)},x=1bp,y=-1bp]
\node[anchor=north west,inner sep=0pt,text=answer,align=left,text width=232.0bp] at (15,15) {\fontsize{7.5bp}{9.5bp}\selectfont 析：\allowbreak{}分母\ensuremath{10}、\allowbreak{}\ensuremath{100}、\allowbreak{}\ensuremath{1000}对应\ensuremath{1}、\allowbreak{}\ensuremath{2}、\allowbreak{}\ensuremath{3}位小数，\allowbreak{}按位数转换连线。\allowbreak{}\par};
\end{scope}\end{tikzpicture}\newpage
\definecolor{answer}{rgb}{0.850980,0.047059,0.070588}
\null\begin{tikzpicture}[remember picture,overlay]\begin{scope}[shift={(current page.north west)},x=1bp,y=-1bp]
\node[anchor=north west,inner sep=0pt,text=answer,align=left,text width=232.0bp] at (15,15) {\fontsize{7.5bp}{9.5bp}\selectfont 析：\allowbreak{}一位小数为\ensuremath{5.1}至\ensuremath{5.9}共\ensuremath{9}个；\allowbreak{}不限位数则有无数个。\allowbreak{}\par};
\end{scope}\end{tikzpicture}\newpage
\definecolor{answer}{rgb}{0.850980,0.047059,0.070588}
\null\begin{tikzpicture}[remember picture,overlay]\begin{scope}[shift={(current page.north west)},x=1bp,y=-1bp]
\node[anchor=north west,inner sep=0pt,text=answer,align=left,text width=232.0bp] at (15,15) {\fontsize{7.5bp}{9.5bp}\selectfont 析：\allowbreak{}三个\ensuremath{8}依次在千位、\allowbreak{}百位、\allowbreak{}十分位，\allowbreak{}分别表示\ensuremath{8000}、\allowbreak{}\ensuremath{800}、\allowbreak{}\ensuremath{0.8}。\allowbreak{}\par};
\end{scope}\end{tikzpicture}\newpage
\definecolor{answer}{rgb}{0.850980,0.047059,0.070588}
\null\begin{tikzpicture}[remember picture,overlay]\begin{scope}[shift={(current page.north west)},x=1bp,y=-1bp]
\node[anchor=north west,inner sep=0pt,text=answer,align=left,text width=232.0bp] at (15,15) {\fontsize{7.5bp}{9.5bp}\selectfont 析：\allowbreak{}整数部分按整数读，\allowbreak{}小数点读“点”，\allowbreak{}小数部分逐位读。\allowbreak{}\par};
\end{scope}\end{tikzpicture}\newpage
\definecolor{answer}{rgb}{0.850980,0.047059,0.070588}
\null\begin{tikzpicture}[remember picture,overlay]\begin{scope}[shift={(current page.north west)},x=1bp,y=-1bp]
\node[anchor=north west,inner sep=0pt,text=answer,align=left,text width=232.0bp] at (15,15) {\fontsize{7.5bp}{9.5bp}\selectfont 析：\allowbreak{}看\ensuremath{6}的数位：\allowbreak{}\ensuremath{2.61}为十分位，\allowbreak{}\ensuremath{231.36}为百分位，\allowbreak{}\ensuremath{36.51}为个位，\allowbreak{}\ensuremath{4.786}为千分位。\allowbreak{}\par};
\end{scope}\end{tikzpicture}\newpage
\definecolor{answer}{rgb}{0.850980,0.047059,0.070588}
\null\begin{tikzpicture}[remember picture,overlay]\begin{scope}[shift={(current page.north west)},x=1bp,y=-1bp]
\node[anchor=north west,inner sep=0pt,text=answer,align=left,text width=232.0bp] at (15,15) {\fontsize{7.5bp}{9.5bp}\selectfont 析：\allowbreak{}\ensuremath{21.22}读二十一点二二；\allowbreak{}六点三八写\ensuremath{6.38}。\allowbreak{}\par};
\end{scope}\end{tikzpicture}\newpage
\definecolor{answer}{rgb}{0.850980,0.047059,0.070588}
\null\begin{tikzpicture}[remember picture,overlay]\begin{scope}[shift={(current page.north west)},x=1bp,y=-1bp]
\node[anchor=north west,inner sep=0pt,text=answer,align=left,text width=232.0bp] at (15,15) {\fontsize{7.5bp}{9.5bp}\selectfont 析：\allowbreak{}相邻计数单位进率为\ensuremath{10},\allowbreak{}小数部分也如此，\allowbreak{}判√。\allowbreak{}\par};
\end{scope}\end{tikzpicture}\newpage
\definecolor{answer}{rgb}{0.850980,0.047059,0.070588}
\null\begin{tikzpicture}[remember picture,overlay]\begin{scope}[shift={(current page.north west)},x=1bp,y=-1bp]
\node[anchor=north west,inner sep=0pt,text=answer,align=left,text width=232.0bp] at (15,15) {\fontsize{7.5bp}{9.5bp}\selectfont 析：\allowbreak{}一百点零八写\ensuremath{100.08}，\allowbreak{}共有\ensuremath{3}个\ensuremath{0}，\allowbreak{}判√。\allowbreak{}\par};
\end{scope}\end{tikzpicture}\newpage
\definecolor{answer}{rgb}{0.850980,0.047059,0.070588}
\null\begin{tikzpicture}[remember picture,overlay]\begin{scope}[shift={(current page.north west)},x=1bp,y=-1bp]
\node[anchor=north west,inner sep=0pt,text=answer,align=left,text width=232.0bp] at (15,15) {\fontsize{7.5bp}{9.5bp}\selectfont 析：\allowbreak{}例如\ensuremath{0.5} < \ensuremath{1},\allowbreak{}判×。\allowbreak{}\par};
\end{scope}\end{tikzpicture}\newpage
\definecolor{answer}{rgb}{0.850980,0.047059,0.070588}
\null\begin{tikzpicture}[remember picture,overlay]\begin{scope}[shift={(current page.north west)},x=1bp,y=-1bp]
\node[anchor=north west,inner sep=0pt,text=answer,align=left,text width=232.0bp] at (15,15) {\fontsize{7.5bp}{9.5bp}\selectfont 析：\allowbreak{}\ensuremath{58}个\ensuremath{0.01}是\ensuremath{0.58},\allowbreak{}判×。\allowbreak{}\par};
\end{scope}\end{tikzpicture}\newpage
\definecolor{answer}{rgb}{0.850980,0.047059,0.070588}
\null\begin{tikzpicture}[remember picture,overlay]\begin{scope}[shift={(current page.north west)},x=1bp,y=-1bp]
\node[anchor=north west,inner sep=0pt,text=answer,align=left,text width=232.0bp] at (15,15) {\fontsize{7.5bp}{9.5bp}\selectfont 析：\allowbreak{}小数点右一位为十分位；\allowbreak{}左三位为百位，\allowbreak{}分别按\ensuremath{0.1}和\ensuremath{100}计数。\allowbreak{}\par};
\end{scope}\end{tikzpicture}\newpage
\definecolor{answer}{rgb}{0.850980,0.047059,0.070588}
\null\begin{tikzpicture}[remember picture,overlay]\begin{scope}[shift={(current page.north west)},x=1bp,y=-1bp]
\node[anchor=north west,inner sep=0pt,text=answer,align=left,text width=232.0bp] at (15,15) {\fontsize{7.5bp}{9.5bp}\selectfont 析：\allowbreak{}数大小不决定计数单位，\allowbreak{}如\ensuremath{2.01} > \ensuremath{1.9}但\ensuremath{0.01} < \ensuremath{0.1},\allowbreak{}判×。\allowbreak{}\par};
\end{scope}\end{tikzpicture}\newpage
\definecolor{answer}{rgb}{0.850980,0.047059,0.070588}
\null\begin{tikzpicture}[remember picture,overlay]\begin{scope}[shift={(current page.north west)},x=1bp,y=-1bp]
\node[anchor=north west,inner sep=0pt,text=answer,align=left,text width=232.0bp] at (15,15) {\fontsize{7.5bp}{9.5bp}\selectfont 析：\allowbreak{}\ensuremath{80} + \ensuremath{0.6} + \ensuremath{0.05} = \ensuremath{80.65},\allowbreak{}是两位小数，\allowbreak{}计数单位\ensuremath{0.01}。\allowbreak{}\par};
\end{scope}\end{tikzpicture}\newpage
\definecolor{answer}{rgb}{0.850980,0.047059,0.070588}
\null\begin{tikzpicture}[remember picture,overlay]\begin{scope}[shift={(current page.north west)},x=1bp,y=-1bp]
\node[anchor=north west,inner sep=0pt,text=answer,align=left,text width=232.0bp] at (15,15) {\fontsize{7.5bp}{9.5bp}\selectfont 析：\allowbreak{}\ensuremath{50} + \ensuremath{0.8} + \ensuremath{0.02} = \ensuremath{50.82},\allowbreak{}选A。\allowbreak{}\par};
\end{scope}\end{tikzpicture}\newpage
\definecolor{answer}{rgb}{0.850980,0.047059,0.070588}
\null\begin{tikzpicture}[remember picture,overlay]\begin{scope}[shift={(current page.north west)},x=1bp,y=-1bp]
\node[anchor=north west,inner sep=0pt,text=answer,align=left,text width=232.0bp] at (15,15) {\fontsize{7.5bp}{9.5bp}\selectfont 析：\allowbreak{}\ensuremath{0.4} + \ensuremath{0.002} = \ensuremath{0.402},\allowbreak{}以\ensuremath{0.001}计数，\allowbreak{}共\ensuremath{402}个。\allowbreak{}\par};
\end{scope}\end{tikzpicture}\newpage
\definecolor{answer}{rgb}{0.850980,0.047059,0.070588}
\null\begin{tikzpicture}[remember picture,overlay]\begin{scope}[shift={(current page.north west)},x=1bp,y=-1bp]
\node[anchor=north west,inner sep=0pt,text=answer,align=left,text width=232.0bp] at (15,15) {\fontsize{7.5bp}{9.5bp}\selectfont 析：\allowbreak{}\ensuremath{2.5}有\ensuremath{25}个\ensuremath{0.1},\allowbreak{}选A。\allowbreak{}\par};
\end{scope}\end{tikzpicture}\newpage
\definecolor{answer}{rgb}{0.850980,0.047059,0.070588}
\null\begin{tikzpicture}[remember picture,overlay]\begin{scope}[shift={(current page.north west)},x=1bp,y=-1bp]
\node[anchor=north west,inner sep=0pt,text=answer,align=left,text width=232.0bp] at (15,15) {\fontsize{7.5bp}{9.5bp}\selectfont 析：\allowbreak{}\ensuremath{6}在十分位，\allowbreak{}\ensuremath{4.65}也如此，\allowbreak{}选C。\allowbreak{}\par};
\end{scope}\end{tikzpicture}\newpage
\definecolor{answer}{rgb}{0.850980,0.047059,0.070588}
\null\begin{tikzpicture}[remember picture,overlay]\begin{scope}[shift={(current page.north west)},x=1bp,y=-1bp]
\node[anchor=north west,inner sep=0pt,text=answer,align=left,text width=232.0bp] at (15,15) {\fontsize{7.5bp}{9.5bp}\selectfont 析：\allowbreak{}整数末尾\ensuremath{0}不读，\allowbreak{}小数逐位读；\allowbreak{}把两个\ensuremath{0}分别放在整数末尾和小数部分。\allowbreak{}\par};
\end{scope}\end{tikzpicture}\newpage
\definecolor{answer}{rgb}{0.850980,0.047059,0.070588}
\null\begin{tikzpicture}[remember picture,overlay]\begin{scope}[shift={(current page.north west)},x=1bp,y=-1bp]
\node[anchor=north west,inner sep=0pt,text=answer,align=left,text width=232.0bp] at (15,15) {\fontsize{7.5bp}{9.5bp}\selectfont 析：\allowbreak{}两个\ensuremath{0}都在整数末尾才不读，\allowbreak{}小数部分放\ensuremath{8}或\ensuremath{9}。\allowbreak{}\par};
\end{scope}\end{tikzpicture}\newpage
\definecolor{answer}{rgb}{0.850980,0.047059,0.070588}
\null\begin{tikzpicture}[remember picture,overlay]\begin{scope}[shift={(current page.north west)},x=1bp,y=-1bp]
\node[anchor=north west,inner sep=0pt,text=answer,align=left,text width=232.0bp] at (15,15) {\fontsize{7.5bp}{9.5bp}\selectfont 析：\allowbreak{}添\ensuremath{1}得到\ensuremath{0.61},\allowbreak{}计数单位变为\ensuremath{0.01},\allowbreak{}选B。\allowbreak{}\par};
\end{scope}\end{tikzpicture}\newpage
\definecolor{answer}{rgb}{0.850980,0.047059,0.070588}
\null\begin{tikzpicture}[remember picture,overlay]\begin{scope}[shift={(current page.north west)},x=1bp,y=-1bp]
\node[anchor=north west,inner sep=0pt,text=answer,align=left,text width=232.0bp] at (15,15) {\fontsize{7.5bp}{9.5bp}\selectfont 析：\allowbreak{}整数位放\ensuremath{8}或\ensuremath{9},\allowbreak{}小数部分三个数字逐位读，\allowbreak{}两个\ensuremath{0}都读出。\allowbreak{}\par};
\end{scope}\end{tikzpicture}\newpage
\definecolor{answer}{rgb}{0.850980,0.047059,0.070588}
\null\begin{tikzpicture}[remember picture,overlay]\begin{scope}[shift={(current page.north west)},x=1bp,y=-1bp]
\node[anchor=north west,inner sep=0pt,text=answer,align=left,text width=232.0bp] at (15,15) {\fontsize{7.5bp}{9.5bp}\selectfont 析：\allowbreak{}\ensuremath{20.006}读“二十点零零六”，\allowbreak{}读\ensuremath{2}个零，\allowbreak{}选C。\allowbreak{}\par};
\end{scope}\end{tikzpicture}\newpage
\definecolor{answer}{rgb}{0.850980,0.047059,0.070588}
\null\begin{tikzpicture}[remember picture,overlay]\begin{scope}[shift={(current page.north west)},x=1bp,y=-1bp]
\node[anchor=north west,inner sep=0pt,text=answer,align=left,text width=232.0bp] at (15,15) {\fontsize{7.5bp}{9.5bp}\selectfont 析：\allowbreak{}E靠近\ensuremath{3},\allowbreak{}候选数\ensuremath{2.95}最符合位置，\allowbreak{}选A。\allowbreak{}\par};
\end{scope}\end{tikzpicture}\newpage
\definecolor{answer}{rgb}{0.850980,0.047059,0.070588}
\null\begin{tikzpicture}[remember picture,overlay]\begin{scope}[shift={(current page.north west)},x=1bp,y=-1bp]
\node[anchor=north west,inner sep=0pt,text=answer,align=left,text width=231.0bp] at (15,15) {\fontsize{9.0bp}{11.25bp}\selectfont 朵朵和飞飞:\allowbreak{}\ensuremath{30.3};\allowbreak{} 佳佳:\allowbreak{} \ensuremath{3.03}。\allowbreak{}\par};
\end{scope}\end{tikzpicture}\newpage
\definecolor{answer}{rgb}{0.850980,0.047059,0.070588}
\null\begin{tikzpicture}[remember picture,overlay]\begin{scope}[shift={(current page.north west)},x=1bp,y=-1bp]
\node[anchor=north west,inner sep=0pt,text=answer,align=left,text width=232.0bp] at (15,15) {\fontsize{7.5bp}{9.5bp}\selectfont 析：\allowbreak{}朵朵、\allowbreak{}飞飞都表示\ensuremath{30.3};\allowbreak{}佳佳\ensuremath{303}个\ensuremath{0.01}表示\ensuremath{3.03}。\allowbreak{}\par};
\end{scope}\end{tikzpicture}\newpage
\definecolor{answer}{rgb}{0.850980,0.047059,0.070588}
\null\begin{tikzpicture}[remember picture,overlay]\begin{scope}[shift={(current page.north west)},x=1bp,y=-1bp]
\node[anchor=north west,inner sep=0pt,text=answer,align=left,text width=232.0bp] at (15,15) {\fontsize{8.5bp}{10.62bp}\selectfont 答：\allowbreak{}百分位\ensuremath{9}、\allowbreak{}千分位\ensuremath{4},\allowbreak{}这个数是\ensuremath{94.394}。\allowbreak{}\par};
\end{scope}\end{tikzpicture}\newpage
\definecolor{answer}{rgb}{0.850980,0.047059,0.070588}
\null\begin{tikzpicture}[remember picture,overlay]\begin{scope}[shift={(current page.north west)},x=1bp,y=-1bp]
\node[anchor=north west,inner sep=0pt,text=answer,align=left,text width=232.0bp] at (15,15) {\fontsize{7.5bp}{9.5bp}\selectfont 析：\allowbreak{}先比较整数，\allowbreak{}再逐位比较小数；\allowbreak{}不足位补\ensuremath{0}比较。\allowbreak{}\par};
\end{scope}\end{tikzpicture}\newpage
\definecolor{answer}{rgb}{0.850980,0.047059,0.070588}
\null\begin{tikzpicture}[remember picture,overlay]\begin{scope}[shift={(current page.north west)},x=1bp,y=-1bp]
\node[anchor=north west,inner sep=0pt,text=answer,align=left,text width=232.0bp] at (15,15) {\fontsize{7.5bp}{9.5bp}\selectfont 析：\allowbreak{}\ensuremath{2.27} > \ensuremath{2.24} > \ensuremath{2.20},\allowbreak{} A最高、\allowbreak{}C最低。\allowbreak{}\par};
\end{scope}\end{tikzpicture}\newpage
\definecolor{answer}{rgb}{0.850980,0.047059,0.070588}
\null\begin{tikzpicture}[remember picture,overlay]\begin{scope}[shift={(current page.north west)},x=1bp,y=-1bp]
\node[anchor=north west,inner sep=0pt,text=answer,align=left,text width=232.0bp] at (15,15) {\fontsize{7.5bp}{9.5bp}\selectfont 析：\allowbreak{}比较十分位、\allowbreak{}再比百分位，\allowbreak{}高到低为A、\allowbreak{}B、\allowbreak{}C。\allowbreak{}\par};
\end{scope}\end{tikzpicture}\newpage
\definecolor{answer}{rgb}{0.850980,0.047059,0.070588}
\null\begin{tikzpicture}[remember picture,overlay]\begin{scope}[shift={(current page.north west)},x=1bp,y=-1bp]
\node[anchor=north west,inner sep=0pt,text=answer,align=left,text width=232.0bp] at (15,15) {\fontsize{7.5bp}{9.5bp}\selectfont 析：\allowbreak{}\ensuremath{6}开头大于\ensuremath{0}开头；\allowbreak{}十分位\ensuremath{3}最大，\allowbreak{}故\ensuremath{6.3}最大、\allowbreak{}\ensuremath{0.063}最小。\allowbreak{}\par};
\end{scope}\end{tikzpicture}\newpage
\definecolor{answer}{rgb}{0.850980,0.047059,0.070588}
\null\begin{tikzpicture}[remember picture,overlay]\begin{scope}[shift={(current page.north west)},x=1bp,y=-1bp]
\node[anchor=north west,inner sep=0pt,text=answer,align=left,text width=232.0bp] at (15,15) {\fontsize{7.5bp}{9.5bp}\selectfont 析：\allowbreak{}整数部分相同，\allowbreak{}十分位越大数越大，\allowbreak{}最大填\ensuremath{9}、\allowbreak{}最小填\ensuremath{0}。\allowbreak{}\par};
\end{scope}\end{tikzpicture}\newpage
\definecolor{answer}{rgb}{0.850980,0.047059,0.070588}
\null\begin{tikzpicture}[remember picture,overlay]\begin{scope}[shift={(current page.north west)},x=1bp,y=-1bp]
\node[anchor=north west,inner sep=0pt,text=answer,align=left,text width=232.0bp] at (15,15) {\fontsize{7.5bp}{9.5bp}\selectfont 析：\allowbreak{}\ensuremath{66.3}最大、\allowbreak{}\ensuremath{23.6}最小，\allowbreak{}分别是金橘、\allowbreak{}番荔枝。\allowbreak{}\par};
\end{scope}\end{tikzpicture}\newpage
\definecolor{answer}{rgb}{0.850980,0.047059,0.070588}
\null\begin{tikzpicture}[remember picture,overlay]\begin{scope}[shift={(current page.north west)},x=1bp,y=-1bp]
\node[anchor=north west,inner sep=0pt,text=answer,align=left,text width=232.0bp] at (15,15) {\fontsize{7.5bp}{9.5bp}\selectfont 析：\allowbreak{}\ensuremath{23.6} < \ensuremath{36.4} < \ensuremath{66.3},\allowbreak{}依次番荔枝、\allowbreak{}{\scfallback 杧}果、\allowbreak{}金橘。\allowbreak{}\par};
\end{scope}\end{tikzpicture}\newpage
\definecolor{answer}{rgb}{0.850980,0.047059,0.070588}
\null\begin{tikzpicture}[remember picture,overlay]\begin{scope}[shift={(current page.north west)},x=1bp,y=-1bp]
\node[anchor=north west,inner sep=0pt,text=answer,align=left,text width=232.00001525878906bp] at (15,15) {\fontsize{7.5bp}{9.5bp}\selectfont 析：\allowbreak{}先找整数部分：\allowbreak{}\ensuremath{6} < \ensuremath{6.7} < \ensuremath{7};\allowbreak{}\ensuremath{8} < \ensuremath{8.004} < \ensuremath{9};\allowbreak{}\ensuremath{79} > \ensuremath{78.93} >\ensuremath{78}。\allowbreak{}\par};
\end{scope}\end{tikzpicture}\newpage
\definecolor{answer}{rgb}{0.850980,0.047059,0.070588}
\null\begin{tikzpicture}[remember picture,overlay]\begin{scope}[shift={(current page.north west)},x=1bp,y=-1bp]
\node[anchor=north west,inner sep=0pt,text=answer,align=left,text width=232.0bp] at (15,15) {\fontsize{7.5bp}{9.5bp}\selectfont 析：\allowbreak{}不限位数时可不断细分，\allowbreak{}有无数个，\allowbreak{}选C。\allowbreak{}\par};
\end{scope}\end{tikzpicture}\newpage
\definecolor{answer}{rgb}{0.850980,0.047059,0.070588}
\null\begin{tikzpicture}[remember picture,overlay]\begin{scope}[shift={(current page.north west)},x=1bp,y=-1bp]
\node[anchor=north west,inner sep=0pt,text=answer,align=left,text width=232.0bp] at (15,15) {\fontsize{7.5bp}{9.5bp}\selectfont 析：\allowbreak{}与\ensuremath{6}的差为\ensuremath{0.1}、\allowbreak{}\ensuremath{0.02}、\allowbreak{}\ensuremath{0.01},\allowbreak{}最后最小，\allowbreak{}选C。\allowbreak{}\par};
\end{scope}\end{tikzpicture}\newpage
\definecolor{answer}{rgb}{0.850980,0.047059,0.070588}
\null\begin{tikzpicture}[remember picture,overlay]\begin{scope}[shift={(current page.north west)},x=1bp,y=-1bp]
\node[anchor=north west,inner sep=0pt,text=answer,align=left,text width=232.0bp] at (15,15) {\fontsize{7.5bp}{9.5bp}\selectfont 析：\allowbreak{}\ensuremath{0.020} < \ensuremath{0.025}不合格，\allowbreak{}\ensuremath{0.030}与\ensuremath{0.026}均合格，\allowbreak{}选A。\allowbreak{}\par};
\end{scope}\end{tikzpicture}\newpage
\definecolor{answer}{rgb}{0.850980,0.047059,0.070588}
\null\begin{tikzpicture}[remember picture,overlay]\begin{scope}[shift={(current page.north west)},x=1bp,y=-1bp]
\node[anchor=north west,inner sep=0pt,text=answer,align=left,text width=232.0bp] at (15,15) {\fontsize{7.5bp}{9.5bp}\selectfont 析：\allowbreak{}整数\ensuremath{735}添小数点和两个\ensuremath{0},\allowbreak{}写\ensuremath{735.00}。\allowbreak{}\par};
\end{scope}\end{tikzpicture}\newpage
\definecolor{answer}{rgb}{0.850980,0.047059,0.070588}
\null\begin{tikzpicture}[remember picture,overlay]\begin{scope}[shift={(current page.north west)},x=1bp,y=-1bp]
\node[anchor=north west,inner sep=0pt,text=answer,align=left,text width=232.0bp] at (15,15) {\fontsize{7.5bp}{9.5bp}\selectfont 析：\allowbreak{}仅限小数末尾添\ensuremath{0}或去\ensuremath{0},\allowbreak{}大小不变，\allowbreak{}这是小数的性质。\allowbreak{}\par};
\end{scope}\end{tikzpicture}\newpage
\definecolor{answer}{rgb}{0.850980,0.047059,0.070588}
\null\begin{tikzpicture}[remember picture,overlay]\begin{scope}[shift={(current page.north west)},x=1bp,y=-1bp]
\node[anchor=north west,inner sep=0pt,text=answer,align=left,text width=232.0bp] at (15,15) {\fontsize{7.5bp}{9.5bp}\selectfont 析：\allowbreak{}\ensuremath{0.5}以\ensuremath{0.1}计数共\ensuremath{5}个；\allowbreak{}写\ensuremath{0.50}后以\ensuremath{0.01}计数。\allowbreak{}\par};
\end{scope}\end{tikzpicture}\newpage
\definecolor{answer}{rgb}{0.850980,0.047059,0.070588}
\null\begin{tikzpicture}[remember picture,overlay]\begin{scope}[shift={(current page.north west)},x=1bp,y=-1bp]
\node[anchor=north west,inner sep=0pt,text=answer,align=left,text width=232.0bp] at (15,15) {\fontsize{7.5bp}{9.5bp}\selectfont 析：\allowbreak{}千分之一对应三位小数，\allowbreak{}写\ensuremath{7.100}。\allowbreak{}\par};
\end{scope}\end{tikzpicture}\newpage
\definecolor{answer}{rgb}{0.850980,0.047059,0.070588}
\null\begin{tikzpicture}[remember picture,overlay]\begin{scope}[shift={(current page.north west)},x=1bp,y=-1bp]
\node[anchor=north west,inner sep=0pt,text=answer,align=left,text width=232.00001525878906bp] at (15,15) {\fontsize{7.5bp}{9.5bp}\selectfont 析：\allowbreak{}只去小数末尾\ensuremath{0},\allowbreak{}中间与整数部分的\ensuremath{0}保留。\allowbreak{}依次\ensuremath{3.2}、\allowbreak{}\ensuremath{0.9}、\allowbreak{}\ensuremath{96.06}、\allowbreak{}\ensuremath{100}、\allowbreak{}\ensuremath{4.05}、\allowbreak{}\ensuremath{140.1}。\allowbreak{}\par};
\end{scope}\end{tikzpicture}\newpage
\definecolor{answer}{rgb}{0.850980,0.047059,0.070588}
\null\begin{tikzpicture}[remember picture,overlay]\begin{scope}[shift={(current page.north west)},x=1bp,y=-1bp]
\node[anchor=north west,inner sep=0pt,text=answer,align=left,text width=232.0bp] at (15,15) {\fontsize{7.5bp}{9.5bp}\selectfont 析：\allowbreak{}\ensuremath{0.006}与\ensuremath{0.060}都以\ensuremath{0.001}计数；\allowbreak{}\ensuremath{0.06}与\ensuremath{0.060}大小相同但计数单位不同。\allowbreak{}\par};
\end{scope}\end{tikzpicture}\newpage
\definecolor{answer}{rgb}{0.850980,0.047059,0.070588}
\null\begin{tikzpicture}[remember picture,overlay]\begin{scope}[shift={(current page.north west)},x=1bp,y=-1bp]
\node[anchor=north west,inner sep=0pt,text=answer,align=left,text width=232.0bp] at (15,15) {\fontsize{7.5bp}{9.5bp}\selectfont 析：\allowbreak{}保留原大小，\allowbreak{}末尾添\ensuremath{0}成为两位小数：\allowbreak{}\ensuremath{51.40}。\allowbreak{}\par};
\end{scope}\end{tikzpicture}\newpage
\definecolor{answer}{rgb}{0.850980,0.047059,0.070588}
\null\begin{tikzpicture}[remember picture,overlay]\begin{scope}[shift={(current page.north west)},x=1bp,y=-1bp]
\node[anchor=north west,inner sep=0pt,text=answer,align=left,text width=232.0bp] at (15,15) {\fontsize{7.5bp}{9.5bp}\selectfont 析：\allowbreak{}首数若\ensuremath{51.6}会超过第二数；\allowbreak{}依次只可\ensuremath{5.16} < \ensuremath{41.7} < \ensuremath{65.8} <\ensuremath{245.8}。\allowbreak{}\par};
\end{scope}\end{tikzpicture}\newpage
\definecolor{answer}{rgb}{0.850980,0.047059,0.070588}
\null\begin{tikzpicture}[remember picture,overlay]\begin{scope}[shift={(current page.north west)},x=1bp,y=-1bp]
\node[anchor=north west,inner sep=0pt,text=answer,align=left,text width=232.0bp] at (15,15) {\fontsize{7.5bp}{9.5bp}\selectfont 析：\allowbreak{}\ensuremath{80}是整数，\allowbreak{}添\ensuremath{0}为\ensuremath{800}、\allowbreak{}去\ensuremath{0}为\ensuremath{8},\allowbreak{}会改变大小，\allowbreak{}判×。\allowbreak{}\par};
\end{scope}\end{tikzpicture}\newpage
\definecolor{answer}{rgb}{0.850980,0.047059,0.070588}
\null\begin{tikzpicture}[remember picture,overlay]\begin{scope}[shift={(current page.north west)},x=1bp,y=-1bp]
\node[anchor=north west,inner sep=0pt,text=answer,align=left,text width=232.0bp] at (15,15) {\fontsize{7.5bp}{9.5bp}\selectfont 析：\allowbreak{}位数不决定大小，\allowbreak{}例如\ensuremath{0.99} < \ensuremath{1.001},\allowbreak{}判×。\allowbreak{}\par};
\end{scope}\end{tikzpicture}\newpage
\definecolor{answer}{rgb}{0.850980,0.047059,0.070588}
\null\begin{tikzpicture}[remember picture,overlay]\begin{scope}[shift={(current page.north west)},x=1bp,y=-1bp]
\node[anchor=north west,inner sep=0pt,text=answer,align=left,text width=232.0bp] at (15,15) {\fontsize{7.5bp}{9.5bp}\selectfont 析：\allowbreak{}末尾可添任意多个\ensuremath{0},\allowbreak{}大小不变，\allowbreak{}故无数个，\allowbreak{}判√。\allowbreak{}\par};
\end{scope}\end{tikzpicture}\newpage
\definecolor{answer}{rgb}{0.850980,0.047059,0.070588}
\null\begin{tikzpicture}[remember picture,overlay]\begin{scope}[shift={(current page.north west)},x=1bp,y=-1bp]
\node[anchor=north west,inner sep=0pt,text=answer,align=left,text width=232.0bp] at (15,15) {\fontsize{7.5bp}{9.5bp}\selectfont 析：\allowbreak{}\ensuremath{6.750} > \ensuremath{6.075},\allowbreak{}大小不同，\allowbreak{}判×。\allowbreak{}\par};
\end{scope}\end{tikzpicture}\newpage
\definecolor{answer}{rgb}{0.850980,0.047059,0.070588}
\null\begin{tikzpicture}[remember picture,overlay]\begin{scope}[shift={(current page.north west)},x=1bp,y=-1bp]
\node[anchor=north west,inner sep=0pt,text=answer,align=left,text width=232.0bp] at (15,15) {\fontsize{7.5bp}{9.5bp}\selectfont 析：\allowbreak{}只去末尾\ensuremath{0},\allowbreak{} \ensuremath{52.450} = \ensuremath{52.45},\allowbreak{}判√。\allowbreak{}\par};
\end{scope}\end{tikzpicture}\newpage
\definecolor{answer}{rgb}{0.850980,0.047059,0.070588}
\null\begin{tikzpicture}[remember picture,overlay]\begin{scope}[shift={(current page.north west)},x=1bp,y=-1bp]
\node[anchor=north west,inner sep=0pt,text=answer,align=left,text width=232.0bp] at (15,15) {\fontsize{7.5bp}{9.5bp}\selectfont 析：\allowbreak{}分米换米÷\ensuremath{10},\allowbreak{}厘米换米÷\ensuremath{100},\allowbreak{}再补足两位小数。\allowbreak{}\par};
\end{scope}\end{tikzpicture}\newpage
\definecolor{answer}{rgb}{0.850980,0.047059,0.070588}
\null\begin{tikzpicture}[remember picture,overlay]\begin{scope}[shift={(current page.north west)},x=1bp,y=-1bp]
\node[anchor=north west,inner sep=0pt,text=answer,align=left,text width=232.0bp] at (15,15) {\fontsize{7.5bp}{9.5bp}\selectfont 析：\allowbreak{}三位小数要保留\ensuremath{3}位小数部分，\allowbreak{}写\ensuremath{6.300},\allowbreak{}选A。\allowbreak{}\par};
\end{scope}\end{tikzpicture}\newpage
\definecolor{answer}{rgb}{0.850980,0.047059,0.070588}
\null\begin{tikzpicture}[remember picture,overlay]\begin{scope}[shift={(current page.north west)},x=1bp,y=-1bp]
\node[anchor=north west,inner sep=0pt,text=answer,align=left,text width=232.0bp] at (15,15) {\fontsize{7.5bp}{9.5bp}\selectfont 析：\allowbreak{}\ensuremath{1}在十分位、\allowbreak{}百分位或千分位，\allowbreak{}另两位为\ensuremath{0},\allowbreak{}故\ensuremath{5.100}、\allowbreak{}\ensuremath{5.010}、\allowbreak{}\ensuremath{5.001}。\allowbreak{}\par};
\end{scope}\end{tikzpicture}\newpage
\definecolor{answer}{rgb}{0.850980,0.047059,0.070588}
\null\begin{tikzpicture}[remember picture,overlay]\begin{scope}[shift={(current page.north west)},x=1bp,y=-1bp]
\node[anchor=north west,inner sep=0pt,text=answer,align=left,text width=232.0bp] at (15,15) {\fontsize{7.5bp}{9.5bp}\selectfont 析：\allowbreak{}十分之一对应一位小数，\allowbreak{}写\ensuremath{5.0},\allowbreak{}选A。\allowbreak{}\par};
\end{scope}\end{tikzpicture}\newpage
\definecolor{answer}{rgb}{0.850980,0.047059,0.070588}
\null\begin{tikzpicture}[remember picture,overlay]\begin{scope}[shift={(current page.north west)},x=1bp,y=-1bp]
\node[anchor=north west,inner sep=0pt,text=answer,align=left,text width=232.0bp] at (15,15) {\fontsize{7.5bp}{9.5bp}\selectfont 析：\allowbreak{}\ensuremath{67.00}的两个\ensuremath{0}都在小数末尾，\allowbreak{}可删，\allowbreak{}选C。\allowbreak{}\par};
\end{scope}\end{tikzpicture}\newpage
\definecolor{answer}{rgb}{0.850980,0.047059,0.070588}
\null\begin{tikzpicture}[remember picture,overlay]\begin{scope}[shift={(current page.north west)},x=1bp,y=-1bp]
\node[anchor=north west,inner sep=0pt,text=answer,align=left,text width=232.0bp] at (15,15) {\fontsize{7.5bp}{9.5bp}\selectfont 析：\allowbreak{}两个\ensuremath{0}都在末尾：\allowbreak{}\ensuremath{4.500}、\allowbreak{}\ensuremath{5.400}、\allowbreak{}\ensuremath{45.00}、\allowbreak{}\ensuremath{54.00},\allowbreak{}共\ensuremath{4}个。\allowbreak{}\par};
\end{scope}\end{tikzpicture}\newpage
\definecolor{answer}{rgb}{0.850980,0.047059,0.070588}
\null\begin{tikzpicture}[remember picture,overlay]\begin{scope}[shift={(current page.north west)},x=1bp,y=-1bp]
\node[anchor=north west,inner sep=0pt,text=answer,align=left,text width=232.0bp] at (15,15) {\fontsize{7.5bp}{9.5bp}\selectfont 析：\allowbreak{}均不在小数末尾：\allowbreak{}\ensuremath{0.045}、\allowbreak{}\ensuremath{0.054}、\allowbreak{}\ensuremath{0.405}、\allowbreak{}\ensuremath{0.504}、\allowbreak{}\ensuremath{4.005}、\allowbreak{}\ensuremath{5.004}、\allowbreak{}\ensuremath{40.05}、\allowbreak{}\ensuremath{50.04}、\allowbreak{}\ensuremath{400.5}、\allowbreak{}\ensuremath{500.4},\allowbreak{}共\ensuremath{10}个。\allowbreak{}\par};
\end{scope}\end{tikzpicture}\newpage
\definecolor{answer}{rgb}{0.850980,0.047059,0.070588}
\null\begin{tikzpicture}[remember picture,overlay]\begin{scope}[shift={(current page.north west)},x=1bp,y=-1bp]
\node[anchor=north west,inner sep=0pt,text=answer,align=left,text width=231.99996185302734bp] at (15,15) {\fontsize{7.5bp}{9.5bp}\selectfont 析：\allowbreak{}按最小正一位小数\ensuremath{0.1}理解，\allowbreak{}改写为\ensuremath{0.10}，\allowbreak{}选A。\allowbreak{}\par};
\end{scope}\end{tikzpicture}\newpage
\definecolor{answer}{rgb}{0.850980,0.047059,0.070588}
\null\begin{tikzpicture}[remember picture,overlay]\begin{scope}[shift={(current page.north west)},x=1bp,y=-1bp]
\node[anchor=north west,inner sep=0pt,text=answer,align=left,text width=232.0bp] at (15,15) {\fontsize{7.5bp}{9.5bp}\selectfont 析：\allowbreak{}利用乘除互逆，\allowbreak{}依次\ensuremath{4.56};\allowbreak{}\ensuremath{230};\allowbreak{} ÷\ensuremath{100};\allowbreak{} ÷\ensuremath{1000}。\allowbreak{}\par};
\end{scope}\end{tikzpicture}\newpage
\definecolor{answer}{rgb}{0.850980,0.047059,0.070588}
\null\begin{tikzpicture}[remember picture,overlay]\begin{scope}[shift={(current page.north west)},x=1bp,y=-1bp]
\node[anchor=north west,inner sep=0pt,text=answer,align=left,text width=232.0bp] at (15,15) {\fontsize{7.5bp}{9.5bp}\selectfont 析：\allowbreak{}×\ensuremath{10},\allowbreak{}小数点右移\ensuremath{1}位，\allowbreak{}\ensuremath{70.9}。\allowbreak{}\par};
\end{scope}\end{tikzpicture}\newpage
\definecolor{answer}{rgb}{0.850980,0.047059,0.070588}
\null\begin{tikzpicture}[remember picture,overlay]\begin{scope}[shift={(current page.north west)},x=1bp,y=-1bp]
\node[anchor=north west,inner sep=0pt,text=answer,align=left,text width=232.0bp] at (15,15) {\fontsize{7.5bp}{9.5bp}\selectfont 析：\allowbreak{}÷\ensuremath{10},\allowbreak{}左移\ensuremath{1}位，\allowbreak{}\ensuremath{0.709}。\allowbreak{}\par};
\end{scope}\end{tikzpicture}\newpage
\definecolor{answer}{rgb}{0.850980,0.047059,0.070588}
\null\begin{tikzpicture}[remember picture,overlay]\begin{scope}[shift={(current page.north west)},x=1bp,y=-1bp]
\node[anchor=north west,inner sep=0pt,text=answer,align=left,text width=232.0bp] at (15,15) {\fontsize{7.5bp}{9.5bp}\selectfont 析：\allowbreak{}×\ensuremath{1000}右移\ensuremath{3}位为\ensuremath{7090};\allowbreak{}÷\ensuremath{100}左移\ensuremath{2}位为\ensuremath{0.0709}。\allowbreak{}\par};
\end{scope}\end{tikzpicture}\newpage
\definecolor{answer}{rgb}{0.850980,0.047059,0.070588}
\null\begin{tikzpicture}[remember picture,overlay]\begin{scope}[shift={(current page.north west)},x=1bp,y=-1bp]
\node[anchor=north west,inner sep=0pt,text=answer,align=left,text width=232.0bp] at (15,15) {\fontsize{7.5bp}{9.5bp}\selectfont 析：\allowbreak{}去点变\ensuremath{709},\allowbreak{}相当于右移\ensuremath{2}位，\allowbreak{}扩大\ensuremath{100}倍。\allowbreak{}\par};
\end{scope}\end{tikzpicture}\newpage
\definecolor{answer}{rgb}{0.850980,0.047059,0.070588}
\null\begin{tikzpicture}[remember picture,overlay]\begin{scope}[shift={(current page.north west)},x=1bp,y=-1bp]
\node[anchor=north west,inner sep=0pt,text=answer,align=left,text width=232.0bp] at (15,15) {\fontsize{7.5bp}{9.5bp}\selectfont 析：\allowbreak{}乘\ensuremath{100},\allowbreak{}小数点右移两位，\allowbreak{}\ensuremath{25.6}。\allowbreak{}\par};
\end{scope}\end{tikzpicture}\newpage
\definecolor{answer}{rgb}{0.850980,0.047059,0.070588}
\null\begin{tikzpicture}[remember picture,overlay]\begin{scope}[shift={(current page.north west)},x=1bp,y=-1bp]
\node[anchor=north west,inner sep=0pt,text=answer,align=left,text width=232.0bp] at (15,15) {\fontsize{7.5bp}{9.5bp}\selectfont 析：\allowbreak{}左移两位，\allowbreak{}填÷\ensuremath{100}。\allowbreak{}\par};
\end{scope}\end{tikzpicture}\newpage
\definecolor{answer}{rgb}{0.850980,0.047059,0.070588}
\null\begin{tikzpicture}[remember picture,overlay]\begin{scope}[shift={(current page.north west)},x=1bp,y=-1bp]
\node[anchor=north west,inner sep=0pt,text=answer,align=left,text width=232.0bp] at (15,15) {\fontsize{7.5bp}{9.5bp}\selectfont 析：\allowbreak{}原数\ensuremath{5060} ÷ \ensuremath{1000} = \ensuremath{5.06},\allowbreak{}选A。\allowbreak{}\par};
\end{scope}\end{tikzpicture}\newpage
\definecolor{answer}{rgb}{0.850980,0.047059,0.070588}
\null\begin{tikzpicture}[remember picture,overlay]\begin{scope}[shift={(current page.north west)},x=1bp,y=-1bp]
\node[anchor=north west,inner sep=0pt,text=answer,align=left,text width=232.0bp] at (15,15) {\fontsize{7.5bp}{9.5bp}\selectfont 析：\allowbreak{}除\ensuremath{1000}左移三位，\allowbreak{}\ensuremath{0.512}。\allowbreak{}\par};
\end{scope}\end{tikzpicture}\newpage
\definecolor{answer}{rgb}{0.850980,0.047059,0.070588}
\null\begin{tikzpicture}[remember picture,overlay]\begin{scope}[shift={(current page.north west)},x=1bp,y=-1bp]
\node[anchor=north west,inner sep=0pt,text=answer,align=left,text width=232.0bp] at (15,15) {\fontsize{7.5bp}{9.5bp}\selectfont 析：\allowbreak{}左移一位，\allowbreak{}填÷\ensuremath{10}。\allowbreak{}\par};
\end{scope}\end{tikzpicture}\newpage
\definecolor{answer}{rgb}{0.850980,0.047059,0.070588}
\null\begin{tikzpicture}[remember picture,overlay]\begin{scope}[shift={(current page.north west)},x=1bp,y=-1bp]
\node[anchor=north west,inner sep=0pt,text=answer,align=left,text width=232.0bp] at (15,15) {\fontsize{7.5bp}{9.5bp}\selectfont 析：\allowbreak{}合起来左移\ensuremath{1}位，\allowbreak{}原数\ensuremath{2.45} × \ensuremath{10} = \ensuremath{24.5},\allowbreak{}选A。\allowbreak{}\par};
\end{scope}\end{tikzpicture}\newpage
\definecolor{answer}{rgb}{0.850980,0.047059,0.070588}
\null\begin{tikzpicture}[remember picture,overlay]\begin{scope}[shift={(current page.north west)},x=1bp,y=-1bp]
\node[anchor=north west,inner sep=0pt,text=answer,align=left,text width=232.0bp] at (15,15) {\fontsize{7.5bp}{9.5bp}\selectfont 析：\allowbreak{}\ensuremath{0.35}去点变\ensuremath{35},\allowbreak{}右移两位，\allowbreak{}扩大\ensuremath{100}倍，\allowbreak{}选A。\allowbreak{}\par};
\end{scope}\end{tikzpicture}\newpage
\definecolor{answer}{rgb}{0.850980,0.047059,0.070588}
\null\begin{tikzpicture}[remember picture,overlay]\begin{scope}[shift={(current page.north west)},x=1bp,y=-1bp]
\node[anchor=north west,inner sep=0pt,text=answer,align=left,text width=232.0bp] at (15,15) {\fontsize{7.5bp}{9.5bp}\selectfont 析：\allowbreak{}合起来右移\ensuremath{2}位，\allowbreak{}扩大\ensuremath{100}倍，\allowbreak{}选B。\allowbreak{}\par};
\end{scope}\end{tikzpicture}\newpage
\definecolor{answer}{rgb}{0.850980,0.047059,0.070588}
\null\begin{tikzpicture}[remember picture,overlay]\begin{scope}[shift={(current page.north west)},x=1bp,y=-1bp]
\node[anchor=north west,inner sep=0pt,text=answer,align=left,text width=232.0bp] at (15,15) {\fontsize{7.5bp}{9.5bp}\selectfont 析：\allowbreak{}\ensuremath{800.02}整数末尾\ensuremath{0}不读，\allowbreak{}小数的\ensuremath{0}读出，\allowbreak{}共\ensuremath{1}个。\allowbreak{}\par};
\end{scope}\end{tikzpicture}\newpage
\definecolor{answer}{rgb}{0.850980,0.047059,0.070588}
\null\begin{tikzpicture}[remember picture,overlay]\begin{scope}[shift={(current page.north west)},x=1bp,y=-1bp]
\node[anchor=north west,inner sep=0pt,text=answer,align=left,text width=232.0bp] at (15,15) {\fontsize{7.5bp}{9.5bp}\selectfont 析：\allowbreak{}\ensuremath{8.0002}小数部分逐位读，\allowbreak{}\ensuremath{3}个\ensuremath{0}都读。\allowbreak{}\par};
\end{scope}\end{tikzpicture}\newpage
\definecolor{answer}{rgb}{0.850980,0.047059,0.070588}
\null\begin{tikzpicture}[remember picture,overlay]\begin{scope}[shift={(current page.north west)},x=1bp,y=-1bp]
\node[anchor=north west,inner sep=0pt,text=answer,align=left,text width=232.0bp] at (15,15) {\fontsize{7.5bp}{9.5bp}\selectfont 析：\allowbreak{}\ensuremath{8000.2}的\ensuremath{0}都在整数末尾，\allowbreak{}一个也不读。\allowbreak{}\par};
\end{scope}\end{tikzpicture}\newpage
\definecolor{answer}{rgb}{0.850980,0.047059,0.070588}
\null\begin{tikzpicture}[remember picture,overlay]\begin{scope}[shift={(current page.north west)},x=1bp,y=-1bp]
\node[anchor=north west,inner sep=0pt,text=answer,align=left,text width=232.0bp] at (15,15) {\fontsize{7.5bp}{9.5bp}\selectfont 析：\allowbreak{}甲=\ensuremath{7.04} ÷ \ensuremath{100} = \ensuremath{0.0704};\allowbreak{}乙=\ensuremath{7.04} × \ensuremath{1000} = \ensuremath{7040}。\allowbreak{}\par};
\end{scope}\end{tikzpicture}\newpage
\definecolor{answer}{rgb}{0.850980,0.047059,0.070588}
\null\begin{tikzpicture}[remember picture,overlay]\begin{scope}[shift={(current page.north west)},x=1bp,y=-1bp]
\node[anchor=north west,inner sep=0pt,text=answer,align=left,text width=222.0bp] at (15,15) {\fontsize{8.5bp}{10.62bp}\selectfont 答：\allowbreak{}分别需要\ensuremath{12}米、\allowbreak{}\ensuremath{120}米胶带。\allowbreak{}\par};
\end{scope}\end{tikzpicture}\newpage
\definecolor{answer}{rgb}{0.850980,0.047059,0.070588}
\null\begin{tikzpicture}[remember picture,overlay]\begin{scope}[shift={(current page.north west)},x=1bp,y=-1bp]
\node[anchor=north west,inner sep=0pt,text=answer,align=left,text width=232.0bp] at (15,15) {\fontsize{8.5bp}{10.62bp}\selectfont 答：\allowbreak{}大约能生产\ensuremath{7}吨塑料颗粒。\allowbreak{}\par};
\end{scope}\end{tikzpicture}\newpage
\definecolor{answer}{rgb}{0.850980,0.047059,0.070588}
\null\begin{tikzpicture}[remember picture,overlay]\begin{scope}[shift={(current page.north west)},x=1bp,y=-1bp]
\node[anchor=north west,inner sep=0pt,text=answer,align=left,text width=232.0bp] at (15,15) {\fontsize{8.5bp}{10.62bp}\selectfont 示例问题：\allowbreak{}回收\ensuremath{100}吨废旧衣物，\allowbreak{}大约能生产多少吨再生纤维？\par};
\end{scope}\end{tikzpicture}\newpage
\definecolor{answer}{rgb}{0.850980,0.047059,0.070588}
\null\begin{tikzpicture}[remember picture,overlay]\begin{scope}[shift={(current page.north west)},x=1bp,y=-1bp]
\node[anchor=north west,inner sep=0pt,text=answer,align=left,text width=222.0bp] at (15,15) {\fontsize{8.5bp}{10.62bp}\selectfont 答：\allowbreak{}分别约排放\ensuremath{7800}克、\allowbreak{}\ensuremath{78}千克二氧化碳。\allowbreak{}\par};
\end{scope}\end{tikzpicture}\newpage
\definecolor{answer}{rgb}{0.850980,0.047059,0.070588}
\null\begin{tikzpicture}[remember picture,overlay]\begin{scope}[shift={(current page.north west)},x=1bp,y=-1bp]
\node[anchor=north west,inner sep=0pt,text=answer,align=left,text width=232.0bp] at (15,15) {\fontsize{8.5bp}{10.62bp}\selectfont 答：\allowbreak{}大约能生产\ensuremath{80}吨再生纤维。\allowbreak{}\par};
\end{scope}\end{tikzpicture}\newpage
\definecolor{answer}{rgb}{0.850980,0.047059,0.070588}
\null\begin{tikzpicture}[remember picture,overlay]\begin{scope}[shift={(current page.north west)},x=1bp,y=-1bp]
\node[anchor=north west,inner sep=0pt,text=answer,align=left,text width=232.0bp] at (15,15) {\fontsize{7.5bp}{9.5bp}\selectfont 析：\allowbreak{}按进率÷\ensuremath{10}、\allowbreak{}÷\ensuremath{100}或÷\ensuremath{1000},\allowbreak{}混合单位先化同一单位。\allowbreak{}\par};
\end{scope}\end{tikzpicture}\newpage
\definecolor{answer}{rgb}{0.850980,0.047059,0.070588}
\null\begin{tikzpicture}[remember picture,overlay]\begin{scope}[shift={(current page.north west)},x=1bp,y=-1bp]
\node[anchor=north west,inner sep=0pt,text=answer,align=left,text width=232bp] at (15,15) {\fontsize{7.5bp}{9.5bp}\selectfont 析：\allowbreak{}千米/米、\allowbreak{}吨/千克、\allowbreak{}千克/克进率\ensuremath{1000};\allowbreak{}米/厘米进率\ensuremath{100};\allowbreak{}元/分进率\ensuremath{100}。\allowbreak{}\par};
\end{scope}\end{tikzpicture}\newpage
\definecolor{answer}{rgb}{0.850980,0.047059,0.070588}
\null\begin{tikzpicture}[remember picture,overlay]\begin{scope}[shift={(current page.north west)},x=1bp,y=-1bp]
\node[anchor=north west,inner sep=0pt,text=answer,align=left,text width=232.0bp] at (15,15) {\fontsize{8.5bp}{10.62bp}\selectfont 比较后依次填：\allowbreak{}>、\allowbreak{}>、\allowbreak{}=、\allowbreak{}>、\allowbreak{}<。\allowbreak{}\par};
\end{scope}\end{tikzpicture}\newpage
\definecolor{answer}{rgb}{0.850980,0.047059,0.070588}
\null\begin{tikzpicture}[remember picture,overlay]\begin{scope}[shift={(current page.north west)},x=1bp,y=-1bp]
\node[anchor=north west,inner sep=0pt,text=answer,align=left,text width=232.0bp] at (15,15) {\fontsize{8.5bp}{10.62bp}\selectfont 答：\allowbreak{}一天用药\ensuremath{30}克，\allowbreak{}合\ensuremath{0.03}千克。\allowbreak{}\par};
\end{scope}\end{tikzpicture}\newpage
\definecolor{answer}{rgb}{0.850980,0.047059,0.070588}
\null\begin{tikzpicture}[remember picture,overlay]\begin{scope}[shift={(current page.north west)},x=1bp,y=-1bp]
\node[anchor=north west,inner sep=0pt,text=answer,align=left,text width=232.0bp] at (15,15) {\fontsize{8.5bp}{10.62bp}\selectfont 答：\allowbreak{}单程需要\ensuremath{5}分钟，\allowbreak{}往返共\ensuremath{10}分钟。\allowbreak{}\par};
\end{scope}\end{tikzpicture}\newpage
\definecolor{answer}{rgb}{0.850980,0.047059,0.070588}
\null\begin{tikzpicture}[remember picture,overlay]\begin{scope}[shift={(current page.north west)},x=1bp,y=-1bp]
\node[anchor=north west,inner sep=0pt,text=answer,align=left,text width=232.0bp] at (15,15) {\fontsize{8.5bp}{10.62bp}\selectfont 答：\allowbreak{}每天至少摄入\ensuremath{0.3}千克；\allowbreak{}\ensuremath{100}天至少摄入\ensuremath{30}千克蔬菜。\allowbreak{}\par};
\end{scope}\end{tikzpicture}\newpage
\definecolor{answer}{rgb}{0.850980,0.047059,0.070588}
\null\begin{tikzpicture}[remember picture,overlay]\begin{scope}[shift={(current page.north west)},x=1bp,y=-1bp]
\node[anchor=north west,inner sep=0pt,text=answer,align=left,text width=232.0bp] at (15,15) {\fontsize{8.5bp}{10.62bp}\selectfont 答：\allowbreak{}\ensuremath{100}天最多摄入油\ensuremath{3000}克，\allowbreak{}合\ensuremath{3}千克。\allowbreak{}\par};
\end{scope}\end{tikzpicture}\newpage
\definecolor{answer}{rgb}{0.850980,0.047059,0.070588}
\null\begin{tikzpicture}[remember picture,overlay]\begin{scope}[shift={(current page.north west)},x=1bp,y=-1bp]
\node[anchor=north west,inner sep=0pt,text=answer,align=left,text width=232.0bp] at (15,15) {\fontsize{7.5bp}{9.5bp}\selectfont 析：\allowbreak{}第一式可千米/米或米/毫米；\allowbreak{}第二式可厘米/米或毫米/分米，\allowbreak{}分别按\ensuremath{1000}、\allowbreak{}\ensuremath{100}换算。\allowbreak{}\par};
\end{scope}\end{tikzpicture}\newpage
\definecolor{answer}{rgb}{0.850980,0.047059,0.070588}
\null\begin{tikzpicture}[remember picture,overlay]\begin{scope}[shift={(current page.north west)},x=1bp,y=-1bp]
\node[anchor=north west,inner sep=0pt,text=answer,align=left,text width=232.0bp] at (15,15) {\fontsize{7.5bp}{9.5bp}\selectfont 析：\allowbreak{}\ensuremath{6.732}保留一位看\ensuremath{3}不进位，\allowbreak{}为\ensuremath{6.7};\allowbreak{}保留两位看\ensuremath{2},\allowbreak{}为\ensuremath{6.73}。\allowbreak{}\par};
\end{scope}\end{tikzpicture}\newpage
\definecolor{answer}{rgb}{0.850980,0.047059,0.070588}
\null\begin{tikzpicture}[remember picture,overlay]\begin{scope}[shift={(current page.north west)},x=1bp,y=-1bp]
\node[anchor=north west,inner sep=0pt,text=answer,align=left,text width=232.0bp] at (15,15) {\fontsize{7.5bp}{9.5bp}\selectfont 析：\allowbreak{}千分位至少\ensuremath{5}才进位，\allowbreak{}最大\ensuremath{9}、\allowbreak{}最小\ensuremath{5}。\allowbreak{}\par};
\end{scope}\end{tikzpicture}\newpage
\definecolor{answer}{rgb}{0.850980,0.047059,0.070588}
\null\begin{tikzpicture}[remember picture,overlay]\begin{scope}[shift={(current page.north west)},x=1bp,y=-1bp]
\node[anchor=north west,inner sep=0pt,text=answer,align=left,text width=232.0bp] at (15,15) {\fontsize{7.5bp}{9.5bp}\selectfont 析：\allowbreak{}保留两位看千分位\ensuremath{9},\allowbreak{}进\ensuremath{1}得\ensuremath{71.68}。\allowbreak{}\par};
\end{scope}\end{tikzpicture}\newpage
\definecolor{answer}{rgb}{0.850980,0.047059,0.070588}
\null\begin{tikzpicture}[remember picture,overlay]\begin{scope}[shift={(current page.north west)},x=1bp,y=-1bp]
\node[anchor=north west,inner sep=0pt,text=answer,align=left,text width=232.0bp] at (15,15) {\fontsize{7.5bp}{9.5bp}\selectfont 析：\allowbreak{}保留整数看十分位\ensuremath{5},\allowbreak{}进\ensuremath{1}得\ensuremath{19}厘米。\allowbreak{}\par};
\end{scope}\end{tikzpicture}\newpage
\definecolor{answer}{rgb}{0.850980,0.047059,0.070588}
\null\begin{tikzpicture}[remember picture,overlay]\begin{scope}[shift={(current page.north west)},x=1bp,y=-1bp]
\node[anchor=north west,inner sep=0pt,text=answer,align=left,text width=232.0bp] at (15,15) {\fontsize{7.5bp}{9.5bp}\selectfont 析：\allowbreak{}精确到十分位；\allowbreak{}原值至少\ensuremath{5.75}万元但小于\ensuremath{5.85}万元。\allowbreak{}\par};
\end{scope}\end{tikzpicture}\newpage
\definecolor{answer}{rgb}{0.850980,0.047059,0.070588}
\null\begin{tikzpicture}[remember picture,overlay]\begin{scope}[shift={(current page.north west)},x=1bp,y=-1bp]
\node[anchor=north west,inner sep=0pt,text=answer,align=left,text width=232.0bp] at (15,15) {\fontsize{7.5bp}{9.5bp}\selectfont 析：\allowbreak{}\ensuremath{1} < \ensuremath{1.731} < \ensuremath{2},\allowbreak{}距\ensuremath{2}较近。\allowbreak{}\par};
\end{scope}\end{tikzpicture}\newpage
\definecolor{answer}{rgb}{0.850980,0.047059,0.070588}
\null\begin{tikzpicture}[remember picture,overlay]\begin{scope}[shift={(current page.north west)},x=1bp,y=-1bp]
\node[anchor=north west,inner sep=0pt,text=answer,align=left,text width=232.0bp] at (15,15) {\fontsize{7.5bp}{9.5bp}\selectfont 析：\allowbreak{}\ensuremath{6} < \ensuremath{6.77} < \ensuremath{7},\allowbreak{}距\ensuremath{7}较近。\allowbreak{}\par};
\end{scope}\end{tikzpicture}\newpage
\definecolor{answer}{rgb}{0.850980,0.047059,0.070588}
\null\begin{tikzpicture}[remember picture,overlay]\begin{scope}[shift={(current page.north west)},x=1bp,y=-1bp]
\node[anchor=north west,inner sep=0pt,text=answer,align=left,text width=232.0bp] at (15,15) {\fontsize{7.5bp}{9.5bp}\selectfont 析：\allowbreak{}近似数也可比准确数大，\allowbreak{}如\ensuremath{3.8}约\ensuremath{4},\allowbreak{}判×。\allowbreak{}\par};
\end{scope}\end{tikzpicture}\newpage
\definecolor{answer}{rgb}{0.850980,0.047059,0.070588}
\null\begin{tikzpicture}[remember picture,overlay]\begin{scope}[shift={(current page.north west)},x=1bp,y=-1bp]
\node[anchor=north west,inner sep=0pt,text=answer,align=left,text width=232.0bp] at (15,15) {\fontsize{7.5bp}{9.5bp}\selectfont 析：\allowbreak{}可取\ensuremath{3.15}至\ensuremath{3.24},\allowbreak{}最小\ensuremath{3.15},\allowbreak{}判√。\allowbreak{}\par};
\end{scope}\end{tikzpicture}\newpage
\definecolor{answer}{rgb}{0.850980,0.047059,0.070588}
\null\begin{tikzpicture}[remember picture,overlay]\begin{scope}[shift={(current page.north west)},x=1bp,y=-1bp]
\node[anchor=north west,inner sep=0pt,text=answer,align=left,text width=232.0bp] at (15,15) {\fontsize{7.5bp}{9.5bp}\selectfont 析：\allowbreak{}\ensuremath{5.0}的\ensuremath{0}表示精确到十分位，\allowbreak{}不能省，\allowbreak{}判√。\allowbreak{}\par};
\end{scope}\end{tikzpicture}\newpage
\definecolor{answer}{rgb}{0.850980,0.047059,0.070588}
\null\begin{tikzpicture}[remember picture,overlay]\begin{scope}[shift={(current page.north west)},x=1bp,y=-1bp]
\node[anchor=north west,inner sep=0pt,text=answer,align=left,text width=228.3333740234375bp] at (15,15) {\fontsize{9.0bp}{11.25bp}\selectfont 例:\allowbreak{} 准确数:\allowbreak{} 一周\ensuremath{7}天、\allowbreak{}一年\ensuremath{12}个月。\allowbreak{}近似数:\allowbreak{} 一支笔长约\ensuremath{15}厘米,\allowbreak{}一棵树高约\ensuremath{8}米。\allowbreak{}\par};
\end{scope}\end{tikzpicture}\newpage
\definecolor{answer}{rgb}{0.850980,0.047059,0.070588}
\null\begin{tikzpicture}[remember picture,overlay]\begin{scope}[shift={(current page.north west)},x=1bp,y=-1bp]
\node[anchor=north west,inner sep=0pt,text=answer,align=left,text width=232.0bp] at (15,15) {\fontsize{7.5bp}{9.5bp}\selectfont 析：\allowbreak{}准确数由确定计数得出；\allowbreak{}近似数来自测量或估计，\allowbreak{}所举例子符合定义即可。\allowbreak{}\par};
\end{scope}\end{tikzpicture}\newpage
\definecolor{answer}{rgb}{0.850980,0.047059,0.070588}
\null\begin{tikzpicture}[remember picture,overlay]\begin{scope}[shift={(current page.north west)},x=1bp,y=-1bp]
\node[anchor=north west,inner sep=0pt,text=answer,align=left,text width=232.0bp] at (15,15) {\fontsize{7.5bp}{9.5bp}\selectfont 析：\allowbreak{}保留两位小数即精确到百分位。\allowbreak{}\par};
\end{scope}\end{tikzpicture}\newpage
\definecolor{answer}{rgb}{0.850980,0.047059,0.070588}
\null\begin{tikzpicture}[remember picture,overlay]\begin{scope}[shift={(current page.north west)},x=1bp,y=-1bp]
\node[anchor=north west,inner sep=0pt,text=answer,align=left,text width=232.0bp] at (15,15) {\fontsize{7.5bp}{9.5bp}\selectfont 析：\allowbreak{}改“万”÷\ensuremath{10000}得\ensuremath{8800}万；\allowbreak{}改“亿”÷\ensuremath{100000000}得\ensuremath{53}亿。\allowbreak{}\par};
\end{scope}\end{tikzpicture}\newpage
\definecolor{answer}{rgb}{0.850980,0.047059,0.070588}
\null\begin{tikzpicture}[remember picture,overlay]\begin{scope}[shift={(current page.north west)},x=1bp,y=-1bp]
\node[anchor=north west,inner sep=0pt,text=answer,align=left,text width=232.0bp] at (15,15) {\fontsize{7.5bp}{9.5bp}\selectfont 析：\allowbreak{}改成“万”作单位，\allowbreak{}÷\ensuremath{10000},\allowbreak{} \ensuremath{166.49}万。\allowbreak{}\par};
\end{scope}\end{tikzpicture}\newpage
\definecolor{answer}{rgb}{0.850980,0.047059,0.070588}
\null\begin{tikzpicture}[remember picture,overlay]\begin{scope}[shift={(current page.north west)},x=1bp,y=-1bp]
\node[anchor=north west,inner sep=0pt,text=answer,align=left,text width=232.0bp] at (15,15) {\fontsize{7.5bp}{9.5bp}\selectfont 析：\allowbreak{}\ensuremath{706500000}为\ensuremath{7.065}亿，\allowbreak{}千分位\ensuremath{5}进\ensuremath{1},\allowbreak{}约\ensuremath{7.07}亿。\allowbreak{}\par};
\end{scope}\end{tikzpicture}\newpage
\definecolor{answer}{rgb}{0.850980,0.047059,0.070588}
\null\begin{tikzpicture}[remember picture,overlay]\begin{scope}[shift={(current page.north west)},x=1bp,y=-1bp]
\node[anchor=north west,inner sep=0pt,text=answer,align=left,text width=232.0bp] at (15,15) {\fontsize{7.5bp}{9.5bp}\selectfont 析：\allowbreak{}\ensuremath{997000000} ÷ \ensuremath{10000} = \ensuremath{99700}万，\allowbreak{}注意保留两个\ensuremath{9}。\allowbreak{}\par};
\end{scope}\end{tikzpicture}\newpage
\definecolor{answer}{rgb}{0.850980,0.047059,0.070588}
\null\begin{tikzpicture}[remember picture,overlay]\begin{scope}[shift={(current page.north west)},x=1bp,y=-1bp]
\node[anchor=north west,inner sep=0pt,text=answer,align=left,text width=232.0bp] at (15,15) {\fontsize{7.5bp}{9.5bp}\selectfont 析：\allowbreak{}\ensuremath{4614.53}万保留一位，\allowbreak{}百分位\ensuremath{3}不进，\allowbreak{}约\ensuremath{4614.5}万。\allowbreak{}\par};
\end{scope}\end{tikzpicture}\newpage
\definecolor{answer}{rgb}{0.850980,0.047059,0.070588}
\null\begin{tikzpicture}[remember picture,overlay]\begin{scope}[shift={(current page.north west)},x=1bp,y=-1bp]
\node[anchor=north west,inner sep=0pt,text=answer,align=left,text width=232.0bp] at (15,15) {\fontsize{7.5bp}{9.5bp}\selectfont 析：\allowbreak{}小数点左移\ensuremath{4}位，\allowbreak{}\ensuremath{1956.8546}万。\allowbreak{}\par};
\end{scope}\end{tikzpicture}\newpage
\definecolor{answer}{rgb}{0.850980,0.047059,0.070588}
\null\begin{tikzpicture}[remember picture,overlay]\begin{scope}[shift={(current page.north west)},x=1bp,y=-1bp]
\node[anchor=north west,inner sep=0pt,text=answer,align=left,text width=232.0bp] at (15,15) {\fontsize{7.5bp}{9.5bp}\selectfont 析：\allowbreak{}\ensuremath{28.71}亿保留一位，\allowbreak{}百分位\ensuremath{1}不进，\allowbreak{}约\ensuremath{28.7}亿。\allowbreak{}\par};
\end{scope}\end{tikzpicture}\newpage
\definecolor{answer}{rgb}{0.850980,0.047059,0.070588}
\null\begin{tikzpicture}[remember picture,overlay]\begin{scope}[shift={(current page.north west)},x=1bp,y=-1bp]
\node[anchor=north west,inner sep=0pt,text=answer,align=left,text width=232.0bp] at (15,15) {\fontsize{7.5bp}{9.5bp}\selectfont 析：\allowbreak{}÷\ensuremath{10000},\allowbreak{} \ensuremath{80500}万户。\allowbreak{}\par};
\end{scope}\end{tikzpicture}\newpage
\definecolor{answer}{rgb}{0.850980,0.047059,0.070588}
\null\begin{tikzpicture}[remember picture,overlay]\begin{scope}[shift={(current page.north west)},x=1bp,y=-1bp]
\node[anchor=north west,inner sep=0pt,text=answer,align=left,text width=232.0bp] at (15,15) {\fontsize{7.5bp}{9.5bp}\selectfont 析：\allowbreak{}九十二万五千写\ensuremath{925000},\allowbreak{}改为\ensuremath{92.5}万。\allowbreak{}\par};
\end{scope}\end{tikzpicture}\newpage
\definecolor{answer}{rgb}{0.850980,0.047059,0.070588}
\null\begin{tikzpicture}[remember picture,overlay]\begin{scope}[shift={(current page.north west)},x=1bp,y=-1bp]
\node[anchor=north west,inner sep=0pt,text=answer,align=left,text width=232.0bp] at (15,15) {\fontsize{7.5bp}{9.5bp}\selectfont 析：\allowbreak{}改成亿是\ensuremath{60}.□\ensuremath{5},\allowbreak{}填\ensuremath{0}得\ensuremath{60.05}约\ensuremath{60.1};\allowbreak{}填\ensuremath{1}就约\ensuremath{60.2},\allowbreak{}选A。\allowbreak{}\par};
\end{scope}\end{tikzpicture}\newpage
\definecolor{answer}{rgb}{0.850980,0.047059,0.070588}
\null\begin{tikzpicture}[remember picture,overlay]\begin{scope}[shift={(current page.north west)},x=1bp,y=-1bp]
\node[anchor=north west,inner sep=0pt,text=answer,align=left,text width=232.0bp] at (15,15) {\fontsize{7.5bp}{9.5bp}\selectfont 析：\allowbreak{}准确改写应为\ensuremath{2.885}亿；\allowbreak{}题目未要求取近似数，\allowbreak{}判×。\allowbreak{}\par};
\end{scope}\end{tikzpicture}\newpage
\definecolor{answer}{rgb}{0.850980,0.047059,0.070588}
\null\begin{tikzpicture}[remember picture,overlay]\begin{scope}[shift={(current page.north west)},x=1bp,y=-1bp]
\node[anchor=north west,inner sep=0pt,text=answer,align=left,text width=232.0bp] at (15,15) {\fontsize{7.5bp}{9.5bp}\selectfont 析：\allowbreak{}千分位\ensuremath{5}进位后得\ensuremath{3.20},\allowbreak{}末尾\ensuremath{0}表示百分位，\allowbreak{}选C。\allowbreak{}\par};
\end{scope}\end{tikzpicture}\newpage
\definecolor{answer}{rgb}{0.850980,0.047059,0.070588}
\null\begin{tikzpicture}[remember picture,overlay]\begin{scope}[shift={(current page.north west)},x=1bp,y=-1bp]
\node[anchor=north west,inner sep=0pt,text=answer,align=left,text width=232.0bp] at (15,15) {\fontsize{7.5bp}{9.5bp}\selectfont 析：\allowbreak{}最大是\ensuremath{7.54},\allowbreak{}不是\ensuremath{7.59},\allowbreak{}判×。\allowbreak{}\par};
\end{scope}\end{tikzpicture}\newpage
\definecolor{answer}{rgb}{0.850980,0.047059,0.070588}
\null\begin{tikzpicture}[remember picture,overlay]\begin{scope}[shift={(current page.north west)},x=1bp,y=-1bp]
\node[anchor=north west,inner sep=0pt,text=answer,align=left,text width=232.0bp] at (15,15) {\fontsize{7.5bp}{9.5bp}\selectfont 析：\allowbreak{}千分位\ensuremath{8}进\ensuremath{1},\allowbreak{}得\ensuremath{4.40},\allowbreak{}判√。\allowbreak{}\par};
\end{scope}\end{tikzpicture}\newpage
\definecolor{answer}{rgb}{0.850980,0.047059,0.070588}
\null\begin{tikzpicture}[remember picture,overlay]\begin{scope}[shift={(current page.north west)},x=1bp,y=-1bp]
\node[anchor=north west,inner sep=0pt,text=answer,align=left,text width=232.0bp] at (15,15) {\fontsize{7.5bp}{9.5bp}\selectfont 析：\allowbreak{}两位小数可取\ensuremath{35.15}至\ensuremath{35.24},\allowbreak{}图C符合，\allowbreak{}选C。\allowbreak{}\par};
\end{scope}\end{tikzpicture}\newpage
\definecolor{answer}{rgb}{0.850980,0.047059,0.070588}
\null\begin{tikzpicture}[remember picture,overlay]\begin{scope}[shift={(current page.north west)},x=1bp,y=-1bp]
\node[anchor=north west,inner sep=0pt,text=answer,align=left,text width=232.0bp] at (15,15) {\fontsize{7.5bp}{9.5bp}\selectfont 析：\allowbreak{}示例均按÷\ensuremath{10000}改写；\allowbreak{}使用给定资料时保留原数值与来源即可。\allowbreak{}\par};
\end{scope}\end{tikzpicture}\newpage
\definecolor{answer}{rgb}{0.850980,0.047059,0.070588}
\null\begin{tikzpicture}[remember picture,overlay]\begin{scope}[shift={(current page.north west)},x=1bp,y=-1bp]
\node[anchor=north west,inner sep=0pt,text=answer,align=left,text width=502.0bp] at (15,15) {\fontsize{7.5bp}{9.5bp}\selectfont 析：\allowbreak{}计数单位\ensuremath{0.1}、\allowbreak{}\ensuremath{0.01}、\allowbreak{}\ensuremath{0.001};\allowbreak{}读小数时点后逐位读。\allowbreak{}\par};
\end{scope}\end{tikzpicture}\newpage
\definecolor{answer}{rgb}{0.850980,0.047059,0.070588}
\null\begin{tikzpicture}[remember picture,overlay]\begin{scope}[shift={(current page.north west)},x=1bp,y=-1bp]
\node[anchor=north west,inner sep=0pt,text=answer,align=left,text width=481bp] at (15,15) {\fontsize{7.5bp}{9.5bp}\selectfont 析：\allowbreak{}仅去添小数末尾\ensuremath{0}不改大小；\allowbreak{}先比整数，\allowbreak{}再逐位比小数。\allowbreak{}\par};
\end{scope}\end{tikzpicture}\newpage
\definecolor{answer}{rgb}{0.850980,0.047059,0.070588}
\null\begin{tikzpicture}[remember picture,overlay]\begin{scope}[shift={(current page.north west)},x=1bp,y=-1bp]
\node[anchor=north west,inner sep=0pt,text=answer,align=left,text width=502.0bp] at (15,15) {\fontsize{7.5bp}{9.5bp}\selectfont 析：\allowbreak{}低级换高级÷进率；\allowbreak{}高级换低级×进率。\allowbreak{}\par};
\end{scope}\end{tikzpicture}\newpage
\definecolor{answer}{rgb}{0.850980,0.047059,0.070588}
\null\begin{tikzpicture}[remember picture,overlay]\begin{scope}[shift={(current page.north west)},x=1bp,y=-1bp]
\node[anchor=north west,inner sep=0pt,text=answer,align=left,text width=481bp] at (15,15) {\fontsize{7.5bp}{9.5bp}\selectfont 析：\allowbreak{}\ensuremath{778000000}改成\ensuremath{7.78}亿，\allowbreak{}百分位\ensuremath{8}进位，\allowbreak{}约\ensuremath{7.8}亿。\allowbreak{}\par};
\end{scope}\end{tikzpicture}\newpage
\definecolor{answer}{rgb}{0.850980,0.047059,0.070588}
\null\begin{tikzpicture}[remember picture,overlay]\begin{scope}[shift={(current page.north west)},x=1bp,y=-1bp]
\node[anchor=north west,inner sep=0pt,text=answer,align=left,text width=502.0bp] at (15,15) {\fontsize{7.5bp}{9.5bp}\selectfont 析：\allowbreak{}左端对准\ensuremath{0},\allowbreak{}右端\ensuremath{4.6}刻度，\allowbreak{}钥匙长\ensuremath{4.6}厘米。\allowbreak{}\par};
\end{scope}\end{tikzpicture}\newpage
\definecolor{answer}{rgb}{0.850980,0.047059,0.070588}
\null\begin{tikzpicture}[remember picture,overlay]\begin{scope}[shift={(current page.north west)},x=1bp,y=-1bp]
\node[anchor=north west,inner sep=0pt,text=answer,align=left,text width=502.0bp] at (15,15) {\fontsize{7.5bp}{9.5bp}\selectfont 析：\allowbreak{}百位、\allowbreak{}百分位放\ensuremath{4},\allowbreak{}其余放\ensuremath{1},\allowbreak{}得\ensuremath{411.14}。\allowbreak{}\par};
\end{scope}\end{tikzpicture}\newpage
\definecolor{answer}{rgb}{0.850980,0.047059,0.070588}
\null\begin{tikzpicture}[remember picture,overlay]\begin{scope}[shift={(current page.north west)},x=1bp,y=-1bp]
\node[anchor=north west,inner sep=0pt,text=answer,align=left,text width=502.0bp] at (15,15) {\fontsize{7.5bp}{9.5bp}\selectfont 析：\allowbreak{}\ensuremath{7}在百分位，\allowbreak{}以\ensuremath{0.01}计数；\allowbreak{}\ensuremath{6}在千分位，\allowbreak{}以\ensuremath{0.001}计数。\allowbreak{}\par};
\end{scope}\end{tikzpicture}\newpage
\definecolor{answer}{rgb}{0.850980,0.047059,0.070588}
\null\begin{tikzpicture}[remember picture,overlay]\begin{scope}[shift={(current page.north west)},x=1bp,y=-1bp]
\node[anchor=north west,inner sep=0pt,text=answer,align=left,text width=502.0bp] at (15,15) {\fontsize{7.5bp}{9.5bp}\selectfont 析：\allowbreak{}\ensuremath{70} + \ensuremath{0.8} + \ensuremath{0.009} = \ensuremath{70.809},\allowbreak{}读七十点八零九。\allowbreak{}\par};
\end{scope}\end{tikzpicture}\newpage
\definecolor{answer}{rgb}{0.850980,0.047059,0.070588}
\null\begin{tikzpicture}[remember picture,overlay]\begin{scope}[shift={(current page.north west)},x=1bp,y=-1bp]
\node[anchor=north west,inner sep=0pt,text=answer,align=left,text width=502.0bp] at (15,15) {\fontsize{7.5bp}{9.5bp}\selectfont 析：\allowbreak{}\ensuremath{0.6}即\ensuremath{60}个百分之一，\allowbreak{}小于\ensuremath{64}个；\allowbreak{}\ensuremath{2.3} > \ensuremath{1.7};\allowbreak{}下排分别涂一半与十分之四，\allowbreak{}所以均填=。\allowbreak{}\par};
\end{scope}\end{tikzpicture}\newpage
\definecolor{answer}{rgb}{0.850980,0.047059,0.070588}
\null\begin{tikzpicture}[remember picture,overlay]\begin{scope}[shift={(current page.north west)},x=1bp,y=-1bp]
\node[anchor=north west,inner sep=0pt,text=answer,align=left,text width=502.0bp] at (15,15) {\fontsize{7.5bp}{9.5bp}\selectfont 析：\allowbreak{}÷\ensuremath{10000}得\ensuremath{186.3404}万，\allowbreak{}百分位\ensuremath{4}不进，\allowbreak{}约\ensuremath{186.3}万。\allowbreak{}\par};
\end{scope}\end{tikzpicture}\newpage
\definecolor{answer}{rgb}{0.850980,0.047059,0.070588}
\null\begin{tikzpicture}[remember picture,overlay]\begin{scope}[shift={(current page.north west)},x=1bp,y=-1bp]
\node[anchor=north west,inner sep=0pt,text=answer,align=left,text width=502.0bp] at (15,15) {\fontsize{7.5bp}{9.5bp}\selectfont 析：\allowbreak{}\ensuremath{9000300} ÷ \ensuremath{10000} = \ensuremath{900.03}万；\allowbreak{}\ensuremath{724700}为\ensuremath{72.47}万，\allowbreak{}约\ensuremath{72.5}万。\allowbreak{}\par};
\end{scope}\end{tikzpicture}\newpage
\definecolor{answer}{rgb}{0.850980,0.047059,0.070588}
\null\begin{tikzpicture}[remember picture,overlay]\begin{scope}[shift={(current page.north west)},x=1bp,y=-1bp]
\node[anchor=north west,inner sep=0pt,text=answer,align=left,text width=502.0bp] at (15,15) {\fontsize{7.5bp}{9.5bp}\selectfont 析：\allowbreak{}按进率\ensuremath{1000}、\allowbreak{}\ensuremath{100}换算，\allowbreak{}混合单位取整数部分，\allowbreak{}大单位小数部分再换小单位。\allowbreak{}\par};
\end{scope}\end{tikzpicture}\newpage
\definecolor{answer}{rgb}{0.850980,0.047059,0.070588}
\null\begin{tikzpicture}[remember picture,overlay]\begin{scope}[shift={(current page.north west)},x=1bp,y=-1bp]
\node[anchor=north west,inner sep=0pt,text=answer,align=left,text width=502.0bp] at (15,15) {\fontsize{7.5bp}{9.5bp}\selectfont 析：\allowbreak{}右移\ensuremath{2}位×\ensuremath{100};\allowbreak{}左移\ensuremath{1}位÷\ensuremath{10};\allowbreak{}\ensuremath{3.1} × \ensuremath{100} = \ensuremath{310},\allowbreak{} \ensuremath{54.2} ÷ \ensuremath{10} = \ensuremath{5.42}。\allowbreak{}\par};
\end{scope}\end{tikzpicture}\newpage
\definecolor{answer}{rgb}{0.850980,0.047059,0.070588}
\null\begin{tikzpicture}[remember picture,overlay]\begin{scope}[shift={(current page.north west)},x=1bp,y=-1bp]
\node[anchor=north west,inner sep=0pt,text=answer,align=left,text width=502.0bp] at (15,15) {\fontsize{7.5bp}{9.5bp}\selectfont 析：\allowbreak{}\ensuremath{9.95}至\ensuremath{9.99}进位、\allowbreak{}\ensuremath{10.00}至\ensuremath{10.04}舍去后均为\ensuremath{10.0};\allowbreak{}最小\ensuremath{9.95}、\allowbreak{}最大\ensuremath{10.04}。\allowbreak{}\par};
\end{scope}\end{tikzpicture}\newpage
\definecolor{answer}{rgb}{0.850980,0.047059,0.070588}
\null\begin{tikzpicture}[remember picture,overlay]\begin{scope}[shift={(current page.north west)},x=1bp,y=-1bp]
\node[anchor=north west,inner sep=0pt,text=answer,align=left,text width=502.0bp] at (15,15) {\fontsize{7.5bp}{9.5bp}\selectfont 析：\allowbreak{}不限位数，\allowbreak{}\ensuremath{0.5}与\ensuremath{0.7}间有无数个小数，\allowbreak{}判×。\allowbreak{}\par};
\end{scope}\end{tikzpicture}\newpage
\definecolor{answer}{rgb}{0.850980,0.047059,0.070588}
\null\begin{tikzpicture}[remember picture,overlay]\begin{scope}[shift={(current page.north west)},x=1bp,y=-1bp]
\node[anchor=north west,inner sep=0pt,text=answer,align=left,text width=502.0bp] at (15,15) {\fontsize{7.5bp}{9.5bp}\selectfont 析：\allowbreak{}如\ensuremath{0.001} < \ensuremath{0.9},\allowbreak{}位多也可能小，\allowbreak{}判×。\allowbreak{}\par};
\end{scope}\end{tikzpicture}\newpage
\definecolor{answer}{rgb}{0.850980,0.047059,0.070588}
\null\begin{tikzpicture}[remember picture,overlay]\begin{scope}[shift={(current page.north west)},x=1bp,y=-1bp]
\node[anchor=north west,inner sep=0pt,text=answer,align=left,text width=502.0bp] at (15,15) {\fontsize{7.5bp}{9.5bp}\selectfont 析：\allowbreak{}大小虽等，\allowbreak{}分别表示\ensuremath{18}个\ensuremath{0.01}与\ensuremath{180}个\ensuremath{0.001},\allowbreak{}意义不同，\allowbreak{}判×。\allowbreak{}\par};
\end{scope}\end{tikzpicture}\newpage
\definecolor{answer}{rgb}{0.850980,0.047059,0.070588}
\null\begin{tikzpicture}[remember picture,overlay]\begin{scope}[shift={(current page.north west)},x=1bp,y=-1bp]
\node[anchor=north west,inner sep=0pt,text=answer,align=left,text width=502.0bp] at (15,15) {\fontsize{7.5bp}{9.5bp}\selectfont 析：\allowbreak{}千分位\ensuremath{5}进\ensuremath{1},\allowbreak{}应为\ensuremath{1.90}万，\allowbreak{}判×。\allowbreak{}\par};
\end{scope}\end{tikzpicture}\newpage
\definecolor{answer}{rgb}{0.850980,0.047059,0.070588}
\null\begin{tikzpicture}[remember picture,overlay]\begin{scope}[shift={(current page.north west)},x=1bp,y=-1bp]
\node[anchor=north west,inner sep=0pt,text=answer,align=left,text width=502.0bp] at (15,15) {\fontsize{7.5bp}{9.5bp}\selectfont 析：\allowbreak{}去点右移\ensuremath{3}位，\allowbreak{}扩大\ensuremath{1000}倍，\allowbreak{}判√。\allowbreak{}\par};
\end{scope}\end{tikzpicture}\newpage
\definecolor{answer}{rgb}{0.850980,0.047059,0.070588}
\null\begin{tikzpicture}[remember picture,overlay]\begin{scope}[shift={(current page.north west)},x=1bp,y=-1bp]
\node[anchor=north west,inner sep=0pt,text=answer,align=left,text width=502.0bp] at (15,15) {\fontsize{7.5bp}{9.5bp}\selectfont 析：\allowbreak{}A涂\ensuremath{9}份中的\ensuremath{4}份，\allowbreak{}不是一半；\allowbreak{}B、\allowbreak{}C涂一半，\allowbreak{}选A。\allowbreak{}\par};
\end{scope}\end{tikzpicture}\newpage
\definecolor{answer}{rgb}{0.850980,0.047059,0.070588}
\null\begin{tikzpicture}[remember picture,overlay]\begin{scope}[shift={(current page.north west)},x=1bp,y=-1bp]
\node[anchor=north west,inner sep=0pt,text=answer,align=left,text width=502.0bp] at (15,15) {\fontsize{7.5bp}{9.5bp}\selectfont 析：\allowbreak{}合起来右移\ensuremath{1}位，\allowbreak{}扩大\ensuremath{10}倍，\allowbreak{}选B。\allowbreak{}\par};
\end{scope}\end{tikzpicture}\newpage
\definecolor{answer}{rgb}{0.850980,0.047059,0.070588}
\null\begin{tikzpicture}[remember picture,overlay]\begin{scope}[shift={(current page.north west)},x=1bp,y=-1bp]
\node[anchor=north west,inner sep=0pt,text=answer,align=left,text width=502.0bp] at (15,15) {\fontsize{7.5bp}{9.5bp}\selectfont 析：\allowbreak{}三位小数应在\ensuremath{7.165}至\ensuremath{7.174};\allowbreak{}\ensuremath{7.178}会约为\ensuremath{7.18},\allowbreak{}选C。\allowbreak{}\par};
\end{scope}\end{tikzpicture}\newpage
\definecolor{answer}{rgb}{0.850980,0.047059,0.070588}
\null\begin{tikzpicture}[remember picture,overlay]\begin{scope}[shift={(current page.north west)},x=1bp,y=-1bp]
\node[anchor=north west,inner sep=0pt,text=answer,align=left,text width=502.0bp] at (15,15) {\fontsize{7.5bp}{9.5bp}\selectfont 析：\allowbreak{}均约\ensuremath{9.8},\allowbreak{}乐乐只能\ensuremath{9.82},\allowbreak{}佳佳\ensuremath{9.76};\allowbreak{}丽丽第三为\ensuremath{9.75},\allowbreak{}所以乐乐第一，\allowbreak{}选A。\allowbreak{}\par};
\end{scope}\end{tikzpicture}\newpage
\definecolor{answer}{rgb}{0.850980,0.047059,0.070588}
\null\begin{tikzpicture}[remember picture,overlay]\begin{scope}[shift={(current page.north west)},x=1bp,y=-1bp]
\node[anchor=north west,inner sep=0pt,text=answer,align=left,text width=502.0bp] at (15,15) {\fontsize{7.5bp}{9.5bp}\selectfont 析：\allowbreak{}飞飞统一成分，\allowbreak{}乐乐比较涂色，\allowbreak{}天天统一成分米，\allowbreak{}\ensuremath{3}法都合理，\allowbreak{}选C。\allowbreak{}\par};
\end{scope}\end{tikzpicture}\newpage
\definecolor{answer}{rgb}{0.850980,0.047059,0.070588}
\null\begin{tikzpicture}[remember picture,overlay]\begin{scope}[shift={(current page.north west)},x=1bp,y=-1bp]
\node[anchor=north west,inner sep=0pt,text=answer,align=left,text width=502.0bp] at (15,15) {\fontsize{7.5bp}{9.5bp}\selectfont 析：\allowbreak{}两位小数按整数位\ensuremath{1}、\allowbreak{}\ensuremath{2}、\allowbreak{}\ensuremath{4}分类，\allowbreak{}各\ensuremath{2}种，\allowbreak{}共\ensuremath{6}种；\allowbreak{}最大数\ensuremath{42.1},\allowbreak{}最小\ensuremath{1.24}。\allowbreak{}\par};
\end{scope}\end{tikzpicture}\newpage
\definecolor{answer}{rgb}{0.850980,0.047059,0.070588}
\null\begin{tikzpicture}[remember picture,overlay]\begin{scope}[shift={(current page.north west)},x=1bp,y=-1bp]
\node[anchor=north west,inner sep=0pt,text=answer,align=left,text width=503.6666717529297bp] at (15,15) {\fontsize{8.0bp}{10.0bp}\selectfont 按四舍五入规则,\allowbreak{}\ensuremath{3.000}也可取近似数\ensuremath{3.00}。\allowbreak{}\par};
\end{scope}\end{tikzpicture}\newpage
\definecolor{answer}{rgb}{0.850980,0.047059,0.070588}
\null\begin{tikzpicture}[remember picture,overlay]\begin{scope}[shift={(current page.north west)},x=1bp,y=-1bp]
\node[anchor=north west,inner sep=0pt,text=answer,align=left,text width=502.0bp] at (15,15) {\fontsize{7.5bp}{9.5bp}\selectfont 析：\allowbreak{}可取\ensuremath{2.995}、\allowbreak{}\ensuremath{2.996}、\allowbreak{}\ensuremath{2.997}、\allowbreak{}\ensuremath{2.998}、\allowbreak{}\ensuremath{2.999}、\allowbreak{}\ensuremath{3.000}、\allowbreak{}\ensuremath{3.001}、\allowbreak{}\ensuremath{3.002}、\allowbreak{}\ensuremath{3.003}、\allowbreak{}\ensuremath{3.004};\allowbreak{}最大\ensuremath{3.004},\allowbreak{}最小\ensuremath{2.995}。\allowbreak{}按规则\ensuremath{3.000}也可取近似数\ensuremath{3.00}。\allowbreak{}\par};
\end{scope}\end{tikzpicture}\newpage
\definecolor{answer}{rgb}{0.850980,0.047059,0.070588}
\null\begin{tikzpicture}[remember picture,overlay]\begin{scope}[shift={(current page.north west)},x=1bp,y=-1bp]
\node[anchor=north west,inner sep=0pt,text=answer,align=left,text width=502.0bp] at (15,15) {\fontsize{7.5bp}{9.5bp}\selectfont 析：\allowbreak{}A先求\ensuremath{1000} ÷ \ensuremath{100} = \ensuremath{10}；\allowbreak{}B先求\ensuremath{3.25} ÷ \ensuremath{100} = \ensuremath{0.0325}千克，\allowbreak{}两法均能求解；\allowbreak{}本题采用A。\allowbreak{}\par};
\end{scope}\end{tikzpicture}\newpage
\definecolor{answer}{rgb}{0.850980,0.047059,0.070588}
\null\begin{tikzpicture}[remember picture,overlay]\begin{scope}[shift={(current page.north west)},x=1bp,y=-1bp]
\node[anchor=north west,inner sep=0pt,text=answer,align=left,text width=502.0bp] at (15,15) {\fontsize{7.5bp}{9.5bp}\selectfont 析：\allowbreak{}上方个位纵式\ensuremath{8};\allowbreak{}下方十分位横式\ensuremath{6}、\allowbreak{}百分位纵式\ensuremath{3},\allowbreak{}共\ensuremath{9}根，\allowbreak{}表示\ensuremath{8.63}。\allowbreak{}\par};
\end{scope}\end{tikzpicture}\newpage
\definecolor{answer}{rgb}{0.850980,0.047059,0.070588}
\null\begin{tikzpicture}[remember picture,overlay]\begin{scope}[shift={(current page.north west)},x=1bp,y=-1bp]
\node[anchor=north west,inner sep=0pt,text=answer,align=left,text width=457.0bp] at (15,15) {\fontsize{7.5bp}{9.5bp}\selectfont 析：\allowbreak{}范围\ensuremath{1.25}至\ensuremath{1.34};\allowbreak{}恰\ensuremath{6}根只有\ensuremath{1.32}(\ensuremath{1} + \ensuremath{3} + \ensuremath{2}根）及\ensuremath{1.27}(\ensuremath{1} + \ensuremath{2} + \ensuremath{3}根），\allowbreak{}均约\ensuremath{1.3}。\allowbreak{}\par};
\end{scope}\end{tikzpicture}\newpage
\definecolor{answer}{rgb}{0.850980,0.047059,0.070588}
\null\begin{tikzpicture}[remember picture,overlay]\begin{scope}[shift={(current page.north west)},x=1bp,y=-1bp]
\node[anchor=north west,inner sep=0pt,text=answer,align=left,text width=232.0bp] at (15,15) {\fontsize{7.5bp}{9.5bp}\selectfont 析：\allowbreak{}最大\ensuremath{6.42}、\allowbreak{}最小\ensuremath{2.46},\allowbreak{}和\ensuremath{8.88}、\allowbreak{}差\ensuremath{3.96}。\allowbreak{}\par};
\end{scope}\end{tikzpicture}\newpage
\definecolor{answer}{rgb}{0.850980,0.047059,0.070588}
\null\begin{tikzpicture}[remember picture,overlay]\begin{scope}[shift={(current page.north west)},x=1bp,y=-1bp]
\node[anchor=north west,inner sep=0pt,text=answer,align=left,text width=232.0bp] at (15,15) {\fontsize{7.5bp}{9.5bp}\selectfont 析：\allowbreak{}减去较多仍相等，\allowbreak{}原来的甲数大，\allowbreak{}比乙数大\ensuremath{1.2}。\allowbreak{}\par};
\end{scope}\end{tikzpicture}\newpage
\definecolor{answer}{rgb}{0.850980,0.047059,0.070588}
\null\begin{tikzpicture}[remember picture,overlay]\begin{scope}[shift={(current page.north west)},x=1bp,y=-1bp]
\node[anchor=north west,inner sep=0pt,text=answer,align=left,text width=232.0bp] at (15,15) {\fontsize{7.5bp}{9.5bp}\selectfont 析：\allowbreak{}B的\ensuremath{2}和\ensuremath{5}都在十分位可相减；\allowbreak{}A、\allowbreak{}C数位不同，\allowbreak{}选B。\allowbreak{}\par};
\end{scope}\end{tikzpicture}\newpage
\definecolor{answer}{rgb}{0.850980,0.047059,0.070588}
\null\begin{tikzpicture}[remember picture,overlay]\begin{scope}[shift={(current page.north west)},x=1bp,y=-1bp]
\node[anchor=north west,inner sep=0pt,text=answer,align=left,text width=232.0bp] at (15,15) {\fontsize{7.5bp}{9.5bp}\selectfont 析：\allowbreak{}最小\ensuremath{3.55},\allowbreak{}最大\ensuremath{3.64},\allowbreak{}最大差\ensuremath{3.64} \ensuremath{-}\ensuremath{3.55} = \ensuremath{0.09}。\allowbreak{}\par};
\end{scope}\end{tikzpicture}\newpage
\definecolor{answer}{rgb}{0.850980,0.047059,0.070588}
\null\begin{tikzpicture}[remember picture,overlay]\begin{scope}[shift={(current page.north west)},x=1bp,y=-1bp]
\node[anchor=north west,inner sep=0pt,text=answer,align=left,text width=232.0bp] at (15,15) {\fontsize{7.5bp}{9.5bp}\selectfont 析：\allowbreak{}A最小\ensuremath{1.06},\allowbreak{} B最小\ensuremath{1.11};\allowbreak{}C为\ensuremath{0.50}至\ensuremath{0.59},\allowbreak{}都小于\ensuremath{1},\allowbreak{}选C。\allowbreak{}\par};
\end{scope}\end{tikzpicture}\newpage
\definecolor{answer}{rgb}{0.850980,0.047059,0.070588}
\null\begin{tikzpicture}[remember picture,overlay]\begin{scope}[shift={(current page.north west)},x=1bp,y=-1bp]
\node[anchor=north west,inner sep=0pt,text=answer,align=left,text width=232.0bp] at (15,15) {\fontsize{7.5bp}{9.5bp}\selectfont 析：\allowbreak{}每次加\ensuremath{0.2},\allowbreak{}填\ensuremath{1.0}、\allowbreak{}\ensuremath{1.2}。\allowbreak{}\par};
\end{scope}\end{tikzpicture}\newpage
\definecolor{answer}{rgb}{0.850980,0.047059,0.070588}
\null\begin{tikzpicture}[remember picture,overlay]\begin{scope}[shift={(current page.north west)},x=1bp,y=-1bp]
\node[anchor=north west,inner sep=0pt,text=answer,align=left,text width=232.0bp] at (15,15) {\fontsize{7.5bp}{9.5bp}\selectfont 析：\allowbreak{}每次减\ensuremath{0.2},\allowbreak{}填\ensuremath{3.06}、\allowbreak{}\ensuremath{2.86}。\allowbreak{}\par};
\end{scope}\end{tikzpicture}\newpage
\definecolor{answer}{rgb}{0.850980,0.047059,0.070588}
\null\begin{tikzpicture}[remember picture,overlay]\begin{scope}[shift={(current page.north west)},x=1bp,y=-1bp]
\node[anchor=north west,inner sep=0pt,text=answer,align=left,text width=232.0bp] at (15,15) {\fontsize{7.5bp}{9.5bp}\selectfont 析：\allowbreak{}每次加\ensuremath{0.003},\allowbreak{}填\ensuremath{4.365}、\allowbreak{}\ensuremath{4.368}。\allowbreak{}\par};
\end{scope}\end{tikzpicture}\newpage
\definecolor{answer}{rgb}{0.850980,0.047059,0.070588}
\null\begin{tikzpicture}[remember picture,overlay]\begin{scope}[shift={(current page.north west)},x=1bp,y=-1bp]
\node[anchor=north west,inner sep=0pt,text=answer,align=left,text width=232.0bp] at (15,15) {\fontsize{8.5bp}{10.62bp}\selectfont 答：\allowbreak{}两座山的海拔相差\ensuremath{622.2}米。\allowbreak{}\par};
\end{scope}\end{tikzpicture}\newpage
\definecolor{answer}{rgb}{0.850980,0.047059,0.070588}
\null\begin{tikzpicture}[remember picture,overlay]\begin{scope}[shift={(current page.north west)},x=1bp,y=-1bp]
\node[anchor=north west,inner sep=0pt,text=answer,align=left,text width=232.0bp] at (15,15) {\fontsize{7.5bp}{9.5bp}\selectfont 析：\allowbreak{}每次减\ensuremath{0.005},\allowbreak{}填\ensuremath{8.376}、\allowbreak{}\ensuremath{8.371}。\allowbreak{}\par};
\end{scope}\end{tikzpicture}\newpage
\definecolor{answer}{rgb}{0.850980,0.047059,0.070588}
\null\begin{tikzpicture}[remember picture,overlay]\begin{scope}[shift={(current page.north west)},x=1bp,y=-1bp]
\node[anchor=north west,inner sep=0pt,text=answer,align=left,text width=232.0bp] at (15,15) {\fontsize{8.5bp}{10.62bp}\selectfont 答：\allowbreak{}黄海面积约为\ensuremath{38}万平方千米。\allowbreak{}\par};
\end{scope}\end{tikzpicture}\newpage
\definecolor{answer}{rgb}{0.850980,0.047059,0.070588}
\null\begin{tikzpicture}[remember picture,overlay]\begin{scope}[shift={(current page.north west)},x=1bp,y=-1bp]
\node[anchor=north west,inner sep=0pt,text=answer,align=left,text width=232.0bp] at (15,15) {\fontsize{7.5bp}{9.5bp}\selectfont 按运算顺序计算，\allowbreak{}也可利用加减法性质简算。\allowbreak{}\par};
\end{scope}\end{tikzpicture}\newpage
\definecolor{answer}{rgb}{0.850980,0.047059,0.070588}
\null\begin{tikzpicture}[remember picture,overlay]\begin{scope}[shift={(current page.north west)},x=1bp,y=-1bp]
\node[anchor=north west,inner sep=0pt,text=answer,align=left,text width=232.0bp] at (15,15) {\fontsize{7.5bp}{9.5bp}\selectfont 按运算顺序计算，\allowbreak{}也可利用加减法性质简算。\allowbreak{}\par};
\end{scope}\end{tikzpicture}\newpage
\definecolor{answer}{rgb}{0.850980,0.047059,0.070588}
\null\begin{tikzpicture}[remember picture,overlay]\begin{scope}[shift={(current page.north west)},x=1bp,y=-1bp]
\node[anchor=north west,inner sep=0pt,text=answer,align=left,text width=233.0bp] at (15,15) {\fontsize{7.5bp}{9.5bp}\selectfont 析：\allowbreak{}无括号先减后加，\allowbreak{}有括号先算加法；\allowbreak{}结果分别\ensuremath{10.9}、\allowbreak{}\ensuremath{5.9},\allowbreak{}运算顺序和结果均不同。\allowbreak{}\par};
\end{scope}\end{tikzpicture}\newpage
\definecolor{answer}{rgb}{0.850980,0.047059,0.070588}
\null\begin{tikzpicture}[remember picture,overlay]\begin{scope}[shift={(current page.north west)},x=1bp,y=-1bp]
\node[anchor=north west,inner sep=0pt,text=answer,align=left,text width=233.0bp] at (15,15) {\fontsize{7.5bp}{9.5bp}\selectfont 析：\allowbreak{}先减：\allowbreak{}\ensuremath{7.25} \ensuremath{-}\ensuremath{5.56} = \ensuremath{1.69};\allowbreak{}再加：\allowbreak{}\ensuremath{3.64} + \ensuremath{1.69} = \ensuremath{5.33}。\allowbreak{}\par};
\end{scope}\end{tikzpicture}\newpage
\definecolor{answer}{rgb}{0.850980,0.047059,0.070588}
\null\begin{tikzpicture}[remember picture,overlay]\begin{scope}[shift={(current page.north west)},x=1bp,y=-1bp]
\node[anchor=north west,inner sep=0pt,text=answer,align=left,text width=232.0bp] at (15,15) {\fontsize{7.5bp}{9.5bp}\selectfont 按运算顺序计算，\allowbreak{}也可利用加减法性质简算。\allowbreak{}\par};
\end{scope}\end{tikzpicture}\newpage
\definecolor{answer}{rgb}{0.850980,0.047059,0.070588}
\null\begin{tikzpicture}[remember picture,overlay]\begin{scope}[shift={(current page.north west)},x=1bp,y=-1bp]
\node[anchor=north west,inner sep=0pt,text=answer,align=left,text width=232.0bp] at (15,15) {\fontsize{8.5bp}{10.62bp}\selectfont 两份蛋糕原价合计\ensuremath{31.6}元，\allowbreak{}达到满\ensuremath{30}减\ensuremath{5}条件，\allowbreak{}再加配送费。\allowbreak{}\par};
\end{scope}\end{tikzpicture}\newpage
\definecolor{answer}{rgb}{0.850980,0.047059,0.070588}
\null\begin{tikzpicture}[remember picture,overlay]\begin{scope}[shift={(current page.north west)},x=1bp,y=-1bp]
\node[anchor=north west,inner sep=0pt,text=answer,align=left,text width=232.0bp] at (15,15) {\fontsize{8.5bp}{10.62bp}\selectfont 把三车总质量与\ensuremath{20}吨比较，\allowbreak{}或者先求前两辆上桥后的剩余承重。\allowbreak{}\par};
\end{scope}\end{tikzpicture}\newpage
\definecolor{answer}{rgb}{0.850980,0.047059,0.070588}
\null\begin{tikzpicture}[remember picture,overlay]\begin{scope}[shift={(current page.north west)},x=1bp,y=-1bp]
\node[anchor=north west,inner sep=0pt,text=answer,align=left,text width=232.0bp] at (15,15) {\fontsize{8.5bp}{10.62bp}\selectfont 答：\allowbreak{}第三辆超过剩余承重\ensuremath{0.4}吨，\allowbreak{}三辆同时上桥会超重。\allowbreak{}\par};
\end{scope}\end{tikzpicture}\newpage
\definecolor{answer}{rgb}{0.850980,0.047059,0.070588}
\null\begin{tikzpicture}[remember picture,overlay]\begin{scope}[shift={(current page.north west)},x=1bp,y=-1bp]
\node[anchor=north west,inner sep=0pt,text=answer,align=left,text width=242.0bp] at (15,15) {\fontsize{7.5bp}{9.5bp}\selectfont 用去的一半黄豆为\ensuremath{9.64} \ensuremath{-}\ensuremath{5.33}，\allowbreak{}另一半同重；\allowbreak{}剩余总质量扣除另一半黄豆即篮重。\allowbreak{}\par};
\end{scope}\end{tikzpicture}\newpage
\definecolor{answer}{rgb}{0.850980,0.047059,0.070588}
\null\begin{tikzpicture}[remember picture,overlay]\begin{scope}[shift={(current page.north west)},x=1bp,y=-1bp]
\node[anchor=north west,inner sep=0pt,text=answer,align=left,text width=233.0bp] at (15,15) {\fontsize{8.5bp}{10.62bp}\selectfont \ensuremath{10.04}\ensuremath{-}\ensuremath{1.77}求磊磊用时；\allowbreak{}再用乐乐用时减磊磊用时，\allowbreak{}问题是“乐乐比磊磊多用多少秒”。\allowbreak{}\par};
\end{scope}\end{tikzpicture}\newpage
\definecolor{answer}{rgb}{0.850980,0.047059,0.070588}
\null\begin{tikzpicture}[remember picture,overlay]\begin{scope}[shift={(current page.north west)},x=1bp,y=-1bp]
\node[anchor=north west,inner sep=0pt,text=answer,align=left,text width=233.00001525878906bp] at (15,15) {\fontsize{8.5bp}{10.62bp}\selectfont 开放提问要至少经过两步。\allowbreak{}先求磊磊\ensuremath{8.27}秒，\allowbreak{}再求不同人的合计时间。\allowbreak{}\par};
\end{scope}\end{tikzpicture}\newpage
\definecolor{answer}{rgb}{0.850980,0.047059,0.070588}
\null\begin{tikzpicture}[remember picture,overlay]\begin{scope}[shift={(current page.north west)},x=1bp,y=-1bp]
\node[anchor=north west,inner sep=0pt,text=answer,align=left,text width=242.0bp] at (15,15) {\fontsize{8.5bp}{10.62bp}\selectfont 升温用加法、\allowbreak{}降温用减法；\allowbreak{}也可先求全天净降温。\allowbreak{}\par};
\end{scope}\end{tikzpicture}\newpage
\definecolor{answer}{rgb}{0.850980,0.047059,0.070588}
\null\begin{tikzpicture}[remember picture,overlay]\begin{scope}[shift={(current page.north west)},x=1bp,y=-1bp]
\node[anchor=north west,inner sep=0pt,text=answer,align=left,text width=242.0bp] at (15,15) {\fontsize{8.5bp}{10.62bp}\selectfont 答：\allowbreak{}这天晚上的气温是\ensuremath{11.9}℃。\allowbreak{}\par};
\end{scope}\end{tikzpicture}\newpage
\definecolor{answer}{rgb}{0.850980,0.047059,0.070588}
\null\begin{tikzpicture}[remember picture,overlay]\begin{scope}[shift={(current page.north west)},x=1bp,y=-1bp]
\node[anchor=north west,inner sep=0pt,text=answer,align=left,text width=233.0bp] at (15,15) {\fontsize{8.5bp}{10.62bp}\selectfont 答：\allowbreak{}乐乐和磊磊共用\ensuremath{16.67}秒。\allowbreak{}\par};
\end{scope}\end{tikzpicture}\newpage
\definecolor{answer}{rgb}{0.850980,0.047059,0.070588}
\null\begin{tikzpicture}[remember picture,overlay]\begin{scope}[shift={(current page.north west)},x=1bp,y=-1bp]
\node[anchor=north west,inner sep=0pt,text=answer,align=left,text width=232bp] at (15,15) {\fontsize{8.5bp}{10.62bp}\selectfont 答：\allowbreak{}倩倩和磊磊共用\ensuremath{18.31}秒。\allowbreak{}\par};
\end{scope}\end{tikzpicture}\newpage
\definecolor{answer}{rgb}{0.850980,0.047059,0.070588}
\null\begin{tikzpicture}[remember picture,overlay]\begin{scope}[shift={(current page.north west)},x=1bp,y=-1bp]
\node[anchor=north west,inner sep=0pt,text=answer,align=left,text width=232bp] at (15,15) {\fontsize{8.5bp}{10.62bp}\selectfont 答：\allowbreak{}乐乐、\allowbreak{}倩倩和磊磊共用\ensuremath{26.71}秒。\allowbreak{}\par};
\end{scope}\end{tikzpicture}\newpage
\definecolor{answer}{rgb}{0.850980,0.047059,0.070588}
\null\begin{tikzpicture}[remember picture,overlay]\begin{scope}[shift={(current page.north west)},x=1bp,y=-1bp]
\node[anchor=north west,inner sep=0pt,text=answer,align=left,text width=242.0bp] at (15,15) {\fontsize{7.5bp}{9.5bp}\selectfont 析：\allowbreak{}新疆有葡萄、\allowbreak{}哈密瓜、\allowbreak{}香梨。\allowbreak{}白天光照足利于积糖，\allowbreak{}夜凉使呼吸消耗少。\allowbreak{}\par};
\end{scope}\end{tikzpicture}\newpage
\definecolor{answer}{rgb}{0.850980,0.047059,0.070588}
\null\begin{tikzpicture}[remember picture,overlay]\begin{scope}[shift={(current page.north west)},x=1bp,y=-1bp]
\node[anchor=north west,inner sep=0pt,text=answer,align=left,text width=233.0bp] at (15,15) {\fontsize{8.5bp}{10.62bp}\selectfont 第一、\allowbreak{}二段与第二、\allowbreak{}三段相加，\allowbreak{}第二段被算了两次；\allowbreak{}减去总长，\allowbreak{}便得到第二段长。\allowbreak{}\par};
\end{scope}\end{tikzpicture}\newpage
\definecolor{answer}{rgb}{0.850980,0.047059,0.070588}
\null\begin{tikzpicture}[remember picture,overlay]\begin{scope}[shift={(current page.north west)},x=1bp,y=-1bp]
\node[anchor=north west,inner sep=0pt,text=answer,align=left,text width=233.0bp] at (15,15) {\fontsize{7.5bp}{9.5bp}\selectfont 先把\ensuremath{5.08}与\ensuremath{1.92}相加凑成\ensuremath{7}，\allowbreak{}再减\ensuremath{3.5}。\allowbreak{}\par};
\end{scope}\end{tikzpicture}\newpage
\definecolor{answer}{rgb}{0.850980,0.047059,0.070588}
\null\begin{tikzpicture}[remember picture,overlay]\begin{scope}[shift={(current page.north west)},x=1bp,y=-1bp]
\node[anchor=north west,inner sep=0pt,text=answer,align=left,text width=232.0bp] at (15,15) {\fontsize{7.5bp}{9.5bp}\selectfont 析：\allowbreak{}A合并减数凑\ensuremath{40},\allowbreak{} B先算\ensuremath{18.69} \ensuremath{-}\ensuremath{8.69}凑\ensuremath{10};\allowbreak{}C不能这样凑整，\allowbreak{}选C。\allowbreak{}\par};
\end{scope}\end{tikzpicture}\newpage
\definecolor{answer}{rgb}{0.850980,0.047059,0.070588}
\null\begin{tikzpicture}[remember picture,overlay]\begin{scope}[shift={(current page.north west)},x=1bp,y=-1bp]
\node[anchor=north west,inner sep=0pt,text=answer,align=left,text width=233.0bp] at (15,15) {\fontsize{7.5bp}{9.5bp}\selectfont 连续减去两个数，\allowbreak{}可以减去它们的和；\allowbreak{}\ensuremath{2.8}与\ensuremath{1.2}的和是\ensuremath{4}。\allowbreak{}\par};
\end{scope}\end{tikzpicture}\newpage
\definecolor{answer}{rgb}{0.850980,0.047059,0.070588}
\null\begin{tikzpicture}[remember picture,overlay]\begin{scope}[shift={(current page.north west)},x=1bp,y=-1bp]
\node[anchor=north west,inner sep=0pt,text=answer,align=left,text width=232.0bp] at (15,15) {\fontsize{7.5bp}{9.5bp}\selectfont 析：\allowbreak{}\ensuremath{4.99}比\ensuremath{5}少\ensuremath{0.01},\allowbreak{}减\ensuremath{4.99}可先减\ensuremath{5}再加\ensuremath{0.01},\allowbreak{}选C。\allowbreak{}\par};
\end{scope}\end{tikzpicture}\newpage
\definecolor{answer}{rgb}{0.850980,0.047059,0.070588}
\null\begin{tikzpicture}[remember picture,overlay]\begin{scope}[shift={(current page.north west)},x=1bp,y=-1bp]
\node[anchor=north west,inner sep=0pt,text=answer,align=left,text width=233.0bp] at (15,15) {\fontsize{7.5bp}{9.5bp}\selectfont 交换、\allowbreak{}结合加数，\allowbreak{}先把\ensuremath{14.3}与\ensuremath{3.7}、\allowbreak{}\ensuremath{2.88}与\ensuremath{0.12}分别相加凑整。\allowbreak{}\par};
\end{scope}\end{tikzpicture}\newpage
\definecolor{answer}{rgb}{0.850980,0.047059,0.070588}
\null\begin{tikzpicture}[remember picture,overlay]\begin{scope}[shift={(current page.north west)},x=1bp,y=-1bp]
\node[anchor=north west,inner sep=0pt,text=answer,align=left,text width=232.0bp] at (15,15) {\fontsize{7.5bp}{9.5bp}\selectfont 红包总额减去已领的两份金额；\allowbreak{}可先求两份之和，\allowbreak{}也可依次扣除。\allowbreak{}\par};
\end{scope}\end{tikzpicture}\newpage
\definecolor{answer}{rgb}{0.850980,0.047059,0.070588}
\null\begin{tikzpicture}[remember picture,overlay]\begin{scope}[shift={(current page.north west)},x=1bp,y=-1bp]
\node[anchor=north west,inner sep=0pt,text=answer,align=left,text width=232bp] at (15,15) {\fontsize{7.5bp}{9.5bp}\selectfont 把\ensuremath{9.99}看成\ensuremath{10}，\allowbreak{}多加了\ensuremath{0.01}，\allowbreak{}最后减回\ensuremath{0.01}。\allowbreak{}\par};
\end{scope}\end{tikzpicture}\newpage
\definecolor{answer}{rgb}{0.850980,0.047059,0.070588}
\null\begin{tikzpicture}[remember picture,overlay]\begin{scope}[shift={(current page.north west)},x=1bp,y=-1bp]
\node[anchor=north west,inner sep=0pt,text=answer,align=left,text width=233.0bp] at (15,15) {\fontsize{7.5bp}{9.5bp}\selectfont 析：\allowbreak{}\ensuremath{5.4}移到前面是交换律；\allowbreak{}把\ensuremath{1.25}与\ensuremath{0.75}先相加是结合律。\allowbreak{}\par};
\end{scope}\end{tikzpicture}\newpage
\definecolor{answer}{rgb}{0.850980,0.047059,0.070588}
\null\begin{tikzpicture}[remember picture,overlay]\begin{scope}[shift={(current page.north west)},x=1bp,y=-1bp]
\node[anchor=north west,inner sep=0pt,text=answer,align=left,text width=233.0bp] at (15,15) {\fontsize{7.5bp}{9.5bp}\selectfont 析：\allowbreak{}错算相当于又减\ensuremath{1.4},\allowbreak{}正确应加\ensuremath{1.4};\allowbreak{}相差\ensuremath{1.4} × \ensuremath{2} = \ensuremath{2.8}。\allowbreak{}\par};
\end{scope}\end{tikzpicture}\newpage
\definecolor{answer}{rgb}{0.850980,0.047059,0.070588}
\null\begin{tikzpicture}[remember picture,overlay]\begin{scope}[shift={(current page.north west)},x=1bp,y=-1bp]
\node[anchor=north west,inner sep=0pt,text=answer,align=left,text width=233.00000762939453bp] at (15,15) {\fontsize{7.5bp}{9.5bp}\selectfont 析：\allowbreak{}依次<、\allowbreak{}=、\allowbreak{}=。\allowbreak{}第一处右式比左式多\ensuremath{7};\allowbreak{}连续减两数等于减其和；\allowbreak{}末式重排加减项。\allowbreak{}\par};
\end{scope}\end{tikzpicture}\newpage
\definecolor{answer}{rgb}{0.850980,0.047059,0.070588}
\null\begin{tikzpicture}[remember picture,overlay]\begin{scope}[shift={(current page.north west)},x=1bp,y=-1bp]
\node[anchor=north west,inner sep=0pt,text=answer,align=left,text width=247.0bp] at (15,15) {\fontsize{7.5bp}{9.5bp}\selectfont 析：\allowbreak{}√。\allowbreak{}沿线段所在直线对折，\allowbreak{}或沿经过线段中点且与它垂直的直线对折，\allowbreak{}都能重合。\allowbreak{}\par};
\end{scope}\end{tikzpicture}\newpage
\definecolor{answer}{rgb}{0.850980,0.047059,0.070588}
\null\begin{tikzpicture}[remember picture,overlay]\begin{scope}[shift={(current page.north west)},x=1bp,y=-1bp]
\node[anchor=north west,inner sep=0pt,text=answer,align=left,text width=247.125bp] at (15,15) {\fontsize{7.51bp}{9.51bp}\selectfont 析：\allowbreak{}×。\allowbreak{}对角线是线段，\allowbreak{}对称轴是直线。\allowbreak{}正方形每条对角线所在的直线都是它的对称轴\par};
\end{scope}\end{tikzpicture}\newpage
\definecolor{answer}{rgb}{0.850980,0.047059,0.070588}
\null\begin{tikzpicture}[remember picture,overlay]\begin{scope}[shift={(current page.north west)},x=1bp,y=-1bp]
\node[anchor=north west,inner sep=0pt,text=answer,align=left,text width=233.0bp] at (15,15) {\fontsize{7.5bp}{9.5bp}\selectfont 析：\allowbreak{}按行依次√××、\allowbreak{}√×√；\allowbreak{}判断能否沿直线对折重合，\allowbreak{}不能只看旋转后是否相同。\allowbreak{}\par};
\end{scope}\end{tikzpicture}\newpage
\definecolor{answer}{rgb}{0.850980,0.047059,0.070588}
\null\begin{tikzpicture}[remember picture,overlay]\begin{scope}[shift={(current page.north west)},x=1bp,y=-1bp]
\node[anchor=north west,inner sep=0pt,text=answer,align=left,text width=233.0bp] at (15,15) {\fontsize{7.5bp}{9.5bp}\selectfont 析：\allowbreak{}两图都只有竖直中线能使两半重合，\allowbreak{}各\ensuremath{1}条对称轴。\allowbreak{}\par};
\end{scope}\end{tikzpicture}\newpage
\definecolor{answer}{rgb}{0.850980,0.047059,0.070588}
\null\begin{tikzpicture}[remember picture,overlay]\begin{scope}[shift={(current page.north west)},x=1bp,y=-1bp]
\node[anchor=north west,inner sep=0pt,text=answer,align=left,text width=247.0bp] at (15,15) {\fontsize{7.5bp}{9.5bp}\selectfont 析：\allowbreak{}①③只能旋转对应，\allowbreak{}不能沿直线对折重合，\allowbreak{}故不是轴对称图形。\allowbreak{}\par};
\end{scope}\end{tikzpicture}\newpage
\definecolor{answer}{rgb}{0.850980,0.047059,0.070588}
\null\begin{tikzpicture}[remember picture,overlay]\begin{scope}[shift={(current page.north west)},x=1bp,y=-1bp]
\node[anchor=north west,inner sep=0pt,text=answer,align=left,text width=233.0bp] at (15,15) {\fontsize{7.5bp}{9.5bp}\selectfont 析：\allowbreak{}×。\allowbreak{}轴对称图形可有\ensuremath{1}条、\allowbreak{}多条或无数条对称轴，\allowbreak{}如圆。\allowbreak{}\par};
\end{scope}\end{tikzpicture}\newpage
\definecolor{answer}{rgb}{0.850980,0.047059,0.070588}
\null\begin{tikzpicture}[remember picture,overlay]\begin{scope}[shift={(current page.north west)},x=1bp,y=-1bp]
\node[anchor=north west,inner sep=0pt,text=answer,align=left,text width=233.0bp] at (15,15) {\fontsize{7.5bp}{9.5bp}\selectfont 析：\allowbreak{}等腰三角形\ensuremath{1}条；\allowbreak{}正方形\ensuremath{4}条；\allowbreak{}正五边形\ensuremath{5}条；\allowbreak{}正六边形\ensuremath{6}条。\allowbreak{}\par};
\end{scope}\end{tikzpicture}\newpage
\definecolor{answer}{rgb}{0.850980,0.047059,0.070588}
\null\begin{tikzpicture}[remember picture,overlay]\begin{scope}[shift={(current page.north west)},x=1bp,y=-1bp]
\node[anchor=north west,inner sep=0pt,text=answer,align=left,text width=232.7bp] at (15,15) {\fontsize{8.57bp}{10.71bp}\selectfont 例如(实线为木棒,\allowbreak{} 虚线为对称轴) :\allowbreak{}\par};
\end{scope}\end{tikzpicture}\newpage
\definecolor{answer}{rgb}{0.850980,0.047059,0.070588}
\null\begin{tikzpicture}[remember picture,overlay]\begin{scope}[shift={(current page.north west)},x=1bp,y=-1bp]
\node[anchor=north west,inner sep=0pt,text=answer,align=left,text width=242.0bp] at (15,15) {\fontsize{7.5bp}{9.5bp}\selectfont 长棒长度相当于两根短棒，\allowbreak{}要用完\ensuremath{3}根长棒和\ensuremath{4}根短棒，\allowbreak{}并使整个图形对折重合。\allowbreak{}\par};
\end{scope}\end{tikzpicture}\newpage
\definecolor{answer}{rgb}{0.850980,0.047059,0.070588}
\null\begin{tikzpicture}[remember picture,overlay]\begin{scope}[shift={(current page.north west)},x=1bp,y=-1bp]
\node[anchor=north west,inner sep=0pt,text=answer,align=left,text width=232.0bp] at (15,15) {\fontsize{7.5bp}{9.5bp}\selectfont 析：\allowbreak{}对应屋顶顶点右移\ensuremath{6}格、\allowbreak{}下移\ensuremath{5}格；\allowbreak{}也可先下后右。\allowbreak{}\par};
\end{scope}\end{tikzpicture}\newpage
\definecolor{answer}{rgb}{0.850980,0.047059,0.070588}
\null\begin{tikzpicture}[remember picture,overlay]\begin{scope}[shift={(current page.north west)},x=1bp,y=-1bp]
\node[anchor=north west,inner sep=0pt,text=answer,align=left,text width=233.0bp] at (15,15) {\fontsize{7.5bp}{9.5bp}\selectfont 析：\allowbreak{}找同一对应顶点数格距：\allowbreak{}直角形下\ensuremath{5}格；\allowbreak{}三角形右\ensuremath{5}格。\allowbreak{}\par};
\end{scope}\end{tikzpicture}\newpage
\definecolor{answer}{rgb}{0.850980,0.047059,0.070588}
\null\begin{tikzpicture}[remember picture,overlay]\begin{scope}[shift={(current page.north west)},x=1bp,y=-1bp]
\node[anchor=north west,inner sep=0pt,text=answer,align=left,text width=233.00000762939453bp] at (15,15) {\fontsize{7.5bp}{9.5bp}\selectfont 析：\allowbreak{}依箭头数对应顶点：\allowbreak{}小船下\ensuremath{5}再右\ensuremath{5};\allowbreak{}小灯左\ensuremath{5}再上\ensuremath{6},\allowbreak{}形状朝向不变。\allowbreak{}\par};
\end{scope}\end{tikzpicture}\newpage
\definecolor{answer}{rgb}{0.850980,0.047059,0.070588}
\null\begin{tikzpicture}[remember picture,overlay]\begin{scope}[shift={(current page.north west)},x=1bp,y=-1bp]
\node[anchor=north west,inner sep=0pt,text=answer,align=left,text width=232.0bp] at (15,15) {\fontsize{7.5bp}{9.5bp}\selectfont 析：\allowbreak{}D两图形大小、\allowbreak{}形状、\allowbreak{}朝向都相同，\allowbreak{}平移可重合，\allowbreak{}选D。\allowbreak{}\par};
\end{scope}\end{tikzpicture}\newpage
\definecolor{answer}{rgb}{0.850980,0.047059,0.070588}
\null\begin{tikzpicture}[remember picture,overlay]\begin{scope}[shift={(current page.north west)},x=1bp,y=-1bp]
\node[anchor=north west,inner sep=0pt,text=answer,align=left,text width=233.0bp] at (15,15) {\fontsize{7.5bp}{9.5bp}\selectfont 析：\allowbreak{}③车挡住出口所在行，\allowbreak{}下移\ensuremath{2}格即可让①车驶出。\allowbreak{}\par};
\end{scope}\end{tikzpicture}\newpage
\definecolor{answer}{rgb}{0.850980,0.047059,0.070588}
\null\begin{tikzpicture}[remember picture,overlay]\begin{scope}[shift={(current page.north west)},x=1bp,y=-1bp]
\node[anchor=north west,inner sep=0pt,text=answer,align=left,text width=232.0bp] at (15,15) {\fontsize{7.5bp}{9.5bp}\selectfont 析：\allowbreak{}从冰壶同一对应点数至P,\allowbreak{}向左\ensuremath{6}格，\allowbreak{}选B。\allowbreak{}\par};
\end{scope}\end{tikzpicture}\newpage
\definecolor{answer}{rgb}{0.850980,0.047059,0.070588}
\null\begin{tikzpicture}[remember picture,overlay]\begin{scope}[shift={(current page.north west)},x=1bp,y=-1bp]
\node[anchor=north west,inner sep=0pt,text=answer,align=left,text width=227.71429443359375bp] at (15,15) {\fontsize{8.57bp}{10.71bp}\selectfont 例如:\allowbreak{} 先向右平移\ensuremath{7}格,\allowbreak{}再向下平移\ensuremath{1}格。\allowbreak{}\par};
\end{scope}\end{tikzpicture}\newpage
\definecolor{answer}{rgb}{0.850980,0.047059,0.070588}
\null\begin{tikzpicture}[remember picture,overlay]\begin{scope}[shift={(current page.north west)},x=1bp,y=-1bp]
\node[anchor=north west,inner sep=0pt,text=answer,align=left,text width=232.0bp] at (15,15) {\fontsize{7.5bp}{9.5bp}\selectfont 析：\allowbreak{}取星形中心为对应点：\allowbreak{}先右移\ensuremath{7}格，\allowbreak{}再下移\ensuremath{1}格。\allowbreak{}\par};
\end{scope}\end{tikzpicture}\newpage
\definecolor{answer}{rgb}{0.850980,0.047059,0.070588}
\null\begin{tikzpicture}[remember picture,overlay]\begin{scope}[shift={(current page.north west)},x=1bp,y=-1bp]
\node[anchor=north west,inner sep=0pt,text=answer,align=left,text width=232.0bp] at (15,15) {\fontsize{8.5bp}{10.62bp}\selectfont 阶梯横边合计等于底边\ensuremath{9}厘米，\allowbreak{}下降边合计等于左边\ensuremath{3}厘米。\allowbreak{}\par};
\end{scope}\end{tikzpicture}\newpage
\definecolor{answer}{rgb}{0.850980,0.047059,0.070588}
\null\begin{tikzpicture}[remember picture,overlay]\begin{scope}[shift={(current page.north west)},x=1bp,y=-1bp]
\node[anchor=north west,inner sep=0pt,text=answer,align=left,text width=233.00001525878906bp] at (15,15) {\fontsize{7.5bp}{9.5bp}\selectfont 析：\allowbreak{}按对应顶点计格：\allowbreak{}①右\ensuremath{10}下\ensuremath{1};\allowbreak{}②右\ensuremath{6}上\ensuremath{1};\allowbreak{}③右\ensuremath{8};\allowbreak{}④右\ensuremath{6}下\ensuremath{2};\allowbreak{}⑤左\ensuremath{6}下\ensuremath{2}。\allowbreak{}\par};
\end{scope}\end{tikzpicture}\newpage
\definecolor{answer}{rgb}{0.850980,0.047059,0.070588}
\null\begin{tikzpicture}[remember picture,overlay]\begin{scope}[shift={(current page.north west)},x=1bp,y=-1bp]
\node[anchor=north west,inner sep=0pt,text=answer,align=left,text width=232.0bp] at (15,15) {\fontsize{8.5bp}{10.62bp}\selectfont 答：\allowbreak{}这个图形的周长是\ensuremath{24}厘米。\allowbreak{}\par};
\end{scope}\end{tikzpicture}\newpage
\definecolor{answer}{rgb}{0.850980,0.047059,0.070588}
\null\begin{tikzpicture}[remember picture,overlay]\begin{scope}[shift={(current page.north west)},x=1bp,y=-1bp]
\node[anchor=north west,inner sep=0pt,text=answer,align=left,text width=232.0bp] at (15,15) {\fontsize{7.5bp}{9.5bp}\selectfont 先画直角三角形，\allowbreak{}利用轴对称和平移构造重复图案。\allowbreak{}\par};
\end{scope}\end{tikzpicture}\newpage
\definecolor{answer}{rgb}{0.850980,0.047059,0.070588}
\null\begin{tikzpicture}[remember picture,overlay]\begin{scope}[shift={(current page.north west)},x=1bp,y=-1bp]
\node[anchor=north west,inner sep=0pt,text=answer,align=left,text width=247.0bp] at (15,15) {\fontsize{7.5bp}{9.5bp}\selectfont 析：\allowbreak{}法\ensuremath{1}向右平移；\allowbreak{}法\ensuremath{2}沿A、\allowbreak{}B之间的竖直中线对折，\allowbreak{}左右对称的{\scfallback 囍}字也能重合。\allowbreak{}\par};
\end{scope}\end{tikzpicture}\newpage
\definecolor{answer}{rgb}{0.850980,0.047059,0.070588}
\null\begin{tikzpicture}[remember picture,overlay]\begin{scope}[shift={(current page.north west)},x=1bp,y=-1bp]
\node[anchor=north west,inner sep=0pt,text=answer,align=left,text width=233.0bp] at (15,15) {\fontsize{7.5bp}{9.5bp}\selectfont 析：\allowbreak{}顺时针A→D→C→B;\allowbreak{}A顺\ensuremath{90}\ensuremath{^\circ}到D,\allowbreak{}逆\ensuremath{90}\ensuremath{^\circ}到B。\allowbreak{}\par};
\end{scope}\end{tikzpicture}\newpage
\definecolor{answer}{rgb}{0.850980,0.047059,0.070588}
\null\begin{tikzpicture}[remember picture,overlay]\begin{scope}[shift={(current page.north west)},x=1bp,y=-1bp]
\node[anchor=north west,inner sep=0pt,text=answer,align=left,text width=247.0bp] at (15,15) {\fontsize{7.5bp}{9.5bp}\selectfont 析：\allowbreak{}A与D位于O的对角两侧；\allowbreak{}绕O顺或逆时针\ensuremath{180}\ensuremath{^\circ}均可。\allowbreak{}\par};
\end{scope}\end{tikzpicture}\newpage
\definecolor{answer}{rgb}{0.850980,0.047059,0.070588}
\null\begin{tikzpicture}[remember picture,overlay]\begin{scope}[shift={(current page.north west)},x=1bp,y=-1bp]
\node[anchor=north west,inner sep=0pt,text=answer,align=left,text width=233.0bp] at (15,15) {\fontsize{7.5bp}{9.5bp}\selectfont 析：\allowbreak{}\ensuremath{6}时到\ensuremath{9}时共\ensuremath{3}小时；\allowbreak{}时针每小时\ensuremath{30}\ensuremath{^\circ},\allowbreak{}共\ensuremath{90}\ensuremath{^\circ}。\allowbreak{}\par};
\end{scope}\end{tikzpicture}\newpage
\definecolor{answer}{rgb}{0.850980,0.047059,0.070588}
\null\begin{tikzpicture}[remember picture,overlay]\begin{scope}[shift={(current page.north west)},x=1bp,y=-1bp]
\node[anchor=north west,inner sep=0pt,text=answer,align=left,text width=232bp] at (15,15) {\fontsize{7.5bp}{9.5bp}\selectfont 析：\allowbreak{}表盘\ensuremath{0}在上、\allowbreak{}\ensuremath{1}在右；\allowbreak{}顺时针\ensuremath{90}\ensuremath{^\circ}指到\ensuremath{1},\allowbreak{}填\ensuremath{1}千克。\allowbreak{}\par};
\end{scope}\end{tikzpicture}\newpage
\definecolor{answer}{rgb}{0.850980,0.047059,0.070588}
\null\begin{tikzpicture}[remember picture,overlay]\begin{scope}[shift={(current page.north west)},x=1bp,y=-1bp]
\node[anchor=north west,inner sep=0pt,text=answer,align=left,text width=233.0bp] at (15,15) {\fontsize{7.5bp}{9.5bp}\selectfont 析：\allowbreak{}B不动，\allowbreak{}A从上方到右方：\allowbreak{}绕B顺时针\ensuremath{90}\ensuremath{^\circ}。\allowbreak{}\par};
\end{scope}\end{tikzpicture}\newpage
\definecolor{answer}{rgb}{0.850980,0.047059,0.070588}
\null\begin{tikzpicture}[remember picture,overlay]\begin{scope}[shift={(current page.north west)},x=1bp,y=-1bp]
\node[anchor=north west,inner sep=0pt,text=answer,align=left,text width=233.0bp] at (15,15) {\fontsize{7.5bp}{9.5bp}\selectfont 析：\allowbreak{}B不动，\allowbreak{}A从上方到下方，\allowbreak{}按箭头逆时针\ensuremath{180}\ensuremath{^\circ}。\allowbreak{}\par};
\end{scope}\end{tikzpicture}\newpage
\definecolor{answer}{rgb}{0.850980,0.047059,0.070588}
\null\begin{tikzpicture}[remember picture,overlay]\begin{scope}[shift={(current page.north west)},x=1bp,y=-1bp]
\node[anchor=north west,inner sep=0pt,text=answer,align=left,text width=233.0bp] at (15,15) {\fontsize{7.5bp}{9.5bp}\selectfont 析：\allowbreak{}B不动，\allowbreak{}A从下方到右方：\allowbreak{}绕B逆时针\ensuremath{90}\ensuremath{^\circ}。\allowbreak{}\par};
\end{scope}\end{tikzpicture}\newpage
\definecolor{answer}{rgb}{0.850980,0.047059,0.070588}
\null\begin{tikzpicture}[remember picture,overlay]\begin{scope}[shift={(current page.north west)},x=1bp,y=-1bp]
\node[anchor=north west,inner sep=0pt,text=answer,align=left,text width=233.0bp] at (15,15) {\fontsize{7.5bp}{9.5bp}\selectfont 析：\allowbreak{}先顺转\ensuremath{90}\ensuremath{^\circ}再右移\ensuremath{2}格，\allowbreak{}恰与梯形拼成大三角形，\allowbreak{}填三角。\allowbreak{}\par};
\end{scope}\end{tikzpicture}\newpage
\definecolor{answer}{rgb}{0.850980,0.047059,0.070588}
\null\begin{tikzpicture}[remember picture,overlay]\begin{scope}[shift={(current page.north west)},x=1bp,y=-1bp]
\node[anchor=north west,inner sep=0pt,text=answer,align=left,text width=242.125bp] at (15,15) {\fontsize{7.54bp}{9.54bp}\selectfont 析：\allowbreak{}以不动的大三角形为参照，\allowbreak{}①向左\ensuremath{1}格，\allowbreak{}再向上\ensuremath{8}格\par};
\end{scope}\end{tikzpicture}\newpage
\definecolor{answer}{rgb}{0.850980,0.047059,0.070588}
\null\begin{tikzpicture}[remember picture,overlay]\begin{scope}[shift={(current page.north west)},x=1bp,y=-1bp]
\node[anchor=north west,inner sep=0pt,text=answer,align=left,text width=242.125bp] at (15,15) {\fontsize{7.53bp}{9.53bp}\selectfont 析：\allowbreak{}②绕O逆时针旋转\ensuremath{90}\ensuremath{^\circ}，\allowbreak{}再向左\ensuremath{3}格\par};
\end{scope}\end{tikzpicture}\newpage
\definecolor{answer}{rgb}{0.850980,0.047059,0.070588}
\null\begin{tikzpicture}[remember picture,overlay]\begin{scope}[shift={(current page.north west)},x=1bp,y=-1bp]
\node[anchor=north west,inner sep=0pt,text=answer,align=left,text width=242.125bp] at (15,15) {\fontsize{7.53bp}{9.53bp}\selectfont 析：\allowbreak{}③向左\ensuremath{4}格，\allowbreak{}再向上\ensuremath{8}格\par};
\end{scope}\end{tikzpicture}\newpage
\definecolor{answer}{rgb}{0.850980,0.047059,0.070588}
\null\begin{tikzpicture}[remember picture,overlay]\begin{scope}[shift={(current page.north west)},x=1bp,y=-1bp]
\node[anchor=north west,inner sep=0pt,text=answer,align=left,text width=233.0bp] at (15,15) {\fontsize{7.5bp}{9.5bp}\selectfont 析：\allowbreak{}A与M同向，\allowbreak{}可平移；\allowbreak{}C、\allowbreak{}D、\allowbreak{}E可旋转重合；\allowbreak{}B、\allowbreak{}F 是镜像。\allowbreak{}\par};
\end{scope}\end{tikzpicture}\newpage
\definecolor{answer}{rgb}{0.850980,0.047059,0.070588}
\null\begin{tikzpicture}[remember picture,overlay]\begin{scope}[shift={(current page.north west)},x=1bp,y=-1bp]
\node[anchor=north west,inner sep=0pt,text=answer,align=left,text width=233.0bp] at (15,15) {\fontsize{7.5bp}{9.5bp}\selectfont 析：\allowbreak{}先顺时针转\ensuremath{90}\ensuremath{^\circ},\allowbreak{}移到最右两列，\allowbreak{}再下落，\allowbreak{}正好补齐下两层。\allowbreak{}\par};
\end{scope}\end{tikzpicture}\newpage
\definecolor{answer}{rgb}{0.850980,0.047059,0.070588}
\null\begin{tikzpicture}[remember picture,overlay]\begin{scope}[shift={(current page.north west)},x=1bp,y=-1bp]
\node[anchor=north west,inner sep=0pt,text=answer,align=left,text width=233.0bp] at (15,15) {\fontsize{7.5bp}{9.5bp}\selectfont 析：\allowbreak{}平均数概括一组数据的总体情况，\allowbreak{}不能当成某个个体的值。\allowbreak{}\par};
\end{scope}\end{tikzpicture}\newpage
\definecolor{answer}{rgb}{0.850980,0.047059,0.070588}
\null\begin{tikzpicture}[remember picture,overlay]\begin{scope}[shift={(current page.north west)},x=1bp,y=-1bp]
\node[anchor=north west,inner sep=0pt,text=answer,align=left,text width=242.0bp] at (15,15) {\fontsize{7.5bp}{9.5bp}\selectfont 析：\allowbreak{}√。\allowbreak{}(\ensuremath{12} + \ensuremath{15} + \ensuremath{18}) ÷ \ensuremath{3} = \ensuremath{15}岁。\allowbreak{}\par};
\end{scope}\end{tikzpicture}\newpage
\definecolor{answer}{rgb}{0.850980,0.047059,0.070588}
\null\begin{tikzpicture}[remember picture,overlay]\begin{scope}[shift={(current page.north west)},x=1bp,y=-1bp]
\node[anchor=north west,inner sep=0pt,text=answer,align=left,text width=242.0bp] at (15,15) {\fontsize{7.5bp}{9.5bp}\selectfont 析：\allowbreak{}×。\allowbreak{}平均每天\ensuremath{5}个，\allowbreak{}不表示每天都恰写\ensuremath{5}个。\allowbreak{}\par};
\end{scope}\end{tikzpicture}\newpage
\definecolor{answer}{rgb}{0.850980,0.047059,0.070588}
\null\begin{tikzpicture}[remember picture,overlay]\begin{scope}[shift={(current page.north west)},x=1bp,y=-1bp]
\node[anchor=north west,inner sep=0pt,text=answer,align=left,text width=242.0bp] at (15,15) {\fontsize{7.5bp}{9.5bp}\selectfont 析：\allowbreak{}×。\allowbreak{}平均\ensuremath{160}厘米，\allowbreak{}可有人\ensuremath{155}厘米，\allowbreak{}由较高者补足差额。\allowbreak{}\par};
\end{scope}\end{tikzpicture}\newpage
\definecolor{answer}{rgb}{0.850980,0.047059,0.070588}
\null\begin{tikzpicture}[remember picture,overlay]\begin{scope}[shift={(current page.north west)},x=1bp,y=-1bp]
\node[anchor=north west,inner sep=0pt,text=answer,align=left,text width=233.0bp] at (15,15) {\fontsize{7.5bp}{9.5bp}\selectfont 析：\allowbreak{}法\ensuremath{1}(\ensuremath{48}+\ensuremath{42})÷\ensuremath{3} = \ensuremath{30};\allowbreak{}法\ensuremath{2}甲乙均\ensuremath{24},\allowbreak{}丙多\ensuremath{18},\allowbreak{}平均补\ensuremath{6},\allowbreak{} \ensuremath{24}+\ensuremath{6} =\ensuremath{30}。\allowbreak{}\par};
\end{scope}\end{tikzpicture}\newpage
\definecolor{answer}{rgb}{0.850980,0.047059,0.070588}
\null\begin{tikzpicture}[remember picture,overlay]\begin{scope}[shift={(current page.north west)},x=1bp,y=-1bp]
\node[anchor=north west,inner sep=0pt,text=answer,align=left,text width=242.0bp] at (15,15) {\fontsize{7.5bp}{9.5bp}\selectfont 析：\allowbreak{}周一至周五共\ensuremath{5}天，\allowbreak{}总本数应除以\ensuremath{5},\allowbreak{}选A;\allowbreak{}选项续在下一页。\allowbreak{}\par};
\end{scope}\end{tikzpicture}\newpage
\definecolor{answer}{rgb}{0.850980,0.047059,0.070588}
\null\begin{tikzpicture}[remember picture,overlay]\begin{scope}[shift={(current page.north west)},x=1bp,y=-1bp]
\node[anchor=north west,inner sep=0pt,text=answer,align=left,text width=233.0bp] at (15,15) {\fontsize{7.5bp}{9.5bp}\selectfont 析：\allowbreak{}×。\allowbreak{}最小数≤平均数≤最大数；\allowbreak{}各数全相等时取等号。\allowbreak{}\par};
\end{scope}\end{tikzpicture}\newpage
\definecolor{answer}{rgb}{0.850980,0.047059,0.070588}
\null\begin{tikzpicture}[remember picture,overlay]\begin{scope}[shift={(current page.north west)},x=1bp,y=-1bp]
\node[anchor=north west,inner sep=0pt,text=answer,align=left,text width=247.0bp] at (15,15) {\fontsize{8.5bp}{10.62bp}\selectfont 平均每包质量为样本\ensuremath{6}包总质量除以\ensuremath{6},\allowbreak{}也可把\ensuremath{11}和\ensuremath{13}移多补少为\ensuremath{12}。\allowbreak{}\par};
\end{scope}\end{tikzpicture}\newpage
\definecolor{answer}{rgb}{0.850980,0.047059,0.070588}
\null\begin{tikzpicture}[remember picture,overlay]\begin{scope}[shift={(current page.north west)},x=1bp,y=-1bp]
\node[anchor=north west,inner sep=0pt,text=answer,align=left,text width=233.125bp] at (15,15) {\fontsize{7.52bp}{9.52bp}\selectfont 析：\allowbreak{}把\ensuremath{35}中的\ensuremath{5}移给\ensuremath{25}，\allowbreak{}三天相当于各\ensuremath{30}个，\allowbreak{}选B\par};
\end{scope}\end{tikzpicture}\newpage
\definecolor{answer}{rgb}{0.850980,0.047059,0.070588}
\null\begin{tikzpicture}[remember picture,overlay]\begin{scope}[shift={(current page.north west)},x=1bp,y=-1bp]
\node[anchor=north west,inner sep=0pt,text=answer,align=left,text width=247.0bp] at (15,15) {\fontsize{8.5bp}{10.62bp}\selectfont 答：\allowbreak{}\ensuremath{11}克与\ensuremath{13}克配成平均\ensuremath{12}克，\allowbreak{}原来\ensuremath{12}克不变，\allowbreak{}因此\ensuremath{6}包平均每包\ensuremath{12}克。\allowbreak{}\par};
\end{scope}\end{tikzpicture}\newpage
\definecolor{answer}{rgb}{0.850980,0.047059,0.070588}
\null\begin{tikzpicture}[remember picture,overlay]\begin{scope}[shift={(current page.north west)},x=1bp,y=-1bp]
\node[anchor=north west,inner sep=0pt,text=answer,align=left,text width=232.39999999999998bp] at (15,15) {\fontsize{8.5bp}{10.62bp}\selectfont 把\ensuremath{12}看作\ensuremath{10}，\allowbreak{}按每小包大约\ensuremath{10}克估计\par};
\end{scope}\end{tikzpicture}\newpage
\definecolor{answer}{rgb}{0.850980,0.047059,0.070588}
\null\begin{tikzpicture}[remember picture,overlay]\begin{scope}[shift={(current page.north west)},x=1bp,y=-1bp]
\node[anchor=north west,inner sep=0pt,text=answer,align=left,text width=233.0bp] at (15,15) {\fontsize{7.5bp}{9.5bp}\selectfont 析：\allowbreak{}加进的数等于原平均数\ensuremath{12},\allowbreak{}平均数才不变，\allowbreak{}选B。\allowbreak{}\par};
\end{scope}\end{tikzpicture}\newpage
\definecolor{answer}{rgb}{0.850980,0.047059,0.070588}
\null\begin{tikzpicture}[remember picture,overlay]\begin{scope}[shift={(current page.north west)},x=1bp,y=-1bp]
\node[anchor=north west,inner sep=0pt,text=answer,align=left,text width=247.0bp] at (15,15) {\fontsize{8.5bp}{10.62bp}\selectfont 四周总数应是\ensuremath{8} × \ensuremath{4},\allowbreak{}扣掉前三周\ensuremath{4}、\allowbreak{}\ensuremath{5}、\allowbreak{}\ensuremath{7}个；\allowbreak{}也可求前三周比平均数的不足。\allowbreak{}\par};
\end{scope}\end{tikzpicture}\newpage
\definecolor{answer}{rgb}{0.850980,0.047059,0.070588}
\null\begin{tikzpicture}[remember picture,overlay]\begin{scope}[shift={(current page.north west)},x=1bp,y=-1bp]
\node[anchor=north west,inner sep=0pt,text=answer,align=left,text width=242.0bp] at (15,15) {\fontsize{8.5bp}{10.62bp}\selectfont 两队各\ensuremath{3}人，\allowbreak{}人数相同，\allowbreak{}可以总分除以总人数，\allowbreak{}也可平均两个队的平均分。\allowbreak{}\par};
\end{scope}\end{tikzpicture}\newpage
\definecolor{answer}{rgb}{0.850980,0.047059,0.070588}
\null\begin{tikzpicture}[remember picture,overlay]\begin{scope}[shift={(current page.north west)},x=1bp,y=-1bp]
\node[anchor=north west,inner sep=0pt,text=answer,align=left,text width=233.0bp] at (15,15) {\fontsize{8.5bp}{10.62bp}\selectfont 两组人数不同，\allowbreak{}要比较平均数；\allowbreak{}不能用总次数直接比较。\allowbreak{}\par};
\end{scope}\end{tikzpicture}\newpage
\definecolor{answer}{rgb}{0.850980,0.047059,0.070588}
\null\begin{tikzpicture}[remember picture,overlay]\begin{scope}[shift={(current page.north west)},x=1bp,y=-1bp]
\node[anchor=north west,inner sep=0pt,text=answer,align=left,text width=242.0bp] at (15,15) {\fontsize{8.5bp}{10.62bp}\selectfont 答：\allowbreak{}两队队员的平均成绩是\ensuremath{78}分。\allowbreak{}\par};
\end{scope}\end{tikzpicture}\newpage
\definecolor{answer}{rgb}{0.850980,0.047059,0.070588}
\null\begin{tikzpicture}[remember picture,overlay]\begin{scope}[shift={(current page.north west)},x=1bp,y=-1bp]
\node[anchor=north west,inner sep=0pt,text=answer,align=left,text width=242.0bp] at (15,15) {\fontsize{8.5bp}{10.62bp}\selectfont 答：\allowbreak{}把高分组的\ensuremath{3}分补给低分组，\allowbreak{}两组都是\ensuremath{78}分。\allowbreak{}\par};
\end{scope}\end{tikzpicture}\newpage
\definecolor{answer}{rgb}{0.850980,0.047059,0.070588}
\null\begin{tikzpicture}[remember picture,overlay]\begin{scope}[shift={(current page.north west)},x=1bp,y=-1bp]
\node[anchor=north west,inner sep=0pt,text=answer,align=left,text width=233.0bp] at (15,15) {\fontsize{8.5bp}{10.62bp}\selectfont 答：\allowbreak{}二组平均数更高，\allowbreak{}整体跳绳水平更好。\allowbreak{}\par};
\end{scope}\end{tikzpicture}\newpage
\definecolor{answer}{rgb}{0.850980,0.047059,0.070588}
\null\begin{tikzpicture}[remember picture,overlay]\begin{scope}[shift={(current page.north west)},x=1bp,y=-1bp]
\node[anchor=north west,inner sep=0pt,text=answer,align=left,text width=233.0bp] at (15,15) {\fontsize{8.5bp}{10.62bp}\selectfont 答：\allowbreak{}二组平均\ensuremath{140}个，\allowbreak{}大于一组\ensuremath{137}个，\allowbreak{}所以二组更好。\allowbreak{}\par};
\end{scope}\end{tikzpicture}\newpage
\definecolor{answer}{rgb}{0.850980,0.047059,0.070588}
\null\begin{tikzpicture}[remember picture,overlay]\begin{scope}[shift={(current page.north west)},x=1bp,y=-1bp]
\node[anchor=north west,inner sep=0pt,text=answer,align=left,text width=233.0bp] at (15,15) {\fontsize{8.5bp}{10.62bp}\selectfont 已知路程\ensuremath{400}米、\allowbreak{}时间\ensuremath{8}分钟，\allowbreak{}速度用路程除以时间。\allowbreak{}\par};
\end{scope}\end{tikzpicture}\newpage
\definecolor{answer}{rgb}{0.850980,0.047059,0.070588}
\null\begin{tikzpicture}[remember picture,overlay]\begin{scope}[shift={(current page.north west)},x=1bp,y=-1bp]
\node[anchor=north west,inner sep=0pt,text=answer,align=left,text width=233.0bp] at (15,15) {\fontsize{8.5bp}{10.62bp}\selectfont 答：\allowbreak{}下坡时平均每分钟走\ensuremath{50}米。\allowbreak{}\par};
\end{scope}\end{tikzpicture}\newpage
\definecolor{answer}{rgb}{0.850980,0.047059,0.070588}
\null\begin{tikzpicture}[remember picture,overlay]\begin{scope}[shift={(current page.north west)},x=1bp,y=-1bp]
\node[anchor=north west,inner sep=0pt,text=answer,align=left,text width=242.0bp] at (15,15) {\fontsize{7.5bp}{9.5bp}\selectfont 析：\allowbreak{}各去\ensuremath{1}个最高分和\ensuremath{1}个最低分，\allowbreak{}\ensuremath{6} \ensuremath{-}\ensuremath{2} = \ensuremath{4}，\allowbreak{}故除\ensuremath{4}：\allowbreak{}于巧\ensuremath{95}、\allowbreak{}李丹\ensuremath{94}、\allowbreak{}王伟\ensuremath{93}、\allowbreak{}刘鑫\ensuremath{92}分。\allowbreak{}\par};
\end{scope}\end{tikzpicture}\newpage
\definecolor{answer}{rgb}{0.850980,0.047059,0.070588}
\null\begin{tikzpicture}[remember picture,overlay]\begin{scope}[shift={(current page.north west)},x=1bp,y=-1bp]
\node[anchor=north west,inner sep=0pt,text=answer,align=left,text width=242.0bp] at (15,15) {\fontsize{7.5bp}{9.5bp}\selectfont 析：\allowbreak{}各去掉一个最高分和一个最低分后再求平均分，\allowbreak{}\ensuremath{95} > \ensuremath{94} > \ensuremath{93} > \ensuremath{92}，\allowbreak{}于巧第一名。\allowbreak{}\par};
\end{scope}\end{tikzpicture}\newpage
\definecolor{answer}{rgb}{0.850980,0.047059,0.070588}
\null\begin{tikzpicture}[remember picture,overlay]\begin{scope}[shift={(current page.north west)},x=1bp,y=-1bp]
\node[anchor=north west,inner sep=0pt,text=answer,align=left,text width=233.0bp] at (15,15) {\fontsize{8.5bp}{10.62bp}\selectfont 平均速度等于总路程除以总时间；\allowbreak{}也可先验证上、\allowbreak{}下坡速度一样。\allowbreak{}\par};
\end{scope}\end{tikzpicture}\newpage
\definecolor{answer}{rgb}{0.850980,0.047059,0.070588}
\null\begin{tikzpicture}[remember picture,overlay]\begin{scope}[shift={(current page.north west)},x=1bp,y=-1bp]
\node[anchor=north west,inner sep=0pt,text=answer,align=left,text width=233.0bp] at (15,15) {\fontsize{8.5bp}{10.62bp}\selectfont 答：\allowbreak{}上、\allowbreak{}下坡平均每分钟走\ensuremath{50}米。\allowbreak{}\par};
\end{scope}\end{tikzpicture}\newpage
\definecolor{answer}{rgb}{0.850980,0.047059,0.070588}
\null\begin{tikzpicture}[remember picture,overlay]\begin{scope}[shift={(current page.north west)},x=1bp,y=-1bp]
\node[anchor=north west,inner sep=0pt,text=answer,align=left,text width=232bp] at (15,15) {\fontsize{8.5bp}{10.62bp}\selectfont 答：\allowbreak{}下坡与上坡速度都是\ensuremath{50}米/分，\allowbreak{}因此全程平均每分钟走\ensuremath{50}米。\allowbreak{}\par};
\end{scope}\end{tikzpicture}\newpage
\definecolor{answer}{rgb}{0.850980,0.047059,0.070588}
\null\begin{tikzpicture}[remember picture,overlay]\begin{scope}[shift={(current page.north west)},x=1bp,y=-1bp]
\node[anchor=north west,inner sep=0pt,text=answer,align=left,text width=233.00001525878906bp] at (15,15) {\fontsize{7.5bp}{9.5bp}\selectfont 析：\allowbreak{}每人只归一段：\allowbreak{}\ensuremath{100}、\allowbreak{}\ensuremath{90}～\ensuremath{99}、\allowbreak{}\ensuremath{80}～\ensuremath{89}、\allowbreak{}\ensuremath{79}及以下分别\ensuremath{1}、\allowbreak{}\ensuremath{7}、\allowbreak{}\ensuremath{5}、\allowbreak{}\ensuremath{3}人，\allowbreak{}合\ensuremath{16}人。\allowbreak{}\par};
\end{scope}\end{tikzpicture}\newpage
\definecolor{answer}{rgb}{0.850980,0.047059,0.070588}
\null\begin{tikzpicture}[remember picture,overlay]\begin{scope}[shift={(current page.north west)},x=1bp,y=-1bp]
\node[anchor=north west,inner sep=0pt,text=answer,align=left,text width=233.0bp] at (15,15) {\fontsize{7.5bp}{9.5bp}\selectfont 析：\allowbreak{}比较频数\ensuremath{1}、\allowbreak{}\ensuremath{7}、\allowbreak{}\ensuremath{5}、\allowbreak{}\ensuremath{3},\allowbreak{} \ensuremath{7}最多，\allowbreak{}对应\ensuremath{90}～\ensuremath{99}分。\allowbreak{}\par};
\end{scope}\end{tikzpicture}\newpage
\definecolor{answer}{rgb}{0.850980,0.047059,0.070588}
\null\begin{tikzpicture}[remember picture,overlay]\begin{scope}[shift={(current page.north west)},x=1bp,y=-1bp]
\node[anchor=north west,inner sep=0pt,text=answer,align=left,text width=234.85714721679688bp] at (15,15) {\fontsize{10.0bp}{12.5bp}\selectfont 例如：\allowbreak{}全组有\ensuremath{16}人，\allowbreak{}\ensuremath{80}分及以上有\ensuremath{13}人。\allowbreak{}\par};
\end{scope}\end{tikzpicture}\newpage
\definecolor{answer}{rgb}{0.850980,0.047059,0.070588}
\null\begin{tikzpicture}[remember picture,overlay]\begin{scope}[shift={(current page.north west)},x=1bp,y=-1bp]
\node[anchor=north west,inner sep=0pt,text=answer,align=left,text width=233.0bp] at (15,15) {\fontsize{7.5bp}{9.5bp}\selectfont 析：\allowbreak{}\ensuremath{80}分及以上包括前三段，\allowbreak{}\ensuremath{1} + \ensuremath{7} + \ensuremath{5} = \ensuremath{13}人，\allowbreak{}共\ensuremath{16}人。\allowbreak{}\par};
\end{scope}\end{tikzpicture}\newpage
\definecolor{answer}{rgb}{0.850980,0.047059,0.070588}
\null\begin{tikzpicture}[remember picture,overlay]\begin{scope}[shift={(current page.north west)},x=1bp,y=-1bp]
\node[anchor=north west,inner sep=0pt,text=answer,align=left,text width=242.0bp] at (15,15) {\fontsize{7.5bp}{9.5bp}\selectfont 析：\allowbreak{}法\ensuremath{1}逐个画记分类；\allowbreak{}法\ensuremath{2}先排序分段。\allowbreak{}男\ensuremath{3}、\allowbreak{}\ensuremath{7}、\allowbreak{}\ensuremath{8}、\allowbreak{}\ensuremath{4};\allowbreak{}女\ensuremath{6}、\allowbreak{}\ensuremath{9}、\allowbreak{}\ensuremath{6}、\allowbreak{}\ensuremath{1};\allowbreak{}合计\ensuremath{9}、\allowbreak{}\ensuremath{16}、\allowbreak{}\ensuremath{14}、\allowbreak{}\ensuremath{5}人。\allowbreak{}\par};
\end{scope}\end{tikzpicture}\newpage
\definecolor{answer}{rgb}{0.850980,0.047059,0.070588}
\null\begin{tikzpicture}[remember picture,overlay]\begin{scope}[shift={(current page.north west)},x=1bp,y=-1bp]
\node[anchor=north west,inner sep=0pt,text=answer,align=left,text width=242.0bp] at (15,15) {\fontsize{7.5bp}{9.5bp}\selectfont 析：\allowbreak{}按性别和区间读表：\allowbreak{}男\ensuremath{31}～\ensuremath{40}为\ensuremath{7}人，\allowbreak{}\ensuremath{41}～\ensuremath{50}为\ensuremath{8}人；\allowbreak{}女\ensuremath{30}及以下\ensuremath{6}人。\allowbreak{}\par};
\end{scope}\end{tikzpicture}\newpage
\definecolor{answer}{rgb}{0.850980,0.047059,0.070588}
\null\begin{tikzpicture}[remember picture,overlay]\begin{scope}[shift={(current page.north west)},x=1bp,y=-1bp]
\node[anchor=north west,inner sep=0pt,text=answer,align=left,text width=233.0bp] at (15,15) {\fontsize{7.5bp}{9.5bp}\selectfont 析：\allowbreak{}按四段分类人数\ensuremath{1}、\allowbreak{}\ensuremath{5}、\allowbreak{}\ensuremath{3}、\allowbreak{}\ensuremath{3},\allowbreak{}总计\ensuremath{12}人。\allowbreak{}\par};
\end{scope}\end{tikzpicture}\newpage
\definecolor{answer}{rgb}{0.850980,0.047059,0.070588}
\null\begin{tikzpicture}[remember picture,overlay]\begin{scope}[shift={(current page.north west)},x=1bp,y=-1bp]
\node[anchor=north west,inner sep=0pt,text=answer,align=left,text width=233.0bp] at (15,15) {\fontsize{7.5bp}{9.5bp}\selectfont 析：\allowbreak{}比较\ensuremath{1}、\allowbreak{}\ensuremath{5}、\allowbreak{}\ensuremath{3}、\allowbreak{}\ensuremath{3},\allowbreak{} \ensuremath{90}～\ensuremath{99}分的\ensuremath{5}人最多。\allowbreak{}\par};
\end{scope}\end{tikzpicture}\newpage
\definecolor{answer}{rgb}{0.850980,0.047059,0.070588}
\null\begin{tikzpicture}[remember picture,overlay]\begin{scope}[shift={(current page.north west)},x=1bp,y=-1bp]
\node[anchor=north west,inner sep=0pt,text=answer,align=left,text width=242.0bp] at (15,15) {\fontsize{7.5bp}{9.5bp}\selectfont 析：\allowbreak{}合计中\ensuremath{16}最大，\allowbreak{}对应\ensuremath{31}～\ensuremath{40}千克；\allowbreak{}\ensuremath{51}及以上\ensuremath{4} + \ensuremath{1} = \ensuremath{5}人。\allowbreak{}\par};
\end{scope}\end{tikzpicture}\newpage
\definecolor{answer}{rgb}{0.850980,0.047059,0.070588}
\null\begin{tikzpicture}[remember picture,overlay]\begin{scope}[shift={(current page.north west)},x=1bp,y=-1bp]
\node[anchor=north west,inner sep=0pt,text=answer,align=left,text width=234.85714721679688bp] at (15,15) {\fontsize{10.0bp}{12.5bp}\selectfont 例如：\allowbreak{}全组有\ensuremath{12}人，\allowbreak{}\ensuremath{80}分及以上有\ensuremath{9}人。\allowbreak{}\par};
\end{scope}\end{tikzpicture}\newpage
\definecolor{answer}{rgb}{0.850980,0.047059,0.070588}
\null\begin{tikzpicture}[remember picture,overlay]\begin{scope}[shift={(current page.north west)},x=1bp,y=-1bp]
\node[anchor=north west,inner sep=0pt,text=answer,align=left,text width=233.0bp] at (15,15) {\fontsize{7.5bp}{9.5bp}\selectfont 析：\allowbreak{}全组\ensuremath{12}人，\allowbreak{}\ensuremath{80}分及以上\ensuremath{1} + \ensuremath{5} + \ensuremath{3} = \ensuremath{9}人。\allowbreak{}\par};
\end{scope}\end{tikzpicture}\newpage
\definecolor{answer}{rgb}{0.850980,0.047059,0.070588}
\null\begin{tikzpicture}[remember picture,overlay]\begin{scope}[shift={(current page.north west)},x=1bp,y=-1bp]
\node[anchor=north west,inner sep=0pt,text=answer,align=left,text width=233.0bp] at (15,15) {\fontsize{7.5bp}{9.5bp}\selectfont 析：\allowbreak{}按\ensuremath{99}及以下、\allowbreak{}\ensuremath{100}～\ensuremath{119}、\allowbreak{}\ensuremath{120}～\ensuremath{139}、\allowbreak{}\ensuremath{140}～\ensuremath{149}、\allowbreak{}\ensuremath{150}～\ensuremath{159}、\allowbreak{}\ensuremath{160}及以上计数：\allowbreak{}\ensuremath{2}、\allowbreak{}\ensuremath{4}、\allowbreak{}\ensuremath{13}、\allowbreak{}\ensuremath{7}、\allowbreak{}\ensuremath{6}、\allowbreak{}\ensuremath{4},\allowbreak{}合\ensuremath{36}人。\allowbreak{}\par};
\end{scope}\end{tikzpicture}\newpage
\definecolor{answer}{rgb}{0.850980,0.047059,0.070588}
\null\begin{tikzpicture}[remember picture,overlay]\begin{scope}[shift={(current page.north west)},x=1bp,y=-1bp]
\node[anchor=north west,inner sep=0pt,text=answer,align=left,text width=242.0bp] at (15,15) {\fontsize{7.5bp}{9.5bp}\selectfont 析：\allowbreak{}按≥\ensuremath{40}、\allowbreak{}\ensuremath{30}～\ensuremath{39}、\allowbreak{}\ensuremath{20}～\ensuremath{29}、\allowbreak{}≤\ensuremath{19}分段；\allowbreak{}一组\ensuremath{8}、\allowbreak{}\ensuremath{7}、\allowbreak{}\ensuremath{4}、\allowbreak{}\ensuremath{1},\allowbreak{}二组\ensuremath{7}、\allowbreak{}\ensuremath{8}、\allowbreak{}\ensuremath{4}、\allowbreak{}\ensuremath{1},\allowbreak{}各\ensuremath{20}人。\allowbreak{}\par};
\end{scope}\end{tikzpicture}\newpage
\definecolor{answer}{rgb}{0.850980,0.047059,0.070588}
\null\begin{tikzpicture}[remember picture,overlay]\begin{scope}[shift={(current page.north west)},x=1bp,y=-1bp]
\node[anchor=north west,inner sep=0pt,text=answer,align=left,text width=233.42857360839844bp] at (15,15) {\fontsize{10.0bp}{12.5bp}\selectfont \ensuremath{120}～\ensuremath{139}个范围内的人数最多，\allowbreak{}有\ensuremath{13}人。\allowbreak{}\par};
\end{scope}\end{tikzpicture}\newpage
\definecolor{answer}{rgb}{0.850980,0.047059,0.070588}
\null\begin{tikzpicture}[remember picture,overlay]\begin{scope}[shift={(current page.north west)},x=1bp,y=-1bp]
\node[anchor=north west,inner sep=0pt,text=answer,align=left,text width=233.0bp] at (15,15) {\fontsize{7.5bp}{9.5bp}\selectfont 析：\allowbreak{}频数\ensuremath{13}最大，\allowbreak{}对应\ensuremath{120}～\ensuremath{139}个，\allowbreak{}共\ensuremath{13}人。\allowbreak{}\par};
\end{scope}\end{tikzpicture}\newpage
\definecolor{answer}{rgb}{0.850980,0.047059,0.070588}
\null\begin{tikzpicture}[remember picture,overlay]\begin{scope}[shift={(current page.north west)},x=1bp,y=-1bp]
\node[anchor=north west,inner sep=0pt,text=answer,align=left,text width=241.28570556640625bp] at (15,15) {\fontsize{10.0bp}{12.5bp}\selectfont 例如：\allowbreak{}两组均有\ensuremath{20}人，\allowbreak{}不及格的各有\ensuremath{1}人；\allowbreak{}一组优秀的比二组多\ensuremath{1}人，\allowbreak{}二组良好的比一组多\ensuremath{1}人。\allowbreak{}\par};
\end{scope}\end{tikzpicture}\newpage
\definecolor{answer}{rgb}{0.850980,0.047059,0.070588}
\null\begin{tikzpicture}[remember picture,overlay]\begin{scope}[shift={(current page.north west)},x=1bp,y=-1bp]
\node[anchor=north west,inner sep=0pt,text=answer,align=left,text width=242.0bp] at (15,15) {\fontsize{7.5bp}{9.5bp}\selectfont 析：\allowbreak{}两组各\ensuremath{20}人，\allowbreak{}及格各\ensuremath{4}、\allowbreak{}不及格各\ensuremath{1};\allowbreak{}一组优秀多\ensuremath{1}人，\allowbreak{}二组良好多\ensuremath{1}人。\allowbreak{}\par};
\end{scope}\end{tikzpicture}\newpage
\definecolor{answer}{rgb}{0.850980,0.047059,0.070588}
\null\begin{tikzpicture}[remember picture,overlay]\begin{scope}[shift={(current page.north west)},x=1bp,y=-1bp]
\node[anchor=north west,inner sep=0pt,text=answer,align=left,text width=233.0bp] at (15,15) {\fontsize{8.5bp}{10.62bp}\selectfont 可以合并相邻组，\allowbreak{}也可比较两个组，\allowbreak{}以下三种提问均有效。\allowbreak{}\par};
\end{scope}\end{tikzpicture}\newpage
\definecolor{answer}{rgb}{0.850980,0.047059,0.070588}
\null\begin{tikzpicture}[remember picture,overlay]\begin{scope}[shift={(current page.north west)},x=1bp,y=-1bp]
\node[anchor=north west,inner sep=0pt,text=answer,align=left,text width=233.0bp] at (15,15) {\fontsize{8.5bp}{10.62bp}\selectfont 答：\allowbreak{}\ensuremath{120}～\ensuremath{139}个组比\ensuremath{140}～\ensuremath{149}个组多\ensuremath{6}人。\allowbreak{}\par};
\end{scope}\end{tikzpicture}\newpage
\definecolor{answer}{rgb}{0.850980,0.047059,0.070588}
\null\begin{tikzpicture}[remember picture,overlay]\begin{scope}[shift={(current page.north west)},x=1bp,y=-1bp]
\node[anchor=north west,inner sep=0pt,text=answer,align=left,text width=233.0bp] at (15,15) {\fontsize{7.5bp}{9.5bp}\selectfont 析：\allowbreak{}每人只计一次，\allowbreak{}四段人数\ensuremath{2}、\allowbreak{}\ensuremath{11}、\allowbreak{}\ensuremath{4}、\allowbreak{}\ensuremath{4},\allowbreak{}合计\ensuremath{21}。\allowbreak{}\par};
\end{scope}\end{tikzpicture}\newpage
\definecolor{answer}{rgb}{0.850980,0.047059,0.070588}
\null\begin{tikzpicture}[remember picture,overlay]\begin{scope}[shift={(current page.north west)},x=1bp,y=-1bp]
\node[anchor=north west,inner sep=0pt,text=answer,align=left,text width=233.0bp] at (15,15) {\fontsize{7.5bp}{9.5bp}\selectfont 析：\allowbreak{}四段频数中\ensuremath{11}最大，\allowbreak{}对应\ensuremath{4.0}～\ensuremath{4.5}。\allowbreak{}\par};
\end{scope}\end{tikzpicture}\newpage
\definecolor{answer}{rgb}{0.850980,0.047059,0.070588}
\null\begin{tikzpicture}[remember picture,overlay]\begin{scope}[shift={(current page.north west)},x=1bp,y=-1bp]
\node[anchor=north west,inner sep=0pt,text=answer,align=left,text width=234.14285278320312bp] at (15,15) {\fontsize{8.57bp}{10.71bp}\selectfont 例如：\allowbreak{}注意用眼卫生，\allowbreak{}保持正确读写姿势，\allowbreak{}适当休息，\allowbreak{}多参加户外活动。\allowbreak{}\par};
\end{scope}\end{tikzpicture}\newpage
\definecolor{answer}{rgb}{0.850980,0.047059,0.070588}
\null\begin{tikzpicture}[remember picture,overlay]\begin{scope}[shift={(current page.north west)},x=1bp,y=-1bp]
\node[anchor=north west,inner sep=0pt,text=answer,align=left,text width=233.0bp] at (15,15) {\fontsize{7.5bp}{9.5bp}\selectfont 析：\allowbreak{}注意用眼卫生、\allowbreak{}保持读写姿势，\allowbreak{}适当休息，\allowbreak{}多做户外活动。\allowbreak{}\par};
\end{scope}\end{tikzpicture}\newpage
\definecolor{answer}{rgb}{0.850980,0.047059,0.070588}
\null\begin{tikzpicture}[remember picture,overlay]\begin{scope}[shift={(current page.north west)},x=1bp,y=-1bp]
\node[anchor=north west,inner sep=0pt,text=answer,align=left,text width=242.0bp] at (15,15) {\fontsize{7.5bp}{9.5bp}\selectfont 析：\allowbreak{}看叉号最高一列：\allowbreak{}男生\ensuremath{155}厘米\ensuremath{9}人；\allowbreak{}女生\ensuremath{150}厘米\ensuremath{8}人。\allowbreak{}\par};
\end{scope}\end{tikzpicture}\newpage
\definecolor{answer}{rgb}{0.850980,0.047059,0.070588}
\null\begin{tikzpicture}[remember picture,overlay]\begin{scope}[shift={(current page.north west)},x=1bp,y=-1bp]
\node[anchor=north west,inner sep=0pt,text=answer,align=left,text width=233.0bp] at (15,15) {\fontsize{7.5bp}{9.5bp}\selectfont 析：\allowbreak{}按四段计数：\allowbreak{}四(\ensuremath{1})\ensuremath{4}、\allowbreak{}\ensuremath{3}、\allowbreak{}\ensuremath{3}、\allowbreak{}\ensuremath{2};\allowbreak{}四(\ensuremath{2})\ensuremath{2}、\allowbreak{}\ensuremath{3}、\allowbreak{}\ensuremath{4}、\allowbreak{}\ensuremath{3},\allowbreak{}各\ensuremath{12}人。\allowbreak{}\par};
\end{scope}\end{tikzpicture}\newpage
\definecolor{answer}{rgb}{0.850980,0.047059,0.070588}
\null\begin{tikzpicture}[remember picture,overlay]\begin{scope}[shift={(current page.north west)},x=1bp,y=-1bp]
\node[anchor=north west,inner sep=0pt,text=answer,align=left,text width=242.0bp] at (15,15) {\fontsize{7.5bp}{9.5bp}\selectfont 析：\allowbreak{}男生\ensuremath{155}、\allowbreak{}\ensuremath{150}厘米多备；\allowbreak{}女生\ensuremath{150}、\allowbreak{}\ensuremath{145}厘米多备；\allowbreak{}其他尺码也适量准备。\allowbreak{}\par};
\end{scope}\end{tikzpicture}\newpage
\definecolor{answer}{rgb}{0.850980,0.047059,0.070588}
\null\begin{tikzpicture}[remember picture,overlay]\begin{scope}[shift={(current page.north west)},x=1bp,y=-1bp]
\node[anchor=north west,inner sep=0pt,text=answer,align=left,text width=233.42857360839844bp] at (15,15) {\fontsize{8.57bp}{10.71bp}\selectfont 四（\ensuremath{2})班更好：\allowbreak{}\ensuremath{80}分及以上人数更多，\allowbreak{}\ensuremath{69}分及以下人数更少。\allowbreak{}\par};
\end{scope}\end{tikzpicture}\newpage
\definecolor{answer}{rgb}{0.850980,0.047059,0.070588}
\null\begin{tikzpicture}[remember picture,overlay]\begin{scope}[shift={(current page.north west)},x=1bp,y=-1bp]
\node[anchor=north west,inner sep=0pt,text=answer,align=left,text width=233.0bp] at (15,15) {\fontsize{7.5bp}{9.5bp}\selectfont 析：\allowbreak{}四(\ensuremath{2})班更好：\allowbreak{}≥\ensuremath{80}分\ensuremath{7}人对\ensuremath{5}人，\allowbreak{}≤\ensuremath{69}分\ensuremath{2}人对\ensuremath{4}人。\allowbreak{}\par};
\end{scope}\end{tikzpicture}\newpage
\definecolor{answer}{rgb}{0.850980,0.047059,0.070588}
\null\begin{tikzpicture}[remember picture,overlay]\begin{scope}[shift={(current page.north west)},x=1bp,y=-1bp]
\node[anchor=north west,inner sep=0pt,text=answer,align=left,text width=242.0bp] at (15,15) {\fontsize{8.5bp}{10.62bp}\selectfont 每人至少借一种书，\allowbreak{}\ensuremath{45} + \ensuremath{42}包含全部\ensuremath{78}人并重复计算两种都借者。\allowbreak{}\par};
\end{scope}\end{tikzpicture}\newpage
\definecolor{answer}{rgb}{0.850980,0.047059,0.070588}
\null\begin{tikzpicture}[remember picture,overlay]\begin{scope}[shift={(current page.north west)},x=1bp,y=-1bp]
\node[anchor=north west,inner sep=0pt,text=answer,align=left,text width=241.999755859375bp] at (15,15) {\fontsize{7.5bp}{9.5bp}\selectfont 析：\allowbreak{}相交处\ensuremath{10}人，\allowbreak{}外圈为全班\ensuremath{40}人；\allowbreak{}至少喜欢一科\ensuremath{25} + \ensuremath{20} \ensuremath{-}\ensuremath{10} =\ensuremath{35}人，\allowbreak{}两科都不喜欢\ensuremath{40} \ensuremath{-}\ensuremath{35} = \ensuremath{5}人，\allowbreak{}选A。\allowbreak{}\par};
\end{scope}\end{tikzpicture}\newpage
\definecolor{answer}{rgb}{0.850980,0.047059,0.070588}
\null\begin{tikzpicture}[remember picture,overlay]\begin{scope}[shift={(current page.north west)},x=1bp,y=-1bp]
\node[anchor=north west,inner sep=0pt,text=answer,align=left,text width=233.0bp] at (15,15) {\fontsize{8.5bp}{10.62bp}\selectfont 两项满分名单相加会重复计\ensuremath{3}人。\allowbreak{}\ensuremath{45}是全班受测人数，\allowbreak{}不能当作满分名单的合计人数。\allowbreak{}\par};
\end{scope}\end{tikzpicture}\newpage
\definecolor{answer}{rgb}{0.850980,0.047059,0.070588}
\null\begin{tikzpicture}[remember picture,overlay]\begin{scope}[shift={(current page.north west)},x=1bp,y=-1bp]
\node[anchor=north west,inner sep=0pt,text=answer,align=left,text width=232.99999237060547bp] at (15,15) {\fontsize{8.5bp}{10.62bp}\selectfont 答：\allowbreak{}两项检测中获得满分的共\ensuremath{13}人。\allowbreak{}\par};
\end{scope}\end{tikzpicture}\newpage
\definecolor{answer}{rgb}{0.850980,0.047059,0.070588}
\null\begin{tikzpicture}[remember picture,overlay]\begin{scope}[shift={(current page.north west)},x=1bp,y=-1bp]
\node[anchor=north west,inner sep=0pt,text=answer,align=left,text width=242.0bp] at (15,15) {\fontsize{8.5bp}{10.62bp}\selectfont 把四年级学生分为只会书法、\allowbreak{}只会绘画、\allowbreak{}两项都不会、\allowbreak{}两项都会四类，\allowbreak{}求最后一类。\allowbreak{}\par};
\end{scope}\end{tikzpicture}\newpage
\definecolor{answer}{rgb}{0.850980,0.047059,0.070588}
\null\begin{tikzpicture}[remember picture,overlay]\begin{scope}[shift={(current page.north west)},x=1bp,y=-1bp]
\node[anchor=north west,inner sep=0pt,text=answer,align=left,text width=233.0bp] at (15,15) {\fontsize{8.5bp}{10.62bp}\selectfont 两项人数相加时，\allowbreak{}都学的人被算两次；\allowbreak{}减去至少学一项的\ensuremath{48}人即重复部分。\allowbreak{}\par};
\end{scope}\end{tikzpicture}\newpage
\definecolor{answer}{rgb}{0.850980,0.047059,0.070588}
\null\begin{tikzpicture}[remember picture,overlay]\begin{scope}[shift={(current page.north west)},x=1bp,y=-1bp]
\node[anchor=north west,inner sep=0pt,text=answer,align=left,text width=507.0bp] at (15,15) {\fontsize{7.5bp}{9.5bp}\selectfont 析：\allowbreak{}把\ensuremath{36}看成\ensuremath{36} × \ensuremath{1} :\allowbreak{} \ensuremath{36} × (\ensuremath{99} + \ensuremath{1}) = \ensuremath{3600},\allowbreak{}用乘法分配律；\allowbreak{}也可先乘\ensuremath{3564}再加\ensuremath{36}。\allowbreak{}\par};
\end{scope}\end{tikzpicture}\newpage
\definecolor{answer}{rgb}{0.850980,0.047059,0.070588}
\null\begin{tikzpicture}[remember picture,overlay]\begin{scope}[shift={(current page.north west)},x=1bp,y=-1bp]
\node[anchor=north west,inner sep=0pt,text=answer,align=left,text width=507.0bp] at (15,15) {\fontsize{7.5bp}{9.5bp}\selectfont 析：\allowbreak{}被减数=差+减数，\allowbreak{}\ensuremath{6.34} + \ensuremath{0.27} = \ensuremath{6.61}。\allowbreak{}\par};
\end{scope}\end{tikzpicture}\newpage
\definecolor{answer}{rgb}{0.850980,0.047059,0.070588}
\null\begin{tikzpicture}[remember picture,overlay]\begin{scope}[shift={(current page.north west)},x=1bp,y=-1bp]
\node[anchor=north west,inner sep=0pt,text=answer,align=left,text width=507.0bp] at (15,15) {\fontsize{7.5bp}{9.5bp}\selectfont 析：\allowbreak{}乙\ensuremath{16.9} \ensuremath{-}\ensuremath{7.26} = \ensuremath{9.64};\allowbreak{}甲乙合计\ensuremath{16.9} + \ensuremath{9.64} = \ensuremath{26.54}。\allowbreak{}\par};
\end{scope}\end{tikzpicture}\newpage
\definecolor{answer}{rgb}{0.850980,0.047059,0.070588}
\null\begin{tikzpicture}[remember picture,overlay]\begin{scope}[shift={(current page.north west)},x=1bp,y=-1bp]
\node[anchor=north west,inner sep=0pt,text=answer,align=left,text width=507.0bp] at (15,15) {\fontsize{7.5bp}{9.5bp}\selectfont 析：\allowbreak{}相邻数增加\ensuremath{1.25};\allowbreak{}前空\ensuremath{4.21} \ensuremath{-}\ensuremath{1.25} = \ensuremath{2.96},\allowbreak{}后空\ensuremath{6.71} + \ensuremath{1.25} = \ensuremath{7.96}。\allowbreak{}\par};
\end{scope}\end{tikzpicture}\newpage
\definecolor{answer}{rgb}{0.850980,0.047059,0.070588}
\null\begin{tikzpicture}[remember picture,overlay]\begin{scope}[shift={(current page.north west)},x=1bp,y=-1bp]
\node[anchor=north west,inner sep=0pt,text=answer,align=left,text width=507.0bp] at (15,15) {\fontsize{7.5bp}{9.5bp}\selectfont 析：\allowbreak{}括号内先乘再减，\allowbreak{}最后除以\ensuremath{49},\allowbreak{}填乘、\allowbreak{}除。\allowbreak{}\par};
\end{scope}\end{tikzpicture}\newpage
\definecolor{answer}{rgb}{0.850980,0.047059,0.070588}
\null\begin{tikzpicture}[remember picture,overlay]\begin{scope}[shift={(current page.north west)},x=1bp,y=-1bp]
\node[anchor=north west,inner sep=0pt,text=answer,align=left,text width=507.0bp] at (15,15) {\fontsize{7.5bp}{9.5bp}\selectfont 析：\allowbreak{}保留两位看千分位\ensuremath{4},\allowbreak{}得\ensuremath{9.96};\allowbreak{}十分位看百分位\ensuremath{6}进位，\allowbreak{}得\ensuremath{10.0},\allowbreak{}末尾\ensuremath{0}保留精度。\allowbreak{}\par};
\end{scope}\end{tikzpicture}\newpage
\definecolor{answer}{rgb}{0.850980,0.047059,0.070588}
\null\begin{tikzpicture}[remember picture,overlay]\begin{scope}[shift={(current page.north west)},x=1bp,y=-1bp]
\node[anchor=north west,inner sep=0pt,text=answer,align=left,text width=507.0bp] at (15,15) {\fontsize{7.5bp}{9.5bp}\selectfont 析：\allowbreak{}法\ensuremath{1}还原被减数\ensuremath{56} + \ensuremath{2.1} = \ensuremath{58.1},\allowbreak{}再减\ensuremath{5.6}得\ensuremath{52.5};\allowbreak{}法\ensuremath{2}补回多减的\ensuremath{50.4},\allowbreak{}也得\ensuremath{52.5}。\allowbreak{}\par};
\end{scope}\end{tikzpicture}\newpage
\definecolor{answer}{rgb}{0.850980,0.047059,0.070588}
\null\begin{tikzpicture}[remember picture,overlay]\begin{scope}[shift={(current page.north west)},x=1bp,y=-1bp]
\node[anchor=north west,inner sep=0pt,text=answer,align=left,text width=507.0bp] at (15,15) {\fontsize{7.5bp}{9.5bp}\selectfont 析：\allowbreak{}依次为(\ensuremath{3000} ÷ \ensuremath{50} \ensuremath{-}\ensuremath{22}) × \ensuremath{2};\allowbreak{}\ensuremath{3000} ÷ (\ensuremath{50} \ensuremath{-}\ensuremath{22} × \ensuremath{2});\allowbreak{} \ensuremath{3000} ÷ [(\ensuremath{25} \ensuremath{-}\ensuremath{22}) × \ensuremath{8}]。\allowbreak{}\par};
\end{scope}\end{tikzpicture}\newpage
\definecolor{answer}{rgb}{0.850980,0.047059,0.070588}
\null\begin{tikzpicture}[remember picture,overlay]\begin{scope}[shift={(current page.north west)},x=1bp,y=-1bp]
\node[anchor=north west,inner sep=0pt,text=answer,align=left,text width=507.125bp] at (15,15) {\fontsize{7.5bp}{9.5bp}\selectfont 析：\allowbreak{}×。\allowbreak{}\ensuremath{168} \ensuremath{-}\ensuremath{48} + \ensuremath{52} = \ensuremath{168} + (\ensuremath{52} \ensuremath{-}\ensuremath{48}) = \ensuremath{168} + \ensuremath{4} = \ensuremath{172}；\allowbreak{}右式\ensuremath{168} \ensuremath{-}(\ensuremath{48} + \ensuremath{52}) = \ensuremath{68}，\allowbreak{}所以不相等\par};
\end{scope}\end{tikzpicture}\newpage
\definecolor{answer}{rgb}{0.850980,0.047059,0.070588}
\null\begin{tikzpicture}[remember picture,overlay]\begin{scope}[shift={(current page.north west)},x=1bp,y=-1bp]
\node[anchor=north west,inner sep=0pt,text=answer,align=left,text width=507.0bp] at (15,15) {\fontsize{7.5bp}{9.5bp}\selectfont 析：\allowbreak{}×。\allowbreak{}乘法交换、\allowbreak{}结合不改变积，\allowbreak{}两式相等。\allowbreak{}\par};
\end{scope}\end{tikzpicture}\newpage
\definecolor{answer}{rgb}{0.850980,0.047059,0.070588}
\null\begin{tikzpicture}[remember picture,overlay]\begin{scope}[shift={(current page.north west)},x=1bp,y=-1bp]
\node[anchor=north west,inner sep=0pt,text=answer,align=left,text width=507.0bp] at (15,15) {\fontsize{7.5bp}{9.5bp}\selectfont 析：\allowbreak{}×。\allowbreak{}左\ensuremath{3}右\ensuremath{2}合计左移\ensuremath{1}位，\allowbreak{}缩小到原来的十分之一。\allowbreak{}\par};
\end{scope}\end{tikzpicture}\newpage
\definecolor{answer}{rgb}{0.850980,0.047059,0.070588}
\null\begin{tikzpicture}[remember picture,overlay]\begin{scope}[shift={(current page.north west)},x=1bp,y=-1bp]
\node[anchor=north west,inner sep=0pt,text=answer,align=left,text width=507.0bp] at (15,15) {\fontsize{7.5bp}{9.5bp}\selectfont 析：\allowbreak{}√。\allowbreak{}连续除以\ensuremath{3}、\allowbreak{}\ensuremath{80}等于除以它们的积。\allowbreak{}\par};
\end{scope}\end{tikzpicture}\newpage
\definecolor{answer}{rgb}{0.850980,0.047059,0.070588}
\null\begin{tikzpicture}[remember picture,overlay]\begin{scope}[shift={(current page.north west)},x=1bp,y=-1bp]
\node[anchor=north west,inner sep=0pt,text=answer,align=left,text width=507.0bp] at (15,15) {\fontsize{7.5bp}{9.5bp}\selectfont 析：\allowbreak{}×。\allowbreak{}先算\ensuremath{56} ÷ \ensuremath{56} = \ensuremath{1},\allowbreak{}再算\ensuremath{4} + \ensuremath{1} + \ensuremath{4} = \ensuremath{9}。\allowbreak{}\par};
\end{scope}\end{tikzpicture}\newpage
\definecolor{answer}{rgb}{0.850980,0.047059,0.070588}
\null\begin{tikzpicture}[remember picture,overlay]\begin{scope}[shift={(current page.north west)},x=1bp,y=-1bp]
\node[anchor=north west,inner sep=0pt,text=answer,align=left,text width=507.0bp] at (15,15) {\fontsize{7.5bp}{9.5bp}\selectfont 析：\allowbreak{}C只交换两个因数的位置：\allowbreak{}\ensuremath{a} × \ensuremath{c} = \ensuremath{c} × \ensuremath{a},\allowbreak{}是乘法交换律。\allowbreak{}\par};
\end{scope}\end{tikzpicture}\newpage
\definecolor{answer}{rgb}{0.850980,0.047059,0.070588}
\null\begin{tikzpicture}[remember picture,overlay]\begin{scope}[shift={(current page.north west)},x=1bp,y=-1bp]
\node[anchor=north west,inner sep=0pt,text=answer,align=left,text width=489bp] at (15,15) {\fontsize{7.5bp}{9.5bp}\selectfont 析：\allowbreak{}\ensuremath{4}个十=\ensuremath{40},\allowbreak{} \ensuremath{6}个百分之一=\ensuremath{0.06},\allowbreak{} \ensuremath{7}个千分之一=\ensuremath{0.007},\allowbreak{}合\ensuremath{40.067},\allowbreak{}选B。\allowbreak{}\par};
\end{scope}\end{tikzpicture}\newpage
\definecolor{answer}{rgb}{0.850980,0.047059,0.070588}
\null\begin{tikzpicture}[remember picture,overlay]\begin{scope}[shift={(current page.north west)},x=1bp,y=-1bp]
\node[anchor=north west,inner sep=0pt,text=answer,align=left,text width=489bp] at (15,15) {\fontsize{7.5bp}{9.5bp}\selectfont 析：\allowbreak{}仅小数末尾\ensuremath{0}可去，\allowbreak{}\ensuremath{9.10} = \ensuremath{9.1};\allowbreak{}其他\ensuremath{0}去掉会改值，\allowbreak{}选A。\allowbreak{}\par};
\end{scope}\end{tikzpicture}\newpage
\definecolor{answer}{rgb}{0.850980,0.047059,0.070588}
\null\begin{tikzpicture}[remember picture,overlay]\begin{scope}[shift={(current page.north west)},x=1bp,y=-1bp]
\node[anchor=north west,inner sep=0pt,text=answer,align=left,text width=489bp] at (15,15) {\fontsize{7.5bp}{9.5bp}\selectfont 析：\allowbreak{}先求\ensuremath{168} + \ensuremath{84},\allowbreak{}再用\ensuremath{336}减这个和，\allowbreak{}最后除\ensuremath{21} :\allowbreak{}选B。\allowbreak{}\par};
\end{scope}\end{tikzpicture}\newpage
\definecolor{answer}{rgb}{0.850980,0.047059,0.070588}
\null\begin{tikzpicture}[remember picture,overlay]\begin{scope}[shift={(current page.north west)},x=1bp,y=-1bp]
\node[anchor=north west,inner sep=0pt,text=answer,align=left,text width=237.0bp] at (15,15) {\fontsize{7.5bp}{9.5bp}\selectfont 把\ensuremath{4.59}与\ensuremath{5.41}、\allowbreak{}\ensuremath{36.43}与\ensuremath{13.57}配对，\allowbreak{}分别凑成\ensuremath{10}、\allowbreak{}\ensuremath{50}。\allowbreak{}\par};
\end{scope}\end{tikzpicture}\newpage
\definecolor{answer}{rgb}{0.850980,0.047059,0.070588}
\null\begin{tikzpicture}[remember picture,overlay]\begin{scope}[shift={(current page.north west)},x=1bp,y=-1bp]
\node[anchor=north west,inner sep=0pt,text=answer,align=left,text width=241.0bp] at (15,15) {\fontsize{7.5bp}{9.5bp}\selectfont 先算\ensuremath{645} \ensuremath{-}\ensuremath{145}，\allowbreak{}再算\ensuremath{180} + \ensuremath{220}，\allowbreak{}分别凑成\ensuremath{500}、\allowbreak{}\ensuremath{400}。\allowbreak{}\par};
\end{scope}\end{tikzpicture}\newpage
\definecolor{answer}{rgb}{0.850980,0.047059,0.070588}
\null\begin{tikzpicture}[remember picture,overlay]\begin{scope}[shift={(current page.north west)},x=1bp,y=-1bp]
\node[anchor=north west,inner sep=0pt,text=answer,align=left,text width=237.0bp] at (15,15) {\fontsize{7.5bp}{9.5bp}\selectfont 把\ensuremath{99}看成\ensuremath{100} \ensuremath{-}\ensuremath{1}，\allowbreak{}用乘法分配律，\allowbreak{}乘\ensuremath{100}后减去一份\ensuremath{43}。\allowbreak{}\par};
\end{scope}\end{tikzpicture}\newpage
\definecolor{answer}{rgb}{0.850980,0.047059,0.070588}
\null\begin{tikzpicture}[remember picture,overlay]\begin{scope}[shift={(current page.north west)},x=1bp,y=-1bp]
\node[anchor=north west,inner sep=0pt,text=answer,align=left,text width=241.0bp] at (15,15) {\fontsize{7.5bp}{9.5bp}\selectfont 交换并结合因数，\allowbreak{}把\ensuremath{25}与\ensuremath{4}配成\ensuremath{100}，\allowbreak{}把\ensuremath{15}与\ensuremath{6}配成\ensuremath{90}。\allowbreak{}\par};
\end{scope}\end{tikzpicture}\newpage
\definecolor{answer}{rgb}{0.850980,0.047059,0.070588}
\null\begin{tikzpicture}[remember picture,overlay]\begin{scope}[shift={(current page.north west)},x=1bp,y=-1bp]
\node[anchor=north west,inner sep=0pt,text=answer,align=left,text width=237.0bp] at (15,15) {\fontsize{7.5bp}{9.5bp}\selectfont 先交换\ensuremath{56} × \ensuremath{135}的因数，\allowbreak{}把末尾\ensuremath{135}看成\ensuremath{135} × \ensuremath{1}，\allowbreak{}再用乘法分配律合并成\ensuremath{135} × \ensuremath{100}。\allowbreak{}\par};
\end{scope}\end{tikzpicture}\newpage
\definecolor{answer}{rgb}{0.850980,0.047059,0.070588}
\null\begin{tikzpicture}[remember picture,overlay]\begin{scope}[shift={(current page.north west)},x=1bp,y=-1bp]
\node[anchor=north west,inner sep=0pt,text=answer,align=left,text width=457.0bp] at (15,15) {\fontsize{8.5bp}{10.62bp}\selectfont 积是\ensuremath{31} × \ensuremath{40}，\allowbreak{}商是\ensuremath{4800} ÷ \ensuremath{60}；\allowbreak{}先算乘、\allowbreak{}除，\allowbreak{}再把积与商相加。\allowbreak{}\par};
\end{scope}\end{tikzpicture}\newpage
\definecolor{answer}{rgb}{0.850980,0.047059,0.070588}
\null\begin{tikzpicture}[remember picture,overlay]\begin{scope}[shift={(current page.north west)},x=1bp,y=-1bp]
\node[anchor=north west,inner sep=0pt,text=answer,align=left,text width=497.0bp] at (15,15) {\fontsize{8.5bp}{10.62bp}\selectfont 两人各工作\ensuremath{6}小时，\allowbreak{}可把每小时加工数相加后乘\ensuremath{6},\allowbreak{}也可分别求总量。\allowbreak{}\par};
\end{scope}\end{tikzpicture}\newpage
\definecolor{answer}{rgb}{0.850980,0.047059,0.070588}
\null\begin{tikzpicture}[remember picture,overlay]\begin{scope}[shift={(current page.north west)},x=1bp,y=-1bp]
\node[anchor=north west,inner sep=0pt,text=answer,align=left,text width=497.0bp] at (15,15) {\fontsize{8.5bp}{10.62bp}\selectfont 李叔叔每小时比王叔叔少\ensuremath{2}个，\allowbreak{}两人工作时间相同，\allowbreak{}用效率差乘\ensuremath{6}小时。\allowbreak{}\par};
\end{scope}\end{tikzpicture}\newpage
\definecolor{answer}{rgb}{0.850980,0.047059,0.070588}
\null\begin{tikzpicture}[remember picture,overlay]\begin{scope}[shift={(current page.north west)},x=1bp,y=-1bp]
\node[anchor=north west,inner sep=0pt,text=answer,align=left,text width=497.0bp] at (15,15) {\fontsize{8.5bp}{10.62bp}\selectfont 上山和下山用时不同，\allowbreak{}必须分别求路程；\allowbreak{}也可把共同的\ensuremath{8}分钟与缆车多用的\ensuremath{2}分钟分开比较。\allowbreak{}\par};
\end{scope}\end{tikzpicture}\newpage
\definecolor{answer}{rgb}{0.850980,0.047059,0.070588}
\null\begin{tikzpicture}[remember picture,overlay]\begin{scope}[shift={(current page.north west)},x=1bp,y=-1bp]
\node[anchor=north west,inner sep=0pt,text=answer,align=left,text width=497.0bp] at (15,15) {\fontsize{8.5bp}{10.62bp}\selectfont 编程组是科技组的\ensuremath{4}倍多\ensuremath{5}人，\allowbreak{}先求\ensuremath{77}人，\allowbreak{}再和创客组\ensuremath{90}人比较。\allowbreak{}\par};
\end{scope}\end{tikzpicture}\newpage
\definecolor{answer}{rgb}{0.850980,0.047059,0.070588}
\null\begin{tikzpicture}[remember picture,overlay]\begin{scope}[shift={(current page.north west)},x=1bp,y=-1bp]
\node[anchor=north west,inner sep=0pt,text=answer,align=left,text width=507.0bp] at (15,15) {\fontsize{7.5bp}{9.5bp}\selectfont 析：\allowbreak{}求几箱是数量；\allowbreak{}每箱多少钱是单价；\allowbreak{}一共花多少钱是总价。\allowbreak{}\par};
\end{scope}\end{tikzpicture}\newpage
\definecolor{answer}{rgb}{0.850980,0.047059,0.070588}
\null\begin{tikzpicture}[remember picture,overlay]\begin{scope}[shift={(current page.north west)},x=1bp,y=-1bp]
\node[anchor=north west,inner sep=0pt,text=answer,align=left,text width=507.0bp] at (15,15) {\fontsize{7.5bp}{9.5bp}\selectfont 析：\allowbreak{}相背而行，\allowbreak{}已知速度、\allowbreak{}时间求路程：\allowbreak{}(\ensuremath{65} + \ensuremath{55}) × \ensuremath{3} = \ensuremath{360}米。\allowbreak{}\par};
\end{scope}\end{tikzpicture}\newpage
\definecolor{answer}{rgb}{0.850980,0.047059,0.070588}
\null\begin{tikzpicture}[remember picture,overlay]\begin{scope}[shift={(current page.north west)},x=1bp,y=-1bp]
\node[anchor=north west,inner sep=0pt,text=answer,align=left,text width=507.0bp] at (15,15) {\fontsize{7.5bp}{9.5bp}\selectfont 析：\allowbreak{}单价写\ensuremath{12}元/千克；\allowbreak{}总价=单价×数量，\allowbreak{}\ensuremath{5}千克需\ensuremath{12} × \ensuremath{5} = \ensuremath{60}元。\allowbreak{}\par};
\end{scope}\end{tikzpicture}\newpage
\definecolor{answer}{rgb}{0.850980,0.047059,0.070588}
\null\begin{tikzpicture}[remember picture,overlay]\begin{scope}[shift={(current page.north west)},x=1bp,y=-1bp]
\node[anchor=north west,inner sep=0pt,text=answer,align=left,text width=507.0bp] at (15,15) {\fontsize{7.5bp}{9.5bp}\selectfont 析：\allowbreak{}速度=路程÷时间，\allowbreak{}\ensuremath{3600} ÷ \ensuremath{40} = \ensuremath{90}千米/时。\allowbreak{}\par};
\end{scope}\end{tikzpicture}\newpage
\definecolor{answer}{rgb}{0.850980,0.047059,0.070588}
\null\begin{tikzpicture}[remember picture,overlay]\begin{scope}[shift={(current page.north west)},x=1bp,y=-1bp]
\node[anchor=north west,inner sep=0pt,text=answer,align=left,text width=507.0bp] at (15,15) {\fontsize{7.5bp}{9.5bp}\selectfont 析：\allowbreak{}单价\ensuremath{9}元/支，\allowbreak{}数量\ensuremath{3}支，\allowbreak{}总价\ensuremath{9} × \ensuremath{3} = \ensuremath{27}元。\allowbreak{}\par};
\end{scope}\end{tikzpicture}\newpage
\definecolor{answer}{rgb}{0.850980,0.047059,0.070588}
\null\begin{tikzpicture}[remember picture,overlay]\begin{scope}[shift={(current page.north west)},x=1bp,y=-1bp]
\node[anchor=north west,inner sep=0pt,text=answer,align=left,text width=507.0bp] at (15,15) {\fontsize{7.5bp}{9.5bp}\selectfont 析：\allowbreak{}行程表按路程=速度×时间，\allowbreak{}得\ensuremath{150}万千米、\allowbreak{}\ensuremath{2}分、\allowbreak{}\ensuremath{150}千米/时；\allowbreak{}购物表按总价=单价×数量，\allowbreak{}得\ensuremath{1800}、\allowbreak{}\ensuremath{320}、\allowbreak{}\ensuremath{8}。\allowbreak{}\par};
\end{scope}\end{tikzpicture}\newpage
\definecolor{answer}{rgb}{0.850980,0.047059,0.070588}
\null\begin{tikzpicture}[remember picture,overlay]\begin{scope}[shift={(current page.north west)},x=1bp,y=-1bp]
\node[anchor=north west,inner sep=0pt,text=answer,align=left,text width=507.0bp] at (15,15) {\fontsize{7.5bp}{9.5bp}\selectfont 析：\allowbreak{}法\ensuremath{1}(\ensuremath{420} \ensuremath{-}\ensuremath{80} × \ensuremath{3}) ÷ \ensuremath{2} = \ensuremath{90};\allowbreak{}法\ensuremath{2}全按\ensuremath{80}共\ensuremath{400}千米，\allowbreak{}后\ensuremath{2}时多\ensuremath{20},\allowbreak{}每时多\ensuremath{10},\allowbreak{}速度\ensuremath{90}千米/时。\allowbreak{}\par};
\end{scope}\end{tikzpicture}\newpage
\definecolor{answer}{rgb}{0.850980,0.047059,0.070588}
\null\begin{tikzpicture}[remember picture,overlay]\begin{scope}[shift={(current page.north west)},x=1bp,y=-1bp]
\node[anchor=north west,inner sep=0pt,text=answer,align=left,text width=507.0bp] at (15,15) {\fontsize{7.5bp}{9.5bp}\selectfont 析：\allowbreak{}往返路程相同；\allowbreak{}返回用\ensuremath{5}时，\allowbreak{}\ensuremath{90} × \ensuremath{4} ÷ \ensuremath{5} = \ensuremath{72}千米/时，\allowbreak{}选B。\allowbreak{}\par};
\end{scope}\end{tikzpicture}\newpage
\definecolor{answer}{rgb}{0.850980,0.047059,0.070588}
\null\begin{tikzpicture}[remember picture,overlay]\begin{scope}[shift={(current page.north west)},x=1bp,y=-1bp]
\node[anchor=north west,inner sep=0pt,text=answer,align=left,text width=507.0bp] at (15,15) {\fontsize{7.5bp}{9.5bp}\selectfont 析：\allowbreak{}B已知单价\ensuremath{98}元和数量\ensuremath{3}套，\allowbreak{}求总价；\allowbreak{}A求路程，\allowbreak{}C求单价，\allowbreak{}选B。\allowbreak{}\par};
\end{scope}\end{tikzpicture}\newpage
\definecolor{answer}{rgb}{0.850980,0.047059,0.070588}
\null\begin{tikzpicture}[remember picture,overlay]\begin{scope}[shift={(current page.north west)},x=1bp,y=-1bp]
\node[anchor=north west,inner sep=0pt,text=answer,align=left,text width=507.0bp] at (15,15) {\fontsize{7.5bp}{9.5bp}\selectfont 析：\allowbreak{}图中男生=女生一半+\ensuremath{6}人，\allowbreak{}右括号\ensuremath{33}是男女合计；\allowbreak{}②④正确，\allowbreak{}选B。\allowbreak{}\par};
\end{scope}\end{tikzpicture}\newpage
\definecolor{answer}{rgb}{0.850980,0.047059,0.070588}
\null\begin{tikzpicture}[remember picture,overlay]\begin{scope}[shift={(current page.north west)},x=1bp,y=-1bp]
\node[anchor=north west,inner sep=0pt,text=answer,align=left,text width=507.0bp] at (15,15) {\fontsize{7.5bp}{9.5bp}\selectfont 析：\allowbreak{}速度\ensuremath{40} ÷ \ensuremath{8}米/分；\allowbreak{}时间\ensuremath{75} ÷ (\ensuremath{40} ÷ \ensuremath{8}),\allowbreak{}选C。\allowbreak{}\par};
\end{scope}\end{tikzpicture}\newpage
\definecolor{answer}{rgb}{0.850980,0.047059,0.070588}
\null\begin{tikzpicture}[remember picture,overlay]\begin{scope}[shift={(current page.north west)},x=1bp,y=-1bp]
\node[anchor=north west,inner sep=0pt,text=answer,align=left,text width=507.0bp] at (15,15) {\fontsize{7.5bp}{9.5bp}\selectfont 析：\allowbreak{}问用去多少，\allowbreak{}不问剩余：\allowbreak{}\ensuremath{2.6} + \ensuremath{5.4} = \ensuremath{8}米，\allowbreak{}选B。\allowbreak{}\par};
\end{scope}\end{tikzpicture}\newpage
\definecolor{answer}{rgb}{0.850980,0.047059,0.070588}
\null\begin{tikzpicture}[remember picture,overlay]\begin{scope}[shift={(current page.north west)},x=1bp,y=-1bp]
\node[anchor=north west,inner sep=0pt,text=answer,align=left,text width=507.0bp] at (15,15) {\fontsize{7.5bp}{9.5bp}\selectfont 析：\allowbreak{}×。\allowbreak{}数量=总价÷单价，\allowbreak{}不是相乘。\allowbreak{}\par};
\end{scope}\end{tikzpicture}\newpage
\definecolor{answer}{rgb}{0.850980,0.047059,0.070588}
\null\begin{tikzpicture}[remember picture,overlay]\begin{scope}[shift={(current page.north west)},x=1bp,y=-1bp]
\node[anchor=north west,inner sep=0pt,text=answer,align=left,text width=507.0bp] at (15,15) {\fontsize{7.5bp}{9.5bp}\selectfont 析：\allowbreak{}×。\allowbreak{}先扣球拍\ensuremath{35}元，\allowbreak{}再平均给\ensuremath{5}个毽子：\allowbreak{}(\ensuremath{60} \ensuremath{-}\ensuremath{35}) ÷ \ensuremath{5} = \ensuremath{5}元。\allowbreak{}\par};
\end{scope}\end{tikzpicture}\newpage
\definecolor{answer}{rgb}{0.850980,0.047059,0.070588}
\null\begin{tikzpicture}[remember picture,overlay]\begin{scope}[shift={(current page.north west)},x=1bp,y=-1bp]
\node[anchor=north west,inner sep=0pt,text=answer,align=left,text width=507.0bp] at (15,15) {\fontsize{7.5bp}{9.5bp}\selectfont 析：\allowbreak{}×。\allowbreak{}大豆比玉米多\ensuremath{0.4}吨，\allowbreak{}玉米应\ensuremath{2.34} \ensuremath{-}\ensuremath{0.4} = \ensuremath{1.94}吨。\allowbreak{}\par};
\end{scope}\end{tikzpicture}\newpage
\definecolor{answer}{rgb}{0.850980,0.047059,0.070588}
\null\begin{tikzpicture}[remember picture,overlay]\begin{scope}[shift={(current page.north west)},x=1bp,y=-1bp]
\node[anchor=north west,inner sep=0pt,text=answer,align=left,text width=497.125bp] at (15,15) {\fontsize{7.5bp}{9.5bp}\selectfont 析：\allowbreak{}两站路程\ensuremath{275} × \ensuremath{5} + \ensuremath{250} × \ensuremath{5} = (\ensuremath{275} + \ensuremath{250}) × \ensuremath{5} = \ensuremath{525} × \ensuremath{5} = \ensuremath{2625}千米；\allowbreak{}多行路程(\ensuremath{275} \ensuremath{-}\ensuremath{250}) × \ensuremath{5}。\allowbreak{}连线不变\par};
\end{scope}\end{tikzpicture}\newpage
\definecolor{answer}{rgb}{0.850980,0.047059,0.070588}
\null\begin{tikzpicture}[remember picture,overlay]\begin{scope}[shift={(current page.north west)},x=1bp,y=-1bp]
\node[anchor=north west,inner sep=0pt,text=answer,align=left,text width=497.0bp] at (15,15) {\fontsize{8.5bp}{10.62bp}\selectfont 同走\ensuremath{1105}千米，\allowbreak{}轿车速度快，\allowbreak{}先分别求时间；\allowbreak{}也可求轿车到达时客车还差多少路程。\allowbreak{}\par};
\end{scope}\end{tikzpicture}\newpage
\definecolor{answer}{rgb}{0.850980,0.047059,0.070588}
\null\begin{tikzpicture}[remember picture,overlay]\begin{scope}[shift={(current page.north west)},x=1bp,y=-1bp]
\node[anchor=north west,inner sep=0pt,text=answer,align=left,text width=497.0bp] at (15,15) {\fontsize{8.5bp}{10.62bp}\selectfont 扣掉\ensuremath{2}张桌子的\ensuremath{160}元后，\allowbreak{}椅子共\ensuremath{170}元，\allowbreak{}平均分给\ensuremath{5}把。\allowbreak{}\par};
\end{scope}\end{tikzpicture}\newpage
\definecolor{answer}{rgb}{0.850980,0.047059,0.070588}
\null\begin{tikzpicture}[remember picture,overlay]\begin{scope}[shift={(current page.north west)},x=1bp,y=-1bp]
\node[anchor=north west,inner sep=0pt,text=answer,align=left,text width=497.0bp] at (15,15) {\fontsize{8.5bp}{10.62bp}\selectfont \ensuremath{5}个足球共\ensuremath{400}元，\allowbreak{}余\ensuremath{258}元为\ensuremath{6}个排球总价；\allowbreak{}也可假设每个排球\ensuremath{40}元再补差。\allowbreak{}\par};
\end{scope}\end{tikzpicture}\newpage
\definecolor{answer}{rgb}{0.850980,0.047059,0.070588}
\null\begin{tikzpicture}[remember picture,overlay]\begin{scope}[shift={(current page.north west)},x=1bp,y=-1bp]
\node[anchor=north west,inner sep=0pt,text=answer,align=left,text width=497.0bp] at (15,15) {\fontsize{7.5bp}{9.5bp}\selectfont 兔子速度\ensuremath{16}米/秒，\allowbreak{}长颈鹿每秒少\ensuremath{1}米，\allowbreak{}为\ensuremath{15}米/秒；\allowbreak{}路程除以速度求时间。\allowbreak{}\par};
\end{scope}\end{tikzpicture}\newpage
\definecolor{answer}{rgb}{0.850980,0.047059,0.070588}
\null\begin{tikzpicture}[remember picture,overlay]\begin{scope}[shift={(current page.north west)},x=1bp,y=-1bp]
\node[anchor=north west,inner sep=0pt,text=answer,align=left,text width=497.0bp] at (15,15) {\fontsize{7.5bp}{9.5bp}\selectfont 析：\allowbreak{}数量=总价÷单价：\allowbreak{}\ensuremath{390} ÷ \ensuremath{13} = \ensuremath{30}、\allowbreak{}\ensuremath{216} ÷ \ensuremath{18} = \ensuremath{12}、\allowbreak{}\ensuremath{408} ÷ \ensuremath{12} = \ensuremath{34}、\allowbreak{}\ensuremath{435} ÷ \ensuremath{29} = \ensuremath{15}本。\allowbreak{}\par};
\end{scope}\end{tikzpicture}\newpage
\definecolor{answer}{rgb}{0.850980,0.047059,0.070588}
\null\begin{tikzpicture}[remember picture,overlay]\begin{scope}[shift={(current page.north west)},x=1bp,y=-1bp]
\node[anchor=north west,inner sep=0pt,text=answer,align=left,text width=497.0bp] at (15,15) {\fontsize{8.5bp}{10.62bp}\selectfont “畅销”比较卖出的本数，\allowbreak{}而非总价。\allowbreak{}四种书依次卖出\ensuremath{30}、\allowbreak{}\ensuremath{12}、\allowbreak{}\ensuremath{34}、\allowbreak{}\ensuremath{15}本，\allowbreak{}\ensuremath{34}本最多。\allowbreak{}\par};
\end{scope}\end{tikzpicture}\newpage
\definecolor{answer}{rgb}{0.850980,0.047059,0.070588}
\null\begin{tikzpicture}[remember picture,overlay]\begin{scope}[shift={(current page.north west)},x=1bp,y=-1bp]
\node[anchor=north west,inner sep=0pt,text=answer,align=left,text width=484.0bp] at (15,15) {\fontsize{8.5bp}{10.62bp}\selectfont 先把四种读物一天的销量相加，\allowbreak{}得到\ensuremath{91}本；\allowbreak{}估算一个月按\ensuremath{30}天计算。\allowbreak{}\par};
\end{scope}\end{tikzpicture}\newpage
\definecolor{answer}{rgb}{0.850980,0.047059,0.070588}
\null\begin{tikzpicture}[remember picture,overlay]\begin{scope}[shift={(current page.north west)},x=1bp,y=-1bp]
\node[anchor=north west,inner sep=0pt,text=answer,align=left,text width=484.0bp] at (15,15) {\fontsize{8.5bp}{10.62bp}\selectfont \ensuremath{29}元接近\ensuremath{30}元，\allowbreak{}先用\ensuremath{600}除以估计单价；\allowbreak{}再用相邻本数检验。\allowbreak{}\par};
\end{scope}\end{tikzpicture}\newpage
\definecolor{answer}{rgb}{0.850980,0.047059,0.070588}
\null\begin{tikzpicture}[remember picture,overlay]\begin{scope}[shift={(current page.north west)},x=1bp,y=-1bp]
\node[anchor=north west,inner sep=0pt,text=answer,align=left,text width=484.0bp] at (15,15) {\fontsize{8.5bp}{10.62bp}\selectfont 求两本书的总价，\allowbreak{}用两本书的单价相加。\allowbreak{}\par};
\end{scope}\end{tikzpicture}\newpage
\definecolor{answer}{rgb}{0.850980,0.047059,0.070588}
\null\begin{tikzpicture}[remember picture,overlay]\begin{scope}[shift={(current page.north west)},x=1bp,y=-1bp]
\node[anchor=north west,inner sep=0pt,text=answer,align=left,text width=478.0bp] at (15,15) {\fontsize{8.5bp}{10.62bp}\selectfont “贵多少”比较两本书的单价，\allowbreak{}用较大价格减较小价格。\allowbreak{}\par};
\end{scope}\end{tikzpicture}\newpage
\definecolor{answer}{rgb}{0.850980,0.047059,0.070588}
\null\begin{tikzpicture}[remember picture,overlay]\begin{scope}[shift={(current page.north west)},x=1bp,y=-1bp]
\node[anchor=north west,inner sep=0pt,text=answer,align=left,text width=478.0bp] at (15,15) {\fontsize{8.5bp}{10.62bp}\selectfont 三种书各买一本，\allowbreak{}总价是三个单价的和。\allowbreak{}\par};
\end{scope}\end{tikzpicture}\newpage
\definecolor{answer}{rgb}{0.850980,0.047059,0.070588}
\null\begin{tikzpicture}[remember picture,overlay]\begin{scope}[shift={(current page.north west)},x=1bp,y=-1bp]
\node[anchor=north west,inner sep=0pt,text=answer,align=left,text width=473.0bp] at (15,15) {\fontsize{8.5bp}{10.62bp}\selectfont 第二根比第三根长\ensuremath{0.99}米，\allowbreak{}比第一根长\ensuremath{2.14}米。\allowbreak{}先加求第二根，\allowbreak{}再减求第一根。\allowbreak{}\par};
\end{scope}\end{tikzpicture}\newpage
\definecolor{answer}{rgb}{0.850980,0.047059,0.070588}
\null\begin{tikzpicture}[remember picture,overlay]\begin{scope}[shift={(current page.north west)},x=1bp,y=-1bp]
\node[anchor=north west,inner sep=0pt,text=answer,align=left,text width=484.0bp] at (15,15) {\fontsize{8.5bp}{10.62bp}\selectfont 装小盒的彩塑总数是\ensuremath{200}减\ensuremath{88},\allowbreak{}再平均分到\ensuremath{28}个小盒。\allowbreak{}\par};
\end{scope}\end{tikzpicture}\newpage
\definecolor{answer}{rgb}{0.850980,0.047059,0.070588}
\null\begin{tikzpicture}[remember picture,overlay]\begin{scope}[shift={(current page.north west)},x=1bp,y=-1bp]
\node[anchor=north west,inner sep=0pt,text=answer,align=left,text width=484.0bp] at (15,15) {\fontsize{8.5bp}{10.62bp}\selectfont 乙每分钟比甲少爬\ensuremath{15}厘米；\allowbreak{}同向各爬一圈，\allowbreak{}比较爬一圈的时间。\allowbreak{}\par};
\end{scope}\end{tikzpicture}\newpage
\definecolor{answer}{rgb}{0.850980,0.047059,0.070588}
\null\begin{tikzpicture}[remember picture,overlay]\begin{scope}[shift={(current page.north west)},x=1bp,y=-1bp]
\node[anchor=north west,inner sep=0pt,text=answer,align=left,text width=484.0bp] at (15,15) {\fontsize{8.5bp}{10.62bp}\selectfont 相向而行时，\allowbreak{}两只蚂蚁行的路程之和达到一圈\ensuremath{900}厘米才相遇。\allowbreak{}\par};
\end{scope}\end{tikzpicture}\newpage
\definecolor{answer}{rgb}{0.850980,0.047059,0.070588}
\null\begin{tikzpicture}[remember picture,overlay]\begin{scope}[shift={(current page.north west)},x=1bp,y=-1bp]
\node[anchor=north west,inner sep=0pt,text=answer,align=left,text width=484.0bp] at (15,15) {\fontsize{7.5bp}{9.5bp}\selectfont 析：\allowbreak{}横杠方向相同且不相交，\allowbreak{}互相平行；\allowbreak{}横杠与竖杠相交成直角，\allowbreak{}互相垂直。\allowbreak{}\par};
\end{scope}\end{tikzpicture}\newpage
\definecolor{answer}{rgb}{0.850980,0.047059,0.070588}
\null\begin{tikzpicture}[remember picture,overlay]\begin{scope}[shift={(current page.north west)},x=1bp,y=-1bp]
\node[anchor=north west,inner sep=0pt,text=answer,align=left,text width=484.0bp] at (15,15) {\fontsize{7.5bp}{9.5bp}\selectfont 析：\allowbreak{}过直线外一点只能画\ensuremath{1}条垂线；\allowbreak{}能画无数条不同方向的相交直线。\allowbreak{}\par};
\end{scope}\end{tikzpicture}\newpage
\definecolor{answer}{rgb}{0.850980,0.047059,0.070588}
\null\begin{tikzpicture}[remember picture,overlay]\begin{scope}[shift={(current page.north west)},x=1bp,y=-1bp]
\node[anchor=north west,inner sep=0pt,text=answer,align=left,text width=483.9988708496094bp] at (15,15) {\fontsize{7.5bp}{9.5bp}\selectfont 析：\allowbreak{}同一平面内，\allowbreak{}两条直线都平行于同一条直线，\allowbreak{}方向相同且彼此不相交，\allowbreak{}所以互相平行。\allowbreak{}\par};
\end{scope}\end{tikzpicture}\newpage
\definecolor{answer}{rgb}{0.850980,0.047059,0.070588}
\null\begin{tikzpicture}[remember picture,overlay]\begin{scope}[shift={(current page.north west)},x=1bp,y=-1bp]
\node[anchor=north west,inner sep=0pt,text=answer,align=left,text width=484.0bp] at (15,15) {\fontsize{7.5bp}{9.5bp}\selectfont 析：\allowbreak{}大角等于另两角之和，\allowbreak{}占\ensuremath{180}\ensuremath{^\circ}的一半，\allowbreak{}即\ensuremath{90}\ensuremath{^\circ},\allowbreak{}所以是直角三角形。\allowbreak{}\par};
\end{scope}\end{tikzpicture}\newpage
\definecolor{answer}{rgb}{0.850980,0.047059,0.070588}
\null\begin{tikzpicture}[remember picture,overlay]\begin{scope}[shift={(current page.north west)},x=1bp,y=-1bp]
\node[anchor=north west,inner sep=0pt,text=answer,align=left,text width=484.0bp] at (15,15) {\fontsize{7.5bp}{8.5bp}\selectfont 析：\allowbreak{}图\ensuremath{1}两组对边平行；\allowbreak{}横边与竖边配成\ensuremath{4}组垂直。\allowbreak{}每条对角线有整条和两半段，\allowbreak{}可配成\ensuremath{9}组，\allowbreak{}共\ensuremath{13}组。\allowbreak{}图\ensuremath{2}每个方向的中间线段，\allowbreak{}分别与外侧整边、\allowbreak{}两半段平行，\allowbreak{}每个方向\ensuremath{3}组，\allowbreak{}共\ensuremath{9}组；\allowbreak{}无垂直。\allowbreak{}\par};
\end{scope}\end{tikzpicture}\newpage
\definecolor{answer}{rgb}{0.850980,0.047059,0.070588}
\null\begin{tikzpicture}[remember picture,overlay]\begin{scope}[shift={(current page.north west)},x=1bp,y=-1bp]
\node[anchor=north west,inner sep=0pt,text=answer,align=left,text width=484.0bp] at (15,15) {\fontsize{7.5bp}{9.5bp}\selectfont 析：\allowbreak{}旋转只改变位置、\allowbreak{}方向，\allowbreak{}不改变形状和大小。\allowbreak{}\par};
\end{scope}\end{tikzpicture}\newpage
\definecolor{answer}{rgb}{0.850980,0.047059,0.070588}
\null\begin{tikzpicture}[remember picture,overlay]\begin{scope}[shift={(current page.north west)},x=1bp,y=-1bp]
\node[anchor=north west,inner sep=0pt,text=answer,align=left,text width=484.0bp] at (15,15) {\fontsize{7.5bp}{9.5bp}\selectfont 析：\allowbreak{}沿对称轴对折，\allowbreak{}两侧对应部分完全重合。\allowbreak{}\par};
\end{scope}\end{tikzpicture}\newpage
\definecolor{answer}{rgb}{0.850980,0.047059,0.070588}
\null\begin{tikzpicture}[remember picture,overlay]\begin{scope}[shift={(current page.north west)},x=1bp,y=-1bp]
\node[anchor=north west,inner sep=0pt,text=answer,align=left,text width=484.0bp] at (15,15) {\fontsize{7.5bp}{9.5bp}\selectfont 析：\allowbreak{}①方向盘绕轴转动、\allowbreak{}④车轮绕轴转动，\allowbreak{}都是旋转；\allowbreak{}②升降国旗和③水平拖地都是平移。\allowbreak{}\par};
\end{scope}\end{tikzpicture}\newpage
\definecolor{answer}{rgb}{0.850980,0.047059,0.070588}
\null\begin{tikzpicture}[remember picture,overlay]\begin{scope}[shift={(current page.north west)},x=1bp,y=-1bp]
\node[anchor=north west,inner sep=0pt,text=answer,align=left,text width=484.0bp] at (15,15) {\fontsize{7.5bp}{9.5bp}\selectfont 析：\allowbreak{}甲到乙绕O逆时针转\ensuremath{90}\ensuremath{^\circ};\allowbreak{}A到B绕O顺时针转\ensuremath{90}\ensuremath{^\circ},\allowbreak{}中心O不动。\allowbreak{}\par};
\end{scope}\end{tikzpicture}\newpage
\definecolor{answer}{rgb}{0.850980,0.047059,0.070588}
\null\begin{tikzpicture}[remember picture,overlay]\begin{scope}[shift={(current page.north west)},x=1bp,y=-1bp]
\node[anchor=north west,inner sep=0pt,text=answer,align=left,text width=484.125bp] at (15,15) {\fontsize{7.54bp}{9.54bp}\selectfont 析：\allowbreak{}长方形有\ensuremath{2}条对称轴；\allowbreak{}正方形再加两条对角线所在的直线，\allowbreak{}共\ensuremath{4}条对称轴；\allowbreak{}正五角星有\ensuremath{5}条对称轴\par};
\end{scope}\end{tikzpicture}\newpage
\definecolor{answer}{rgb}{0.850980,0.047059,0.070588}
\null\begin{tikzpicture}[remember picture,overlay]\begin{scope}[shift={(current page.north west)},x=1bp,y=-1bp]
\node[anchor=north west,inner sep=0pt,text=answer,align=left,text width=484.0bp] at (15,15) {\fontsize{7.5bp}{9.5bp}\selectfont 析：\allowbreak{}从上方OFF 到右方\ensuremath{1}档，\allowbreak{}顺时针转四分之一周，\allowbreak{}即\ensuremath{90}\ensuremath{^\circ}。\allowbreak{}\par};
\end{scope}\end{tikzpicture}\newpage
\definecolor{answer}{rgb}{0.850980,0.047059,0.070588}
\null\begin{tikzpicture}[remember picture,overlay]\begin{scope}[shift={(current page.north west)},x=1bp,y=-1bp]
\node[anchor=north west,inner sep=0pt,text=answer,align=left,text width=484.0bp] at (15,15) {\fontsize{7.5bp}{9.5bp}\selectfont 析：\allowbreak{}等边周长\ensuremath{28} × \ensuremath{3} = \ensuremath{84}米；\allowbreak{}腰(\ensuremath{25} \ensuremath{-}\ensuremath{7}) ÷ \ensuremath{2} = \ensuremath{9}米，\allowbreak{}也可\ensuremath{7} + (\ensuremath{25} \ensuremath{-}\ensuremath{7} × \ensuremath{3}) ÷ \ensuremath{2} = \ensuremath{9}米。\allowbreak{}\par};
\end{scope}\end{tikzpicture}\newpage
\definecolor{answer}{rgb}{0.850980,0.047059,0.070588}
\null\begin{tikzpicture}[remember picture,overlay]\begin{scope}[shift={(current page.north west)},x=1bp,y=-1bp]
\node[anchor=north west,inner sep=0pt,text=answer,align=left,text width=484.0bp] at (15,15) {\fontsize{7.5bp}{9.5bp}\selectfont 析：\allowbreak{}\ensuremath{1}个不能成四边形，\allowbreak{}\ensuremath{2}个拼成平行四边形；\allowbreak{}\ensuremath{3}个可拼上底\ensuremath{1}份、\allowbreak{}下底\ensuremath{2}份的梯形，\allowbreak{}至少\ensuremath{3}个。\allowbreak{}\par};
\end{scope}\end{tikzpicture}\newpage
\definecolor{answer}{rgb}{0.850980,0.047059,0.070588}
\null\begin{tikzpicture}[remember picture,overlay]\begin{scope}[shift={(current page.north west)},x=1bp,y=-1bp]
\node[anchor=north west,inner sep=0pt,text=answer,align=left,text width=484.0bp] at (15,15) {\fontsize{7.5bp}{9.5bp}\selectfont 析：\allowbreak{}底\ensuremath{8}配高\ensuremath{5},\allowbreak{}底\ensuremath{6}配高\ensuremath{7};\allowbreak{}周长\ensuremath{6} × \ensuremath{2} + \ensuremath{8} × \ensuremath{2} = \ensuremath{28}厘米，\allowbreak{}也可(\ensuremath{6} + \ensuremath{8}) × \ensuremath{2} = \ensuremath{28}厘米。\allowbreak{}\par};
\end{scope}\end{tikzpicture}\newpage
\definecolor{answer}{rgb}{0.850980,0.047059,0.070588}
\null\begin{tikzpicture}[remember picture,overlay]\begin{scope}[shift={(current page.north west)},x=1bp,y=-1bp]
\node[anchor=north west,inner sep=0pt,text=answer,align=left,text width=484.0bp] at (15,15) {\fontsize{7.5bp}{9.5bp}\selectfont 析：\allowbreak{}上底\ensuremath{3}厘米全部缩掉后变成一个顶点，\allowbreak{}两腰与下底围成三角形。\allowbreak{}\par};
\end{scope}\end{tikzpicture}\newpage
\definecolor{answer}{rgb}{0.850980,0.047059,0.070588}
\null\begin{tikzpicture}[remember picture,overlay]\begin{scope}[shift={(current page.north west)},x=1bp,y=-1bp]
\node[anchor=north west,inner sep=0pt,text=answer,align=left,text width=484.0bp] at (15,15) {\fontsize{7.5bp}{9.5bp}\selectfont 析：\allowbreak{}单个三角形\ensuremath{4}个；\allowbreak{}相邻\ensuremath{2}个组成的小平行四边形\ensuremath{3}个，\allowbreak{}加整个\ensuremath{1}个，\allowbreak{}共\ensuremath{4}个；\allowbreak{}连续\ensuremath{3}个小三角形组成梯形，\allowbreak{}共\ensuremath{2}个。\allowbreak{}\par};
\end{scope}\end{tikzpicture}\newpage
\definecolor{answer}{rgb}{0.850980,0.047059,0.070588}
\null\begin{tikzpicture}[remember picture,overlay]\begin{scope}[shift={(current page.north west)},x=1bp,y=-1bp]
\node[anchor=north west,inner sep=0pt,text=answer,align=left,text width=484.0bp] at (15,15) {\fontsize{7.5bp}{9.5bp}\selectfont 析：\allowbreak{}三角形由三条线段围成，\allowbreak{}直线无限延伸，\allowbreak{}原说法错。\allowbreak{}\par};
\end{scope}\end{tikzpicture}\newpage
\definecolor{answer}{rgb}{0.850980,0.047059,0.070588}
\null\begin{tikzpicture}[remember picture,overlay]\begin{scope}[shift={(current page.north west)},x=1bp,y=-1bp]
\node[anchor=north west,inner sep=0pt,text=answer,align=left,text width=484.0bp] at (15,15) {\fontsize{7.5bp}{9.5bp}\selectfont 析：\allowbreak{}\ensuremath{3}时\ensuremath{30}分时针在\ensuremath{3}与\ensuremath{4}中间，\allowbreak{}夹角\ensuremath{75}\ensuremath{^\circ},\allowbreak{}不是直角。\allowbreak{}\par};
\end{scope}\end{tikzpicture}\newpage
\definecolor{answer}{rgb}{0.850980,0.047059,0.070588}
\null\begin{tikzpicture}[remember picture,overlay]\begin{scope}[shift={(current page.north west)},x=1bp,y=-1bp]
\node[anchor=north west,inner sep=0pt,text=answer,align=left,text width=484.0bp] at (15,15) {\fontsize{7.5bp}{9.5bp}\selectfont 析：\allowbreak{}每条直线上的相邻两角合\ensuremath{180}°；\allowbreak{}依次用\ensuremath{180}\ensuremath{^\circ}\ensuremath{-}\ensuremath{90}\ensuremath{^\circ}= \ensuremath{90}\ensuremath{^\circ}，\allowbreak{}另外三个角也都是直角。\allowbreak{}\par};
\end{scope}\end{tikzpicture}\newpage
\definecolor{answer}{rgb}{0.850980,0.047059,0.070588}
\null\begin{tikzpicture}[remember picture,overlay]\begin{scope}[shift={(current page.north west)},x=1bp,y=-1bp]
\node[anchor=north west,inner sep=0pt,text=answer,align=left,text width=484.0bp] at (15,15) {\fontsize{7.5bp}{9.5bp}\selectfont 析：\allowbreak{}伸缩门利用平行四边形容易变形的特点，\allowbreak{}不是稳定性。\allowbreak{}\par};
\end{scope}\end{tikzpicture}\newpage
\definecolor{answer}{rgb}{0.850980,0.047059,0.070588}
\null\begin{tikzpicture}[remember picture,overlay]\begin{scope}[shift={(current page.north west)},x=1bp,y=-1bp]
\node[anchor=north west,inner sep=0pt,text=answer,align=left,text width=484.0bp] at (15,15) {\fontsize{7.5bp}{9.5bp}\selectfont 析：\allowbreak{}三角形至少有\ensuremath{2}个锐角，\allowbreak{}最多只有\ensuremath{1}个直角或钝角。\allowbreak{}\par};
\end{scope}\end{tikzpicture}\newpage
\definecolor{answer}{rgb}{0.850980,0.047059,0.070588}
\null\begin{tikzpicture}[remember picture,overlay]\begin{scope}[shift={(current page.north west)},x=1bp,y=-1bp]
\node[anchor=north west,inner sep=0pt,text=answer,align=left,text width=484.0bp] at (15,15) {\fontsize{7.5bp}{9.5bp}\selectfont 析：\allowbreak{}图案沿水平中线对折，\allowbreak{}上下能完全重合。\allowbreak{}\par};
\end{scope}\end{tikzpicture}\newpage
\definecolor{answer}{rgb}{0.850980,0.047059,0.070588}
\null\begin{tikzpicture}[remember picture,overlay]\begin{scope}[shift={(current page.north west)},x=1bp,y=-1bp]
\node[anchor=north west,inner sep=0pt,text=answer,align=left,text width=484.0bp] at (15,15) {\fontsize{7.5bp}{9.5bp}\selectfont 析：\allowbreak{}逆时针转\ensuremath{90}\ensuremath{^\circ}后小旗应朝另一侧，\allowbreak{}图示方向不符。\allowbreak{}\par};
\end{scope}\end{tikzpicture}\newpage
\definecolor{answer}{rgb}{0.850980,0.047059,0.070588}
\null\begin{tikzpicture}[remember picture,overlay]\begin{scope}[shift={(current page.north west)},x=1bp,y=-1bp]
\node[anchor=north west,inner sep=0pt,text=answer,align=left,text width=484.0bp] at (15,15) {\fontsize{7.5bp}{9.5bp}\selectfont 析：\allowbreak{}最短两边\ensuremath{6} + \ensuremath{6}等于\ensuremath{12},\allowbreak{}不能围成三角形；\allowbreak{}必须大于第三边。\allowbreak{}\par};
\end{scope}\end{tikzpicture}\newpage
\definecolor{answer}{rgb}{0.850980,0.047059,0.070588}
\null\begin{tikzpicture}[remember picture,overlay]\begin{scope}[shift={(current page.north west)},x=1bp,y=-1bp]
\node[anchor=north west,inner sep=0pt,text=answer,align=left,text width=484.0bp] at (15,15) {\fontsize{7.5bp}{9.5bp}\selectfont 析：\allowbreak{}同平面内，\allowbreak{}相交与平行互斥；\allowbreak{}垂直属于相交，\allowbreak{}故选B。\allowbreak{}\par};
\end{scope}\end{tikzpicture}\newpage
\definecolor{answer}{rgb}{0.850980,0.047059,0.070588}
\null\begin{tikzpicture}[remember picture,overlay]\begin{scope}[shift={(current page.north west)},x=1bp,y=-1bp]
\node[anchor=north west,inner sep=0pt,text=answer,align=left,text width=484.0bp] at (15,15) {\fontsize{7.5bp}{9.5bp}\selectfont 析：\allowbreak{}两条平行线间的垂直距离处处相等，\allowbreak{}两根等长绳子表示同一距离，\allowbreak{}选C。\allowbreak{}\par};
\end{scope}\end{tikzpicture}\newpage
\definecolor{answer}{rgb}{0.850980,0.047059,0.070588}
\null\begin{tikzpicture}[remember picture,overlay]\begin{scope}[shift={(current page.north west)},x=1bp,y=-1bp]
\node[anchor=north west,inner sep=0pt,text=answer,align=left,text width=484.0bp] at (15,15) {\fontsize{7.5bp}{9.5bp}\selectfont 析：\allowbreak{}AB边上的高应从C出发且垂直AB,\allowbreak{}图中为CD,\allowbreak{}选C。\allowbreak{}\par};
\end{scope}\end{tikzpicture}\newpage
\definecolor{answer}{rgb}{0.850980,0.047059,0.070588}
\null\begin{tikzpicture}[remember picture,overlay]\begin{scope}[shift={(current page.north west)},x=1bp,y=-1bp]
\node[anchor=north west,inner sep=0pt,text=answer,align=left,text width=484.0bp] at (15,15) {\fontsize{7.5bp}{9.5bp}\selectfont 析：\allowbreak{}可平切成两个平行四边形，\allowbreak{}或斜切成两个梯形；\allowbreak{}不能一个平行四边形加一个梯形，\allowbreak{}选B。\allowbreak{}\par};
\end{scope}\end{tikzpicture}\newpage
\definecolor{answer}{rgb}{0.850980,0.047059,0.070588}
\null\begin{tikzpicture}[remember picture,overlay]\begin{scope}[shift={(current page.north west)},x=1bp,y=-1bp]
\node[anchor=north west,inner sep=0pt,text=answer,align=left,text width=484.0bp] at (15,15) {\fontsize{7.5bp}{9.5bp}\selectfont 析：\allowbreak{}按角分的三类三角形互不重叠且覆盖全部，\allowbreak{}图A正确；\allowbreak{}正方形应在长方形内部，\allowbreak{}图B错，\allowbreak{}等边应在等腰内部，\allowbreak{}续页C也错。\allowbreak{}\par};
\end{scope}\end{tikzpicture}\newpage
\definecolor{answer}{rgb}{0.850980,0.047059,0.070588}
\null\begin{tikzpicture}[remember picture,overlay]\begin{scope}[shift={(current page.north west)},x=1bp,y=-1bp]
\node[anchor=north west,inner sep=0pt,text=answer,align=left,text width=484.0bp] at (15,15) {\fontsize{7.5bp}{9.5bp}\selectfont 析：\allowbreak{}顺时针转\ensuremath{90}\ensuremath{^\circ}从上方A到右侧B,\allowbreak{}选C。\allowbreak{}\par};
\end{scope}\end{tikzpicture}\newpage
\definecolor{answer}{rgb}{0.850980,0.047059,0.070588}
\null\begin{tikzpicture}[remember picture,overlay]\begin{scope}[shift={(current page.north west)},x=1bp,y=-1bp]
\node[anchor=north west,inner sep=0pt,text=answer,align=left,text width=484.0bp] at (15,15) {\fontsize{7.5bp}{9.5bp}\selectfont 析：\allowbreak{}等腰钝角三角形的两底角和小于\ensuremath{90}\ensuremath{^\circ},\allowbreak{}每个底角小于\ensuremath{45}\ensuremath{^\circ};\allowbreak{} \ensuremath{45}\ensuremath{^\circ}不行，\allowbreak{}选C。\allowbreak{}\par};
\end{scope}\end{tikzpicture}\newpage
\definecolor{answer}{rgb}{0.850980,0.047059,0.070588}
\null\begin{tikzpicture}[remember picture,overlay]\begin{scope}[shift={(current page.north west)},x=1bp,y=-1bp]
\node[anchor=north west,inner sep=0pt,text=answer,align=left,text width=484.0bp] at (15,15) {\fontsize{7.5bp}{9.5bp}\selectfont 析：\allowbreak{}旋转保持各对应点到中心的距离相等，\allowbreak{}C正确；\allowbreak{}旋转一般改变方向，\allowbreak{}平移和旋转均不改变大小。\allowbreak{}\par};
\end{scope}\end{tikzpicture}\newpage
\definecolor{answer}{rgb}{0.850980,0.047059,0.070588}
\null\begin{tikzpicture}[remember picture,overlay]\begin{scope}[shift={(current page.north west)},x=1bp,y=-1bp]
\node[anchor=north west,inner sep=0pt,text=answer,align=left,text width=484.0bp] at (15,15) {\fontsize{7.5bp}{9.5bp}\selectfont 析：\allowbreak{}高垂直于对应底；\allowbreak{}梯形可从上底顶点或内部点作高，\allowbreak{}垂足可在下底延长线上。\allowbreak{}三图按原PDF 约\ensuremath{1.1}、\allowbreak{}\ensuremath{1.2}、\allowbreak{}\ensuremath{1.7}厘米。\allowbreak{}\par};
\end{scope}\end{tikzpicture}\newpage
\definecolor{answer}{rgb}{0.850980,0.047059,0.070588}
\null\begin{tikzpicture}[remember picture,overlay]\begin{scope}[shift={(current page.north west)},x=1bp,y=-1bp]
\node[anchor=north west,inner sep=0pt,text=answer,align=left,text width=483.6666564941406bp] at (15,15) {\fontsize{8.66bp}{10.82bp}\selectfont 注:\allowbreak{} 按本PDF原尺寸测量,\allowbreak{}印刷或缩放后长度可能不同。\allowbreak{}\par};
\end{scope}\end{tikzpicture}\newpage
\definecolor{answer}{rgb}{0.850980,0.047059,0.070588}
\null\begin{tikzpicture}[remember picture,overlay]\begin{scope}[shift={(current page.north west)},x=1bp,y=-1bp]
\node[anchor=north west,inner sep=0pt,text=answer,align=left,text width=321.7bp] at (15,15) {\fontsize{10.67bp}{13.34bp}\selectfont 向左岸作垂线,\allowbreak{} 沿这条线段运送。\allowbreak{}理由:\allowbreak{} 点到直线的垂直线段最短。\allowbreak{}\par};
\end{scope}\end{tikzpicture}\newpage
\definecolor{answer}{rgb}{0.850980,0.047059,0.070588}
\null\begin{tikzpicture}[remember picture,overlay]\begin{scope}[shift={(current page.north west)},x=1bp,y=-1bp]
\node[anchor=north west,inner sep=0pt,text=answer,align=left,text width=484.0bp] at (15,15) {\fontsize{8.5bp}{10.62bp}\selectfont 两把曲尺都紧靠同一条直边，\allowbreak{}测得的相同刻度表示同样的垂直高度。\allowbreak{}此题用几何关系说明，\allowbreak{}无需数值脱式计算。\allowbreak{}\par};
\end{scope}\end{tikzpicture}\newpage
\definecolor{answer}{rgb}{0.850980,0.047059,0.070588}
\null\begin{tikzpicture}[remember picture,overlay]\begin{scope}[shift={(current page.north west)},x=1bp,y=-1bp]
\node[anchor=north west,inner sep=0pt,text=answer,align=left,text width=484.0bp] at (15,15) {\fontsize{8.5bp}{10.62bp}\selectfont 一周\ensuremath{360}\ensuremath{^\circ}对应\ensuremath{12}千克，\allowbreak{}\ensuremath{90}\ensuremath{^\circ}是四分之一周；\allowbreak{}也可直接读右侧刻度。\allowbreak{}\par};
\end{scope}\end{tikzpicture}\newpage
\definecolor{answer}{rgb}{0.850980,0.047059,0.070588}
\null\begin{tikzpicture}[remember picture,overlay]\begin{scope}[shift={(current page.north west)},x=1bp,y=-1bp]
\node[anchor=north west,inner sep=0pt,text=answer,align=left,text width=484.0bp] at (15,15) {\fontsize{8.5bp}{10.62bp}\selectfont 拿走物品使示数减小，\allowbreak{}应逆时针旋转；\allowbreak{}拿走\ensuremath{6}千克占量程\ensuremath{12}千克的一半，\allowbreak{}指针转半周。\allowbreak{}\par};
\end{scope}\end{tikzpicture}\newpage
\definecolor{answer}{rgb}{0.850980,0.047059,0.070588}
\null\begin{tikzpicture}[remember picture,overlay]\begin{scope}[shift={(current page.north west)},x=1bp,y=-1bp]
\node[anchor=north west,inner sep=0pt,text=answer,align=left,text width=484.0bp] at (15,15) {\fontsize{8.5bp}{10.62bp}\selectfont 等边三角形三边相等，\allowbreak{}用铁丝总长除以\ensuremath{3}。\allowbreak{}\par};
\end{scope}\end{tikzpicture}\newpage
\definecolor{answer}{rgb}{0.850980,0.047059,0.070588}
\null\begin{tikzpicture}[remember picture,overlay]\begin{scope}[shift={(current page.north west)},x=1bp,y=-1bp]
\node[anchor=north west,inner sep=0pt,text=answer,align=left,text width=484.0bp] at (15,15) {\fontsize{8.5bp}{10.62bp}\selectfont 等腰三角形两腰相等，\allowbreak{}用周长减底边，\allowbreak{}再平均分为\ensuremath{2}份。\allowbreak{}\par};
\end{scope}\end{tikzpicture}\newpage
\definecolor{answer}{rgb}{0.850980,0.047059,0.070588}
\null\begin{tikzpicture}[remember picture,overlay]\begin{scope}[shift={(current page.north west)},x=1bp,y=-1bp]
\node[anchor=north west,inner sep=0pt,text=answer,align=left,text width=484.0bp] at (15,15) {\fontsize{8.5bp}{10.62bp}\selectfont 铁丝总长不变。\allowbreak{}等边三角形的\ensuremath{3}条边总长，\allowbreak{}改平均分给正方形\ensuremath{4}条边。\allowbreak{}\par};
\end{scope}\end{tikzpicture}\newpage
\definecolor{answer}{rgb}{0.850980,0.047059,0.070588}
\null\begin{tikzpicture}[remember picture,overlay]\begin{scope}[shift={(current page.north west)},x=1bp,y=-1bp]
\node[anchor=north west,inner sep=0pt,text=answer,align=left,text width=484.0bp] at (15,15) {\fontsize{8.5bp}{10.62bp}\selectfont 所有踏面的水平长度之和\ensuremath{3.4}米，\allowbreak{}所有立面的高度之和\ensuremath{2.2}米；\allowbreak{}宽\ensuremath{1.5}米不参与长度相加。\allowbreak{}\par};
\end{scope}\end{tikzpicture}\newpage
\definecolor{answer}{rgb}{0.850980,0.047059,0.070588}
\null\begin{tikzpicture}[remember picture,overlay]\begin{scope}[shift={(current page.north west)},x=1bp,y=-1bp]
\node[anchor=north west,inner sep=0pt,text=answer,align=left,text width=484.0bp] at (15,15) {\fontsize{8.5bp}{10.62bp}\selectfont 两条由原长方形长边得到的边长\ensuremath{40}厘米；\allowbreak{}平行四边形的一组邻边和等于周长的一半。\allowbreak{}\par};
\end{scope}\end{tikzpicture}\newpage
\definecolor{answer}{rgb}{0.850980,0.047059,0.070588}
\null\begin{tikzpicture}[remember picture,overlay]\begin{scope}[shift={(current page.north west)},x=1bp,y=-1bp]
\node[anchor=north west,inner sep=0pt,text=answer,align=left,text width=483.9999694824219bp] at (15,15) {\fontsize{8.5bp}{10.62bp}\selectfont 两底相差\ensuremath{6}厘米。\allowbreak{}上底\ensuremath{1}份、\allowbreak{}下底\ensuremath{4}份，\allowbreak{}相差\ensuremath{3}份；\allowbreak{}求下底时用每份数×份数。\allowbreak{}\par};
\end{scope}\end{tikzpicture}\newpage
\definecolor{answer}{rgb}{0.850980,0.047059,0.070588}
\null\begin{tikzpicture}[remember picture,overlay]\begin{scope}[shift={(current page.north west)},x=1bp,y=-1bp]
\node[anchor=north west,inner sep=0pt,text=answer,align=left,text width=484.0bp] at (15,15) {\fontsize{8.5bp}{10.62bp}\selectfont 沿一条腰拼接，\allowbreak{}重合的两条腰都成为内部线段，\allowbreak{}不计入拼成图形的周长。\allowbreak{}\par};
\end{scope}\end{tikzpicture}\newpage
\definecolor{answer}{rgb}{0.850980,0.047059,0.070588}
\null\begin{tikzpicture}[remember picture,overlay]\begin{scope}[shift={(current page.north west)},x=1bp,y=-1bp]
\node[anchor=north west,inner sep=0pt,text=answer,align=left,text width=484.0bp] at (15,15) {\fontsize{8.5bp}{10.62bp}\selectfont 铁丝周长不变。\allowbreak{}图中等腰梯形上底仍为\ensuremath{8}厘米，\allowbreak{}用总周长减两腰和上底。\allowbreak{}\par};
\end{scope}\end{tikzpicture}\newpage
\definecolor{answer}{rgb}{0.850980,0.047059,0.070588}
\null\begin{tikzpicture}[remember picture,overlay]\begin{scope}[shift={(current page.north west)},x=1bp,y=-1bp]
\node[anchor=north west,inner sep=0pt,text=answer,align=left,text width=483.9988708496094bp] at (15,15) {\fontsize{7.5bp}{9.5bp}\selectfont 析：\allowbreak{}两人分别有\ensuremath{3}次、\allowbreak{}\ensuremath{4}次记录，\allowbreak{}不能直接比总分\ensuremath{33}与\ensuremath{40}；\allowbreak{}平均\ensuremath{33} ÷ \ensuremath{3} = \ensuremath{11}分、\allowbreak{}\ensuremath{40} ÷ \ensuremath{4} = \ensuremath{10}分，\allowbreak{}\ensuremath{1}号较高；\allowbreak{}也可移多补少。\allowbreak{}\par};
\end{scope}\end{tikzpicture}\newpage
\definecolor{answer}{rgb}{0.850980,0.047059,0.070588}
\null\begin{tikzpicture}[remember picture,overlay]\begin{scope}[shift={(current page.north west)},x=1bp,y=-1bp]
\node[anchor=north west,inner sep=0pt,text=answer,align=left,text width=484.0bp] at (15,15) {\fontsize{7.5bp}{9.5bp}\selectfont 析：\allowbreak{}总和\ensuremath{450}÷\ensuremath{5}=\ensuremath{90}(步)；\allowbreak{}另法以\ensuremath{90}为基准，\allowbreak{}多\ensuremath{1}、\allowbreak{}多\ensuremath{3}与少\ensuremath{1}、\allowbreak{}少\ensuremath{3}抵消。\allowbreak{}\ensuremath{48} × \ensuremath{90} = \ensuremath{4320}(厘米)\par};
\end{scope}\end{tikzpicture}\newpage
\definecolor{answer}{rgb}{0.850980,0.047059,0.070588}
\null\begin{tikzpicture}[remember picture,overlay]\begin{scope}[shift={(current page.north west)},x=1bp,y=-1bp]
\node[anchor=north west,inner sep=0pt,text=answer,align=left,text width=484.0bp] at (15,15) {\fontsize{7.5bp}{9.5bp}\selectfont 析：\allowbreak{}甲乙总和\ensuremath{75} × \ensuremath{2} = \ensuremath{150},\allowbreak{}三数平均(\ensuremath{150} + \ensuremath{90}) ÷ \ensuremath{3} = \ensuremath{80},\allowbreak{}不是\ensuremath{85}。\allowbreak{}\par};
\end{scope}\end{tikzpicture}\newpage
\definecolor{answer}{rgb}{0.850980,0.047059,0.070588}
\null\begin{tikzpicture}[remember picture,overlay]\begin{scope}[shift={(current page.north west)},x=1bp,y=-1bp]
\node[anchor=north west,inner sep=0pt,text=answer,align=left,text width=484.0bp] at (15,15) {\fontsize{7.5bp}{9.5bp}\selectfont 析：\allowbreak{}问平均每组，\allowbreak{}不是每人；\allowbreak{}(\ensuremath{12} + \ensuremath{8}) ÷ \ensuremath{2} = \ensuremath{10}本。\allowbreak{}\par};
\end{scope}\end{tikzpicture}\newpage
\definecolor{answer}{rgb}{0.850980,0.047059,0.070588}
\null\begin{tikzpicture}[remember picture,overlay]\begin{scope}[shift={(current page.north west)},x=1bp,y=-1bp]
\node[anchor=north west,inner sep=0pt,text=answer,align=left,text width=484.0bp] at (15,15) {\fontsize{7.5bp}{9.5bp}\selectfont 析：\allowbreak{}平均\ensuremath{35}千克反映整体水平，\allowbreak{}不表示每个人都重\ensuremath{35}千克。\allowbreak{}\par};
\end{scope}\end{tikzpicture}\newpage
\definecolor{answer}{rgb}{0.850980,0.047059,0.070588}
\null\begin{tikzpicture}[remember picture,overlay]\begin{scope}[shift={(current page.north west)},x=1bp,y=-1bp]
\node[anchor=north west,inner sep=0pt,text=answer,align=left,text width=484.0bp] at (15,15) {\fontsize{7.5bp}{9.5bp}\selectfont 析：\allowbreak{}平均\ensuremath{201}厘米时，\allowbreak{}某位队员仍可能\ensuremath{180}厘米，\allowbreak{}其他较高队员可把平均数拉高，\allowbreak{}选A。\allowbreak{}\par};
\end{scope}\end{tikzpicture}\newpage
\definecolor{answer}{rgb}{0.850980,0.047059,0.070588}
\null\begin{tikzpicture}[remember picture,overlay]\begin{scope}[shift={(current page.north west)},x=1bp,y=-1bp]
\node[anchor=north west,inner sep=0pt,text=answer,align=left,text width=484.0bp] at (15,15) {\fontsize{7.5bp}{9.5bp}\selectfont 析：\allowbreak{}四数总和\ensuremath{14} × \ensuremath{4} = \ensuremath{56},\allowbreak{}已知和\ensuremath{16} + \ensuremath{12} + \ensuremath{14} = \ensuremath{42},\allowbreak{} \ensuremath{a} = \ensuremath{56} \ensuremath{-}\ensuremath{42} = \ensuremath{14},\allowbreak{}选B。\allowbreak{}\par};
\end{scope}\end{tikzpicture}\newpage
\definecolor{answer}{rgb}{0.850980,0.047059,0.070588}
\null\begin{tikzpicture}[remember picture,overlay]\begin{scope}[shift={(current page.north west)},x=1bp,y=-1bp]
\node[anchor=north west,inner sep=0pt,text=answer,align=left,text width=484.0bp] at (15,15) {\fontsize{7.5bp}{9.5bp}\selectfont 析：\allowbreak{}总价\ensuremath{60} + \ensuremath{100} = \ensuremath{160}元，\allowbreak{}总质量\ensuremath{3} + \ensuremath{5} = \ensuremath{8}千克，\allowbreak{}平均\ensuremath{160} ÷ \ensuremath{8} = \ensuremath{20}元/千克，\allowbreak{}选B。\allowbreak{}\par};
\end{scope}\end{tikzpicture}\newpage
\definecolor{answer}{rgb}{0.850980,0.047059,0.070588}
\null\begin{tikzpicture}[remember picture,overlay]\begin{scope}[shift={(current page.north west)},x=1bp,y=-1bp]
\node[anchor=north west,inner sep=0pt,text=answer,align=left,text width=494bp] at (15,15) {\fontsize{8.5bp}{10.62bp}\selectfont 比较\ensuremath{8}、\allowbreak{}\ensuremath{15}、\allowbreak{}\ensuremath{13}、\allowbreak{}\ensuremath{12},\allowbreak{}最大是\ensuremath{15},\allowbreak{}最小是\ensuremath{8}。\allowbreak{}\par};
\end{scope}\end{tikzpicture}\newpage
\definecolor{answer}{rgb}{0.850980,0.047059,0.070588}
\null\begin{tikzpicture}[remember picture,overlay]\begin{scope}[shift={(current page.north west)},x=1bp,y=-1bp]
\node[anchor=north west,inner sep=0pt,text=answer,align=left,text width=494.0bp] at (15,15) {\fontsize{8.5bp}{10.62bp}\selectfont \ensuremath{15}对应书法，\allowbreak{}\ensuremath{8}对应绘画。\allowbreak{}\par};
\end{scope}\end{tikzpicture}\newpage
\definecolor{answer}{rgb}{0.850980,0.047059,0.070588}
\null\begin{tikzpicture}[remember picture,overlay]\begin{scope}[shift={(current page.north west)},x=1bp,y=-1bp]
\node[anchor=north west,inner sep=0pt,text=answer,align=left,text width=483.99952697753906bp] at (15,15) {\fontsize{7.5bp}{9.5bp}\selectfont 析：\allowbreak{}按\ensuremath{50}及以上、\allowbreak{}\ensuremath{40}～\ensuremath{49}、\allowbreak{}\ensuremath{30}～\ensuremath{39}、\allowbreak{}\ensuremath{20}～\ensuremath{29}、\allowbreak{}\ensuremath{10}～\ensuremath{19}、\allowbreak{}\ensuremath{10}以下分类，\allowbreak{}各数只记一次。\allowbreak{}一组\ensuremath{1}、\allowbreak{}\ensuremath{4}、\allowbreak{}\ensuremath{4}、\allowbreak{}\ensuremath{1}、\allowbreak{}\ensuremath{1}、\allowbreak{}\ensuremath{1}；\allowbreak{}二组\ensuremath{1}、\allowbreak{}\ensuremath{4}、\allowbreak{}\ensuremath{4}、\allowbreak{}\ensuremath{1}、\allowbreak{}\ensuremath{2}、\allowbreak{}\ensuremath{0}；\allowbreak{}总计\ensuremath{2}、\allowbreak{}\ensuremath{8}、\allowbreak{}\ensuremath{8}、\allowbreak{}\ensuremath{2}、\allowbreak{}\ensuremath{3}、\allowbreak{}\ensuremath{1}。\allowbreak{}核验两组各\ensuremath{12}人，\allowbreak{}共\ensuremath{24}人。\allowbreak{}\par};
\end{scope}\end{tikzpicture}\newpage
\definecolor{answer}{rgb}{0.850980,0.047059,0.070588}
\null\begin{tikzpicture}[remember picture,overlay]\begin{scope}[shift={(current page.north west)},x=1bp,y=-1bp]
\node[anchor=north west,inner sep=0pt,text=answer,align=left,text width=494.0bp] at (15,15) {\fontsize{8.5bp}{10.62bp}\selectfont 开放题，\allowbreak{}可比较达到同一等级门槛的人数，\allowbreak{}或比较两组是否全部达标。\allowbreak{}\par};
\end{scope}\end{tikzpicture}\newpage
\definecolor{answer}{rgb}{0.850980,0.047059,0.070588}
\null\begin{tikzpicture}[remember picture,overlay]\begin{scope}[shift={(current page.north west)},x=1bp,y=-1bp]
\node[anchor=north west,inner sep=0pt,text=answer,align=left,text width=484.0bp] at (15,15) {\fontsize{7.5bp}{9.5bp}\selectfont 析：\allowbreak{}方法\ensuremath{1}把四位评委分数相加再÷\ensuremath{4} :\allowbreak{}低年级缺\ensuremath{57}、\allowbreak{}\ensuremath{82}、\allowbreak{}\ensuremath{72},\allowbreak{}高年级缺\ensuremath{86}、\allowbreak{}\ensuremath{90}、\allowbreak{}\ensuremath{68}、\allowbreak{}\ensuremath{57};\allowbreak{}方法\ensuremath{2}可用这些基准分检查多分与少分是否抵消。\allowbreak{}\par};
\end{scope}\end{tikzpicture}\newpage
\definecolor{answer}{rgb}{0.850980,0.047059,0.070588}
\null\begin{tikzpicture}[remember picture,overlay]\begin{scope}[shift={(current page.north west)},x=1bp,y=-1bp]
\node[anchor=north west,inner sep=0pt,text=answer,align=left,text width=484.0bp] at (15,15) {\fontsize{7.5bp}{9.5bp}\selectfont 析：\allowbreak{}按优秀\ensuremath{90}～\ensuremath{100}、\allowbreak{}良好\ensuremath{80}～\ensuremath{89}、\allowbreak{}合格\ensuremath{60}～\ensuremath{79}、\allowbreak{}不合格≤\ensuremath{59}分类。\allowbreak{}低年级\ensuremath{10}、\allowbreak{}\ensuremath{2}、\allowbreak{}\ensuremath{3}、\allowbreak{}\ensuremath{3}、\allowbreak{}\ensuremath{2};\allowbreak{}高年级\ensuremath{10}、\allowbreak{}\ensuremath{3}、\allowbreak{}\ensuremath{2}、\allowbreak{}\ensuremath{3}、\allowbreak{}\ensuremath{2};\allowbreak{}合计\ensuremath{20}、\allowbreak{}\ensuremath{5}、\allowbreak{}\ensuremath{5}、\allowbreak{}\ensuremath{6}、\allowbreak{}\ensuremath{4}。\allowbreak{}\par};
\end{scope}\end{tikzpicture}\newpage
\definecolor{answer}{rgb}{0.850980,0.047059,0.070588}
\null\begin{tikzpicture}[remember picture,overlay]\begin{scope}[shift={(current page.north west)},x=1bp,y=-1bp]
\node[anchor=north west,inner sep=0pt,text=answer,align=left,text width=484.0bp] at (15,15) {\fontsize{8.5bp}{10.62bp}\selectfont 比较时先确定评价角度。\allowbreak{}可看各等级作品数，\allowbreak{}也可把每件作品平均分求和后再求平均。\allowbreak{}\par};
\end{scope}\end{tikzpicture}\newpage
\definecolor{answer}{rgb}{0.850980,0.047059,0.070588}
\null\begin{tikzpicture}[remember picture,overlay]\begin{scope}[shift={(current page.north west)},x=1bp,y=-1bp]
\node[anchor=north west,inner sep=0pt,text=answer,align=left,text width=484.0bp] at (15,15) {\fontsize{8.5bp}{10.62bp}\selectfont 平均每天借书人数等于总人数除以天数；\allowbreak{}第\ensuremath{4}天加入后应重新除以\ensuremath{4}。\allowbreak{}\par};
\end{scope}\end{tikzpicture}\newpage
\definecolor{answer}{rgb}{0.850980,0.047059,0.070588}
\null\begin{tikzpicture}[remember picture,overlay]\begin{scope}[shift={(current page.north west)},x=1bp,y=-1bp]
\node[anchor=north west,inner sep=0pt,text=answer,align=left,text width=484.0bp] at (15,15) {\fontsize{8.5bp}{10.62bp}\selectfont 计划总销量是每天\ensuremath{150}本乘\ensuremath{5}天，\allowbreak{}扣掉前三天\ensuremath{420}本，\allowbreak{}再平均分到剩下\ensuremath{2}天。\allowbreak{}\par};
\end{scope}\end{tikzpicture}\newpage
\definecolor{answer}{rgb}{0.850980,0.047059,0.070588}
\null\begin{tikzpicture}[remember picture,overlay]\begin{scope}[shift={(current page.north west)},x=1bp,y=-1bp]
\node[anchor=north west,inner sep=0pt,text=answer,align=left,text width=484.0bp] at (15,15) {\fontsize{8.5bp}{10.62bp}\selectfont \ensuremath{4}次测量之间只有\ensuremath{3}个一周间隔；\allowbreak{}平均每周生长量用末次与首次高度差除以\ensuremath{3},\allowbreak{}预测第\ensuremath{5}周再加一周生长量。\allowbreak{}\par};
\end{scope}\end{tikzpicture}\newpage
\definecolor{answer}{rgb}{0.850980,0.047059,0.070588}
\null\begin{tikzpicture}[remember picture,overlay]\begin{scope}[shift={(current page.north west)},x=1bp,y=-1bp]
\node[anchor=north west,inner sep=0pt,text=answer,align=left,text width=484.0bp] at (15,15) {\fontsize{8.5bp}{10.62bp}\selectfont 黄瓜四周从\ensuremath{6}到\ensuremath{18}厘米，\allowbreak{}菠菜从\ensuremath{7}到\ensuremath{25}厘米；\allowbreak{}均比较\ensuremath{3}个一周间隔。\allowbreak{}\par};
\end{scope}\end{tikzpicture}\newpage
\definecolor{answer}{rgb}{0.850980,0.047059,0.070588}
\null\begin{tikzpicture}[remember picture,overlay]\begin{scope}[shift={(current page.north west)},x=1bp,y=-1bp]
\node[anchor=north west,inner sep=0pt,text=answer,align=left,text width=484.0bp] at (15,15) {\fontsize{8.5bp}{10.62bp}\selectfont 答：\allowbreak{}菠菜生长较快；\allowbreak{}可能与品种、\allowbreak{}光照、\allowbreak{}水分、\allowbreak{}养分等有关，\allowbreak{}表中数据不能单独确定具体原因。\allowbreak{}\par};
\end{scope}\end{tikzpicture}\newpage
\definecolor{answer}{rgb}{0.850980,0.047059,0.070588}
\null\begin{tikzpicture}[remember picture,overlay]\begin{scope}[shift={(current page.north west)},x=1bp,y=-1bp]
\node[anchor=north west,inner sep=0pt,text=answer,align=left,text width=484bp] at (15,15) {\fontsize{8.5bp}{10.62bp}\selectfont 开放题，\allowbreak{}选“生长之星”还应综合衡量健康和发育情况。\allowbreak{}\par};
\end{scope}\end{tikzpicture}\newpage
\definecolor{answer}{rgb}{0.850980,0.047059,0.070588}
\null\begin{tikzpicture}[remember picture,overlay]\begin{scope}[shift={(current page.north west)},x=1bp,y=-1bp]
\node[anchor=north west,inner sep=0pt,text=answer,align=left,text width=484.0bp] at (15,15) {\fontsize{8.5bp}{10.62bp}\selectfont 综合观察健康和发育还看叶片数量与颜色、\allowbreak{}茎的粗细、\allowbreak{}根系是否发达、\allowbreak{}有无病虫害等。\allowbreak{}\par};
\end{scope}\end{tikzpicture}\newpage
\definecolor{answer}{rgb}{0.850980,0.047059,0.070588}
\null\begin{tikzpicture}[remember picture,overlay]\begin{scope}[shift={(current page.north west)},x=1bp,y=-1bp]
\node[anchor=north west,inner sep=0pt,text=answer,align=left,text width=484.0bp] at (15,15) {\fontsize{8.5bp}{10.62bp}\selectfont 答：\allowbreak{}这些特征能帮助判断植株是否健壮，\allowbreak{}不能只看高度或增长量。\allowbreak{}\par};
\end{scope}\end{tikzpicture}\newpage
\definecolor{answer}{rgb}{0.850980,0.047059,0.070588}
\null\begin{tikzpicture}[remember picture,overlay]\begin{scope}[shift={(current page.north west)},x=1bp,y=-1bp]
\node[anchor=north west,inner sep=0pt,text=answer,align=left,text width=87bp] at (15,15) {\fontsize{8.0bp}{10.0bp}\selectfont 法\ensuremath{3}：\allowbreak{}先求第三段\par};
\end{scope}\end{tikzpicture}\newpage
\definecolor{answer}{rgb}{0.850980,0.047059,0.070588}
\null\begin{tikzpicture}[remember picture,overlay]\begin{scope}[shift={(current page.north west)},x=1bp,y=-1bp]
\node[anchor=north west,inner sep=0pt,text=answer,align=left,text width=232bp] at (15,15) {\fontsize{8.5bp}{10.62bp}\selectfont 法\ensuremath{1} :\allowbreak{}连加配对凑整：\allowbreak{}选小数部分能互补的一对数，\allowbreak{}再添加另一个整数。\allowbreak{}\par};
\end{scope}\end{tikzpicture}\newpage
\definecolor{answer}{rgb}{0.850980,0.047059,0.070588}
\null\begin{tikzpicture}[remember picture,overlay]\begin{scope}[shift={(current page.north west)},x=1bp,y=-1bp]
\node[anchor=north west,inner sep=0pt,text=answer,align=left,text width=232bp] at (15,15) {\fontsize{8.5bp}{10.62bp}\selectfont 法\ensuremath{3} :\allowbreak{}加减混合先抵消小数部分：\allowbreak{}\ensuremath{6.78}与\ensuremath{0.78}的小数部分相同，\allowbreak{}相减后得整数\ensuremath{6}。\allowbreak{}\par};
\end{scope}\end{tikzpicture}\newpage
\definecolor{answer}{rgb}{0.850980,0.047059,0.070588}
\null\begin{tikzpicture}[remember picture,overlay]\begin{scope}[shift={(current page.north west)},x=1bp,y=-1bp]
\node[anchor=north west,inner sep=0pt,text=answer,align=left,text width=235bp] at (15,15) {\fontsize{7.5bp}{9.5bp}\selectfont 析：\allowbreak{}\ensuremath{a}、\allowbreak{}\ensuremath{b}、\allowbreak{}\ensuremath{c}连续，\allowbreak{}中间数\ensuremath{b}两侧的差恰好抵消，\allowbreak{}平均数为\ensuremath{b}。\allowbreak{}\par};
\end{scope}\end{tikzpicture}\newpage
\definecolor{answer}{rgb}{0.850980,0.047059,0.070588}
\null\begin{tikzpicture}[remember picture,overlay]\begin{scope}[shift={(current page.north west)},x=1bp,y=-1bp]
\node[anchor=north west,inner sep=0pt,text=answer,align=left,text width=122bp] at (15,15) {\fontsize{7.5bp}{9.5bp}\selectfont 通常算法与运算性质简算均可，\allowbreak{}保留两种思路。\allowbreak{}\par};
\end{scope}\end{tikzpicture}\newpage
\definecolor{answer}{rgb}{0.850980,0.047059,0.070588}
\null\begin{tikzpicture}[remember picture,overlay]\begin{scope}[shift={(current page.north west)},x=1bp,y=-1bp]
\node[anchor=north west,inner sep=0pt,text=answer,align=left,text width=243bp] at (15,15) {\fontsize{7.5bp}{9.5bp}\selectfont \ensuremath{12.56} + \ensuremath{8.44} = \ensuremath{21}；\allowbreak{}连续减去两个数，\allowbreak{}可以减去它们的和。\allowbreak{}\par};
\end{scope}\end{tikzpicture}\newpage
\definecolor{answer}{rgb}{0.850980,0.047059,0.070588}
\null\begin{tikzpicture}[remember picture,overlay]\begin{scope}[shift={(current page.north west)},x=1bp,y=-1bp]
\node[anchor=north west,inner sep=0pt,text=answer,align=left,text width=217bp] at (15,15) {\fontsize{7.5bp}{9.5bp}\selectfont 按顺序计算，\allowbreak{}或先把\ensuremath{125}与\ensuremath{25}相除，\allowbreak{}再乘括号内的差。\allowbreak{}\par};
\end{scope}\end{tikzpicture}\newpage
\definecolor{answer}{rgb}{0.772549,0.000000,0.109804}
\null\begin{tikzpicture}[remember picture,overlay]\begin{scope}[shift={(current page.north west)},x=1bp,y=-1bp]
\node[anchor=north west,inner sep=0pt,text=answer,align=left,text width=246bp] at (15,15) {\fontsize{10.2bp}{12.75bp}\selectfont 作法：\allowbreak{}先画两条平行线之间的一条垂线段。\allowbreak{}以这条垂线段为边画正方形。\allowbreak{}它的边长就是两条平行线间的距离。\allowbreak{}\par};
\end{scope}\end{tikzpicture}\newpage
\definecolor{answer}{rgb}{0.850980,0.047059,0.070588}
\null\begin{tikzpicture}[remember picture,overlay]\begin{scope}[shift={(current page.north west)},x=1bp,y=-1bp]
\node[anchor=north west,inner sep=0pt,text=answer,align=left,text width=122bp] at (15,15) {\fontsize{7.5bp}{9.5bp}\selectfont 通常算法与运算性质简算均可，\allowbreak{}保留两种思路。\allowbreak{}\par};
\end{scope}\end{tikzpicture}\newpage
\definecolor{answer}{rgb}{0.850980,0.047059,0.070588}
\null\begin{tikzpicture}[remember picture,overlay]\begin{scope}[shift={(current page.north west)},x=1bp,y=-1bp]
\node[anchor=north west,inner sep=0pt,text=answer,align=left,text width=126bp] at (15,15) {\fontsize{7.5bp}{9.5bp}\selectfont 通常算法与运算性质简算均可，\allowbreak{}保留两种思路。\allowbreak{}\par};
\end{scope}\end{tikzpicture}\newpage
\definecolor{answer}{rgb}{0.850980,0.047059,0.070588}
\null\begin{tikzpicture}[remember picture,overlay]\begin{scope}[shift={(current page.north west)},x=1bp,y=-1bp]
\node[anchor=north west,inner sep=0pt,text=answer,align=left,text width=122bp] at (15,15) {\fontsize{7.5bp}{9.5bp}\selectfont 通常算法与运算性质简算均可，\allowbreak{}保留两种思路。\allowbreak{}\par};
\end{scope}\end{tikzpicture}\newpage
\definecolor{answer}{rgb}{0.850980,0.047059,0.070588}
\null\begin{tikzpicture}[remember picture,overlay]\begin{scope}[shift={(current page.north west)},x=1bp,y=-1bp]
\node[anchor=north west,inner sep=0pt,outer sep=0pt,text=answer,align=left,text width=500bp] at (15,15) {\fontsize{7.5bp}{12bp}\selectfont 析：交换$2.9$与$4.7$的位置，先算$5.3+4.7=10$，再算$10+2.9=12.9$，运用加法交换律。\par};
\end{scope}\end{tikzpicture}\newpage
\definecolor{answer}{rgb}{0.850980,0.047059,0.070588}
\null\begin{tikzpicture}[remember picture,overlay]\begin{scope}[shift={(current page.north west)},x=1bp,y=-1bp]
\node[anchor=north west,inner sep=0pt,text=answer,align=left,text width=483.0bp] at (15,15) {\fontsize{7.5bp}{9.5bp}\selectfont 析：\allowbreak{}保留一位看百分位；\allowbreak{}\ensuremath{7.049}可舍成\ensuremath{7.0},\allowbreak{} \ensuremath{7.050}会进位为\ensuremath{7.1},\allowbreak{}所以最大\ensuremath{7.049}。\allowbreak{}\par};
\end{scope}\end{tikzpicture}\newpage
\definecolor{answer}{rgb}{0.850980,0.047059,0.070588}
\null\begin{tikzpicture}[remember picture,overlay]\begin{scope}[shift={(current page.north west)},x=1bp,y=-1bp]
\node[anchor=north west,inner sep=0pt,text=answer,align=left,text width=483.0bp] at (15,15) {\fontsize{7.5bp}{9.5bp}\selectfont 析：\allowbreak{}交换律交换两加数，\allowbreak{}\ensuremath{a} + \ensuremath{b} = \ensuremath{b} + \ensuremath{a};\allowbreak{}分配律把同一因数分别乘两加数，\allowbreak{}(\ensuremath{a} + \ensuremath{b}) × \ensuremath{c} = \ensuremath{a} × \ensuremath{c} + \ensuremath{b} × \ensuremath{c}。\allowbreak{}\par};
\end{scope}\end{tikzpicture}\newpage
\definecolor{answer}{rgb}{0.850980,0.047059,0.070588}
\null\begin{tikzpicture}[remember picture,overlay]\begin{scope}[shift={(current page.north west)},x=1bp,y=-1bp]
\node[anchor=north west,inner sep=0pt,text=answer,align=left,text width=483.0bp] at (15,15) {\fontsize{7.5bp}{9.5bp}\selectfont 析：\allowbreak{}\ensuremath{0}～\ensuremath{0.1}分\ensuremath{10}小格，\allowbreak{}每格\ensuremath{0.01};\allowbreak{} \ensuremath{0.04}、\allowbreak{}\ensuremath{0.15}、\allowbreak{}\ensuremath{0.28}分别在\ensuremath{0}、\allowbreak{}\ensuremath{0.1}、\allowbreak{}\ensuremath{0.2}右第\ensuremath{4}、\allowbreak{}\ensuremath{5}、\allowbreak{}\ensuremath{8}格，\allowbreak{}\ensuremath{0.4}是第\ensuremath{40}格。\allowbreak{}\par};
\end{scope}\end{tikzpicture}\newpage
\definecolor{answer}{rgb}{0.850980,0.047059,0.070588}
\null\begin{tikzpicture}[remember picture,overlay]\begin{scope}[shift={(current page.north west)},x=1bp,y=-1bp]
\node[anchor=north west,inner sep=0pt,text=answer,align=left,text width=483.0bp] at (15,15) {\fontsize{7.5bp}{9.5bp}\selectfont 析：\allowbreak{}保留两位看千分位\ensuremath{8},\allowbreak{}满\ensuremath{5}进\ensuremath{1},\allowbreak{} \ensuremath{299.988}约为\ensuremath{299.99}。\allowbreak{}\par};
\end{scope}\end{tikzpicture}\newpage
\definecolor{answer}{rgb}{0.850980,0.047059,0.070588}
\null\begin{tikzpicture}[remember picture,overlay]\begin{scope}[shift={(current page.north west)},x=1bp,y=-1bp]
\node[anchor=north west,inner sep=0pt,text=answer,align=left,text width=483.0bp] at (15,15) {\fontsize{7.5bp}{9.5bp}\selectfont 析：\allowbreak{}除以一亿得\ensuremath{14.295}亿；\allowbreak{}保留两位看千分位\ensuremath{5},\allowbreak{}进位得\ensuremath{14.30}亿，\allowbreak{}末尾\ensuremath{0}须保留。\allowbreak{}\par};
\end{scope}\end{tikzpicture}\newpage
\definecolor{answer}{rgb}{0.850980,0.047059,0.070588}
\null\begin{tikzpicture}[remember picture,overlay]\begin{scope}[shift={(current page.north west)},x=1bp,y=-1bp]
\node[anchor=north west,inner sep=0pt,text=answer,align=left,text width=483.0bp] at (15,15) {\fontsize{7.5bp}{9.5bp}\selectfont 析：\allowbreak{}小数末尾添\ensuremath{0}或去\ensuremath{0}只改变小数位数，\allowbreak{}不改变大小。\allowbreak{}\par};
\end{scope}\end{tikzpicture}\newpage
\definecolor{answer}{rgb}{0.850980,0.047059,0.070588}
\null\begin{tikzpicture}[remember picture,overlay]\begin{scope}[shift={(current page.north west)},x=1bp,y=-1bp]
\node[anchor=north west,inner sep=0pt,text=answer,align=left,text width=483.0bp] at (15,15) {\fontsize{7.5bp}{9.5bp}\selectfont 析：\allowbreak{}\ensuremath{125}乘\ensuremath{8}与\ensuremath{4}的积不能改为两个积相加；\allowbreak{}分配律对应括号里的加减法。\allowbreak{}\par};
\end{scope}\end{tikzpicture}\newpage
\definecolor{answer}{rgb}{0.850980,0.047059,0.070588}
\null\begin{tikzpicture}[remember picture,overlay]\begin{scope}[shift={(current page.north west)},x=1bp,y=-1bp]
\node[anchor=north west,inner sep=0pt,text=answer,align=left,text width=483.0bp] at (15,15) {\fontsize{7.5bp}{9.5bp}\selectfont 析：\allowbreak{}小数点对齐，\allowbreak{}个位、\allowbreak{}十分位、\allowbreak{}百分位等相同数位就对齐。\allowbreak{}\par};
\end{scope}\end{tikzpicture}\newpage
\definecolor{answer}{rgb}{0.850980,0.047059,0.070588}
\null\begin{tikzpicture}[remember picture,overlay]\begin{scope}[shift={(current page.north west)},x=1bp,y=-1bp]
\node[anchor=north west,inner sep=0pt,text=answer,align=left,text width=483.0bp] at (15,15) {\fontsize{7.5bp}{9.5bp}\selectfont 析：\allowbreak{}同一平面内，\allowbreak{}都垂直于同一条直线的两条直线互相平行。\allowbreak{}\par};
\end{scope}\end{tikzpicture}\newpage
\definecolor{answer}{rgb}{0.850980,0.047059,0.070588}
\null\begin{tikzpicture}[remember picture,overlay]\begin{scope}[shift={(current page.north west)},x=1bp,y=-1bp]
\node[anchor=north west,inner sep=0pt,text=answer,align=left,text width=483.0bp] at (15,15) {\fontsize{7.5bp}{9.5bp}\selectfont 析：\allowbreak{}C两式都是乘除同级运算，\allowbreak{}均从左到右；\allowbreak{}A、\allowbreak{}B两组的先后运算不同，\allowbreak{}选C。\allowbreak{}\par};
\end{scope}\end{tikzpicture}\newpage
\definecolor{answer}{rgb}{0.850980,0.047059,0.070588}
\null\begin{tikzpicture}[remember picture,overlay]\begin{scope}[shift={(current page.north west)},x=1bp,y=-1bp]
\node[anchor=north west,inner sep=0pt,text=answer,align=left,text width=483.0bp] at (15,15) {\fontsize{7.5bp}{9.5bp}\selectfont 析：\allowbreak{}\ensuremath{75} × \ensuremath{99}可拆为\ensuremath{75} × (\ensuremath{100} \ensuremath{-}\ensuremath{1}),\allowbreak{}再用分配律得到\ensuremath{75} × \ensuremath{100} \ensuremath{-}\ensuremath{75},\allowbreak{}选B。\allowbreak{}\par};
\end{scope}\end{tikzpicture}\newpage
\definecolor{answer}{rgb}{0.850980,0.047059,0.070588}
\null\begin{tikzpicture}[remember picture,overlay]\begin{scope}[shift={(current page.north west)},x=1bp,y=-1bp]
\node[anchor=north west,inner sep=0pt,text=answer,align=left,text width=483.0bp] at (15,15) {\fontsize{7.5bp}{9.5bp}\selectfont 析：\allowbreak{}相向而行用速度和，\allowbreak{}\ensuremath{480} ÷ (\ensuremath{65} + \ensuremath{55}) = \ensuremath{4}小时，\allowbreak{}选A。\allowbreak{}\par};
\end{scope}\end{tikzpicture}\newpage
\definecolor{answer}{rgb}{0.850980,0.047059,0.070588}
\null\begin{tikzpicture}[remember picture,overlay]\begin{scope}[shift={(current page.north west)},x=1bp,y=-1bp]
\node[anchor=north west,inner sep=0pt,text=answer,align=left,text width=483.0bp] at (15,15) {\fontsize{7.5bp}{9.5bp}\selectfont 析：\allowbreak{}每个等宽长条占\ensuremath{0.2},\allowbreak{}涂了\ensuremath{2}整条和半条，\allowbreak{}共\ensuremath{0.4} + \ensuremath{0.1} = \ensuremath{0.5},\allowbreak{}选C。\allowbreak{}\par};
\end{scope}\end{tikzpicture}\newpage
\definecolor{answer}{rgb}{0.850980,0.047059,0.070588}
\null\begin{tikzpicture}[remember picture,overlay]\begin{scope}[shift={(current page.north west)},x=1bp,y=-1bp]
\node[anchor=north west,inner sep=0pt,text=answer,align=left,text width=483.0bp] at (15,15) {\fontsize{7.5bp}{9.5bp}\selectfont 析：\allowbreak{}六个月是两个季度，\allowbreak{}平均每季度应把六个月总费用除以\ensuremath{2},\allowbreak{}选C。\allowbreak{}\par};
\end{scope}\end{tikzpicture}\newpage
\definecolor{answer}{rgb}{0.850980,0.047059,0.070588}
\null\begin{tikzpicture}[remember picture,overlay]\begin{scope}[shift={(current page.north west)},x=1bp,y=-1bp]
\node[anchor=north west,inner sep=0pt,text=answer,align=left,text width=483.0bp] at (15,15) {\fontsize{7.5bp}{9.5bp}\selectfont 析：\allowbreak{}端点不计入，\allowbreak{}只有\ensuremath{0.7}、\allowbreak{}\ensuremath{0.8}两个一位小数，\allowbreak{}选B。\allowbreak{}\par};
\end{scope}\end{tikzpicture}\newpage
\definecolor{answer}{rgb}{0.850980,0.047059,0.070588}
\null\begin{tikzpicture}[remember picture,overlay]\begin{scope}[shift={(current page.north west)},x=1bp,y=-1bp]
\node[anchor=north west,inner sep=0pt,text=answer,align=left,text width=483.0bp] at (15,15) {\fontsize{8.5bp}{10.62bp}\selectfont 丽丽比佳佳高\ensuremath{0.05}米要加；\allowbreak{}飞飞比丽丽矮\ensuremath{0.1}米要减。\allowbreak{}\par};
\end{scope}\end{tikzpicture}\newpage
\definecolor{answer}{rgb}{0.850980,0.047059,0.070588}
\null\begin{tikzpicture}[remember picture,overlay]\begin{scope}[shift={(current page.north west)},x=1bp,y=-1bp]
\node[anchor=north west,inner sep=0pt,text=answer,align=left,text width=483.0bp] at (15,15) {\fontsize{7.5bp}{9.5bp}\selectfont 析：\allowbreak{}同平面内不相交就互相平行；\allowbreak{}两线垂直时交点叫垂足。\allowbreak{}\par};
\end{scope}\end{tikzpicture}\newpage
\definecolor{answer}{rgb}{0.850980,0.047059,0.070588}
\null\begin{tikzpicture}[remember picture,overlay]\begin{scope}[shift={(current page.north west)},x=1bp,y=-1bp]
\node[anchor=north west,inner sep=0pt,text=answer,align=left,text width=483.0bp] at (15,15) {\fontsize{7.5bp}{9.5bp}\selectfont 析：\allowbreak{}等腰梯形两腰相等，\allowbreak{}过两底中点有\ensuremath{1}条对称轴。\allowbreak{}\par};
\end{scope}\end{tikzpicture}\newpage
\definecolor{answer}{rgb}{0.850980,0.047059,0.070588}
\null\begin{tikzpicture}[remember picture,overlay]\begin{scope}[shift={(current page.north west)},x=1bp,y=-1bp]
\node[anchor=north west,inner sep=0pt,text=answer,align=left,text width=483.125bp] at (15,15) {\fontsize{7.52bp}{9.52bp}\selectfont 析：\allowbreak{}第三角\ensuremath{180}\ensuremath{^\circ}\ensuremath{-}(\ensuremath{35}\ensuremath{^\circ}+ \ensuremath{65}\ensuremath{^\circ}) = \ensuremath{180}\ensuremath{^\circ}\ensuremath{-}\ensuremath{100}\ensuremath{^\circ}= \ensuremath{80}\ensuremath{^\circ}；\allowbreak{}三个角都小于\ensuremath{90}\ensuremath{^\circ}，\allowbreak{}为锐角三角形\par};
\end{scope}\end{tikzpicture}\newpage
\definecolor{answer}{rgb}{0.850980,0.047059,0.070588}
\null\begin{tikzpicture}[remember picture,overlay]\begin{scope}[shift={(current page.north west)},x=1bp,y=-1bp]
\node[anchor=north west,inner sep=0pt,text=answer,align=left,text width=483.0bp] at (15,15) {\fontsize{7.5bp}{9.5bp}\selectfont 析：\allowbreak{}旋转后形状和大小不变，\allowbreak{}位置、\allowbreak{}方向改变。\allowbreak{}\par};
\end{scope}\end{tikzpicture}\newpage
\definecolor{answer}{rgb}{0.850980,0.047059,0.070588}
\null\begin{tikzpicture}[remember picture,overlay]\begin{scope}[shift={(current page.north west)},x=1bp,y=-1bp]
\node[anchor=north west,inner sep=0pt,text=answer,align=left,text width=483.0bp] at (15,15) {\fontsize{7.5bp}{9.5bp}\selectfont 析：\allowbreak{}三角形有\ensuremath{3}条高；\allowbreak{}直角三角形两条直角边互相垂直，\allowbreak{}正好也是其中两条高。\allowbreak{}\par};
\end{scope}\end{tikzpicture}\newpage
\definecolor{answer}{rgb}{0.850980,0.047059,0.070588}
\null\begin{tikzpicture}[remember picture,overlay]\begin{scope}[shift={(current page.north west)},x=1bp,y=-1bp]
\node[anchor=north west,inner sep=0pt,text=answer,align=left,text width=483.0bp] at (15,15) {\fontsize{7.5bp}{9.5bp}\selectfont 析：\allowbreak{}顶角\ensuremath{80}\ensuremath{^\circ}时底角(\ensuremath{180} \ensuremath{-}\ensuremath{80}) ÷ \ensuremath{2} = \ensuremath{50}\ensuremath{^\circ};\allowbreak{}底角\ensuremath{80}\ensuremath{^\circ}时顶角\ensuremath{180} \ensuremath{-}\ensuremath{80} \ensuremath{-}\ensuremath{80} = \ensuremath{20}\ensuremath{^\circ}。\allowbreak{}\par};
\end{scope}\end{tikzpicture}\newpage
\definecolor{answer}{rgb}{0.850980,0.047059,0.070588}
\null\begin{tikzpicture}[remember picture,overlay]\begin{scope}[shift={(current page.north west)},x=1bp,y=-1bp]
\node[anchor=north west,inner sep=0pt,text=answer,align=left,text width=483.0bp] at (15,15) {\fontsize{7.5bp}{9.5bp}\selectfont 析：\allowbreak{}最大角\ensuremath{89}\ensuremath{^\circ},\allowbreak{}其余也小于\ensuremath{90}\ensuremath{^\circ},\allowbreak{}是锐角三角形；\allowbreak{}最大角\ensuremath{91}\ensuremath{^\circ},\allowbreak{}是钝角三角形。\allowbreak{}\par};
\end{scope}\end{tikzpicture}\newpage
\definecolor{answer}{rgb}{0.850980,0.047059,0.070588}
\null\begin{tikzpicture}[remember picture,overlay]\begin{scope}[shift={(current page.north west)},x=1bp,y=-1bp]
\node[anchor=north west,inner sep=0pt,text=answer,align=left,text width=483.0bp] at (15,15) {\fontsize{7.5bp}{9.5bp}\selectfont 析：\allowbreak{}每边\ensuremath{24} ÷ \ensuremath{3} = \ensuremath{8}厘米，\allowbreak{}每角\ensuremath{180} ÷ \ensuremath{3} = \ensuremath{60}\ensuremath{^\circ}。\allowbreak{}\par};
\end{scope}\end{tikzpicture}\newpage
\definecolor{answer}{rgb}{0.850980,0.047059,0.070588}
\null\begin{tikzpicture}[remember picture,overlay]\begin{scope}[shift={(current page.north west)},x=1bp,y=-1bp]
\node[anchor=north west,inner sep=0pt,text=answer,align=left,text width=483.0bp] at (15,15) {\fontsize{7.5bp}{9.5bp}\selectfont 析：\allowbreak{}三角形有\ensuremath{3}条高，\allowbreak{}但平行四边形可作无数条高，\allowbreak{}“四条”错。\allowbreak{}\par};
\end{scope}\end{tikzpicture}\newpage
\definecolor{answer}{rgb}{0.850980,0.047059,0.070588}
\null\begin{tikzpicture}[remember picture,overlay]\begin{scope}[shift={(current page.north west)},x=1bp,y=-1bp]
\node[anchor=north west,inner sep=0pt,text=answer,align=left,text width=483.0bp] at (15,15) {\fontsize{7.5bp}{9.5bp}\selectfont 析：\allowbreak{}梯形须有且只有一组对边平行；\allowbreak{}只说“有一组”也包括平行四边形。\allowbreak{}\par};
\end{scope}\end{tikzpicture}\newpage
\definecolor{answer}{rgb}{0.850980,0.047059,0.070588}
\null\begin{tikzpicture}[remember picture,overlay]\begin{scope}[shift={(current page.north west)},x=1bp,y=-1bp]
\node[anchor=north west,inner sep=0pt,text=answer,align=left,text width=483.0bp] at (15,15) {\fontsize{7.5bp}{9.5bp}\selectfont 析：\allowbreak{}等边一定等腰，\allowbreak{}但一般等腰三角形的底不必等于腰，\allowbreak{}所以整句错。\allowbreak{}\par};
\end{scope}\end{tikzpicture}\newpage
\definecolor{answer}{rgb}{0.850980,0.047059,0.070588}
\null\begin{tikzpicture}[remember picture,overlay]\begin{scope}[shift={(current page.north west)},x=1bp,y=-1bp]
\node[anchor=north west,inner sep=0pt,text=answer,align=left,text width=483.0bp] at (15,15) {\fontsize{7.5bp}{9.5bp}\selectfont 析：\allowbreak{}准确说法是对应点的平移距离相等，\allowbreak{}原句笼统说两图距离相等不准确。\allowbreak{}\par};
\end{scope}\end{tikzpicture}\newpage
\definecolor{answer}{rgb}{0.850980,0.047059,0.070588}
\null\begin{tikzpicture}[remember picture,overlay]\begin{scope}[shift={(current page.north west)},x=1bp,y=-1bp]
\node[anchor=north west,inner sep=0pt,text=answer,align=left,text width=483.0bp] at (15,15) {\fontsize{7.5bp}{9.5bp}\selectfont 析：\allowbreak{}两全同直角三角形可沿斜边拼长方形，\allowbreak{}也可沿直角边拼平行四边形。\allowbreak{}\par};
\end{scope}\end{tikzpicture}\newpage
\definecolor{answer}{rgb}{0.850980,0.047059,0.070588}
\null\begin{tikzpicture}[remember picture,overlay]\begin{scope}[shift={(current page.north west)},x=1bp,y=-1bp]
\node[anchor=north west,inner sep=0pt,text=answer,align=left,text width=483.0bp] at (15,15) {\fontsize{7.5bp}{9.5bp}\selectfont 析：\allowbreak{}任意大小、\allowbreak{}形状的三角形，\allowbreak{}内角和都为\ensuremath{180}\ensuremath{^\circ}。\allowbreak{}\par};
\end{scope}\end{tikzpicture}\newpage
\definecolor{answer}{rgb}{0.850980,0.047059,0.070588}
\null\begin{tikzpicture}[remember picture,overlay]\begin{scope}[shift={(current page.north west)},x=1bp,y=-1bp]
\node[anchor=north west,inner sep=0pt,text=answer,align=left,text width=483.0bp] at (15,15) {\fontsize{7.5bp}{9.5bp}\selectfont 析：\allowbreak{}\ensuremath{3}、\allowbreak{}\ensuremath{4}、\allowbreak{}\ensuremath{5}互不相等，\allowbreak{}按边是“不等边”三角形，\allowbreak{}选C。\allowbreak{}\par};
\end{scope}\end{tikzpicture}\newpage
\definecolor{answer}{rgb}{0.850980,0.047059,0.070588}
\null\begin{tikzpicture}[remember picture,overlay]\begin{scope}[shift={(current page.north west)},x=1bp,y=-1bp]
\node[anchor=north west,inner sep=0pt,text=answer,align=left,text width=483.0bp] at (15,15) {\fontsize{7.5bp}{9.5bp}\selectfont 析：\allowbreak{}最短两边和须大于最长边；\allowbreak{}\ensuremath{5} + \ensuremath{6} > \ensuremath{7}可以，\allowbreak{}\ensuremath{1} + \ensuremath{2} = \ensuremath{3}、\allowbreak{}\ensuremath{3} + \ensuremath{4} < \ensuremath{8}不行，\allowbreak{}选A。\allowbreak{}\par};
\end{scope}\end{tikzpicture}\newpage
\definecolor{answer}{rgb}{0.850980,0.047059,0.070588}
\null\begin{tikzpicture}[remember picture,overlay]\begin{scope}[shift={(current page.north west)},x=1bp,y=-1bp]
\node[anchor=north west,inner sep=0pt,text=answer,align=left,text width=483.0bp] at (15,15) {\fontsize{7.5bp}{9.5bp}\selectfont 析：\allowbreak{}椭圆、\allowbreak{}等腰三角形、\allowbreak{}箭头图有对称轴，\allowbreak{}共\ensuremath{3}个，\allowbreak{}选B。\allowbreak{}\par};
\end{scope}\end{tikzpicture}\newpage
\definecolor{answer}{rgb}{0.850980,0.047059,0.070588}
\null\begin{tikzpicture}[remember picture,overlay]\begin{scope}[shift={(current page.north west)},x=1bp,y=-1bp]
\node[anchor=north west,inner sep=0pt,text=answer,align=left,text width=483.0bp] at (15,15) {\fontsize{7.5bp}{9.5bp}\selectfont 析：\allowbreak{}两个钝角之和已超过\ensuremath{180}\ensuremath{^\circ},\allowbreak{}所以三角形最多有\ensuremath{1}个钝角，\allowbreak{}选A。\allowbreak{}\par};
\end{scope}\end{tikzpicture}\newpage
\definecolor{answer}{rgb}{0.850980,0.047059,0.070588}
\null\begin{tikzpicture}[remember picture,overlay]\begin{scope}[shift={(current page.north west)},x=1bp,y=-1bp]
\node[anchor=north west,inner sep=0pt,text=answer,align=left,text width=483.0bp] at (15,15) {\fontsize{7.5bp}{9.5bp}\selectfont 析：\allowbreak{}两条平行底边间距离处处相等，\allowbreak{}能作无数条等长高，\allowbreak{}选A;\allowbreak{}梯形有且只有一组对边平行。\allowbreak{}\par};
\end{scope}\end{tikzpicture}\newpage
\definecolor{answer}{rgb}{0.850980,0.047059,0.070588}
\null\begin{tikzpicture}[remember picture,overlay]\begin{scope}[shift={(current page.north west)},x=1bp,y=-1bp]
\node[anchor=north west,inner sep=0pt,text=answer,align=left,text width=483.0bp] at (15,15) {\fontsize{7.5bp}{9.5bp}\selectfont 析：\allowbreak{}平均数把一组总量平均分配，\allowbreak{}能反映整体水平。\allowbreak{}\par};
\end{scope}\end{tikzpicture}\newpage
\definecolor{answer}{rgb}{0.850980,0.047059,0.070588}
\null\begin{tikzpicture}[remember picture,overlay]\begin{scope}[shift={(current page.north west)},x=1bp,y=-1bp]
\node[anchor=north west,inner sep=0pt,text=answer,align=left,text width=483.0bp] at (15,15) {\fontsize{7.5bp}{9.5bp}\selectfont 析：\allowbreak{}平均数用全部数据的和除以个数，\allowbreak{}每个数据都必须参加。\allowbreak{}\par};
\end{scope}\end{tikzpicture}\newpage
\definecolor{answer}{rgb}{0.850980,0.047059,0.070588}
\null\begin{tikzpicture}[remember picture,overlay]\begin{scope}[shift={(current page.north west)},x=1bp,y=-1bp]
\node[anchor=north west,inner sep=0pt,text=answer,align=left,text width=483.0bp] at (15,15) {\fontsize{7.5bp}{9.5bp}\selectfont 析：\allowbreak{}班级平均身高较高，\allowbreak{}不能说明每个学生都比另一班某人高。\allowbreak{}\par};
\end{scope}\end{tikzpicture}\newpage
\definecolor{answer}{rgb}{0.850980,0.047059,0.070588}
\null\begin{tikzpicture}[remember picture,overlay]\begin{scope}[shift={(current page.north west)},x=1bp,y=-1bp]
\node[anchor=north west,inner sep=0pt,text=answer,align=left,text width=483.0bp] at (15,15) {\fontsize{7.5bp}{9.5bp}\selectfont 析：\allowbreak{}总数是平均每人数乘人数；\allowbreak{}平均相同而人数不同，\allowbreak{}总棵数未必相同。\allowbreak{}\par};
\end{scope}\end{tikzpicture}\newpage
\definecolor{answer}{rgb}{0.850980,0.047059,0.070588}
\null\begin{tikzpicture}[remember picture,overlay]\begin{scope}[shift={(current page.north west)},x=1bp,y=-1bp]
\node[anchor=north west,inner sep=0pt,text=answer,align=left,text width=483.0bp] at (15,15) {\fontsize{7.5bp}{9.5bp}\selectfont 析：\allowbreak{}\ensuremath{5}个\ensuremath{20}与\ensuremath{1}个\ensuremath{30}总和\ensuremath{130},\allowbreak{}平均不可能超过最大数\ensuremath{30},\allowbreak{}更不可能\ensuremath{35}。\allowbreak{}\par};
\end{scope}\end{tikzpicture}\newpage
\definecolor{answer}{rgb}{0.850980,0.047059,0.070588}
\null\begin{tikzpicture}[remember picture,overlay]\begin{scope}[shift={(current page.north west)},x=1bp,y=-1bp]
\node[anchor=north west,inner sep=0pt,text=answer,align=left,text width=483.0bp] at (15,15) {\fontsize{7.5bp}{9.5bp}\selectfont 析：\allowbreak{}八数和\ensuremath{128},\allowbreak{} \ensuremath{128} ÷ \ensuremath{8} = \ensuremath{16},\allowbreak{}选B。\allowbreak{}\par};
\end{scope}\end{tikzpicture}\newpage
\definecolor{answer}{rgb}{0.850980,0.047059,0.070588}
\null\begin{tikzpicture}[remember picture,overlay]\begin{scope}[shift={(current page.north west)},x=1bp,y=-1bp]
\node[anchor=north west,inner sep=0pt,text=answer,align=left,text width=482.9988708496094bp] at (15,15) {\fontsize{7.5bp}{9.5bp}\selectfont 析：\allowbreak{}\ensuremath{100} ÷ \ensuremath{5} = \ensuremath{20}道是实际每天；\allowbreak{}计划每天比实际少\ensuremath{2}道，\allowbreak{}\ensuremath{100} ÷ \ensuremath{5} \ensuremath{-}\ensuremath{2} = \ensuremath{18}道，\allowbreak{}选A。\allowbreak{}\par};
\end{scope}\end{tikzpicture}\newpage
\definecolor{answer}{rgb}{0.850980,0.047059,0.070588}
\null\begin{tikzpicture}[remember picture,overlay]\begin{scope}[shift={(current page.north west)},x=1bp,y=-1bp]
\node[anchor=north west,inner sep=0pt,text=answer,align=left,text width=482.9988708496094bp] at (15,15) {\fontsize{7.5bp}{9.5bp}\selectfont 析：\allowbreak{}按每人\ensuremath{29}千克计，\allowbreak{}男生总重多\ensuremath{1}×男生人数，\allowbreak{}女生总重少\ensuremath{1}×女生人数；\allowbreak{}男生人数多，\allowbreak{}全班平均>\ensuremath{29}千克，\allowbreak{}选B。\allowbreak{}\par};
\end{scope}\end{tikzpicture}\newpage
\definecolor{answer}{rgb}{0.850980,0.047059,0.070588}
\null\begin{tikzpicture}[remember picture,overlay]\begin{scope}[shift={(current page.north west)},x=1bp,y=-1bp]
\node[anchor=north west,inner sep=0pt,text=answer,align=left,text width=483.0bp] at (15,15) {\fontsize{7.5bp}{9.5bp}\selectfont 析：\allowbreak{}六个月是\ensuremath{2}个季度，\allowbreak{}应把六个月总费用除以\ensuremath{2},\allowbreak{}选C。\allowbreak{}\par};
\end{scope}\end{tikzpicture}\newpage
\definecolor{answer}{rgb}{0.850980,0.047059,0.070588}
\null\begin{tikzpicture}[remember picture,overlay]\begin{scope}[shift={(current page.north west)},x=1bp,y=-1bp]
\node[anchor=north west,inner sep=0pt,text=answer,align=left,text width=483.0bp] at (15,15) {\fontsize{7.5bp}{9.5bp}\selectfont 析：\allowbreak{}男女各\ensuremath{6}人，\allowbreak{}人数相等，\allowbreak{}平均(\ensuremath{72} + \ensuremath{82}) ÷ \ensuremath{2} = \ensuremath{77}岁，\allowbreak{}选C。\allowbreak{}\par};
\end{scope}\end{tikzpicture}\newpage
\definecolor{answer}{rgb}{0.850980,0.047059,0.070588}
\null\begin{tikzpicture}[remember picture,overlay]\begin{scope}[shift={(current page.north west)},x=1bp,y=-1bp]
\node[anchor=north west,inner sep=0pt,text=answer,align=left,text width=483.0bp] at (15,15) {\fontsize{7.5bp}{9.5bp}\selectfont 析：\allowbreak{}按\ensuremath{100}、\allowbreak{}\ensuremath{90}～\ensuremath{99}、\allowbreak{}\ensuremath{80}～\ensuremath{89}、\allowbreak{}\ensuremath{70}～\ensuremath{79}、\allowbreak{}\ensuremath{60}～\ensuremath{69}、\allowbreak{}\ensuremath{60}以下分类，\allowbreak{}人数依次\ensuremath{2}、\allowbreak{}\ensuremath{16}、\allowbreak{}\ensuremath{8}、\allowbreak{}\ensuremath{2}、\allowbreak{}\ensuremath{1}、\allowbreak{}\ensuremath{1},\allowbreak{}合计\ensuremath{30}。\allowbreak{}\par};
\end{scope}\end{tikzpicture}\newpage
\definecolor{answer}{rgb}{0.850980,0.047059,0.070588}
\null\begin{tikzpicture}[remember picture,overlay]\begin{scope}[shift={(current page.north west)},x=1bp,y=-1bp]
\node[anchor=north west,inner sep=0pt,text=answer,align=left,text width=483.0bp] at (15,15) {\fontsize{7.5bp}{9.5bp}\selectfont 析：\allowbreak{}最多\ensuremath{90}～\ensuremath{99}分段\ensuremath{16}人；\allowbreak{}优秀加\ensuremath{100}分的\ensuremath{2}人，\allowbreak{}共\ensuremath{18};\allowbreak{}良好\ensuremath{8};\allowbreak{}中等合并\ensuremath{2} + \ensuremath{1} = \ensuremath{3};\allowbreak{} \ensuremath{60}分以下\ensuremath{1}人。\allowbreak{}\par};
\end{scope}\end{tikzpicture}\newpage
\definecolor{answer}{rgb}{0.850980,0.047059,0.070588}
\null\begin{tikzpicture}[remember picture,overlay]\begin{scope}[shift={(current page.north west)},x=1bp,y=-1bp]
\node[anchor=north west,inner sep=0pt,text=answer,align=left,text width=483.0bp] at (15,15) {\fontsize{8.5bp}{10.62bp}\selectfont 男女平均身高对应人数不同，\allowbreak{}不能直接把\ensuremath{145}与\ensuremath{141}平均；\allowbreak{}应先求全队总身高，\allowbreak{}再除以总人数。\allowbreak{}\par};
\end{scope}\end{tikzpicture}\newpage
\definecolor{answer}{rgb}{0.850980,0.047059,0.070588}
\null\begin{tikzpicture}[remember picture,overlay]\begin{scope}[shift={(current page.north west)},x=1bp,y=-1bp]
\node[anchor=north west,inner sep=0pt,text=answer,align=left,text width=483.0bp] at (15,15) {\fontsize{8.5bp}{10.62bp}\selectfont 上、\allowbreak{}下山路程相同，\allowbreak{}总路程是上山路程的\ensuremath{2}倍，\allowbreak{}总时间为\ensuremath{18}加\ensuremath{12};\allowbreak{}两段用时不同，\allowbreak{}不能直接平均两个速度。\allowbreak{}\par};
\end{scope}\end{tikzpicture}\newpage
\definecolor{answer}{rgb}{0.850980,0.047059,0.070588}
\null\begin{tikzpicture}[remember picture,overlay]\begin{scope}[shift={(current page.north west)},x=1bp,y=-1bp]
\node[anchor=north west,inner sep=0pt,text=answer,align=left,text width=453.36700439453125bp] at (15,15) {\fontsize{8.5bp}{10.62bp}\selectfont 先用去掉\ensuremath{88}个者前后的平均差推人数，\allowbreak{}再使其余较慢同学都只做最低\ensuremath{70}个，\allowbreak{}把剩余份额尽量集中给最快的同学。\allowbreak{}\par};
\end{scope}\end{tikzpicture}\newpage
\definecolor{answer}{rgb}{0.850980,0.047059,0.070588}
\null\begin{tikzpicture}[remember picture,overlay]\begin{scope}[shift={(current page.north west)},x=1bp,y=-1bp]
\node[anchor=north west,inner sep=0pt,text=answer,align=left,text width=338.5bp] at (15,15) {\fontsize{7.5bp}{9.5bp}\selectfont 析：\allowbreak{}同级从左到右；\allowbreak{}先乘除后加减；\allowbreak{}有括号先小括号，\allowbreak{}再中括号。\allowbreak{}\par};
\end{scope}\end{tikzpicture}\newpage
\definecolor{answer}{rgb}{0.850980,0.047059,0.070588}
\null\begin{tikzpicture}[remember picture,overlay]\begin{scope}[shift={(current page.north west)},x=1bp,y=-1bp]
\node[anchor=north west,inner sep=0pt,text=answer,align=left,text width=366.906494140625bp] at (15,15) {\fontsize{7.5bp}{9.5bp}\selectfont 析：\allowbreak{}分配律：\allowbreak{}(\ensuremath{a} + \ensuremath{b}) × \ensuremath{c} = \ensuremath{a} × \ensuremath{c} + \ensuremath{b} × \ensuremath{c};\allowbreak{}结合律：\allowbreak{}(\ensuremath{a} + \ensuremath{b}) + \ensuremath{c} = \ensuremath{a} + (\ensuremath{b} + \ensuremath{c})。\allowbreak{}\par};
\end{scope}\end{tikzpicture}\newpage
\definecolor{answer}{rgb}{0.850980,0.047059,0.070588}
\null\begin{tikzpicture}[remember picture,overlay]\begin{scope}[shift={(current page.north west)},x=1bp,y=-1bp]
\node[anchor=north west,inner sep=0pt,outer sep=0pt,text=answer,align=left,text width=500bp] at (15,15) {\fontsize{7.5bp}{12bp}\selectfont 析：交换$6.8$与$5.5$的位置，先算$4.5+5.5=10$，再算$10+6.8=16.8$，运用加法交换律。\par};
\end{scope}\end{tikzpicture}\newpage
\definecolor{answer}{rgb}{0.850980,0.047059,0.070588}
\null\begin{tikzpicture}[remember picture,overlay]\begin{scope}[shift={(current page.north west)},x=1bp,y=-1bp]
\node[anchor=north west,inner sep=0pt,text=answer,align=left,text width=322.2799987792969bp] at (15,15) {\fontsize{7.5bp}{9.5bp}\selectfont 析：\allowbreak{}\ensuremath{0.85}是\ensuremath{85}个\ensuremath{0.01};\allowbreak{} \ensuremath{1}是\ensuremath{100}个\ensuremath{0.01},\allowbreak{}还差\ensuremath{100} \ensuremath{-}\ensuremath{85} = \ensuremath{15}个。\allowbreak{}\par};
\end{scope}\end{tikzpicture}\newpage
\definecolor{answer}{rgb}{0.850980,0.047059,0.070588}
\null\begin{tikzpicture}[remember picture,overlay]\begin{scope}[shift={(current page.north west)},x=1bp,y=-1bp]
\node[anchor=north west,inner sep=0pt,text=answer,align=left,text width=379.4859924316406bp] at (15,15) {\fontsize{7.5bp}{9.5bp}\selectfont 析：\allowbreak{}\ensuremath{4.495}四舍五入为\ensuremath{4.50};\allowbreak{}再小的\ensuremath{4.494}只能得\ensuremath{4.49},\allowbreak{}所以最小\ensuremath{4.495}。\allowbreak{}\par};
\end{scope}\end{tikzpicture}\newpage
\definecolor{answer}{rgb}{0.850980,0.047059,0.070588}
\null\begin{tikzpicture}[remember picture,overlay]\begin{scope}[shift={(current page.north west)},x=1bp,y=-1bp]
\node[anchor=north west,inner sep=0pt,text=answer,align=left,text width=372.9530029296875bp] at (15,15) {\fontsize{7.5bp}{9.5bp}\selectfont 析：\allowbreak{}原数：\allowbreak{}\ensuremath{0.458} ÷ \ensuremath{100} = \ensuremath{0.00458};\allowbreak{}缩小到百分之一：\allowbreak{}\ensuremath{0.5} ÷ \ensuremath{100} = \ensuremath{0.005}。\allowbreak{}\par};
\end{scope}\end{tikzpicture}\newpage
\definecolor{answer}{rgb}{0.850980,0.047059,0.070588}
\null\begin{tikzpicture}[remember picture,overlay]\begin{scope}[shift={(current page.north west)},x=1bp,y=-1bp]
\node[anchor=north west,inner sep=0pt,text=answer,align=left,text width=348.12799072265625bp] at (15,15) {\fontsize{7.5bp}{9.5bp}\selectfont 析：\allowbreak{}每格\ensuremath{0.1};\allowbreak{} \ensuremath{1.45}在\ensuremath{1.4}与\ensuremath{1.5}正中，\allowbreak{}\ensuremath{2.350} = \ensuremath{2.35},\allowbreak{}在\ensuremath{2.3}与\ensuremath{2.4}正中。\allowbreak{}\par};
\end{scope}\end{tikzpicture}\newpage
\definecolor{answer}{rgb}{0.850980,0.047059,0.070588}
\null\begin{tikzpicture}[remember picture,overlay]\begin{scope}[shift={(current page.north west)},x=1bp,y=-1bp]
\node[anchor=north west,inner sep=0pt,text=answer,align=left,text width=433.1730041503906bp] at (15,15) {\fontsize{7.5bp}{9.5bp}\selectfont 析：\allowbreak{}\ensuremath{7}个评分去掉最高\ensuremath{98}、\allowbreak{}最低\ensuremath{84}，\allowbreak{}剩\ensuremath{5}个，\allowbreak{}所以除\ensuremath{5}：\allowbreak{}(\ensuremath{92} + \ensuremath{94} + \ensuremath{93} + \ensuremath{95} + \ensuremath{96}) ÷ \ensuremath{5} = \ensuremath{94}分；\allowbreak{}也可移多补少。\allowbreak{}\par};
\end{scope}\end{tikzpicture}\newpage
\definecolor{answer}{rgb}{0.850980,0.047059,0.070588}
\null\begin{tikzpicture}[remember picture,overlay]\begin{scope}[shift={(current page.north west)},x=1bp,y=-1bp]
\node[anchor=north west,inner sep=0pt,text=answer,align=left,text width=371.5150146484375bp] at (15,15) {\fontsize{7.5bp}{9.5bp}\selectfont 析：\allowbreak{}错：\allowbreak{}大小不变，\allowbreak{}计数单位可能变；\allowbreak{}如\ensuremath{0.8}的单位是\ensuremath{0.1},\allowbreak{} \ensuremath{0.80}的单位是\ensuremath{0.01}。\allowbreak{}\par};
\end{scope}\end{tikzpicture}\newpage
\definecolor{answer}{rgb}{0.850980,0.047059,0.070588}
\null\begin{tikzpicture}[remember picture,overlay]\begin{scope}[shift={(current page.north west)},x=1bp,y=-1bp]
\node[anchor=north west,inner sep=0pt,text=answer,align=left,text width=501.125bp] at (15,15) {\fontsize{7.5bp}{9.5bp}\selectfont 析：\allowbreak{}×。\allowbreak{}\ensuremath{2.51} \ensuremath{-}\ensuremath{0.52} + \ensuremath{0.48} = \ensuremath{2.51} \ensuremath{-}(\ensuremath{0.52} \ensuremath{-}\ensuremath{0.48}) = \ensuremath{2.51} \ensuremath{-}\ensuremath{0.04} = \ensuremath{2.47}；\allowbreak{}右式\ensuremath{2.51} \ensuremath{-}(\ensuremath{0.52} + \ensuremath{0.48}) = \ensuremath{1.51}，\allowbreak{}不相等\par};
\end{scope}\end{tikzpicture}\newpage
\definecolor{answer}{rgb}{0.850980,0.047059,0.070588}
\null\begin{tikzpicture}[remember picture,overlay]\begin{scope}[shift={(current page.north west)},x=1bp,y=-1bp]
\node[anchor=north west,inner sep=0pt,text=answer,align=left,text width=400.989013671875bp] at (15,15) {\fontsize{7.5bp}{9.5bp}\selectfont 析：\allowbreak{}错：\allowbreak{}括号内是乘积，\allowbreak{}应用结合律写\ensuremath{25} × \ensuremath{8} × \ensuremath{4},\allowbreak{}不能改成两个积相加。\allowbreak{}\par};
\end{scope}\end{tikzpicture}\newpage
\definecolor{answer}{rgb}{0.850980,0.047059,0.070588}
\null\begin{tikzpicture}[remember picture,overlay]\begin{scope}[shift={(current page.north west)},x=1bp,y=-1bp]
\node[anchor=north west,inner sep=0pt,text=answer,align=left,text width=338.6090087890625bp] at (15,15) {\fontsize{7.5bp}{9.5bp}\selectfont 析：\allowbreak{}对：\allowbreak{}\ensuremath{0.8}与\ensuremath{0.800}大小相等，\allowbreak{}计数单位分别是\ensuremath{0.1}和\ensuremath{0.001}。\allowbreak{}\par};
\end{scope}\end{tikzpicture}\newpage
\definecolor{answer}{rgb}{0.850980,0.047059,0.070588}
\null\begin{tikzpicture}[remember picture,overlay]\begin{scope}[shift={(current page.north west)},x=1bp,y=-1bp]
\node[anchor=north west,inner sep=0pt,text=answer,align=left,text width=331.0bp] at (15,15) {\fontsize{7.5bp}{9.5bp}\selectfont 析：\allowbreak{}错：\allowbreak{}先乘\ensuremath{10}再除\ensuremath{100},\allowbreak{}等于除\ensuremath{10},\allowbreak{}应缩小到原来的十分之一。\allowbreak{}\par};
\end{scope}\end{tikzpicture}\newpage
\definecolor{answer}{rgb}{0.850980,0.047059,0.070588}
\null\begin{tikzpicture}[remember picture,overlay]\begin{scope}[shift={(current page.north west)},x=1bp,y=-1bp]
\node[anchor=north west,inner sep=0pt,text=answer,align=left,text width=323.5bp] at (15,15) {\fontsize{7.5bp}{9.5bp}\selectfont 析：\allowbreak{}B :\allowbreak{}先算小括号的减法，\allowbreak{}再按从左到右先乘后除。\allowbreak{}\par};
\end{scope}\end{tikzpicture}\newpage
\definecolor{answer}{rgb}{0.850980,0.047059,0.070588}
\null\begin{tikzpicture}[remember picture,overlay]\begin{scope}[shift={(current page.north west)},x=1bp,y=-1bp]
\node[anchor=north west,inner sep=0pt,text=answer,align=left,text width=316.968994140625bp] at (15,15) {\fontsize{7.5bp}{9.5bp}\selectfont 析：\allowbreak{}A :\allowbreak{} \ensuremath{85} ÷ \ensuremath{100} = \ensuremath{0.85},\allowbreak{}向右移两位才成为\ensuremath{85}。\allowbreak{}\par};
\end{scope}\end{tikzpicture}\newpage
\definecolor{answer}{rgb}{0.850980,0.047059,0.070588}
\null\begin{tikzpicture}[remember picture,overlay]\begin{scope}[shift={(current page.north west)},x=1bp,y=-1bp]
\node[anchor=north west,inner sep=0pt,text=answer,align=left,text width=324.1629943847656bp] at (15,15) {\fontsize{7.5bp}{9.5bp}\selectfont 析：\allowbreak{}A :\allowbreak{}与\ensuremath{10}的差依次是\ensuremath{0.002}、\allowbreak{}\ensuremath{0.01}、\allowbreak{}\ensuremath{0.02},\allowbreak{} \ensuremath{0.002}最小。\allowbreak{}\par};
\end{scope}\end{tikzpicture}\newpage
\definecolor{answer}{rgb}{0.850980,0.047059,0.070588}
\null\begin{tikzpicture}[remember picture,overlay]\begin{scope}[shift={(current page.north west)},x=1bp,y=-1bp]
\node[anchor=north west,inner sep=0pt,text=answer,align=left,text width=336.8370056152344bp] at (15,15) {\fontsize{7.5bp}{9.5bp}\selectfont 析：\allowbreak{}A :\allowbreak{}连续减去两数，\allowbreak{}等于减去它们的和；\allowbreak{}\ensuremath{2.86} + \ensuremath{1.14} = \ensuremath{4},\allowbreak{}便于计算。\allowbreak{}\par};
\end{scope}\end{tikzpicture}\newpage
\definecolor{answer}{rgb}{0.850980,0.047059,0.070588}
\null\begin{tikzpicture}[remember picture,overlay]\begin{scope}[shift={(current page.north west)},x=1bp,y=-1bp]
\node[anchor=north west,inner sep=0pt,text=answer,align=left,text width=346.0bp] at (15,15) {\fontsize{7.5bp}{9.5bp}\selectfont 析：\allowbreak{}C :\allowbreak{} \ensuremath{0.31}、\allowbreak{}\ensuremath{0.301}、\allowbreak{}\ensuremath{0.3001}等都在两数之间，\allowbreak{}小数位数不限，\allowbreak{}所以有无数个。\allowbreak{}\par};
\end{scope}\end{tikzpicture}\newpage
\definecolor{answer}{rgb}{0.850980,0.047059,0.070588}
\null\begin{tikzpicture}[remember picture,overlay]\begin{scope}[shift={(current page.north west)},x=1bp,y=-1bp]
\node[anchor=north west,inner sep=0pt,text=answer,align=left,text width=501.0bp] at (15,15) {\fontsize{8.5bp}{10.62bp}\selectfont 先从总价中扣掉\ensuremath{50}本笔记本的价钱，\allowbreak{}再平均分给\ensuremath{50}支笔。\allowbreak{}\par};
\end{scope}\end{tikzpicture}\newpage
\definecolor{answer}{rgb}{0.850980,0.047059,0.070588}
\null\begin{tikzpicture}[remember picture,overlay]\begin{scope}[shift={(current page.north west)},x=1bp,y=-1bp]
\node[anchor=north west,inner sep=0pt,text=answer,align=left,text width=501.0bp] at (15,15) {\fontsize{8.5bp}{10.62bp}\selectfont 两人相向而行，\allowbreak{}速度和乘共同时间就是学校到图书馆的距离。\allowbreak{}\par};
\end{scope}\end{tikzpicture}\newpage
\definecolor{answer}{rgb}{0.850980,0.047059,0.070588}
\null\begin{tikzpicture}[remember picture,overlay]\begin{scope}[shift={(current page.north west)},x=1bp,y=-1bp]
\node[anchor=north west,inner sep=0pt,text=answer,align=left,text width=501.0bp] at (15,15) {\fontsize{8.5bp}{10.62bp}\selectfont 先由第一棵比第二棵矮求第二棵高度，\allowbreak{}再求第三棵。\allowbreak{}\par};
\end{scope}\end{tikzpicture}\newpage
\definecolor{answer}{rgb}{0.850980,0.047059,0.070588}
\null\begin{tikzpicture}[remember picture,overlay]\begin{scope}[shift={(current page.north west)},x=1bp,y=-1bp]
\node[anchor=north west,inner sep=0pt,text=answer,align=left,text width=501.0bp] at (15,15) {\fontsize{8.5bp}{10.62bp}\selectfont 同路程分别除以两种速度，\allowbreak{}求出到达时间，\allowbreak{}再相减。\allowbreak{}\par};
\end{scope}\end{tikzpicture}\newpage
\definecolor{answer}{rgb}{0.850980,0.047059,0.070588}
\null\begin{tikzpicture}[remember picture,overlay]\begin{scope}[shift={(current page.north west)},x=1bp,y=-1bp]
\node[anchor=north west,inner sep=0pt,text=answer,align=left,text width=501.0bp] at (15,15) {\fontsize{8.5bp}{10.62bp}\selectfont 从一周总销量中减前三天销量，\allowbreak{}再除以后四天的天数。\allowbreak{}\par};
\end{scope}\end{tikzpicture}\newpage
\definecolor{answer}{rgb}{0.850980,0.047059,0.070588}
\null\begin{tikzpicture}[remember picture,overlay]\begin{scope}[shift={(current page.north west)},x=1bp,y=-1bp]
\node[anchor=north west,inner sep=0pt,text=answer,align=left,text width=501.0bp] at (15,15) {\fontsize{8.5bp}{10.62bp}\selectfont 一个来回是\ensuremath{2}个单程，\allowbreak{}\ensuremath{6}个来回共有\ensuremath{12}个单程。\allowbreak{}\par};
\end{scope}\end{tikzpicture}\newpage
\definecolor{answer}{rgb}{0.850980,0.047059,0.070588}
\null\begin{tikzpicture}[remember picture,overlay]\begin{scope}[shift={(current page.north west)},x=1bp,y=-1bp]
\node[anchor=north west,inner sep=0pt,text=answer,align=left,text width=323.5bp] at (15,15) {\fontsize{7.5bp}{9.5bp}\selectfont 析：\allowbreak{}同一平面内两直线相交或平行；\allowbreak{}相交成直角时互相垂直。\allowbreak{}\par};
\end{scope}\end{tikzpicture}\newpage
\definecolor{answer}{rgb}{0.850980,0.047059,0.070588}
\null\begin{tikzpicture}[remember picture,overlay]\begin{scope}[shift={(current page.north west)},x=1bp,y=-1bp]
\node[anchor=north west,inner sep=0pt,text=answer,align=left,text width=335.2340087890625bp] at (15,15) {\fontsize{7.5bp}{9.5bp}\selectfont 析：\allowbreak{}过直线外一点，\allowbreak{}平行线与垂线各只能画\ensuremath{1}条。\allowbreak{}\par};
\end{scope}\end{tikzpicture}\newpage
\definecolor{answer}{rgb}{0.850980,0.047059,0.070588}
\null\begin{tikzpicture}[remember picture,overlay]\begin{scope}[shift={(current page.north west)},x=1bp,y=-1bp]
\node[anchor=north west,inner sep=0pt,text=answer,align=left,text width=312.7340087890625bp] at (15,15) {\fontsize{7.5bp}{9.5bp}\selectfont 析：\allowbreak{}三角形有\ensuremath{3}条边、\allowbreak{}\ensuremath{3}个角，\allowbreak{}具有稳定性。\allowbreak{}\par};
\end{scope}\end{tikzpicture}\newpage
\definecolor{answer}{rgb}{0.850980,0.047059,0.070588}
\null\begin{tikzpicture}[remember picture,overlay]\begin{scope}[shift={(current page.north west)},x=1bp,y=-1bp]
\node[anchor=north west,inner sep=0pt,text=answer,align=left,text width=375.4834899902344bp] at (15,15) {\fontsize{7.5bp}{9.5bp}\selectfont 析：\allowbreak{}按角分锐角、\allowbreak{}直角、\allowbreak{}钝角三角形；\allowbreak{}等边三角形三角均\ensuremath{60}\ensuremath{^\circ},\allowbreak{}属于锐角三角形。\allowbreak{}\par};
\end{scope}\end{tikzpicture}\newpage
\definecolor{answer}{rgb}{0.850980,0.047059,0.070588}
\null\begin{tikzpicture}[remember picture,overlay]\begin{scope}[shift={(current page.north west)},x=1bp,y=-1bp]
\node[anchor=north west,inner sep=0pt,text=answer,align=left,text width=320.2340087890625bp] at (15,15) {\fontsize{7.5bp}{9.5bp}\selectfont 析：\allowbreak{}平行四边形对边分别平行且相等；\allowbreak{}梯形只有\ensuremath{1}组对边平行。\allowbreak{}\par};
\end{scope}\end{tikzpicture}\newpage
\definecolor{answer}{rgb}{0.850980,0.047059,0.070588}
\null\begin{tikzpicture}[remember picture,overlay]\begin{scope}[shift={(current page.north west)},x=1bp,y=-1bp]
\node[anchor=north west,inner sep=0pt,text=answer,align=left,text width=308.5bp] at (15,15) {\fontsize{7.5bp}{9.5bp}\selectfont 析：\allowbreak{}平移和旋转只改变位置，\allowbreak{}不改变形状与大小。\allowbreak{}\par};
\end{scope}\end{tikzpicture}\newpage
\definecolor{answer}{rgb}{0.850980,0.047059,0.070588}
\null\begin{tikzpicture}[remember picture,overlay]\begin{scope}[shift={(current page.north west)},x=1bp,y=-1bp]
\node[anchor=north west,inner sep=0pt,text=answer,align=left,text width=293.5bp] at (15,15) {\fontsize{7.5bp}{9.5bp}\selectfont 析：\allowbreak{}沿直线对折两侧完全重合，\allowbreak{}叫轴对称图形。\allowbreak{}\par};
\end{scope}\end{tikzpicture}\newpage
\definecolor{answer}{rgb}{0.850980,0.047059,0.070588}
\null\begin{tikzpicture}[remember picture,overlay]\begin{scope}[shift={(current page.north west)},x=1bp,y=-1bp]
\node[anchor=north west,inner sep=0pt,text=answer,align=left,text width=342.7340087890625bp] at (15,15) {\fontsize{7.5bp}{9.5bp}\selectfont 析：\allowbreak{}两者都是四边形：\allowbreak{}梯形\ensuremath{1}组对边平行，\allowbreak{}平行四边形\ensuremath{2}组对边分别平行。\allowbreak{}\par};
\end{scope}\end{tikzpicture}\newpage
\definecolor{answer}{rgb}{0.850980,0.047059,0.070588}
\null\begin{tikzpicture}[remember picture,overlay]\begin{scope}[shift={(current page.north west)},x=1bp,y=-1bp]
\node[anchor=north west,inner sep=0pt,text=answer,align=left,text width=283.1304931640625bp] at (15,15) {\fontsize{7.5bp}{9.5bp}\selectfont 析：\allowbreak{}对：\allowbreak{}等边三角形三个角都是\ensuremath{60}\ensuremath{^\circ},\allowbreak{}均小于\ensuremath{90}\ensuremath{^\circ}。\allowbreak{}\par};
\end{scope}\end{tikzpicture}\newpage
\definecolor{answer}{rgb}{0.850980,0.047059,0.070588}
\null\begin{tikzpicture}[remember picture,overlay]\begin{scope}[shift={(current page.north west)},x=1bp,y=-1bp]
\node[anchor=north west,inner sep=0pt,text=answer,align=left,text width=331.0bp] at (15,15) {\fontsize{7.5bp}{9.5bp}\selectfont 析：\allowbreak{}错：\allowbreak{}只见一个锐角，\allowbreak{}另两个角未知，\allowbreak{}可能是直角或钝角三角形。\allowbreak{}\par};
\end{scope}\end{tikzpicture}\newpage
\definecolor{answer}{rgb}{0.850980,0.047059,0.070588}
\null\begin{tikzpicture}[remember picture,overlay]\begin{scope}[shift={(current page.north west)},x=1bp,y=-1bp]
\node[anchor=north west,inner sep=0pt,text=answer,align=left,text width=346.0bp] at (15,15) {\fontsize{7.5bp}{9.5bp}\selectfont 析：\allowbreak{}对：\allowbreak{}过直线外一点，\allowbreak{}该直线的平行线、\allowbreak{}垂线各唯一。\allowbreak{}\par};
\end{scope}\end{tikzpicture}\newpage
\definecolor{answer}{rgb}{0.850980,0.047059,0.070588}
\null\begin{tikzpicture}[remember picture,overlay]\begin{scope}[shift={(current page.north west)},x=1bp,y=-1bp]
\node[anchor=north west,inner sep=0pt,text=answer,align=left,text width=353.5bp] at (15,15) {\fontsize{7.5bp}{9.5bp}\selectfont 析：\allowbreak{}错：\allowbreak{}不等边三角形没有对称轴，\allowbreak{}并非所有三角形都是轴对称图形。\allowbreak{}\par};
\end{scope}\end{tikzpicture}\newpage
\definecolor{answer}{rgb}{0.850980,0.047059,0.070588}
\null\begin{tikzpicture}[remember picture,overlay]\begin{scope}[shift={(current page.north west)},x=1bp,y=-1bp]
\node[anchor=north west,inner sep=0pt,text=answer,align=left,text width=501.0bp] at (15,15) {\fontsize{7.5bp}{9.5bp}\selectfont 析：\allowbreak{}B :\allowbreak{}梯形包含在四边形中；\allowbreak{}等边属于等腰，\allowbreak{}正方形属于长方形，\allowbreak{}A、\allowbreak{}C关系画错。\allowbreak{}\par};
\end{scope}\end{tikzpicture}\newpage
\definecolor{answer}{rgb}{0.850980,0.047059,0.070588}
\null\begin{tikzpicture}[remember picture,overlay]\begin{scope}[shift={(current page.north west)},x=1bp,y=-1bp]
\node[anchor=north west,inner sep=0pt,text=answer,align=left,text width=501.0bp] at (15,15) {\fontsize{7.5bp}{9.5bp}\selectfont 析：\allowbreak{}B :\allowbreak{}平行四边形两条平行对边之间可作无数条垂线段；\allowbreak{}三角形只有\ensuremath{3}条高。\allowbreak{}\par};
\end{scope}\end{tikzpicture}\newpage
\definecolor{answer}{rgb}{0.850980,0.047059,0.070588}
\null\begin{tikzpicture}[remember picture,overlay]\begin{scope}[shift={(current page.north west)},x=1bp,y=-1bp]
\node[anchor=north west,inner sep=0pt,text=answer,align=left,text width=501.0bp] at (15,15) {\fontsize{7.5bp}{9.5bp}\selectfont 析：\allowbreak{}B :\allowbreak{}圆的每条直径所在直线都是对称轴，\allowbreak{}有无数条。\allowbreak{}\par};
\end{scope}\end{tikzpicture}\newpage
\definecolor{answer}{rgb}{0.850980,0.047059,0.070588}
\null\begin{tikzpicture}[remember picture,overlay]\begin{scope}[shift={(current page.north west)},x=1bp,y=-1bp]
\node[anchor=north west,inner sep=0pt,text=answer,align=left,text width=501.0bp] at (15,15) {\fontsize{7.5bp}{9.5bp}\selectfont 析：\allowbreak{}C :\allowbreak{}一般平行四边形没有对称轴；\allowbreak{}对边平行且相等、\allowbreak{}对角相等均正确。\allowbreak{}\par};
\end{scope}\end{tikzpicture}\newpage
\definecolor{answer}{rgb}{0.850980,0.047059,0.070588}
\null\begin{tikzpicture}[remember picture,overlay]\begin{scope}[shift={(current page.north west)},x=1bp,y=-1bp]
\node[anchor=north west,inner sep=0pt,text=answer,align=left,text width=501.0bp] at (15,15) {\fontsize{7.5bp}{9.5bp}\selectfont 析：\allowbreak{}B旋转后能与L形缺口完全重合；\allowbreak{}A、\allowbreak{}C不能只靠旋转吻合。\allowbreak{}\par};
\end{scope}\end{tikzpicture}\newpage
\definecolor{answer}{rgb}{0.850980,0.047059,0.070588}
\null\begin{tikzpicture}[remember picture,overlay]\begin{scope}[shift={(current page.north west)},x=1bp,y=-1bp]
\node[anchor=north west,inner sep=0pt,text=answer,align=left,text width=501.0bp] at (15,15) {\fontsize{7.5bp}{9.5bp}\selectfont 析：\allowbreak{}倒推叠放顺序：\allowbreak{}左上\ensuremath{1}、\allowbreak{}左下\ensuremath{2}、\allowbreak{}右下\ensuremath{3}、\allowbreak{}右上\ensuremath{4};\allowbreak{}折好后从前到后依次\ensuremath{1}、\allowbreak{}\ensuremath{2}、\allowbreak{}\ensuremath{3}、\allowbreak{}\ensuremath{4}。\allowbreak{}\par};
\end{scope}\end{tikzpicture}\newpage
\definecolor{answer}{rgb}{0.850980,0.047059,0.070588}
\null\begin{tikzpicture}[remember picture,overlay]\begin{scope}[shift={(current page.north west)},x=1bp,y=-1bp]
\node[anchor=north west,inner sep=0pt,text=answer,align=left,text width=369.3269958496094bp] at (15,15) {\fontsize{7.5bp}{9.5bp}\selectfont 析：\allowbreak{}平均本数=总本数÷人数；\allowbreak{}\ensuremath{873} ÷ \ensuremath{97} = \ensuremath{9}本，\allowbreak{}\ensuremath{968} ÷ \ensuremath{88} = \ensuremath{11}本。\allowbreak{}\par};
\end{scope}\end{tikzpicture}\newpage
\definecolor{answer}{rgb}{0.850980,0.047059,0.070588}
\null\begin{tikzpicture}[remember picture,overlay]\begin{scope}[shift={(current page.north west)},x=1bp,y=-1bp]
\node[anchor=north west,inner sep=0pt,text=answer,align=left,text width=431.2795104980469bp] at (15,15) {\fontsize{7.5bp}{9.5bp}\selectfont 析：\allowbreak{}两班人数不同，\allowbreak{}要合并总量和人数：\allowbreak{}(\ensuremath{270} + \ensuremath{252}) ÷ (\ensuremath{45} + \ensuremath{42}) = \ensuremath{6}棵；\allowbreak{}本题两班人均恰好均为\ensuremath{6}，\allowbreak{}合并仍为\ensuremath{6}。\allowbreak{}\par};
\end{scope}\end{tikzpicture}\newpage
\definecolor{answer}{rgb}{0.850980,0.047059,0.070588}
\null\begin{tikzpicture}[remember picture,overlay]\begin{scope}[shift={(current page.north west)},x=1bp,y=-1bp]
\node[anchor=north west,inner sep=0pt,text=answer,align=left,text width=374.6709899902344bp] at (15,15) {\fontsize{7.5bp}{9.5bp}\selectfont 析：\allowbreak{}读图五次为\ensuremath{70}、\allowbreak{}\ensuremath{80}、\allowbreak{}\ensuremath{60}、\allowbreak{}\ensuremath{100}、\allowbreak{}\ensuremath{80};\allowbreak{}第四次最多\ensuremath{100}个，\allowbreak{}第二、\allowbreak{}五次均\ensuremath{80}个。\allowbreak{}\par};
\end{scope}\end{tikzpicture}\newpage
\definecolor{answer}{rgb}{0.850980,0.047059,0.070588}
\null\begin{tikzpicture}[remember picture,overlay]\begin{scope}[shift={(current page.north west)},x=1bp,y=-1bp]
\node[anchor=north west,inner sep=0pt,text=answer,align=left,text width=327.37298583984375bp] at (15,15) {\fontsize{7.5bp}{9.5bp}\selectfont 析：\allowbreak{}连续三数的平均数就是中间数；\allowbreak{}\ensuremath{16}、\allowbreak{}\ensuremath{17}、\allowbreak{}\ensuremath{18}中最大是\ensuremath{18}。\allowbreak{}\par};
\end{scope}\end{tikzpicture}\newpage
\definecolor{answer}{rgb}{0.850980,0.047059,0.070588}
\null\begin{tikzpicture}[remember picture,overlay]\begin{scope}[shift={(current page.north west)},x=1bp,y=-1bp]
\node[anchor=north west,inner sep=0pt,text=answer,align=left,text width=402.78900146484375bp] at (15,15) {\fontsize{7.5bp}{9.5bp}\selectfont 析：\allowbreak{}对：\allowbreak{}女生均分\ensuremath{86.5}低于全班\ensuremath{87},\allowbreak{}男生均分必须高于\ensuremath{87},\allowbreak{}才能把全班平均提高到\ensuremath{87}。\allowbreak{}\par};
\end{scope}\end{tikzpicture}\newpage
\definecolor{answer}{rgb}{0.850980,0.047059,0.070588}
\null\begin{tikzpicture}[remember picture,overlay]\begin{scope}[shift={(current page.north west)},x=1bp,y=-1bp]
\node[anchor=north west,inner sep=0pt,text=answer,align=left,text width=361.0bp] at (15,15) {\fontsize{7.5bp}{9.5bp}\selectfont 析：\allowbreak{}错：\allowbreak{}平均数不小于最小值，\allowbreak{}不大于最大值，\allowbreak{}不一定大于每个数。\allowbreak{}\par};
\end{scope}\end{tikzpicture}\newpage
\definecolor{answer}{rgb}{0.850980,0.047059,0.070588}
\null\begin{tikzpicture}[remember picture,overlay]\begin{scope}[shift={(current page.north west)},x=1bp,y=-1bp]
\node[anchor=north west,inner sep=0pt,text=answer,align=left,text width=383.1390075683594bp] at (15,15) {\fontsize{7.5bp}{9.5bp}\selectfont 析：\allowbreak{}错：\allowbreak{}平均数不保证恰有一半数据更大；\allowbreak{}若\ensuremath{32}位教师都\ensuremath{35}岁，\allowbreak{}则高于\ensuremath{35}岁的有\ensuremath{0}位。\allowbreak{}\par};
\end{scope}\end{tikzpicture}\newpage
\definecolor{answer}{rgb}{0.850980,0.047059,0.070588}
\null\begin{tikzpicture}[remember picture,overlay]\begin{scope}[shift={(current page.north west)},x=1bp,y=-1bp]
\node[anchor=north west,inner sep=0pt,text=answer,align=left,text width=501.0bp] at (15,15) {\fontsize{7.5bp}{9.5bp}\selectfont 析：\allowbreak{}A :\allowbreak{}平均年龄应在最小\ensuremath{18}与最大\ensuremath{20}岁之间，\allowbreak{}只有\ensuremath{19}岁符合。\allowbreak{}\par};
\end{scope}\end{tikzpicture}\newpage
\definecolor{answer}{rgb}{0.850980,0.047059,0.070588}
\null\begin{tikzpicture}[remember picture,overlay]\begin{scope}[shift={(current page.north west)},x=1bp,y=-1bp]
\node[anchor=north west,inner sep=0pt,text=answer,align=left,text width=501.0bp] at (15,15) {\fontsize{7.5bp}{9.5bp}\selectfont 析：\allowbreak{}C :\allowbreak{}实际\ensuremath{134} + \ensuremath{46} = \ensuremath{180}道，\allowbreak{}\ensuremath{180} ÷ \ensuremath{4} = \ensuremath{45}道/天。\allowbreak{}\par};
\end{scope}\end{tikzpicture}\newpage
\definecolor{answer}{rgb}{0.850980,0.047059,0.070588}
\null\begin{tikzpicture}[remember picture,overlay]\begin{scope}[shift={(current page.north west)},x=1bp,y=-1bp]
\node[anchor=north west,inner sep=0pt,text=answer,align=left,text width=501.0bp] at (15,15) {\fontsize{7.5bp}{9.5bp}\selectfont 析：\allowbreak{}C :\allowbreak{}高于虚线的部分恰好能补齐低于虚线的部分，\allowbreak{}移多补少后四柱等高。\allowbreak{}\par};
\end{scope}\end{tikzpicture}\newpage
\definecolor{answer}{rgb}{0.850980,0.047059,0.070588}
\null\begin{tikzpicture}[remember picture,overlay]\begin{scope}[shift={(current page.north west)},x=1bp,y=-1bp]
\node[anchor=north west,inner sep=0pt,text=answer,align=left,text width=501.0bp] at (15,15) {\fontsize{7.5bp}{9.5bp}\selectfont 析：\allowbreak{}D :\allowbreak{}一共\ensuremath{2} + \ensuremath{3} + \ensuremath{1} = \ensuremath{6}个数，\allowbreak{}两个\ensuremath{15}、\allowbreak{}三个\ensuremath{16}、\allowbreak{}一个\ensuremath{14}全相加，\allowbreak{}再除以\ensuremath{6}。\allowbreak{}\par};
\end{scope}\end{tikzpicture}\newpage
\definecolor{answer}{rgb}{0.850980,0.047059,0.070588}
\null\begin{tikzpicture}[remember picture,overlay]\begin{scope}[shift={(current page.north west)},x=1bp,y=-1bp]
\node[anchor=north west,inner sep=0pt,text=answer,align=left,text width=501.0bp] at (15,15) {\fontsize{7.5bp}{9.5bp}\selectfont 析：\allowbreak{}D :\allowbreak{}平均数大于最小数\ensuremath{c}、\allowbreak{}小于最大数\ensuremath{a},\allowbreak{}不一定等于\ensuremath{b}。\allowbreak{}\par};
\end{scope}\end{tikzpicture}\newpage
\definecolor{answer}{rgb}{0.850980,0.047059,0.070588}
\null\begin{tikzpicture}[remember picture,overlay]\begin{scope}[shift={(current page.north west)},x=1bp,y=-1bp]
\node[anchor=north west,inner sep=0pt,outer sep=0pt,text=answer,align=left,text width=490bp] at (15,15) {\fontsize{8.5bp}{11.5bp}\selectfont 答：不对。\ensuremath{152}下和\ensuremath{150}下只是略高于\ensuremath{149}下，两人多出的次数，连\ensuremath{117}下这名同学不足\ensuremath{149}下的部分都补不齐，\\其余同学的成绩也都低于\ensuremath{149}下，所以平均成绩一定小于\ensuremath{149}下。\par};
\end{scope}\end{tikzpicture}
\newpage
\definecolor{answer}{rgb}{0.850980,0.047059,0.070588}
\null\begin{tikzpicture}[remember picture,overlay]\begin{scope}[shift={(current page.north west)},x=1bp,y=-1bp]
\node[anchor=north west,inner sep=0pt,text=answer,align=left,text width=333.4775085449219bp] at (15,15) {\fontsize{7.5bp}{9.5bp}\selectfont 析：\allowbreak{}先求\ensuremath{840}与\ensuremath{790}的差，\allowbreak{}再乘\ensuremath{2},\allowbreak{}列(\ensuremath{840} \ensuremath{-}\ensuremath{790}) × \ensuremath{2}。\allowbreak{}\par};
\end{scope}\end{tikzpicture}\newpage
\definecolor{answer}{rgb}{0.850980,0.047059,0.070588}
\null\begin{tikzpicture}[remember picture,overlay]\begin{scope}[shift={(current page.north west)},x=1bp,y=-1bp]
\node[anchor=north west,inner sep=0pt,text=answer,align=left,text width=327.510986328125bp] at (15,15) {\fontsize{7.5bp}{9.5bp}\selectfont 析：\allowbreak{}\ensuremath{240}是乘法结果，\allowbreak{}\ensuremath{28} \ensuremath{-}\ensuremath{16}应整体作除数，\allowbreak{}列\ensuremath{15} × \ensuremath{16} ÷ (\ensuremath{28} \ensuremath{-}\ensuremath{16})。\allowbreak{}\par};
\end{scope}\end{tikzpicture}\newpage
\definecolor{answer}{rgb}{0.850980,0.047059,0.070588}
\null\begin{tikzpicture}[remember picture,overlay]\begin{scope}[shift={(current page.north west)},x=1bp,y=-1bp]
\node[anchor=north west,inner sep=0pt,text=answer,align=left,text width=326.5150146484375bp] at (15,15) {\fontsize{7.5bp}{9.5bp}\selectfont 析：\allowbreak{}先加\ensuremath{20} + \ensuremath{40},\allowbreak{}再除\ensuremath{540} ÷ \ensuremath{60},\allowbreak{}最后乘\ensuremath{160} × \ensuremath{9} = \ensuremath{1440}。\allowbreak{}\par};
\end{scope}\end{tikzpicture}\newpage
\definecolor{answer}{rgb}{0.850980,0.047059,0.070588}
\null\begin{tikzpicture}[remember picture,overlay]\begin{scope}[shift={(current page.north west)},x=1bp,y=-1bp]
\node[anchor=north west,inner sep=0pt,text=answer,align=left,text width=325.04901123046875bp] at (15,15) {\fontsize{7.5bp}{9.5bp}\selectfont 析：\allowbreak{}一楼座位\ensuremath{288} × \ensuremath{2} + \ensuremath{50} = \ensuremath{626}个；\allowbreak{}\ensuremath{626} > \ensuremath{600},\allowbreak{}所以能坐下。\allowbreak{}\par};
\end{scope}\end{tikzpicture}\newpage
\definecolor{answer}{rgb}{0.850980,0.047059,0.070588}
\null\begin{tikzpicture}[remember picture,overlay]\begin{scope}[shift={(current page.north west)},x=1bp,y=-1bp]
\node[anchor=north west,inner sep=0pt,text=answer,align=left,text width=330.26348876953125bp] at (15,15) {\fontsize{7.5bp}{9.5bp}\selectfont 析：\allowbreak{}由○÷□=☆知○=□×☆；\allowbreak{}除数非\ensuremath{0},\allowbreak{}所以○÷(□×☆)=\ensuremath{1}。\allowbreak{}\par};
\end{scope}\end{tikzpicture}\newpage
\definecolor{answer}{rgb}{0.850980,0.047059,0.070588}
\null\begin{tikzpicture}[remember picture,overlay]\begin{scope}[shift={(current page.north west)},x=1bp,y=-1bp]
\node[anchor=north west,inner sep=0pt,text=answer,align=left,text width=382.7239990234375bp] at (15,15) {\fontsize{7.5bp}{9.5bp}\selectfont 析：\allowbreak{}B :\allowbreak{} \ensuremath{44}个\ensuremath{25}是\ensuremath{25} × \ensuremath{44},\allowbreak{}除数是\ensuremath{25}的\ensuremath{11}倍，\allowbreak{}应把\ensuremath{25} × \ensuremath{11}整体括起。\allowbreak{}\par};
\end{scope}\end{tikzpicture}\newpage
\definecolor{answer}{rgb}{0.850980,0.047059,0.070588}
\null\begin{tikzpicture}[remember picture,overlay]\begin{scope}[shift={(current page.north west)},x=1bp,y=-1bp]
\node[anchor=north west,inner sep=0pt,text=answer,align=left,text width=428.22698974609375bp] at (15,15) {\fontsize{7.5bp}{9.5bp}\selectfont 析：\allowbreak{}C先算\ensuremath{30} ÷ \ensuremath{3} = \ensuremath{10},\allowbreak{}正确；\allowbreak{}A先算\ensuremath{25} × \ensuremath{18},\allowbreak{} B的乘除应从左到右，\allowbreak{}不能先算\ensuremath{25} × \ensuremath{4}。\allowbreak{}\par};
\end{scope}\end{tikzpicture}\newpage
\definecolor{answer}{rgb}{0.850980,0.047059,0.070588}
\null\begin{tikzpicture}[remember picture,overlay]\begin{scope}[shift={(current page.north west)},x=1bp,y=-1bp]
\node[anchor=north west,inner sep=0pt,text=answer,align=left,text width=294.5780029296875bp] at (15,15) {\fontsize{7.5bp}{9.5bp}\selectfont 析：\allowbreak{}C :\allowbreak{}全天\ensuremath{30} + \ensuremath{24} = \ensuremath{54}棵，\allowbreak{}\ensuremath{54} ÷ \ensuremath{6} = \ensuremath{9}棵。\allowbreak{}\par};
\end{scope}\end{tikzpicture}\newpage
\definecolor{answer}{rgb}{0.850980,0.047059,0.070588}
\null\begin{tikzpicture}[remember picture,overlay]\begin{scope}[shift={(current page.north west)},x=1bp,y=-1bp]
\node[anchor=north west,inner sep=0pt,text=answer,align=left,text width=399.70550537109375bp] at (15,15) {\fontsize{7.5bp}{9.5bp}\selectfont 析：\allowbreak{}A :\allowbreak{}先减去\ensuremath{40}与\ensuremath{30}的和，\allowbreak{}再把所得差除以\ensuremath{11},\allowbreak{}列[\ensuremath{180} \ensuremath{-}(\ensuremath{40} + \ensuremath{30})] ÷ \ensuremath{11}。\allowbreak{}\par};
\end{scope}\end{tikzpicture}\newpage
\definecolor{answer}{rgb}{0.850980,0.047059,0.070588}
\null\begin{tikzpicture}[remember picture,overlay]\begin{scope}[shift={(current page.north west)},x=1bp,y=-1bp]
\node[anchor=north west,inner sep=0pt,text=answer,align=left,text width=307.281005859375bp] at (15,15) {\fontsize{7.5bp}{9.5bp}\selectfont 析：\allowbreak{}B :\allowbreak{} \ensuremath{200} × \ensuremath{5} \ensuremath{-}\ensuremath{50} = \ensuremath{950}个，\allowbreak{}\ensuremath{950} ÷ \ensuremath{25} = \ensuremath{38}盒。\allowbreak{}\par};
\end{scope}\end{tikzpicture}\newpage
\definecolor{answer}{rgb}{0.850980,0.047059,0.070588}
\null\begin{tikzpicture}[remember picture,overlay]\begin{scope}[shift={(current page.north west)},x=1bp,y=-1bp]
\node[anchor=north west,inner sep=0pt,text=answer,align=left,text width=501.0bp] at (15,15) {\fontsize{8.5bp}{10.62bp}\selectfont 总吨数除以每车吨数求总车数，\allowbreak{}再减去已经运输的车数。\allowbreak{}\par};
\end{scope}\end{tikzpicture}\newpage
\definecolor{answer}{rgb}{0.850980,0.047059,0.070588}
\null\begin{tikzpicture}[remember picture,overlay]\begin{scope}[shift={(current page.north west)},x=1bp,y=-1bp]
\node[anchor=north west,inner sep=0pt,text=answer,align=left,text width=435.48248291015625bp] at (15,15) {\fontsize{7.5bp}{9.5bp}\selectfont 析：\allowbreak{}上排第\ensuremath{2}、\allowbreak{}\ensuremath{4}组垂直画○；\allowbreak{}下排第\ensuremath{1}、\allowbreak{}\ensuremath{2}组平行画△。\allowbreak{}垂直须成\ensuremath{90}\ensuremath{^\circ},\allowbreak{}平行须为同平面内不相交的直线，\allowbreak{}曲线不算。\allowbreak{}\par};
\end{scope}\end{tikzpicture}\newpage
\definecolor{answer}{rgb}{0.850980,0.047059,0.070588}
\null\begin{tikzpicture}[remember picture,overlay]\begin{scope}[shift={(current page.north west)},x=1bp,y=-1bp]
\node[anchor=north west,inner sep=0pt,text=answer,align=left,text width=398.5bp] at (15,15) {\fontsize{7.5bp}{9.5bp}\selectfont 析：\allowbreak{}选上排第\ensuremath{2}、\allowbreak{}\ensuremath{3}个和下排第\ensuremath{1}、\allowbreak{}\ensuremath{2}个：\allowbreak{}这些图中既有平行边，\allowbreak{}也有相交成直角的线段。\allowbreak{}\par};
\end{scope}\end{tikzpicture}\newpage
\definecolor{answer}{rgb}{0.850980,0.047059,0.070588}
\null\begin{tikzpicture}[remember picture,overlay]\begin{scope}[shift={(current page.north west)},x=1bp,y=-1bp]
\node[anchor=north west,inner sep=0pt,text=answer,align=left,text width=298.1304931640625bp] at (15,15) {\fontsize{7.5bp}{9.5bp}\selectfont 析：\allowbreak{}错：\allowbreak{}还可能斜交，\allowbreak{}交角不一定\ensuremath{90}\ensuremath{^\circ}。\allowbreak{}\par};
\end{scope}\end{tikzpicture}\newpage
\definecolor{answer}{rgb}{0.850980,0.047059,0.070588}
\null\begin{tikzpicture}[remember picture,overlay]\begin{scope}[shift={(current page.north west)},x=1bp,y=-1bp]
\node[anchor=north west,inner sep=0pt,text=answer,align=left,text width=305.6304931640625bp] at (15,15) {\fontsize{7.5bp}{9.5bp}\selectfont 析：\allowbreak{}错：\allowbreak{}\ensuremath{3}时半分针指\ensuremath{6},\allowbreak{}时针已走到\ensuremath{3}与\ensuremath{4}之间，\allowbreak{}两针较小夹角\ensuremath{75}\ensuremath{^\circ}。\allowbreak{}\par};
\end{scope}\end{tikzpicture}\newpage
\definecolor{answer}{rgb}{0.850980,0.047059,0.070588}
\null\begin{tikzpicture}[remember picture,overlay]\begin{scope}[shift={(current page.north west)},x=1bp,y=-1bp]
\node[anchor=north west,inner sep=0pt,text=answer,align=left,text width=308.5bp] at (15,15) {\fontsize{7.5bp}{9.5bp}\selectfont 析：\allowbreak{}错：\allowbreak{}平行只表示位置关系，\allowbreak{}线段长度可以不同。\allowbreak{}\par};
\end{scope}\end{tikzpicture}\newpage
\definecolor{answer}{rgb}{0.850980,0.047059,0.070588}
\null\begin{tikzpicture}[remember picture,overlay]\begin{scope}[shift={(current page.north west)},x=1bp,y=-1bp]
\node[anchor=north west,inner sep=0pt,text=answer,align=left,text width=353.5bp] at (15,15) {\fontsize{7.5bp}{9.5bp}\selectfont 析：\allowbreak{}错：\allowbreak{}必须在同一平面内；\allowbreak{}空间中不相交的直线不一定平行。\allowbreak{}\par};
\end{scope}\end{tikzpicture}\newpage
\definecolor{answer}{rgb}{0.850980,0.047059,0.070588}
\null\begin{tikzpicture}[remember picture,overlay]\begin{scope}[shift={(current page.north west)},x=1bp,y=-1bp]
\node[anchor=north west,inner sep=0pt,text=answer,align=left,text width=480.0bp] at (15,15) {\fontsize{8.5bp}{10.62bp}\selectfont 指示牌挂歪；\allowbreak{}牌的上边与天花板所在直线不平行。\allowbreak{}\par};
\end{scope}\end{tikzpicture}\newpage
\definecolor{answer}{rgb}{0.850980,0.047059,0.070588}
\null\begin{tikzpicture}[remember picture,overlay]\begin{scope}[shift={(current page.north west)},x=1bp,y=-1bp]
\node[anchor=north west,inner sep=0pt,text=answer,align=left,text width=394.4100036621094bp] at (15,15) {\fontsize{8.5bp}{10.62bp}\selectfont 过A点画与下水道平行的直线；\allowbreak{}下水道沿横向，\allowbreak{}所以过A的水管也应沿横向。\allowbreak{}\par};
\end{scope}\end{tikzpicture}\newpage
\definecolor{answer}{rgb}{0.850980,0.047059,0.070588}
\null\begin{tikzpicture}[remember picture,overlay]\begin{scope}[shift={(current page.north west)},x=1bp,y=-1bp]
\node[anchor=north west,inner sep=0pt,text=answer,align=left,text width=277.0bp] at (15,15) {\fontsize{8.5bp}{10.62bp}\selectfont 点到直线之间所有线段中，\allowbreak{}垂线段最短。\allowbreak{}\par};
\end{scope}\end{tikzpicture}\newpage
\definecolor{answer}{rgb}{0.850980,0.047059,0.070588}
\null\begin{tikzpicture}[remember picture,overlay]\begin{scope}[shift={(current page.north west)},x=1bp,y=-1bp]
\node[anchor=north west,inner sep=0pt,text=answer,align=left,text width=480.0bp] at (15,15) {\fontsize{7.5bp}{9.5bp}\selectfont 直线\ensuremath{a}与\ensuremath{b}垂直形成\ensuremath{90}°，\allowbreak{}∠\ensuremath{1}与∠\ensuremath{2}之和也是\ensuremath{90}°。\allowbreak{}\par};
\end{scope}\end{tikzpicture}\newpage
\definecolor{answer}{rgb}{0.850980,0.047059,0.070588}
\null\begin{tikzpicture}[remember picture,overlay]\begin{scope}[shift={(current page.north west)},x=1bp,y=-1bp]
\node[anchor=north west,inner sep=0pt,text=answer,align=left,text width=337.2149963378906bp] at (15,15) {\fontsize{8.5bp}{10.62bp}\selectfont 长方形一组对边间距等于其中一条邻边长，\allowbreak{}另一组间距用半周长减\ensuremath{12}米。\allowbreak{}\par};
\end{scope}\end{tikzpicture}\newpage
\definecolor{answer}{rgb}{0.850980,0.047059,0.070588}
\null\begin{tikzpicture}[remember picture,overlay]\begin{scope}[shift={(current page.north west)},x=1bp,y=-1bp]
\node[anchor=north west,inner sep=0pt,text=answer,align=left,text width=362.5299987792969bp] at (15,15) {\fontsize{7.5bp}{9.5bp}\selectfont 析：\allowbreak{}总价=单价×数量；\allowbreak{}单价=总价÷数量；\allowbreak{}数量=总价÷单价。\allowbreak{}\par};
\end{scope}\end{tikzpicture}\newpage
\definecolor{answer}{rgb}{0.850980,0.047059,0.070588}
\null\begin{tikzpicture}[remember picture,overlay]\begin{scope}[shift={(current page.north west)},x=1bp,y=-1bp]
\node[anchor=north west,inner sep=0pt,text=answer,align=left,text width=321.31201171875bp] at (15,15) {\fontsize{7.5bp}{9.5bp}\selectfont 析：\allowbreak{}单价\ensuremath{9}元/支读作\ensuremath{9}元每支；\allowbreak{}买\ensuremath{3}支，\allowbreak{}\ensuremath{9} × \ensuremath{3} = \ensuremath{27}元。\allowbreak{}\par};
\end{scope}\end{tikzpicture}\newpage
\definecolor{answer}{rgb}{0.850980,0.047059,0.070588}
\null\begin{tikzpicture}[remember picture,overlay]\begin{scope}[shift={(current page.north west)},x=1bp,y=-1bp]
\node[anchor=north west,inner sep=0pt,text=answer,align=left,text width=362.5299987792969bp] at (15,15) {\fontsize{7.5bp}{9.5bp}\selectfont 析：\allowbreak{}速度×时间=路程；\allowbreak{}路程÷时间=速度；\allowbreak{}路程÷速度=时间。\allowbreak{}\par};
\end{scope}\end{tikzpicture}\newpage
\definecolor{answer}{rgb}{0.850980,0.047059,0.070588}
\null\begin{tikzpicture}[remember picture,overlay]\begin{scope}[shift={(current page.north west)},x=1bp,y=-1bp]
\node[anchor=north west,inner sep=0pt,text=answer,align=left,text width=417.7349853515625bp] at (15,15) {\fontsize{7.5bp}{9.5bp}\selectfont 析：\allowbreak{}①速度和，\allowbreak{}单位千米/时；\allowbreak{}②相遇时间，\allowbreak{}单位小时；\allowbreak{}③尚未相遇时，\allowbreak{}\ensuremath{3}小时后两车相距的路程，\allowbreak{}单位千米。\allowbreak{}\par};
\end{scope}\end{tikzpicture}\newpage
\definecolor{answer}{rgb}{0.850980,0.047059,0.070588}
\null\begin{tikzpicture}[remember picture,overlay]\begin{scope}[shift={(current page.north west)},x=1bp,y=-1bp]
\node[anchor=north west,inner sep=0pt,text=answer,align=left,text width=350.1700134277344bp] at (15,15) {\fontsize{7.5bp}{9.5bp}\selectfont 析：\allowbreak{}返回时间\ensuremath{7}\ensuremath{-}\ensuremath{1}=\ensuremath{6}(小时)，\allowbreak{}速度\ensuremath{420}÷\ensuremath{6}=\ensuremath{70}(千米/时)，\allowbreak{}选C。\allowbreak{}\par};
\end{scope}\end{tikzpicture}\newpage
\definecolor{answer}{rgb}{0.850980,0.047059,0.070588}
\null\begin{tikzpicture}[remember picture,overlay]\begin{scope}[shift={(current page.north west)},x=1bp,y=-1bp]
\node[anchor=north west,inner sep=0pt,text=answer,align=left,text width=501.0bp] at (15,15) {\fontsize{7.5bp}{9.5bp}\selectfont 析：\allowbreak{}C :\allowbreak{}第二个部分积从十位起写，\allowbreak{}表示\ensuremath{243} × \ensuremath{20} = \ensuremath{4860}元，\allowbreak{}即\ensuremath{20}套总价。\allowbreak{}\par};
\end{scope}\end{tikzpicture}\newpage
\definecolor{answer}{rgb}{0.850980,0.047059,0.070588}
\null\begin{tikzpicture}[remember picture,overlay]\begin{scope}[shift={(current page.north west)},x=1bp,y=-1bp]
\node[anchor=north west,inner sep=0pt,text=answer,align=left,text width=501.0bp] at (15,15) {\fontsize{7.5bp}{9.5bp}\selectfont 析：\allowbreak{}速度=路程÷时间；\allowbreak{}依次\ensuremath{240} ÷ \ensuremath{4} = \ensuremath{60}千米/时、\allowbreak{}\ensuremath{264} ÷ \ensuremath{8} = \ensuremath{33}米/秒、\allowbreak{}\ensuremath{420} ÷ \ensuremath{6} = \ensuremath{70}千米/时、\allowbreak{}\ensuremath{112} ÷ \ensuremath{7} = \ensuremath{16}米/秒。\allowbreak{}\par};
\end{scope}\end{tikzpicture}\newpage
\definecolor{answer}{rgb}{0.850980,0.047059,0.070588}
\null\begin{tikzpicture}[remember picture,overlay]\begin{scope}[shift={(current page.north west)},x=1bp,y=-1bp]
\node[anchor=north west,inner sep=0pt,text=answer,align=left,text width=501.0bp] at (15,15) {\fontsize{8.5bp}{10.62bp}\selectfont 已知前一段路程与时间，\allowbreak{}速度为路程除以时间。\allowbreak{}\par};
\end{scope}\end{tikzpicture}\newpage
\definecolor{answer}{rgb}{0.850980,0.047059,0.070588}
\null\begin{tikzpicture}[remember picture,overlay]\begin{scope}[shift={(current page.north west)},x=1bp,y=-1bp]
\node[anchor=north west,inner sep=0pt,text=answer,align=left,text width=438.58245849609375bp] at (15,15) {\fontsize{7.5bp}{9.5bp}\selectfont 析：\allowbreak{}四律依次：\allowbreak{}\ensuremath{a} + \ensuremath{b} = \ensuremath{b} + \ensuremath{a};\allowbreak{} (\ensuremath{a} + \ensuremath{b}) + \ensuremath{c} = \ensuremath{a} + (\ensuremath{b} + \ensuremath{c});\allowbreak{}(\ensuremath{a} × \ensuremath{b}) × \ensuremath{c} = \ensuremath{a} × (\ensuremath{b} × \ensuremath{c});\allowbreak{} (\ensuremath{a} + \ensuremath{b}) × \ensuremath{c} = \ensuremath{a} × \ensuremath{c} + \ensuremath{b} × \ensuremath{c}。\allowbreak{}\par};
\end{scope}\end{tikzpicture}\newpage
\definecolor{answer}{rgb}{0.850980,0.047059,0.070588}
\null\begin{tikzpicture}[remember picture,overlay]\begin{scope}[shift={(current page.north west)},x=1bp,y=-1bp]
\node[anchor=north west,inner sep=0pt,text=answer,align=left,text width=301.0bp] at (15,15) {\fontsize{7.5bp}{9.5bp}\selectfont 析：\allowbreak{}(\ensuremath{100} + \ensuremath{1})分别乘\ensuremath{42}后相加，\allowbreak{}运用乘法分配律。\allowbreak{}\par};
\end{scope}\end{tikzpicture}\newpage
\definecolor{answer}{rgb}{0.850980,0.047059,0.070588}
\null\begin{tikzpicture}[remember picture,overlay]\begin{scope}[shift={(current page.north west)},x=1bp,y=-1bp]
\node[anchor=north west,inner sep=0pt,text=answer,align=left,text width=399.9534912109375bp] at (15,15) {\fontsize{7.5bp}{9.5bp}\selectfont 析：\allowbreak{}两块先分别数再相加，\allowbreak{}与合并后算(\ensuremath{6} + \ensuremath{4}) × \ensuremath{3}相同，\allowbreak{}运用乘法分配律。\allowbreak{}\par};
\end{scope}\end{tikzpicture}\newpage
\definecolor{answer}{rgb}{0.850980,0.047059,0.070588}
\null\begin{tikzpicture}[remember picture,overlay]\begin{scope}[shift={(current page.north west)},x=1bp,y=-1bp]
\node[anchor=north west,inner sep=0pt,text=answer,align=left,text width=362.93701171875bp] at (15,15) {\fontsize{7.5bp}{9.5bp}\selectfont 析：\allowbreak{}B :\allowbreak{}只改变计算分组，\allowbreak{}把\ensuremath{27}与\ensuremath{73}先相加，\allowbreak{}运用加法结合律。\allowbreak{}\par};
\end{scope}\end{tikzpicture}\newpage
\definecolor{answer}{rgb}{0.850980,0.047059,0.070588}
\null\begin{tikzpicture}[remember picture,overlay]\begin{scope}[shift={(current page.north west)},x=1bp,y=-1bp]
\node[anchor=north west,inner sep=0pt,text=answer,align=left,text width=335.4259948730469bp] at (15,15) {\fontsize{7.5bp}{9.5bp}\selectfont 析：\allowbreak{}C :\allowbreak{}把\ensuremath{98}当\ensuremath{100}多减\ensuremath{2},\allowbreak{}应补回\ensuremath{2},\allowbreak{}写\ensuremath{357} \ensuremath{-}\ensuremath{100} + \ensuremath{2}。\allowbreak{}\par};
\end{scope}\end{tikzpicture}\newpage
\definecolor{answer}{rgb}{0.850980,0.047059,0.070588}
\null\begin{tikzpicture}[remember picture,overlay]\begin{scope}[shift={(current page.north west)},x=1bp,y=-1bp]
\node[anchor=north west,inner sep=0pt,text=answer,align=left,text width=352.4179992675781bp] at (15,15) {\fontsize{7.5bp}{9.5bp}\selectfont 析：\allowbreak{}A :\allowbreak{} \ensuremath{99} = \ensuremath{100} \ensuremath{-}\ensuremath{1},\allowbreak{}所以\ensuremath{54} × \ensuremath{99} = \ensuremath{54} × \ensuremath{100} \ensuremath{-}\ensuremath{54}。\allowbreak{}\par};
\end{scope}\end{tikzpicture}\newpage
\definecolor{answer}{rgb}{0.850980,0.047059,0.070588}
\null\begin{tikzpicture}[remember picture,overlay]\begin{scope}[shift={(current page.north west)},x=1bp,y=-1bp]
\node[anchor=north west,inner sep=0pt,text=answer,align=left,text width=368.2770080566406bp] at (15,15) {\fontsize{7.5bp}{9.5bp}\selectfont 析：\allowbreak{}A :\allowbreak{}减去的\ensuremath{29}相当于\ensuremath{29} × \ensuremath{1},\allowbreak{}合并共同因数得\ensuremath{29} × (\ensuremath{33} + \ensuremath{68} \ensuremath{-}\ensuremath{1})。\allowbreak{}\par};
\end{scope}\end{tikzpicture}\newpage
\definecolor{answer}{rgb}{0.850980,0.047059,0.070588}
\null\begin{tikzpicture}[remember picture,overlay]\begin{scope}[shift={(current page.north west)},x=1bp,y=-1bp]
\node[anchor=north west,inner sep=0pt,text=answer,align=left,text width=357.24798583984375bp] at (15,15) {\fontsize{7.5bp}{9.5bp}\selectfont 析：\allowbreak{}A :\allowbreak{}正确有\ensuremath{5}个□和\ensuremath{5}个○，\allowbreak{}错式仅\ensuremath{1}个○，\allowbreak{}差\ensuremath{4}个○，\allowbreak{}应圈\ensuremath{4}个○。\allowbreak{}\par};
\end{scope}\end{tikzpicture}\newpage
\definecolor{answer}{rgb}{0.850980,0.047059,0.070588}
\null\begin{tikzpicture}[remember picture,overlay]\begin{scope}[shift={(current page.north west)},x=1bp,y=-1bp]
\node[anchor=north west,inner sep=0pt,text=answer,align=left,text width=370.989013671875bp] at (15,15) {\fontsize{7.5bp}{9.5bp}\selectfont 析：\allowbreak{}错：\allowbreak{}括号内是乘积，\allowbreak{}应写\ensuremath{25} × \ensuremath{4} × \ensuremath{8},\allowbreak{}不能改为两个积相加。\allowbreak{}\par};
\end{scope}\end{tikzpicture}\newpage
\definecolor{answer}{rgb}{0.850980,0.047059,0.070588}
\null\begin{tikzpicture}[remember picture,overlay]\begin{scope}[shift={(current page.north west)},x=1bp,y=-1bp]
\node[anchor=north west,inner sep=0pt,text=answer,align=left,text width=376.79901123046875bp] at (15,15) {\fontsize{7.5bp}{9.5bp}\selectfont 析：\allowbreak{}错：\allowbreak{}共同因数是\ensuremath{36},\allowbreak{}正确为\ensuremath{36} × (\ensuremath{9} + \ensuremath{1}),\allowbreak{}不是\ensuremath{35} × (\ensuremath{9} + \ensuremath{1})。\allowbreak{}\par};
\end{scope}\end{tikzpicture}\newpage
\definecolor{answer}{rgb}{0.850980,0.047059,0.070588}
\null\begin{tikzpicture}[remember picture,overlay]\begin{scope}[shift={(current page.north west)},x=1bp,y=-1bp]
\node[anchor=north west,inner sep=0pt,text=answer,align=left,text width=327.45599365234375bp] at (15,15) {\fontsize{7.5bp}{9.5bp}\selectfont 析：\allowbreak{}错：\allowbreak{}\ensuremath{15}应分别乘\ensuremath{6}和\ensuremath{2},\allowbreak{}正确为\ensuremath{15} × \ensuremath{6} + \ensuremath{15} × \ensuremath{2}。\allowbreak{}\par};
\end{scope}\end{tikzpicture}\newpage
\definecolor{answer}{rgb}{0.850980,0.047059,0.070588}
\null\begin{tikzpicture}[remember picture,overlay]\begin{scope}[shift={(current page.north west)},x=1bp,y=-1bp]
\node[anchor=north west,inner sep=0pt,text=answer,align=left,text width=336.3399963378906bp] at (15,15) {\fontsize{7.5bp}{9.5bp}\selectfont 析：\allowbreak{}对：\allowbreak{}\ensuremath{99} = \ensuremath{100} \ensuremath{-}\ensuremath{1},\allowbreak{}用分配律得\ensuremath{100} × \ensuremath{78} \ensuremath{-}\ensuremath{78}。\allowbreak{}\par};
\end{scope}\end{tikzpicture}\newpage
\definecolor{answer}{rgb}{0.850980,0.047059,0.070588}
\null\begin{tikzpicture}[remember picture,overlay]\begin{scope}[shift={(current page.north west)},x=1bp,y=-1bp]
\node[anchor=north west,inner sep=0pt,text=answer,align=left,text width=301.0bp] at (15,15) {\fontsize{7.5bp}{9.5bp}\selectfont 析：\allowbreak{}对：\allowbreak{}只改变三个因数相乘的分组，\allowbreak{}运用乘法结合律。\allowbreak{}\par};
\end{scope}\end{tikzpicture}\newpage
\definecolor{answer}{rgb}{0.850980,0.047059,0.070588}
\null\begin{tikzpicture}[remember picture,overlay]\begin{scope}[shift={(current page.north west)},x=1bp,y=-1bp]
\node[anchor=north west,inner sep=0pt,text=answer,align=left,text width=501.375bp] at (15,15) {\fontsize{8.5bp}{10.62bp}\selectfont 每组\ensuremath{8}片，\allowbreak{}每班\ensuremath{6}组，\allowbreak{}共\ensuremath{15}班；\allowbreak{}先求总需求再与\ensuremath{800}片比较。\allowbreak{}\par};
\end{scope}\end{tikzpicture}\newpage
\definecolor{answer}{rgb}{0.850980,0.047059,0.070588}
\null\begin{tikzpicture}[remember picture,overlay]\begin{scope}[shift={(current page.north west)},x=1bp,y=-1bp]
\node[anchor=north west,inner sep=0pt,text=answer,align=left,text width=209.0bp] at (15,15) {\fontsize{8.5bp}{10.62bp}\selectfont 图中两个高度差为\ensuremath{198}厘米和\ensuremath{102}厘米，\allowbreak{}依次从\ensuremath{386}厘米中扣去。\allowbreak{}\par};
\end{scope}\end{tikzpicture}\newpage
\definecolor{answer}{rgb}{0.850980,0.047059,0.070588}
\null\begin{tikzpicture}[remember picture,overlay]\begin{scope}[shift={(current page.north west)},x=1bp,y=-1bp]
\node[anchor=north west,inner sep=0pt,text=answer,align=left,text width=501.0bp] at (15,15) {\fontsize{8.5bp}{10.62bp}\selectfont 同样走\ensuremath{4}分钟，\allowbreak{}速度较快的小英家离学校更远。\allowbreak{}\par};
\end{scope}\end{tikzpicture}\newpage
\definecolor{answer}{rgb}{0.850980,0.047059,0.070588}
\null\begin{tikzpicture}[remember picture,overlay]\begin{scope}[shift={(current page.north west)},x=1bp,y=-1bp]
\node[anchor=north west,inner sep=0pt,text=answer,align=left,text width=391.0bp] at (15,15) {\fontsize{7.5bp}{9.5bp}\selectfont 析：\allowbreak{}红领巾有钝角且两腰相等，\allowbreak{}按角分为钝角三角形，\allowbreak{}按边分为等腰三角形。\allowbreak{}\par};
\end{scope}\end{tikzpicture}\newpage
\definecolor{answer}{rgb}{0.850980,0.047059,0.070588}
\null\begin{tikzpicture}[remember picture,overlay]\begin{scope}[shift={(current page.north west)},x=1bp,y=-1bp]
\node[anchor=north west,inner sep=0pt,text=answer,align=left,text width=320.8599967956543bp] at (15,15) {\fontsize{6.6bp}{8.6bp}\selectfont *整边拼接共\ensuremath{6}种：\allowbreak{}表中\ensuremath{5}种，\allowbreak{}另一个四边形周长\ensuremath{14}厘米。\allowbreak{}\par};
\end{scope}\end{tikzpicture}\newpage
\definecolor{answer}{rgb}{0.850980,0.047059,0.070588}
\null\begin{tikzpicture}[remember picture,overlay]\begin{scope}[shift={(current page.north west)},x=1bp,y=-1bp]
\node[anchor=north west,inner sep=0pt,text=answer,align=left,text width=446.50628662109375bp] at (15,15) {\fontsize{7.52bp}{9.52bp}\selectfont 析：\allowbreak{}（\ensuremath{3}）先合并两个已知角：\allowbreak{}\ensuremath{180}\ensuremath{^\circ}\ensuremath{-}\ensuremath{54}\ensuremath{^\circ}\ensuremath{-}\ensuremath{66}\ensuremath{^\circ}= \ensuremath{180}\ensuremath{^\circ}\ensuremath{-}(\ensuremath{54}\ensuremath{^\circ}+ \ensuremath{66}\ensuremath{^\circ}) = \ensuremath{180}\ensuremath{^\circ}\ensuremath{-}\ensuremath{120}\ensuremath{^\circ}= \ensuremath{60}\ensuremath{^\circ}\par};
\end{scope}\end{tikzpicture}\newpage
\definecolor{answer}{rgb}{0.850980,0.047059,0.070588}
\null\begin{tikzpicture}[remember picture,overlay]\begin{scope}[shift={(current page.north west)},x=1bp,y=-1bp]
\node[anchor=north west,inner sep=0pt,text=answer,align=left,text width=297.7340087890625bp] at (15,15) {\fontsize{7.5bp}{9.5bp}\selectfont 析：\allowbreak{}（\ensuremath{4})第三边须大于\ensuremath{14} \ensuremath{-}\ensuremath{6} = \ensuremath{8}厘米且小于\ensuremath{20}厘米，\allowbreak{}整数最小\ensuremath{9}厘米。\allowbreak{}\par};
\end{scope}\end{tikzpicture}\newpage
\definecolor{answer}{rgb}{0.850980,0.047059,0.070588}
\null\begin{tikzpicture}[remember picture,overlay]\begin{scope}[shift={(current page.north west)},x=1bp,y=-1bp]
\node[anchor=north west,inner sep=0pt,text=answer,align=left,text width=329.0108642578125bp] at (15,15) {\fontsize{7.5bp}{9.5bp}\selectfont 析：\allowbreak{}两个底角相等，\allowbreak{}(\ensuremath{180}\ensuremath{^\circ}\ensuremath{-}\ensuremath{90}\ensuremath{^\circ}) ÷ \ensuremath{2} = \ensuremath{45}\ensuremath{^\circ}。\allowbreak{}\par};
\end{scope}\end{tikzpicture}\newpage
\definecolor{answer}{rgb}{0.850980,0.047059,0.070588}
\null\begin{tikzpicture}[remember picture,overlay]\begin{scope}[shift={(current page.north west)},x=1bp,y=-1bp]
\node[anchor=north west,inner sep=0pt,text=answer,align=left,text width=308.5bp] at (15,15) {\fontsize{7.5bp}{9.5bp}\selectfont 析：\allowbreak{}A :\allowbreak{}从同一个点到指定直线，\allowbreak{}只能作一条垂线段。\allowbreak{}\par};
\end{scope}\end{tikzpicture}\newpage
\definecolor{answer}{rgb}{0.850980,0.047059,0.070588}
\null\begin{tikzpicture}[remember picture,overlay]\begin{scope}[shift={(current page.north west)},x=1bp,y=-1bp]
\node[anchor=north west,inner sep=0pt,text=answer,align=left,text width=406.0bp] at (15,15) {\fontsize{7.5bp}{9.5bp}\selectfont 析：\allowbreak{}A：\allowbreak{}两个形状、\allowbreak{}大小完全相同的三角形沿对应边拼合，\allowbreak{}可使另外两组对应边分别平行，\allowbreak{}拼成平行四边形。\allowbreak{}\par};
\end{scope}\end{tikzpicture}\newpage
\definecolor{answer}{rgb}{0.850980,0.047059,0.070588}
\null\begin{tikzpicture}[remember picture,overlay]\begin{scope}[shift={(current page.north west)},x=1bp,y=-1bp]
\node[anchor=north west,inner sep=0pt,text=answer,align=left,text width=316.0bp] at (15,15) {\fontsize{7.5bp}{9.5bp}\selectfont 析：\allowbreak{}A :\allowbreak{}斜撑形成三角形，\allowbreak{}利用三角形的稳定性。\allowbreak{}\par};
\end{scope}\end{tikzpicture}\newpage
\definecolor{answer}{rgb}{0.850980,0.047059,0.070588}
\null\begin{tikzpicture}[remember picture,overlay]\begin{scope}[shift={(current page.north west)},x=1bp,y=-1bp]
\node[anchor=north west,inner sep=0pt,text=answer,align=left,text width=396.5610046386719bp] at (15,15) {\fontsize{7.5bp}{9.5bp}\selectfont 析：\allowbreak{}A只露一个锐角，\allowbreak{}三类三角形都有锐角，\allowbreak{}不能确定；\allowbreak{}B露直角、\allowbreak{}C露钝角，\allowbreak{}D可由两斜边方向判断。\allowbreak{}\par};
\end{scope}\end{tikzpicture}\newpage
\definecolor{answer}{rgb}{0.850980,0.047059,0.070588}
\null\begin{tikzpicture}[remember picture,overlay]\begin{scope}[shift={(current page.north west)},x=1bp,y=-1bp]
\node[anchor=north west,inner sep=0pt,text=answer,align=left,text width=501.0bp] at (15,15) {\fontsize{7.5bp}{9.5bp}\selectfont 析：\allowbreak{}C :\allowbreak{}两腰各\ensuremath{30}厘米，\allowbreak{}底边\ensuremath{85} \ensuremath{-}\ensuremath{30} × \ensuremath{2} = \ensuremath{25}厘米。\allowbreak{}\par};
\end{scope}\end{tikzpicture}\newpage
\definecolor{answer}{rgb}{0.850980,0.047059,0.070588}
\null\begin{tikzpicture}[remember picture,overlay]\begin{scope}[shift={(current page.north west)},x=1bp,y=-1bp]
\node[anchor=north west,inner sep=0pt,text=answer,align=left,text width=501.0bp] at (15,15) {\fontsize{7.5bp}{9.5bp}\selectfont 析：\allowbreak{}B :\allowbreak{}两底角相等，\allowbreak{}顶角\ensuremath{180}\ensuremath{^\circ}\ensuremath{-}\ensuremath{35}\ensuremath{^\circ}× \ensuremath{2} = \ensuremath{110}\ensuremath{^\circ}。\allowbreak{}\par};
\end{scope}\end{tikzpicture}\newpage
\definecolor{answer}{rgb}{0.850980,0.047059,0.070588}
\null\begin{tikzpicture}[remember picture,overlay]\begin{scope}[shift={(current page.north west)},x=1bp,y=-1bp]
\node[anchor=north west,inner sep=0pt,text=answer,align=left,text width=501.0bp] at (15,15) {\fontsize{7.5bp}{9.5bp}\selectfont 析：\allowbreak{}C :\allowbreak{} \ensuremath{5} + \ensuremath{6} > \ensuremath{10}可围成；\allowbreak{}A两短边和等于最长边，\allowbreak{}B、\allowbreak{}D小于最长边，\allowbreak{}均不行。\allowbreak{}\par};
\end{scope}\end{tikzpicture}\newpage
\definecolor{answer}{rgb}{0.850980,0.047059,0.070588}
\null\begin{tikzpicture}[remember picture,overlay]\begin{scope}[shift={(current page.north west)},x=1bp,y=-1bp]
\node[anchor=north west,inner sep=0pt,text=answer,align=left,text width=501.0bp] at (15,15) {\fontsize{7.5bp}{9.5bp}\selectfont 析：\allowbreak{}错：\allowbreak{}平行四边形容易变形；\allowbreak{}边长确定，\allowbreak{}角度仍可变，\allowbreak{}形状不一定确定。\allowbreak{}\par};
\end{scope}\end{tikzpicture}\newpage
\definecolor{answer}{rgb}{0.850980,0.047059,0.070588}
\null\begin{tikzpicture}[remember picture,overlay]\begin{scope}[shift={(current page.north west)},x=1bp,y=-1bp]
\node[anchor=north west,inner sep=0pt,text=answer,align=left,text width=501.0bp] at (15,15) {\fontsize{7.5bp}{9.5bp}\selectfont 析：\allowbreak{}对：\allowbreak{}沿两直角边中点连线剪下小直角三角形，\allowbreak{}旋转后可与余下梯形拼成平行四边形。\allowbreak{}\par};
\end{scope}\end{tikzpicture}\newpage
\definecolor{answer}{rgb}{0.850980,0.047059,0.070588}
\null\begin{tikzpicture}[remember picture,overlay]\begin{scope}[shift={(current page.north west)},x=1bp,y=-1bp]
\node[anchor=north west,inner sep=0pt,text=answer,align=left,text width=501.0bp] at (15,15) {\fontsize{7.5bp}{9.5bp}\selectfont 析：\allowbreak{}对：\allowbreak{}长方形是特殊平行四边形，\allowbreak{}正方形又是四角为直角的特殊长方形。\allowbreak{}\par};
\end{scope}\end{tikzpicture}\newpage
\definecolor{answer}{rgb}{0.850980,0.047059,0.070588}
\null\begin{tikzpicture}[remember picture,overlay]\begin{scope}[shift={(current page.north west)},x=1bp,y=-1bp]
\node[anchor=north west,inner sep=0pt,text=answer,align=left,text width=436.0bp] at (15,15) {\fontsize{7.5bp}{9.5bp}\selectfont 析：\allowbreak{}依次为三角形、\allowbreak{}平行四边形、\allowbreak{}梯形、\allowbreak{}三角形；\allowbreak{}高须垂直于指定底边所在直线，\allowbreak{}末图需向右延长底边再落垂足。\allowbreak{}\par};
\end{scope}\end{tikzpicture}\newpage
\definecolor{answer}{rgb}{0.850980,0.047059,0.070588}
\null\begin{tikzpicture}[remember picture,overlay]\begin{scope}[shift={(current page.north west)},x=1bp,y=-1bp]
\node[anchor=north west,inner sep=0pt,text=answer,align=left,text width=474.0bp] at (15,15) {\fontsize{8.5bp}{10.62bp}\selectfont 周长用四条边长相加；\allowbreak{}\ensuremath{3}米是高，\allowbreak{}不能当作邻边长来算周长。\allowbreak{}\par};
\end{scope}\end{tikzpicture}\newpage
\definecolor{answer}{rgb}{0.850980,0.047059,0.070588}
\null\begin{tikzpicture}[remember picture,overlay]\begin{scope}[shift={(current page.north west)},x=1bp,y=-1bp]
\node[anchor=north west,inner sep=0pt,text=answer,align=left,text width=501.375bp] at (15,15) {\fontsize{8.5bp}{10.62bp}\selectfont 等腰梯形的两条腰相等；\allowbreak{}先把两底\ensuremath{11}米和\ensuremath{19}米合成\ensuremath{30}米。\allowbreak{}\par};
\end{scope}\end{tikzpicture}\newpage
\definecolor{answer}{rgb}{0.850980,0.047059,0.070588}
\null\begin{tikzpicture}[remember picture,overlay]\begin{scope}[shift={(current page.north west)},x=1bp,y=-1bp]
\node[anchor=north west,inner sep=0pt,text=answer,align=left,text width=501.0bp] at (15,15) {\fontsize{7.5bp}{9.5bp}\selectfont 析：\allowbreak{}固定底边，\allowbreak{}顶点在底边的垂直平分线上；\allowbreak{}高越长，\allowbreak{}两腰越靠近中线，\allowbreak{}顶角越小。\allowbreak{}\par};
\end{scope}\end{tikzpicture}\newpage
\definecolor{answer}{rgb}{0.850980,0.047059,0.070588}
\null\begin{tikzpicture}[remember picture,overlay]\begin{scope}[shift={(current page.north west)},x=1bp,y=-1bp]
\node[anchor=north west,inner sep=0pt,text=answer,align=left,text width=501.0bp] at (15,15) {\fontsize{7.5bp}{9.5bp}\selectfont 析：\allowbreak{}在同一底边中点的垂线上改取顶点并连两端，\allowbreak{}可见高越长顶角越小，\allowbreak{}高越短顶角越大。\allowbreak{}\par};
\end{scope}\end{tikzpicture}\newpage
\definecolor{answer}{rgb}{0.850980,0.047059,0.070588}
\null\begin{tikzpicture}[remember picture,overlay]\begin{scope}[shift={(current page.north west)},x=1bp,y=-1bp]
\node[anchor=north west,inner sep=0pt,text=answer,align=left,text width=382.2250061035156bp] at (15,15) {\fontsize{7.5bp}{9.5bp}\selectfont 析：\allowbreak{}先算\ensuremath{30} × \ensuremath{8}为①④；\allowbreak{}最后不是除法为③；\allowbreak{}④结果\ensuremath{0};\allowbreak{}②③都相当于\ensuremath{210} × \ensuremath{2},\allowbreak{}结果相等。\allowbreak{}\par};
\end{scope}\end{tikzpicture}\newpage
\definecolor{answer}{rgb}{0.850980,0.047059,0.070588}
\null\begin{tikzpicture}[remember picture,overlay]\begin{scope}[shift={(current page.north west)},x=1bp,y=-1bp]
\node[anchor=north west,inner sep=0pt,text=answer,align=left,text width=383.5bp] at (15,15) {\fontsize{7.5bp}{9.5bp}\selectfont 析：\allowbreak{}①⑤相交成直角，\allowbreak{}互相垂直；\allowbreak{}③④两线方向相同，\allowbreak{}延长不相交，\allowbreak{}互相平行；\allowbreak{}②两线方向不同。\allowbreak{}\par};
\end{scope}\end{tikzpicture}\newpage
\definecolor{answer}{rgb}{0.850980,0.047059,0.070588}
\null\begin{tikzpicture}[remember picture,overlay]\begin{scope}[shift={(current page.north west)},x=1bp,y=-1bp]
\node[anchor=north west,inner sep=0pt,text=answer,align=left,text width=370.093505859375bp] at (15,15) {\fontsize{7.5bp}{9.5bp}\selectfont 析：\allowbreak{}比较速度：\allowbreak{}红红\ensuremath{1000} ÷ \ensuremath{20} = \ensuremath{50}米/分，\allowbreak{}明明\ensuremath{900} ÷ \ensuremath{15} = \ensuremath{60}米/分；\allowbreak{}明明更快。\allowbreak{}\par};
\end{scope}\end{tikzpicture}\newpage
\definecolor{answer}{rgb}{0.850980,0.047059,0.070588}
\null\begin{tikzpicture}[remember picture,overlay]\begin{scope}[shift={(current page.north west)},x=1bp,y=-1bp]
\node[anchor=north west,inner sep=0pt,text=answer,align=left,text width=335.8590087890625bp] at (15,15) {\fontsize{7.5bp}{9.5bp}\selectfont 析：\allowbreak{}速度和\ensuremath{288} ÷ \ensuremath{3} = \ensuremath{96}千米/时，\allowbreak{}乙车\ensuremath{96} \ensuremath{-}\ensuremath{45} = \ensuremath{51}千米/时。\allowbreak{}\par};
\end{scope}\end{tikzpicture}\newpage
\definecolor{answer}{rgb}{0.850980,0.047059,0.070588}
\null\begin{tikzpicture}[remember picture,overlay]\begin{scope}[shift={(current page.north west)},x=1bp,y=-1bp]
\node[anchor=north west,inner sep=0pt,text=answer,align=left,text width=351.2569885253906bp] at (15,15) {\fontsize{7.5bp}{9.5bp}\selectfont 析：\allowbreak{}前式提出共同因数\ensuremath{4},\allowbreak{}填\ensuremath{138}、\allowbreak{}\ensuremath{38}、\allowbreak{}\ensuremath{4};\allowbreak{}后式只改分组，\allowbreak{}填\ensuremath{25}、\allowbreak{}\ensuremath{4}。\allowbreak{}\par};
\end{scope}\end{tikzpicture}\newpage
\definecolor{answer}{rgb}{0.850980,0.047059,0.070588}
\null\begin{tikzpicture}[remember picture,overlay]\begin{scope}[shift={(current page.north west)},x=1bp,y=-1bp]
\node[anchor=north west,inner sep=0pt,text=answer,align=left,text width=484.0bp] at (15,15) {\fontsize{7.5bp}{9.5bp}\selectfont 析：\allowbreak{}先求剩余量，\allowbreak{}再平均分\ensuremath{7}天，\allowbreak{}列(\ensuremath{1000}\ensuremath{-}\ensuremath{510})÷\ensuremath{7}，\allowbreak{}选C。\allowbreak{}\par};
\end{scope}\end{tikzpicture}\newpage
\definecolor{answer}{rgb}{0.850980,0.047059,0.070588}
\null\begin{tikzpicture}[remember picture,overlay]\begin{scope}[shift={(current page.north west)},x=1bp,y=-1bp]
\node[anchor=north west,inner sep=0pt,text=answer,align=left,text width=484.0bp] at (15,15) {\fontsize{7.5bp}{9.5bp}\selectfont 析：\allowbreak{}把\ensuremath{3}与\ensuremath{9}、\allowbreak{}\ensuremath{25}与\ensuremath{4}结合，\allowbreak{}既换位置又改结合顺序，\allowbreak{}用交换律和结合律，\allowbreak{}选C。\allowbreak{}\par};
\end{scope}\end{tikzpicture}\newpage
\definecolor{answer}{rgb}{0.850980,0.047059,0.070588}
\null\begin{tikzpicture}[remember picture,overlay]\begin{scope}[shift={(current page.north west)},x=1bp,y=-1bp]
\node[anchor=north west,inner sep=0pt,text=answer,align=left,text width=484.0bp] at (15,15) {\fontsize{7.5bp}{9.5bp}\selectfont 析：\allowbreak{}等边三角形每角\ensuremath{180}\ensuremath{^\circ}÷ \ensuremath{3} = \ensuremath{60}\ensuremath{^\circ};\allowbreak{}其余说法不成立，\allowbreak{}选C。\allowbreak{}\par};
\end{scope}\end{tikzpicture}\newpage
\definecolor{answer}{rgb}{0.850980,0.047059,0.070588}
\null\begin{tikzpicture}[remember picture,overlay]\begin{scope}[shift={(current page.north west)},x=1bp,y=-1bp]
\node[anchor=north west,inner sep=0pt,text=answer,align=left,text width=484.0bp] at (15,15) {\fontsize{7.5bp}{9.5bp}\selectfont 析：\allowbreak{}找钱=\ensuremath{200} \ensuremath{-}\ensuremath{24} × \ensuremath{8} = \ensuremath{8}元，\allowbreak{}需①③④；\allowbreak{}每盒支数无关，\allowbreak{}选C。\allowbreak{}\par};
\end{scope}\end{tikzpicture}\newpage
\definecolor{answer}{rgb}{0.850980,0.047059,0.070588}
\null\begin{tikzpicture}[remember picture,overlay]\begin{scope}[shift={(current page.north west)},x=1bp,y=-1bp]
\node[anchor=north west,inner sep=0pt,text=answer,align=left,text width=484.0bp] at (15,15) {\fontsize{7.5bp}{9.5bp}\selectfont 析：\allowbreak{}平行线中一条垂直于\ensuremath{c},\allowbreak{}另一条也垂直于\ensuremath{c};\allowbreak{} \ensuremath{a} ⊥\ensuremath{c},\allowbreak{}选A。\allowbreak{}\par};
\end{scope}\end{tikzpicture}\newpage
\definecolor{answer}{rgb}{0.850980,0.047059,0.070588}
\null\begin{tikzpicture}[remember picture,overlay]\begin{scope}[shift={(current page.north west)},x=1bp,y=-1bp]
\node[anchor=north west,inner sep=0pt,text=answer,align=left,text width=484.0bp] at (15,15) {\fontsize{8.5bp}{10.62bp}\selectfont 两队共同完成\ensuremath{1200}米，\allowbreak{}\ensuremath{8}天的平均合挖量减去甲队每天挖的米数，\allowbreak{}就是乙队每天挖的米数。\allowbreak{}\par};
\end{scope}\end{tikzpicture}\newpage
\definecolor{answer}{rgb}{0.850980,0.047059,0.070588}
\null\begin{tikzpicture}[remember picture,overlay]\begin{scope}[shift={(current page.north west)},x=1bp,y=-1bp]
\node[anchor=north west,inner sep=0pt,text=answer,align=left,text width=484.0bp] at (15,15) {\fontsize{8.5bp}{10.62bp}\selectfont 两车同时相向而行，\allowbreak{}相遇时两车所走路程的和，\allowbreak{}正好等于A、\allowbreak{}B两地的路程。\allowbreak{}\par};
\end{scope}\end{tikzpicture}\newpage
\definecolor{answer}{rgb}{0.850980,0.047059,0.070588}
\null\begin{tikzpicture}[remember picture,overlay]\begin{scope}[shift={(current page.north west)},x=1bp,y=-1bp]
\node[anchor=north west,inner sep=0pt,text=answer,align=left,text width=484.0bp] at (15,15) {\fontsize{8.5bp}{10.62bp}\selectfont 每天\ensuremath{4}场，\allowbreak{}每场\ensuremath{150}个座位，\allowbreak{}最多售出\ensuremath{600}张票；\allowbreak{}按每张\ensuremath{25}元计算营业额。\allowbreak{}\par};
\end{scope}\end{tikzpicture}\newpage
\definecolor{answer}{rgb}{0.850980,0.047059,0.070588}
\null\begin{tikzpicture}[remember picture,overlay]\begin{scope}[shift={(current page.north west)},x=1bp,y=-1bp]
\node[anchor=north west,inner sep=0pt,text=answer,align=left,text width=484.0bp] at (15,15) {\fontsize{8.5bp}{10.62bp}\selectfont 松树苗比杨树苗多\ensuremath{174}棵，\allowbreak{}用松树苗棵数减\ensuremath{174},\allowbreak{}先求杨树苗的总棵数，\allowbreak{}再平均分成\ensuremath{4}行。\allowbreak{}\par};
\end{scope}\end{tikzpicture}\newpage
\definecolor{answer}{rgb}{0.850980,0.047059,0.070588}
\null\begin{tikzpicture}[remember picture,overlay]\begin{scope}[shift={(current page.north west)},x=1bp,y=-1bp]
\node[anchor=north west,inner sep=0pt,text=answer,align=left,text width=484.0bp] at (15,15) {\fontsize{8.5bp}{10.62bp}\selectfont \ensuremath{1}台电饭煲\ensuremath{190}元，\allowbreak{}\ensuremath{2}台电扇每台\ensuremath{204}元，\allowbreak{}应把电饭煲的钱和\ensuremath{2}台电扇的钱相加。\allowbreak{}\par};
\end{scope}\end{tikzpicture}\newpage
\definecolor{answer}{rgb}{0.850980,0.047059,0.070588}
\null\begin{tikzpicture}[remember picture,overlay]\begin{scope}[shift={(current page.north west)},x=1bp,y=-1bp]
\node[anchor=north west,inner sep=0pt,text=answer,align=left,text width=484.0bp] at (15,15) {\fontsize{8.5bp}{10.62bp}\selectfont 先确定\ensuremath{4}台电饭煲花\ensuremath{760}元，\allowbreak{}余\ensuremath{240}元；\allowbreak{}\ensuremath{3}台音箱需\ensuremath{264}元，\allowbreak{}因此还差\ensuremath{24}元。\allowbreak{}\par};
\end{scope}\end{tikzpicture}\newpage
\definecolor{answer}{rgb}{0.850980,0.047059,0.070588}
\null\begin{tikzpicture}[remember picture,overlay]\begin{scope}[shift={(current page.north west)},x=1bp,y=-1bp]
\node[anchor=north west,inner sep=0pt,text=answer,align=left,text width=484.0bp] at (15,15) {\fontsize{7.5bp}{9.5bp}\selectfont 析：\allowbreak{}按分步顺序嵌套括号：\allowbreak{}①[(\ensuremath{48} \ensuremath{-}\ensuremath{23}) + \ensuremath{4}] × \ensuremath{102};\allowbreak{}②\ensuremath{64} ÷ [\ensuremath{31} \ensuremath{-}(\ensuremath{35} \ensuremath{-}\ensuremath{20})]。\allowbreak{}\par};
\end{scope}\end{tikzpicture}\newpage
\definecolor{answer}{rgb}{0.850980,0.047059,0.070588}
\null\begin{tikzpicture}[remember picture,overlay]\begin{scope}[shift={(current page.north west)},x=1bp,y=-1bp]
\node[anchor=north west,inner sep=0pt,text=answer,align=left,text width=484.0bp] at (15,15) {\fontsize{7.5bp}{9.5bp}\selectfont 析：\allowbreak{}等边三角形每角\ensuremath{60}\ensuremath{^\circ};\allowbreak{}对折平分一个角，\allowbreak{}得\ensuremath{30}\ensuremath{^\circ}、\allowbreak{}\ensuremath{60}\ensuremath{^\circ}两个锐角。\allowbreak{}\par};
\end{scope}\end{tikzpicture}\newpage
\definecolor{answer}{rgb}{0.850980,0.047059,0.070588}
\null\begin{tikzpicture}[remember picture,overlay]\begin{scope}[shift={(current page.north west)},x=1bp,y=-1bp]
\node[anchor=north west,inner sep=0pt,text=answer,align=left,text width=484.0bp] at (15,15) {\fontsize{7.5bp}{9.5bp}\selectfont 析：\allowbreak{}\ensuremath{438} ÷ \ensuremath{3}是前\ensuremath{3}个月平均每月节水量；\allowbreak{}再乘(\ensuremath{12} \ensuremath{-}\ensuremath{3}),\allowbreak{}得后\ensuremath{9}个月共节水量。\allowbreak{}\par};
\end{scope}\end{tikzpicture}\newpage
\definecolor{answer}{rgb}{0.850980,0.047059,0.070588}
\null\begin{tikzpicture}[remember picture,overlay]\begin{scope}[shift={(current page.north west)},x=1bp,y=-1bp]
\node[anchor=north west,inner sep=0pt,text=answer,align=left,text width=484.0bp] at (15,15) {\fontsize{7.5bp}{9.5bp}\selectfont 析：\allowbreak{}□=(\ensuremath{18} + \ensuremath{72}) ÷ \ensuremath{10} = \ensuremath{9};\allowbreak{}□-(□-▲)=▲=\ensuremath{4};\allowbreak{}○=\ensuremath{94} \ensuremath{-}\ensuremath{44} × \ensuremath{2} = \ensuremath{6},\allowbreak{}密码▲○□为\ensuremath{469}。\allowbreak{}\par};
\end{scope}\end{tikzpicture}\newpage
\definecolor{answer}{rgb}{0.850980,0.047059,0.070588}
\null\begin{tikzpicture}[remember picture,overlay]\begin{scope}[shift={(current page.north west)},x=1bp,y=-1bp]
\node[anchor=north west,inner sep=0pt,text=answer,align=left,text width=484.0bp] at (15,15) {\fontsize{7.5bp}{9.5bp}\selectfont 析：\allowbreak{}拆\ensuremath{36}为\ensuremath{30} + \ensuremath{6}对应图②；\allowbreak{}把\ensuremath{36}看成\ensuremath{40} \ensuremath{-}\ensuremath{4}对应图③；\allowbreak{}拆\ensuremath{45}为\ensuremath{40} + \ensuremath{5}对应图①。\allowbreak{}\par};
\end{scope}\end{tikzpicture}\newpage
\definecolor{answer}{rgb}{0.850980,0.047059,0.070588}
\null\begin{tikzpicture}[remember picture,overlay]\begin{scope}[shift={(current page.north west)},x=1bp,y=-1bp]
\node[anchor=north west,inner sep=0pt,text=answer,align=left,text width=484.0bp] at (15,15) {\fontsize{7.5bp}{9.5bp}\selectfont 析：\allowbreak{}同时出发、\allowbreak{}同刻相遇，\allowbreak{}行驶时间相同，\allowbreak{}选C。\allowbreak{}\par};
\end{scope}\end{tikzpicture}\newpage
\definecolor{answer}{rgb}{0.850980,0.047059,0.070588}
\null\begin{tikzpicture}[remember picture,overlay]\begin{scope}[shift={(current page.north west)},x=1bp,y=-1bp]
\node[anchor=north west,inner sep=0pt,text=answer,align=left,text width=484.0bp] at (15,15) {\fontsize{7.5bp}{9.5bp}\selectfont 析：\allowbreak{}B去括号仍先算乘除再相加；\allowbreak{}A、\allowbreak{}C的运算顺序会变，\allowbreak{}选B。\allowbreak{}\par};
\end{scope}\end{tikzpicture}\newpage
\definecolor{answer}{rgb}{0.850980,0.047059,0.070588}
\null\begin{tikzpicture}[remember picture,overlay]\begin{scope}[shift={(current page.north west)},x=1bp,y=-1bp]
\node[anchor=north west,inner sep=0pt,text=answer,align=left,text width=484.0bp] at (15,15) {\fontsize{7.5bp}{9.5bp}\selectfont 析：\allowbreak{}第三边必须小于\ensuremath{75} + \ensuremath{75} = \ensuremath{150}厘米，\allowbreak{}只有\ensuremath{140}厘米符合，\allowbreak{}选A。\allowbreak{}\par};
\end{scope}\end{tikzpicture}\newpage
\definecolor{answer}{rgb}{0.850980,0.047059,0.070588}
\null\begin{tikzpicture}[remember picture,overlay]\begin{scope}[shift={(current page.north west)},x=1bp,y=-1bp]
\node[anchor=north west,inner sep=0pt,text=answer,align=left,text width=484.0bp] at (15,15) {\fontsize{7.5bp}{9.5bp}\selectfont 析：\allowbreak{}○×(\ensuremath{33} + \ensuremath{17} + \ensuremath{50}) = \ensuremath{3000},\allowbreak{}○=\ensuremath{3000} ÷ \ensuremath{100} = \ensuremath{30},\allowbreak{}选B。\allowbreak{}\par};
\end{scope}\end{tikzpicture}\newpage
\definecolor{answer}{rgb}{0.850980,0.047059,0.070588}
\null\begin{tikzpicture}[remember picture,overlay]\begin{scope}[shift={(current page.north west)},x=1bp,y=-1bp]
\node[anchor=north west,inner sep=0pt,text=answer,align=left,text width=484.0bp] at (15,15) {\fontsize{7.5bp}{9.5bp}\selectfont 析：\allowbreak{}沿保留的两斜边向上延长，\allowbreak{}三个角均为锐角，\allowbreak{}选A。\allowbreak{}\par};
\end{scope}\end{tikzpicture}\newpage
\definecolor{answer}{rgb}{0.850980,0.047059,0.070588}
\null\begin{tikzpicture}[remember picture,overlay]\begin{scope}[shift={(current page.north west)},x=1bp,y=-1bp]
\node[anchor=north west,inner sep=0pt,text=answer,align=left,text width=484.0bp] at (15,15) {\fontsize{7.5bp}{9.5bp}\selectfont 析：\allowbreak{}全长=(\ensuremath{75} + \ensuremath{85}) × \ensuremath{15} + \ensuremath{450} = \ensuremath{2850}米，\allowbreak{}选B。\allowbreak{}\par};
\end{scope}\end{tikzpicture}\newpage
\definecolor{answer}{rgb}{0.850980,0.047059,0.070588}
\null\begin{tikzpicture}[remember picture,overlay]\begin{scope}[shift={(current page.north west)},x=1bp,y=-1bp]
\node[anchor=north west,inner sep=0pt,text=answer,align=left,text width=484.0bp] at (15,15) {\fontsize{8.5bp}{10.62bp}\selectfont 两人从同一地点反向跑，\allowbreak{}第一次相遇时，\allowbreak{}两人的路程之和正好是一圈。\allowbreak{}\par};
\end{scope}\end{tikzpicture}\newpage
\definecolor{answer}{rgb}{0.850980,0.047059,0.070588}
\null\begin{tikzpicture}[remember picture,overlay]\begin{scope}[shift={(current page.north west)},x=1bp,y=-1bp]
\node[anchor=north west,inner sep=0pt,text=answer,align=left,text width=484.0bp] at (15,15) {\fontsize{8.5bp}{10.62bp}\selectfont 甲、\allowbreak{}乙同时从两端修路，\allowbreak{}\ensuremath{15}天修完。\allowbreak{}两队\ensuremath{15}天所修的路程相加，\allowbreak{}就是公路总长。\allowbreak{}\par};
\end{scope}\end{tikzpicture}\newpage
\definecolor{answer}{rgb}{0.850980,0.047059,0.070588}
\null\begin{tikzpicture}[remember picture,overlay]\begin{scope}[shift={(current page.north west)},x=1bp,y=-1bp]
\node[anchor=north west,inner sep=0pt,text=answer,align=left,text width=484.0bp] at (15,15) {\fontsize{7.5bp}{9.5bp}\selectfont 析：\allowbreak{}浩浩：\allowbreak{}两底\ensuremath{9}、\allowbreak{}\ensuremath{5},\allowbreak{}另两根\ensuremath{5}作等腰；\allowbreak{}明明两底会相等，\allowbreak{}丽丽两底差\ensuremath{12}大于两腰和\ensuremath{10},\allowbreak{}均不行。\allowbreak{}\par};
\end{scope}\end{tikzpicture}\newpage
\definecolor{answer}{rgb}{0.850980,0.047059,0.070588}
\null\begin{tikzpicture}[remember picture,overlay]\begin{scope}[shift={(current page.north west)},x=1bp,y=-1bp]
\node[anchor=north west,inner sep=0pt,text=answer,align=left,text width=484.0bp] at (15,15) {\fontsize{7.5bp}{9.5bp}\selectfont 析：\allowbreak{}缩小到十分之一即÷\ensuremath{10},\allowbreak{} \ensuremath{0.36}→\ensuremath{0.036};\allowbreak{}扩大\ensuremath{1000}倍即×\ensuremath{1000},\allowbreak{} \ensuremath{0.042}→\ensuremath{42}。\allowbreak{}\par};
\end{scope}\end{tikzpicture}\newpage
\definecolor{answer}{rgb}{0.850980,0.047059,0.070588}
\null\begin{tikzpicture}[remember picture,overlay]\begin{scope}[shift={(current page.north west)},x=1bp,y=-1bp]
\node[anchor=north west,inner sep=0pt,text=answer,align=left,text width=484.0bp] at (15,15) {\fontsize{7.5bp}{9.5bp}\selectfont 析：\allowbreak{}\ensuremath{0.08}计数单位\ensuremath{0.01},\allowbreak{}另两数都是\ensuremath{0.001}且大小不等，\allowbreak{}所以选\ensuremath{0.080}和\ensuremath{0.008}。\allowbreak{}\par};
\end{scope}\end{tikzpicture}\newpage
\definecolor{answer}{rgb}{0.850980,0.047059,0.070588}
\null\begin{tikzpicture}[remember picture,overlay]\begin{scope}[shift={(current page.north west)},x=1bp,y=-1bp]
\node[anchor=north west,inner sep=0pt,text=answer,align=left,text width=484.0bp] at (15,15) {\fontsize{7.5bp}{9.5bp}\selectfont 析：\allowbreak{}保留\ensuremath{1}位看百分位\ensuremath{6},\allowbreak{}进\ensuremath{1}得\ensuremath{9.2};\allowbreak{}保留\ensuremath{2}位看千分位\ensuremath{4},\allowbreak{}舍去得\ensuremath{9.16}。\allowbreak{}\par};
\end{scope}\end{tikzpicture}\newpage
\definecolor{answer}{rgb}{0.850980,0.047059,0.070588}
\null\begin{tikzpicture}[remember picture,overlay]\begin{scope}[shift={(current page.north west)},x=1bp,y=-1bp]
\node[anchor=north west,inner sep=0pt,text=answer,align=left,text width=484.0bp] at (15,15) {\fontsize{7.5bp}{9.5bp}\selectfont 析：\allowbreak{}最高位是十位，\allowbreak{}移到\ensuremath{2}前须左移\ensuremath{2}位，\allowbreak{}得\ensuremath{0.2189}。\allowbreak{}\par};
\end{scope}\end{tikzpicture}\newpage
\definecolor{answer}{rgb}{0.850980,0.047059,0.070588}
\null\begin{tikzpicture}[remember picture,overlay]\begin{scope}[shift={(current page.north west)},x=1bp,y=-1bp]
\node[anchor=north west,inner sep=0pt,text=answer,align=left,text width=484.0bp] at (15,15) {\fontsize{7.5bp}{9.5bp}\selectfont 析：\allowbreak{}\ensuremath{4.55}至\ensuremath{4.64}都约为\ensuremath{4.6},\allowbreak{}共\ensuremath{10}种；\allowbreak{}最大\ensuremath{4.64},\allowbreak{}最小\ensuremath{4.55}。\allowbreak{}\par};
\end{scope}\end{tikzpicture}\newpage
\definecolor{answer}{rgb}{0.850980,0.047059,0.070588}
\null\begin{tikzpicture}[remember picture,overlay]\begin{scope}[shift={(current page.north west)},x=1bp,y=-1bp]
\node[anchor=north west,inner sep=0pt,text=answer,align=left,text width=484.0bp] at (15,15) {\fontsize{7.5bp}{9.5bp}\selectfont 析：\allowbreak{}平方米与平方厘米进率\ensuremath{10000};\allowbreak{}千克与克\ensuremath{1000};\allowbreak{}米与厘米\ensuremath{100};\allowbreak{}元与分\ensuremath{100}。\allowbreak{}得\ensuremath{0.095};\allowbreak{} \ensuremath{3}和\ensuremath{50};\allowbreak{} \ensuremath{7.5};\allowbreak{} \ensuremath{2.09}。\allowbreak{}\par};
\end{scope}\end{tikzpicture}\newpage
\definecolor{answer}{rgb}{0.850980,0.047059,0.070588}
\null\begin{tikzpicture}[remember picture,overlay]\begin{scope}[shift={(current page.north west)},x=1bp,y=-1bp]
\node[anchor=north west,inner sep=0pt,text=answer,align=left,text width=484.0bp] at (15,15) {\fontsize{7.5bp}{9.5bp}\selectfont 析：\allowbreak{}位数多未必大，\allowbreak{}如\ensuremath{0.12}＜\ensuremath{0.9},\allowbreak{}判错。\allowbreak{}\par};
\end{scope}\end{tikzpicture}\newpage
\definecolor{answer}{rgb}{0.850980,0.047059,0.070588}
\null\begin{tikzpicture}[remember picture,overlay]\begin{scope}[shift={(current page.north west)},x=1bp,y=-1bp]
\node[anchor=north west,inner sep=0pt,text=answer,align=left,text width=484.0bp] at (15,15) {\fontsize{7.5bp}{9.5bp}\selectfont 析：\allowbreak{}\ensuremath{28.0}亿末尾\ensuremath{0}表示保留\ensuremath{1}位小数，\allowbreak{}省略会改变近似程度，\allowbreak{}判错。\allowbreak{}\par};
\end{scope}\end{tikzpicture}\newpage
\definecolor{answer}{rgb}{0.850980,0.047059,0.070588}
\null\begin{tikzpicture}[remember picture,overlay]\begin{scope}[shift={(current page.north west)},x=1bp,y=-1bp]
\node[anchor=north west,inner sep=0pt,text=answer,align=left,text width=484.125bp] at (15,15) {\fontsize{7.52bp}{9.52bp}\selectfont 析：\allowbreak{}向右移动\ensuremath{3}位会扩大到原来的\ensuremath{1000}倍；\allowbreak{}缩小到原数的千分之一应向左移动\ensuremath{3}位，\allowbreak{}原句判×\par};
\end{scope}\end{tikzpicture}\newpage
\definecolor{answer}{rgb}{0.850980,0.047059,0.070588}
\null\begin{tikzpicture}[remember picture,overlay]\begin{scope}[shift={(current page.north west)},x=1bp,y=-1bp]
\node[anchor=north west,inner sep=0pt,text=answer,align=left,text width=484.125bp] at (15,15) {\fontsize{7.52bp}{9.52bp}\selectfont 析：\allowbreak{}从右往左，\allowbreak{}第一个\ensuremath{2}在百分位，\allowbreak{}表示\ensuremath{2}个\ensuremath{0.01}；\allowbreak{}第二个\ensuremath{2}在十分位，\allowbreak{}表示\ensuremath{2}个\ensuremath{0.1}，\allowbreak{}原句判×\par};
\end{scope}\end{tikzpicture}\newpage
\definecolor{answer}{rgb}{0.850980,0.047059,0.070588}
\null\begin{tikzpicture}[remember picture,overlay]\begin{scope}[shift={(current page.north west)},x=1bp,y=-1bp]
\node[anchor=north west,inner sep=0pt,text=answer,align=left,text width=484.0bp] at (15,15) {\fontsize{7.5bp}{9.5bp}\selectfont 析：\allowbreak{}\ensuremath{7.8} = \ensuremath{7.80},\allowbreak{}但计数单位分别为\ensuremath{0.1}、\allowbreak{}\ensuremath{0.01},\allowbreak{}意义不同，\allowbreak{}判错。\allowbreak{}\par};
\end{scope}\end{tikzpicture}\newpage
\definecolor{answer}{rgb}{0.850980,0.047059,0.070588}
\null\begin{tikzpicture}[remember picture,overlay]\begin{scope}[shift={(current page.north west)},x=1bp,y=-1bp]
\node[anchor=north west,inner sep=0pt,text=answer,align=left,text width=484.0bp] at (15,15) {\fontsize{7.5bp}{9.5bp}\selectfont 析：\allowbreak{}右移\ensuremath{3}位再左移\ensuremath{2}位，\allowbreak{}合计右移\ensuremath{1}位，\allowbreak{}扩大\ensuremath{10}倍，\allowbreak{}选A。\allowbreak{}\par};
\end{scope}\end{tikzpicture}\newpage
\definecolor{answer}{rgb}{0.850980,0.047059,0.070588}
\null\begin{tikzpicture}[remember picture,overlay]\begin{scope}[shift={(current page.north west)},x=1bp,y=-1bp]
\node[anchor=north west,inner sep=0pt,text=answer,align=left,text width=484.0bp] at (15,15) {\fontsize{7.5bp}{9.5bp}\selectfont 析：\allowbreak{}\ensuremath{5}在百分位，\allowbreak{}表示\ensuremath{5}个\ensuremath{0.01},\allowbreak{}选B。\allowbreak{}\par};
\end{scope}\end{tikzpicture}\newpage
\definecolor{answer}{rgb}{0.850980,0.047059,0.070588}
\null\begin{tikzpicture}[remember picture,overlay]\begin{scope}[shift={(current page.north west)},x=1bp,y=-1bp]
\node[anchor=north west,inner sep=0pt,text=answer,align=left,text width=484.0bp] at (15,15) {\fontsize{7.5bp}{9.5bp}\selectfont 析：\allowbreak{}没有限定位数；\allowbreak{}\ensuremath{0.51}、\allowbreak{}\ensuremath{0.501}、\allowbreak{}\ensuremath{0.5001}等有无数个，\allowbreak{}选C。\allowbreak{}\par};
\end{scope}\end{tikzpicture}\newpage
\definecolor{answer}{rgb}{0.850980,0.047059,0.070588}
\null\begin{tikzpicture}[remember picture,overlay]\begin{scope}[shift={(current page.north west)},x=1bp,y=-1bp]
\node[anchor=north west,inner sep=0pt,text=answer,align=left,text width=484.0bp] at (15,15) {\fontsize{7.5bp}{9.5bp}\selectfont 析：\allowbreak{}千万位填\ensuremath{0}、\allowbreak{}\ensuremath{1}、\allowbreak{}\ensuremath{2}、\allowbreak{}\ensuremath{3}、\allowbreak{}\ensuremath{4}均约成\ensuremath{4}亿，\allowbreak{}共\ensuremath{5}种，\allowbreak{}选B。\allowbreak{}\par};
\end{scope}\end{tikzpicture}\newpage
\definecolor{answer}{rgb}{0.850980,0.047059,0.070588}
\null\begin{tikzpicture}[remember picture,overlay]\begin{scope}[shift={(current page.north west)},x=1bp,y=-1bp]
\node[anchor=north west,inner sep=0pt,text=answer,align=left,text width=484.0bp] at (15,15) {\fontsize{7.5bp}{9.5bp}\selectfont 析：\allowbreak{}A、\allowbreak{}C均表示\ensuremath{0.08};\allowbreak{} B每格\ensuremath{0.02},\allowbreak{}标点第\ensuremath{3}格为\ensuremath{0.06},\allowbreak{}选B。\allowbreak{}\par};
\end{scope}\end{tikzpicture}\newpage
\definecolor{answer}{rgb}{0.850980,0.047059,0.070588}
\null\begin{tikzpicture}[remember picture,overlay]\begin{scope}[shift={(current page.north west)},x=1bp,y=-1bp]
\node[anchor=north west,inner sep=0pt,text=answer,align=left,text width=484.0bp] at (15,15) {\fontsize{7.5bp}{9.5bp}\selectfont 析：\allowbreak{}只去末尾\ensuremath{0},\allowbreak{}中间\ensuremath{0}保留；\allowbreak{}依次得\ensuremath{5.02}、\allowbreak{}\ensuremath{10.18}、\allowbreak{}\ensuremath{2.07}、\allowbreak{}\ensuremath{41}。\allowbreak{}\par};
\end{scope}\end{tikzpicture}\newpage
\definecolor{answer}{rgb}{0.850980,0.047059,0.070588}
\null\begin{tikzpicture}[remember picture,overlay]\begin{scope}[shift={(current page.north west)},x=1bp,y=-1bp]
\node[anchor=north west,inner sep=0pt,text=answer,align=left,text width=484.0bp] at (15,15) {\fontsize{7.5bp}{9.5bp}\selectfont 析：\allowbreak{}利用小数末尾添去\ensuremath{0}大小不变，\allowbreak{}保留恰好\ensuremath{3}位：\allowbreak{}\ensuremath{0.480}、\allowbreak{}\ensuremath{30.000}、\allowbreak{}\ensuremath{0.107}、\allowbreak{}\ensuremath{9.500}。\allowbreak{}\par};
\end{scope}\end{tikzpicture}\newpage
\definecolor{answer}{rgb}{0.850980,0.047059,0.070588}
\null\begin{tikzpicture}[remember picture,overlay]\begin{scope}[shift={(current page.north west)},x=1bp,y=-1bp]
\node[anchor=north west,inner sep=0pt,text=answer,align=left,text width=484.0bp] at (15,15) {\fontsize{7.5bp}{9.5bp}\selectfont 析：\allowbreak{}改写成万用÷\ensuremath{10000},\allowbreak{}亿用÷\ensuremath{100000000},\allowbreak{}再按要求四舍五入；\allowbreak{}依次\ensuremath{4.5}万、\allowbreak{}\ensuremath{107.59}万、\allowbreak{}\ensuremath{7.5}亿、\allowbreak{}\ensuremath{0.59}亿、\allowbreak{}\ensuremath{18.0}亿（末\ensuremath{0}保留）。\allowbreak{}\par};
\end{scope}\end{tikzpicture}\newpage
\definecolor{answer}{rgb}{0.850980,0.047059,0.070588}
\null\begin{tikzpicture}[remember picture,overlay]\begin{scope}[shift={(current page.north west)},x=1bp,y=-1bp]
\node[anchor=north west,inner sep=0pt,text=answer,align=left,text width=484.0bp] at (15,15) {\fontsize{7.5bp}{9.5bp}\selectfont 析：\allowbreak{}\ensuremath{0}放两位整数的个位，\allowbreak{}整数为\ensuremath{10}、\allowbreak{}\ensuremath{20}、\allowbreak{}\ensuremath{50};\allowbreak{}余两数字互换，\allowbreak{}共\ensuremath{6}种，\allowbreak{}按原空从小到大排列。\allowbreak{}\par};
\end{scope}\end{tikzpicture}\newpage
\definecolor{answer}{rgb}{0.850980,0.047059,0.070588}
\null\begin{tikzpicture}[remember picture,overlay]\begin{scope}[shift={(current page.north west)},x=1bp,y=-1bp]
\node[anchor=north west,inner sep=0pt,text=answer,align=left,text width=484.0bp] at (15,15) {\fontsize{8.5bp}{10.62bp}\selectfont 已知续重质量和续重费，\allowbreak{}用\ensuremath{38}元除以\ensuremath{10}千克求续重单价，\allowbreak{}再在表中比较。\allowbreak{}注意这里的\ensuremath{38}元只含续重费用。\allowbreak{}\par};
\end{scope}\end{tikzpicture}\newpage
\definecolor{answer}{rgb}{0.850980,0.047059,0.070588}
\null\begin{tikzpicture}[remember picture,overlay]\begin{scope}[shift={(current page.north west)},x=1bp,y=-1bp]
\node[anchor=north west,inner sep=0pt,text=answer,align=left,text width=437.3385009765625bp] at (15,15) {\fontsize{8.5bp}{10.62bp}\selectfont 收费只按超出\ensuremath{1000}元免费额度的部分计算。\allowbreak{}原来超出\ensuremath{100}元收费\ensuremath{0.2}元，\allowbreak{}现在超出\ensuremath{1000}元，\allowbreak{}是原来的\ensuremath{10}倍。\allowbreak{}\par};
\end{scope}\end{tikzpicture}\newpage
\definecolor{answer}{rgb}{0.850980,0.047059,0.070588}
\null\begin{tikzpicture}[remember picture,overlay]\begin{scope}[shift={(current page.north west)},x=1bp,y=-1bp]
\node[anchor=north west,inner sep=0pt,text=answer,align=left,text width=484.0bp] at (15,15) {\fontsize{7.5bp}{9.5bp}\selectfont 析：\allowbreak{}多用加、\allowbreak{}求相差用减：\allowbreak{}\ensuremath{4.75} + \ensuremath{0.55} = \ensuremath{5.3};\allowbreak{} \ensuremath{7.26} \ensuremath{-}\ensuremath{5.4} = \ensuremath{1.86}。\allowbreak{}\par};
\end{scope}\end{tikzpicture}\newpage
\definecolor{answer}{rgb}{0.850980,0.047059,0.070588}
\null\begin{tikzpicture}[remember picture,overlay]\begin{scope}[shift={(current page.north west)},x=1bp,y=-1bp]
\node[anchor=north west,inner sep=0pt,text=answer,align=left,text width=484.0bp] at (15,15) {\fontsize{7.5bp}{9.5bp}\selectfont 析：\allowbreak{}原量相同，\allowbreak{}余下越少卖得越多；\allowbreak{}第一筐多卖\ensuremath{5.8} \ensuremath{-}\ensuremath{3.75} = \ensuremath{2.05}千克。\allowbreak{}\par};
\end{scope}\end{tikzpicture}\newpage
\definecolor{answer}{rgb}{0.850980,0.047059,0.070588}
\null\begin{tikzpicture}[remember picture,overlay]\begin{scope}[shift={(current page.north west)},x=1bp,y=-1bp]
\node[anchor=north west,inner sep=0pt,text=answer,align=left,text width=484.0bp] at (15,15) {\fontsize{7.5bp}{9.5bp}\selectfont 析：\allowbreak{}①每次减\ensuremath{1.25}，\allowbreak{}得\ensuremath{4.45}、\allowbreak{}\ensuremath{3.2}；\allowbreak{}②相邻差依次\ensuremath{0.4}、\allowbreak{}\ensuremath{0.8}、\allowbreak{}\ensuremath{1.6}、\allowbreak{}\ensuremath{3.2}、\allowbreak{}\ensuremath{6.4}，\allowbreak{}得\ensuremath{3.1}、\allowbreak{}\ensuremath{12.7}；\allowbreak{}也可把前数与自己相加，\allowbreak{}再加\ensuremath{0.1}。\allowbreak{}\par};
\end{scope}\end{tikzpicture}\newpage
\definecolor{answer}{rgb}{0.850980,0.047059,0.070588}
\null\begin{tikzpicture}[remember picture,overlay]\begin{scope}[shift={(current page.north west)},x=1bp,y=-1bp]
\node[anchor=north west,inner sep=0pt,text=answer,align=left,text width=484.0bp] at (15,15) {\fontsize{7.5bp}{9.5bp}\selectfont 析：\allowbreak{}改变加数位置与结合顺序以凑整，\allowbreak{}用加法交换律和结合律；\allowbreak{}移项须连同原减号。\allowbreak{}\par};
\end{scope}\end{tikzpicture}\newpage
\definecolor{answer}{rgb}{0.850980,0.047059,0.070588}
\null\begin{tikzpicture}[remember picture,overlay]\begin{scope}[shift={(current page.north west)},x=1bp,y=-1bp]
\node[anchor=north west,inner sep=0pt,text=answer,align=left,text width=484.0bp] at (15,15) {\fontsize{7.5bp}{9.5bp}\selectfont 析：\allowbreak{}同起点减小加大；\allowbreak{}两加数一增一减\ensuremath{0.01}和不变；\allowbreak{}\ensuremath{8.10}元＜\ensuremath{8.40}元；\allowbreak{}\ensuremath{14.9} + \ensuremath{0.16} = \ensuremath{15.06};\allowbreak{} \ensuremath{1.03}千克＜\ensuremath{1.5}千克；\allowbreak{}减数小，\allowbreak{}差大。\allowbreak{}\par};
\end{scope}\end{tikzpicture}\newpage
\definecolor{answer}{rgb}{0.850980,0.047059,0.070588}
\null\begin{tikzpicture}[remember picture,overlay]\begin{scope}[shift={(current page.north west)},x=1bp,y=-1bp]
\node[anchor=north west,inner sep=0pt,text=answer,align=left,text width=484.0bp] at (15,15) {\fontsize{7.5bp}{9.5bp}\selectfont 析：\allowbreak{}买三赠一，\allowbreak{}得\ensuremath{4}盒只付\ensuremath{3}盒：\allowbreak{}\ensuremath{17.16} + \ensuremath{17.16} + \ensuremath{17.16} = \ensuremath{51.48}元。\allowbreak{}\par};
\end{scope}\end{tikzpicture}\newpage
\definecolor{answer}{rgb}{0.850980,0.047059,0.070588}
\null\begin{tikzpicture}[remember picture,overlay]\begin{scope}[shift={(current page.north west)},x=1bp,y=-1bp]
\node[anchor=north west,inner sep=0pt,text=answer,align=left,text width=484.0bp] at (15,15) {\fontsize{7.5bp}{9.5bp}\selectfont 析：\allowbreak{}原被减数\ensuremath{12.7} + \ensuremath{4.5} = \ensuremath{17.2};\allowbreak{}差不变则两数同减\ensuremath{1.7},\allowbreak{}新减数\ensuremath{4.5} \ensuremath{-}\ensuremath{1.7} = \ensuremath{2.8}(也可\ensuremath{17.2} \ensuremath{-}\ensuremath{1.7} \ensuremath{-}\ensuremath{12.7})。\allowbreak{}\par};
\end{scope}\end{tikzpicture}\newpage
\definecolor{answer}{rgb}{0.850980,0.047059,0.070588}
\null\begin{tikzpicture}[remember picture,overlay]\begin{scope}[shift={(current page.north west)},x=1bp,y=-1bp]
\node[anchor=north west,inner sep=0pt,text=answer,align=left,text width=484.0bp] at (15,15) {\fontsize{7.5bp}{9.5bp}\selectfont 析：\allowbreak{}\ensuremath{3.17}与\ensuremath{2.83}凑成\ensuremath{6},\allowbreak{}可用\ensuremath{7.42} \ensuremath{-}\ensuremath{6}简算，\allowbreak{}选C。\allowbreak{}\par};
\end{scope}\end{tikzpicture}\newpage
\definecolor{answer}{rgb}{0.850980,0.047059,0.070588}
\null\begin{tikzpicture}[remember picture,overlay]\begin{scope}[shift={(current page.north west)},x=1bp,y=-1bp]
\node[anchor=north west,inner sep=0pt,text=answer,align=left,text width=484.0bp] at (15,15) {\fontsize{7.5bp}{9.5bp}\selectfont 析：\allowbreak{}圈中\ensuremath{8}、\allowbreak{}\ensuremath{3}、\allowbreak{}\ensuremath{5}在百分位，\allowbreak{}单位都是\ensuremath{0.01}平方米；\allowbreak{}D误说成整平方米，\allowbreak{}选D。\allowbreak{}\par};
\end{scope}\end{tikzpicture}\newpage
\definecolor{answer}{rgb}{0.850980,0.047059,0.070588}
\null\begin{tikzpicture}[remember picture,overlay]\begin{scope}[shift={(current page.north west)},x=1bp,y=-1bp]
\node[anchor=north west,inner sep=0pt,text=answer,align=left,text width=484.0bp] at (15,15) {\fontsize{7.5bp}{9.5bp}\selectfont 析：\allowbreak{}\ensuremath{0.89} = \ensuremath{1} \ensuremath{-}\ensuremath{0.11},\allowbreak{} \ensuremath{13.47} + \ensuremath{0.89} = \ensuremath{13.47} + \ensuremath{1} \ensuremath{-}\ensuremath{0.11},\allowbreak{}选D。\allowbreak{}\par};
\end{scope}\end{tikzpicture}\newpage
\definecolor{answer}{rgb}{0.850980,0.047059,0.070588}
\null\begin{tikzpicture}[remember picture,overlay]\begin{scope}[shift={(current page.north west)},x=1bp,y=-1bp]
\node[anchor=north west,inner sep=0pt,text=answer,align=left,text width=484.0bp] at (15,15) {\fontsize{7.5bp}{9.5bp}\selectfont 析：\allowbreak{}误多加\ensuremath{6.77} \ensuremath{-}\ensuremath{6.27} = \ensuremath{0.5},\allowbreak{}继续输入-\ensuremath{0.5},\allowbreak{}选D。\allowbreak{}\par};
\end{scope}\end{tikzpicture}\newpage
\definecolor{answer}{rgb}{0.850980,0.047059,0.070588}
\null\begin{tikzpicture}[remember picture,overlay]\begin{scope}[shift={(current page.north west)},x=1bp,y=-1bp]
\node[anchor=north west,inner sep=0pt,text=answer,align=left,text width=484.0bp] at (15,15) {\fontsize{7.5bp}{9.5bp}\selectfont 析：\allowbreak{}糖=\ensuremath{3.62} + \ensuremath{1.54};\allowbreak{}巧克力再少\ensuremath{1.87},\allowbreak{}列\ensuremath{3.62} + \ensuremath{1.54} \ensuremath{-}\ensuremath{1.87},\allowbreak{}选C。\allowbreak{}\par};
\end{scope}\end{tikzpicture}\newpage
\definecolor{answer}{rgb}{0.850980,0.047059,0.070588}
\null\begin{tikzpicture}[remember picture,overlay]\begin{scope}[shift={(current page.north west)},x=1bp,y=-1bp]
\node[anchor=north west,inner sep=0pt,text=answer,align=left,text width=484.0bp] at (15,15) {\fontsize{7.5bp}{9.5bp}\selectfont 析：\allowbreak{}约成整数\ensuremath{3}的一位小数为\ensuremath{2.5}、\allowbreak{}\ensuremath{2.6}、\allowbreak{}…、\allowbreak{}\ensuremath{3.4}(含\ensuremath{3.0}),\allowbreak{}共\ensuremath{10}个刻度；\allowbreak{}最大\ensuremath{3.4}、\allowbreak{}最小\ensuremath{2.5}。\allowbreak{}\par};
\end{scope}\end{tikzpicture}\newpage
\definecolor{answer}{rgb}{0.850980,0.047059,0.070588}
\null\begin{tikzpicture}[remember picture,overlay]\begin{scope}[shift={(current page.north west)},x=1bp,y=-1bp]
\node[anchor=north west,inner sep=0pt,text=answer,align=left,text width=430.0bp] at (15,15) {\fontsize{8.5bp}{10.62bp}\selectfont 第三阶段比第二阶段多\ensuremath{2.4}吨，\allowbreak{}应先求第三阶段\ensuremath{8}吨。\allowbreak{}总量减去第二、\allowbreak{}第三阶段的量，\allowbreak{}就是第一阶段回收量。\allowbreak{}\par};
\end{scope}\end{tikzpicture}\newpage
\definecolor{answer}{rgb}{0.850980,0.047059,0.070588}
\null\begin{tikzpicture}[remember picture,overlay]\begin{scope}[shift={(current page.north west)},x=1bp,y=-1bp]
\node[anchor=north west,inner sep=0pt,text=answer,align=left,text width=403.8399963378906bp] at (15,15) {\fontsize{8.5bp}{10.62bp}\selectfont 三处在同一直线上，\allowbreak{}没有说明豆豆家与图书馆位于电影院的同侧还是两侧，\allowbreak{}因此有两种结果。\allowbreak{}\par};
\end{scope}\end{tikzpicture}\newpage
\definecolor{answer}{rgb}{0.850980,0.047059,0.070588}
\null\begin{tikzpicture}[remember picture,overlay]\begin{scope}[shift={(current page.north west)},x=1bp,y=-1bp]
\node[anchor=north west,inner sep=0pt,text=answer,align=left,text width=484.125bp] at (15,15) {\fontsize{7.52bp}{9.52bp}\selectfont 析：\allowbreak{}仅第\ensuremath{4}幅是轴对称图形。\allowbreak{}第\ensuremath{1}幅中\ensuremath{8}的上环小、\allowbreak{}下环大，\allowbreak{}整个图案不对称\par};
\end{scope}\end{tikzpicture}\newpage
\definecolor{answer}{rgb}{0.850980,0.047059,0.070588}
\null\begin{tikzpicture}[remember picture,overlay]\begin{scope}[shift={(current page.north west)},x=1bp,y=-1bp]
\node[anchor=north west,inner sep=0pt,text=answer,align=left,text width=484.0bp] at (15,15) {\fontsize{7.5bp}{9.5bp}\selectfont 析：\allowbreak{}平移、\allowbreak{}旋转只改位置或方向，\allowbreak{}形状、\allowbreak{}大小均不变。\allowbreak{}\par};
\end{scope}\end{tikzpicture}\newpage
\definecolor{answer}{rgb}{0.850980,0.047059,0.070588}
\null\begin{tikzpicture}[remember picture,overlay]\begin{scope}[shift={(current page.north west)},x=1bp,y=-1bp]
\node[anchor=north west,inner sep=0pt,text=answer,align=left,text width=256.0bp] at (15,15) {\fontsize{7.5bp}{9.5bp}\selectfont 析：\allowbreak{}②选同一顶点：\allowbreak{}先右\ensuremath{3}格，\allowbreak{}再下\ensuremath{5}格。\allowbreak{}\par};
\end{scope}\end{tikzpicture}\newpage
\definecolor{answer}{rgb}{0.850980,0.047059,0.070588}
\null\begin{tikzpicture}[remember picture,overlay]\begin{scope}[shift={(current page.north west)},x=1bp,y=-1bp]
\node[anchor=north west,inner sep=0pt,text=answer,align=left,text width=484.0bp] at (15,15) {\fontsize{7.5bp}{9.5bp}\selectfont 析：\allowbreak{}选同一顶点：\allowbreak{}先左\ensuremath{3}格，\allowbreak{}再上\ensuremath{5}格。\allowbreak{}\par};
\end{scope}\end{tikzpicture}\newpage
\definecolor{answer}{rgb}{0.850980,0.047059,0.070588}
\null\begin{tikzpicture}[remember picture,overlay]\begin{scope}[shift={(current page.north west)},x=1bp,y=-1bp]
\node[anchor=north west,inner sep=0pt,text=answer,align=left,text width=256.0bp] at (15,15) {\fontsize{7.5bp}{9.5bp}\selectfont 析：\allowbreak{}②绕O顺时针每次\ensuremath{90}\ensuremath{^\circ};\allowbreak{}转\ensuremath{4}次共\ensuremath{360}\ensuremath{^\circ}回原位。\allowbreak{}\par};
\end{scope}\end{tikzpicture}\newpage
\definecolor{answer}{rgb}{0.850980,0.047059,0.070588}
\null\begin{tikzpicture}[remember picture,overlay]\begin{scope}[shift={(current page.north west)},x=1bp,y=-1bp]
\node[anchor=north west,inner sep=0pt,text=answer,align=left,text width=222.0bp] at (15,15) {\fontsize{7.5bp}{9.5bp}\selectfont 析：\allowbreak{}①同一顶点比较：\allowbreak{}先下\ensuremath{5}格，\allowbreak{}再右\ensuremath{6}格。\allowbreak{}\par};
\end{scope}\end{tikzpicture}\newpage
\definecolor{answer}{rgb}{0.850980,0.047059,0.070588}
\null\begin{tikzpicture}[remember picture,overlay]\begin{scope}[shift={(current page.north west)},x=1bp,y=-1bp]
\node[anchor=north west,inner sep=0pt,text=answer,align=left,text width=484.0bp] at (15,15) {\fontsize{7.5bp}{9.5bp}\selectfont 析：\allowbreak{}A、\allowbreak{}B、\allowbreak{}D可沿竖直中线对称；\allowbreak{}C的内部笔画不能对应，\allowbreak{}选C。\allowbreak{}\par};
\end{scope}\end{tikzpicture}\newpage
\definecolor{answer}{rgb}{0.850980,0.047059,0.070588}
\null\begin{tikzpicture}[remember picture,overlay]\begin{scope}[shift={(current page.north west)},x=1bp,y=-1bp]
\node[anchor=north west,inner sep=0pt,text=answer,align=left,text width=484.0bp] at (15,15) {\fontsize{7.5bp}{9.5bp}\selectfont 析：\allowbreak{}旋转中心位置不动：\allowbreak{}②中A不动、\allowbreak{}③中B不动、\allowbreak{}①中C不动，\allowbreak{}依次②③①。\allowbreak{}\par};
\end{scope}\end{tikzpicture}\newpage
\definecolor{answer}{rgb}{0.850980,0.047059,0.070588}
\null\begin{tikzpicture}[remember picture,overlay]\begin{scope}[shift={(current page.north west)},x=1bp,y=-1bp]
\node[anchor=north west,inner sep=0pt,text=answer,align=left,text width=484.0bp] at (15,15) {\fontsize{7.5bp}{9.5bp}\selectfont 析：\allowbreak{}绕O逆时针\ensuremath{90}\ensuremath{^\circ}移到左下拼平行四边形；\allowbreak{}等效为顺时针\ensuremath{270}\ensuremath{^\circ}。\allowbreak{}\par};
\end{scope}\end{tikzpicture}\newpage
\definecolor{answer}{rgb}{0.850980,0.047059,0.070588}
\null\begin{tikzpicture}[remember picture,overlay]\begin{scope}[shift={(current page.north west)},x=1bp,y=-1bp]
\node[anchor=north west,inner sep=0pt,text=answer,align=left,text width=484.0bp] at (15,15) {\fontsize{7.5bp}{9.5bp}\selectfont 析：\allowbreak{}绕O顺时针\ensuremath{90}\ensuremath{^\circ}移到右上拼大三角形；\allowbreak{}等效为逆时针\ensuremath{270}\ensuremath{^\circ}。\allowbreak{}\par};
\end{scope}\end{tikzpicture}\newpage
\definecolor{answer}{rgb}{0.850980,0.047059,0.070588}
\null\begin{tikzpicture}[remember picture,overlay]\begin{scope}[shift={(current page.north west)},x=1bp,y=-1bp]
\node[anchor=north west,inner sep=0pt,text=answer,align=left,text width=484.0bp] at (15,15) {\fontsize{7.5bp}{9.5bp}\selectfont 析：\allowbreak{}总数÷个数得到平均数，\allowbreak{}反映整体水平，\allowbreak{}不一定等于某个人的实际值。\allowbreak{}\par};
\end{scope}\end{tikzpicture}\newpage
\definecolor{answer}{rgb}{0.850980,0.047059,0.070588}
\null\begin{tikzpicture}[remember picture,overlay]\begin{scope}[shift={(current page.north west)},x=1bp,y=-1bp]
\node[anchor=north west,inner sep=0pt,text=answer,align=left,text width=484.0bp] at (15,15) {\fontsize{7.5bp}{9.5bp}\selectfont 析：\allowbreak{}连续三个自然数为\ensuremath{14}、\allowbreak{}\ensuremath{15}、\allowbreak{}\ensuremath{16},\allowbreak{}平均数就是中间数\ensuremath{15},\allowbreak{}最大\ensuremath{16}。\allowbreak{}\par};
\end{scope}\end{tikzpicture}\newpage
\definecolor{answer}{rgb}{0.850980,0.047059,0.070588}
\null\begin{tikzpicture}[remember picture,overlay]\begin{scope}[shift={(current page.north west)},x=1bp,y=-1bp]
\node[anchor=north west,inner sep=0pt,text=answer,align=left,text width=484.0bp] at (15,15) {\fontsize{7.5bp}{9.5bp}\selectfont 析：\allowbreak{}平均数在最小\ensuremath{30}与最大\ensuremath{45}之间，\allowbreak{}只有\ensuremath{34}符合，\allowbreak{}选B。\allowbreak{}\par};
\end{scope}\end{tikzpicture}\newpage
\definecolor{answer}{rgb}{0.850980,0.047059,0.070588}
\null\begin{tikzpicture}[remember picture,overlay]\begin{scope}[shift={(current page.north west)},x=1bp,y=-1bp]
\node[anchor=north west,inner sep=0pt,text=answer,align=left,text width=484.0bp] at (15,15) {\fontsize{7.5bp}{9.5bp}\selectfont 析：\allowbreak{}平均\ensuremath{62}岁，\allowbreak{}个人可高于或低于\ensuremath{62}岁；\allowbreak{}\ensuremath{70}岁有可能，\allowbreak{}选A。\allowbreak{}\par};
\end{scope}\end{tikzpicture}\newpage
\definecolor{answer}{rgb}{0.850980,0.047059,0.070588}
\null\begin{tikzpicture}[remember picture,overlay]\begin{scope}[shift={(current page.north west)},x=1bp,y=-1bp]
\node[anchor=north west,inner sep=0pt,text=answer,align=left,text width=484.0bp] at (15,15) {\fontsize{7.5bp}{9.5bp}\selectfont 析：\allowbreak{}连续偶数相差\ensuremath{2},\allowbreak{}平均数\ensuremath{a}为中间数；\allowbreak{}最大\ensuremath{a} + \ensuremath{2}选C,\allowbreak{}最小\ensuremath{a} \ensuremath{-}\ensuremath{2}选B。\allowbreak{}\par};
\end{scope}\end{tikzpicture}\newpage
\definecolor{answer}{rgb}{0.850980,0.047059,0.070588}
\null\begin{tikzpicture}[remember picture,overlay]\begin{scope}[shift={(current page.north west)},x=1bp,y=-1bp]
\node[anchor=north west,inner sep=0pt,text=answer,align=left,text width=431.6485290527344bp] at (15,15) {\fontsize{8.5bp}{10.62bp}\selectfont \ensuremath{5}天平均每天读\ensuremath{20}页，\allowbreak{}先求\ensuremath{5}天共读\ensuremath{100}页，\allowbreak{}再减去已知\ensuremath{4}天共读的\ensuremath{74}页。\allowbreak{}也可按每天\ensuremath{20}页作基准移多补少。\allowbreak{}\par};
\end{scope}\end{tikzpicture}\newpage
\definecolor{answer}{rgb}{0.850980,0.047059,0.070588}
\null\begin{tikzpicture}[remember picture,overlay]\begin{scope}[shift={(current page.north west)},x=1bp,y=-1bp]
\node[anchor=north west,inner sep=0pt,text=answer,align=left,text width=484.0bp] at (15,15) {\fontsize{7.5bp}{9.5bp}\selectfont 析：\allowbreak{}把相同分段对齐合并；\allowbreak{}男\ensuremath{1}、\allowbreak{}\ensuremath{3}、\allowbreak{}\ensuremath{12}、\allowbreak{}\ensuremath{4},\allowbreak{}女\ensuremath{2}、\allowbreak{}\ensuremath{10}、\allowbreak{}\ensuremath{7}、\allowbreak{}\ensuremath{1},\allowbreak{}每行合计\ensuremath{20}人。\allowbreak{}\par};
\end{scope}\end{tikzpicture}\newpage
\definecolor{answer}{rgb}{0.850980,0.047059,0.070588}
\null\begin{tikzpicture}[remember picture,overlay]\begin{scope}[shift={(current page.north west)},x=1bp,y=-1bp]
\node[anchor=north west,inner sep=0pt,text=answer,align=left,text width=484.0bp] at (15,15) {\fontsize{7.5bp}{9.5bp}\selectfont 析：\allowbreak{}各组取人数最多的段：\allowbreak{}女\ensuremath{10}人对应\ensuremath{5}～\ensuremath{10}个；\allowbreak{}男\ensuremath{12}人对应\ensuremath{11}～\ensuremath{15}个。\allowbreak{}\par};
\end{scope}\end{tikzpicture}\newpage
\definecolor{answer}{rgb}{0.850980,0.047059,0.070588}
\null\begin{tikzpicture}[remember picture,overlay]\begin{scope}[shift={(current page.north west)},x=1bp,y=-1bp]
\node[anchor=north west,inner sep=0pt,text=answer,align=left,text width=484.0bp] at (15,15) {\fontsize{7.5bp}{9.5bp}\selectfont 析：\allowbreak{}改写万用÷\ensuremath{10000} :\allowbreak{} \ensuremath{33.76}万保留\ensuremath{1}位为\ensuremath{33.8}万；\allowbreak{}改写亿用÷\ensuremath{100000000},\allowbreak{}得\ensuremath{35}亿。\allowbreak{}\par};
\end{scope}\end{tikzpicture}\newpage
\definecolor{answer}{rgb}{0.850980,0.047059,0.070588}
\null\begin{tikzpicture}[remember picture,overlay]\begin{scope}[shift={(current page.north west)},x=1bp,y=-1bp]
\node[anchor=north west,inner sep=0pt,text=answer,align=left,text width=484.0bp] at (15,15) {\fontsize{7.5bp}{9.5bp}\selectfont 析：\allowbreak{}任一点均可作垂线，\allowbreak{}故无数条；\allowbreak{}各垂线相互平行，\allowbreak{}平行线间距离固定故长度相等。\allowbreak{}\par};
\end{scope}\end{tikzpicture}\newpage
\definecolor{answer}{rgb}{0.850980,0.047059,0.070588}
\null\begin{tikzpicture}[remember picture,overlay]\begin{scope}[shift={(current page.north west)},x=1bp,y=-1bp]
\node[anchor=north west,inner sep=0pt,text=answer,align=left,text width=484.0bp] at (15,15) {\fontsize{7.5bp}{9.5bp}\selectfont 析：\allowbreak{}内角和\ensuremath{180}\ensuremath{^\circ},\allowbreak{}另两角之和＞\ensuremath{90}\ensuremath{^\circ},\allowbreak{}第三角必＜\ensuremath{90}\ensuremath{^\circ}。\allowbreak{}\par};
\end{scope}\end{tikzpicture}\newpage
\definecolor{answer}{rgb}{0.850980,0.047059,0.070588}
\null\begin{tikzpicture}[remember picture,overlay]\begin{scope}[shift={(current page.north west)},x=1bp,y=-1bp]
\node[anchor=north west,inner sep=0pt,text=answer,align=left,text width=484.0bp] at (15,15) {\fontsize{7.5bp}{9.5bp}\selectfont 析：\allowbreak{}三边都是\ensuremath{6},\allowbreak{}属等边三角形；\allowbreak{}周长\ensuremath{6} × \ensuremath{3} = \ensuremath{18}厘米。\allowbreak{}\par};
\end{scope}\end{tikzpicture}\newpage
\definecolor{answer}{rgb}{0.850980,0.047059,0.070588}
\null\begin{tikzpicture}[remember picture,overlay]\begin{scope}[shift={(current page.north west)},x=1bp,y=-1bp]
\node[anchor=north west,inner sep=0pt,text=answer,align=left,text width=484.0bp] at (15,15) {\fontsize{7.5bp}{9.5bp}\selectfont 析：\allowbreak{}两项直接相加重复算\ensuremath{8}人，\allowbreak{}\ensuremath{15} + \ensuremath{16} \ensuremath{-}\ensuremath{8} = \ensuremath{23}人。\allowbreak{}\par};
\end{scope}\end{tikzpicture}\newpage
\definecolor{answer}{rgb}{0.850980,0.047059,0.070588}
\null\begin{tikzpicture}[remember picture,overlay]\begin{scope}[shift={(current page.north west)},x=1bp,y=-1bp]
\node[anchor=north west,inner sep=0pt,text=answer,align=left,text width=484.0bp] at (15,15) {\fontsize{7.5bp}{9.5bp}\selectfont 析：\allowbreak{}\ensuremath{5}在百分位，\allowbreak{}表示\ensuremath{5}个百分之一，\allowbreak{}选B。\allowbreak{}\par};
\end{scope}\end{tikzpicture}\newpage
\definecolor{answer}{rgb}{0.850980,0.047059,0.070588}
\null\begin{tikzpicture}[remember picture,overlay]\begin{scope}[shift={(current page.north west)},x=1bp,y=-1bp]
\node[anchor=north west,inner sep=0pt,text=answer,align=left,text width=484.0bp] at (15,15) {\fontsize{7.5bp}{9.5bp}\selectfont 析：\allowbreak{}须由“五入”得到\ensuremath{1.8},\allowbreak{}只能\ensuremath{1.75}至\ensuremath{1.79},\allowbreak{}共\ensuremath{5}个，\allowbreak{}选A。\allowbreak{}\par};
\end{scope}\end{tikzpicture}\newpage
\definecolor{answer}{rgb}{0.850980,0.047059,0.070588}
\null\begin{tikzpicture}[remember picture,overlay]\begin{scope}[shift={(current page.north west)},x=1bp,y=-1bp]
\node[anchor=north west,inner sep=0pt,text=answer,align=left,text width=484.0bp] at (15,15) {\fontsize{7.5bp}{9.5bp}\selectfont 析：\allowbreak{}第三边须＞\ensuremath{10} \ensuremath{-}\ensuremath{5} = \ensuremath{5}且＜\ensuremath{10} + \ensuremath{5} = \ensuremath{15},\allowbreak{}只有\ensuremath{10}厘米符合，\allowbreak{}选B。\allowbreak{}\par};
\end{scope}\end{tikzpicture}\newpage
\definecolor{answer}{rgb}{0.850980,0.047059,0.070588}
\null\begin{tikzpicture}[remember picture,overlay]\begin{scope}[shift={(current page.north west)},x=1bp,y=-1bp]
\node[anchor=north west,inner sep=0pt,text=answer,align=left,text width=349.1400146484375bp] at (15,15) {\fontsize{7.5bp}{9.5bp}\selectfont 析：\allowbreak{}连续数n、\allowbreak{}n + \ensuremath{1}、\allowbreak{}n + \ensuremath{2}的平均数为中间数n + \ensuremath{1},\allowbreak{}选A。\allowbreak{}\par};
\end{scope}\end{tikzpicture}\newpage
\definecolor{answer}{rgb}{0.850980,0.047059,0.070588}
\null\begin{tikzpicture}[remember picture,overlay]\begin{scope}[shift={(current page.north west)},x=1bp,y=-1bp]
\node[anchor=north west,inner sep=0pt,outer sep=0pt,text=answer,align=left,text width=490bp] at (15,15) {\fontsize{7.5bp}{9.375bp}\selectfont 析：绕\ensuremath{O}点逆时针旋转\ensuremath{90^{\circ}}后，上半部的蓝色区域在右侧，下半部的蓝色区域在左侧，选\ensuremath{C}。\par};
\end{scope}\end{tikzpicture}
\newpage
\definecolor{answer}{rgb}{0.850980,0.047059,0.070588}
\null\begin{tikzpicture}[remember picture,overlay]\begin{scope}[shift={(current page.north west)},x=1bp,y=-1bp]
\node[anchor=north west,inner sep=0pt,text=answer,align=left,text width=484.0bp] at (15,15) {\fontsize{8.5bp}{10.62bp}\selectfont 跑步比骑行多消耗\ensuremath{26.5}千卡，\allowbreak{}用跑步消耗减\ensuremath{26.5}求骑行，\allowbreak{}再求两项的和。\allowbreak{}\par};
\end{scope}\end{tikzpicture}\newpage
\definecolor{answer}{rgb}{0.850980,0.047059,0.070588}
\null\begin{tikzpicture}[remember picture,overlay]\begin{scope}[shift={(current page.north west)},x=1bp,y=-1bp]
\node[anchor=north west,inner sep=0pt,text=answer,align=left,text width=484.0bp] at (15,15) {\fontsize{8.5bp}{10.62bp}\selectfont 两队已掘的长度相加后，\allowbreak{}还要加上没有掘通的\ensuremath{26}米，\allowbreak{}才是隧道总长。\allowbreak{}\par};
\end{scope}\end{tikzpicture}\newpage
\definecolor{answer}{rgb}{0.850980,0.047059,0.070588}
\null\begin{tikzpicture}[remember picture,overlay]\begin{scope}[shift={(current page.north west)},x=1bp,y=-1bp]
\node[anchor=north west,inner sep=0pt,text=answer,align=left,text width=484.0bp] at (15,15) {\fontsize{7.5bp}{9.5bp}\selectfont 析：\allowbreak{}逐个按分数段分类；\allowbreak{}四(\ensuremath{1})依次\ensuremath{1}、\allowbreak{}\ensuremath{1}、\allowbreak{}\ensuremath{1}、\allowbreak{}\ensuremath{7},\allowbreak{}四(\ensuremath{2})为\ensuremath{2}、\allowbreak{}\ensuremath{2}、\allowbreak{}\ensuremath{2}、\allowbreak{}\ensuremath{4},\allowbreak{}每班合计\ensuremath{10}人。\allowbreak{}\par};
\end{scope}\end{tikzpicture}\newpage
\definecolor{answer}{rgb}{0.850980,0.047059,0.070588}
\null\begin{tikzpicture}[remember picture,overlay]\begin{scope}[shift={(current page.north west)},x=1bp,y=-1bp]
\node[anchor=north west,inner sep=0pt,text=answer,align=left,text width=484.0bp] at (15,15) {\fontsize{8.5bp}{10.62bp}\selectfont 两个班都参加\ensuremath{10}人，\allowbreak{}可比较优秀人数、\allowbreak{}高分段人数，\allowbreak{}也可比较\ensuremath{10}个人的平均成绩。\allowbreak{}\par};
\end{scope}\end{tikzpicture}\newpage
\definecolor{answer}{rgb}{0.850980,0.047059,0.070588}
\null\begin{tikzpicture}[remember picture,overlay]\begin{scope}[shift={(current page.north west)},x=1bp,y=-1bp]
\node[anchor=north west,inner sep=0pt,text=answer,align=left,text width=484.0bp] at (15,15) {\fontsize{7.5bp}{9.5bp}\selectfont 析：\allowbreak{}从边的长度/平行关系及角的特征研究，\allowbreak{}概括为边、\allowbreak{}角。\allowbreak{}\par};
\end{scope}\end{tikzpicture}\newpage
\definecolor{answer}{rgb}{0.850980,0.047059,0.070588}
\null\begin{tikzpicture}[remember picture,overlay]\begin{scope}[shift={(current page.north west)},x=1bp,y=-1bp]
\node[anchor=north west,inner sep=0pt,text=answer,align=left,text width=484.0bp] at (15,15) {\fontsize{7.5bp}{9.5bp}\selectfont 析：\allowbreak{}平行四边形\ensuremath{2};\allowbreak{}梯形为外框与左右部分，\allowbreak{}共\ensuremath{3};\allowbreak{}直角三角形\ensuremath{2};\allowbreak{}左、\allowbreak{}右、\allowbreak{}中央锐角三角形共\ensuremath{3}。\allowbreak{}\par};
\end{scope}\end{tikzpicture}\newpage
\definecolor{answer}{rgb}{0.850980,0.047059,0.070588}
\null\begin{tikzpicture}[remember picture,overlay]\begin{scope}[shift={(current page.north west)},x=1bp,y=-1bp]
\node[anchor=north west,inner sep=0pt,text=answer,align=left,text width=484.0bp] at (15,15) {\fontsize{7.5bp}{9.5bp}\selectfont 析：\allowbreak{}\ensuremath{20278} ÷ \ensuremath{10000} = \ensuremath{2.0278}万；\allowbreak{}保留\ensuremath{1}位看百分位\ensuremath{2},\allowbreak{}得\ensuremath{2.0}万，\allowbreak{}末\ensuremath{0}须保留。\allowbreak{}\par};
\end{scope}\end{tikzpicture}\newpage
\definecolor{answer}{rgb}{0.850980,0.047059,0.070588}
\null\begin{tikzpicture}[remember picture,overlay]\begin{scope}[shift={(current page.north west)},x=1bp,y=-1bp]
\node[anchor=north west,inner sep=0pt,text=answer,align=left,text width=484.0bp] at (15,15) {\fontsize{7.5bp}{9.5bp}\selectfont 析：\allowbreak{}先比整数部分，\allowbreak{}再逐位比较：\allowbreak{}\ensuremath{0.58}＜\ensuremath{0.74}＜\ensuremath{4.7}＜\ensuremath{5.08}＜\ensuremath{8.04}。\allowbreak{}\par};
\end{scope}\end{tikzpicture}\newpage
\definecolor{answer}{rgb}{0.850980,0.047059,0.070588}
\null\begin{tikzpicture}[remember picture,overlay]\begin{scope}[shift={(current page.north west)},x=1bp,y=-1bp]
\node[anchor=north west,inner sep=0pt,text=answer,align=left,text width=484.0bp] at (15,15) {\fontsize{7.5bp}{9.5bp}\selectfont 析：\allowbreak{}连续减可减和，\allowbreak{}首空=；\allowbreak{}\ensuremath{35} × \ensuremath{101}分配后多\ensuremath{35}而非\ensuremath{1},\allowbreak{}填＞；\allowbreak{}连续除可除积，\allowbreak{}填=；\allowbreak{}\ensuremath{127} \ensuremath{-}\ensuremath{98}应=\ensuremath{127} \ensuremath{-}\ensuremath{100} + \ensuremath{2},\allowbreak{}填＞。\allowbreak{}\par};
\end{scope}\end{tikzpicture}\newpage
\definecolor{answer}{rgb}{0.850980,0.047059,0.070588}
\null\begin{tikzpicture}[remember picture,overlay]\begin{scope}[shift={(current page.north west)},x=1bp,y=-1bp]
\node[anchor=north west,inner sep=0pt,text=answer,align=left,text width=484.0bp] at (15,15) {\fontsize{7.5bp}{9.5bp}\selectfont 析：\allowbreak{}扩大\ensuremath{1000}倍右移\ensuremath{3}位得\ensuremath{5090};\allowbreak{}缩小为百分之一左移\ensuremath{2}位得\ensuremath{0.037}。\allowbreak{}\par};
\end{scope}\end{tikzpicture}\newpage
\definecolor{answer}{rgb}{0.850980,0.047059,0.070588}
\null\begin{tikzpicture}[remember picture,overlay]\begin{scope}[shift={(current page.north west)},x=1bp,y=-1bp]
\node[anchor=north west,inner sep=0pt,text=answer,align=left,text width=484.0bp] at (15,15) {\fontsize{7.5bp}{9.5bp}\selectfont 析：\allowbreak{}\ensuremath{40}\ensuremath{^\circ}若为底角将有\ensuremath{100}\ensuremath{^\circ}钝角，\allowbreak{}不合“锐角”；\allowbreak{}故是顶角，\allowbreak{}两底角各(\ensuremath{180}\ensuremath{^\circ}\ensuremath{-}\ensuremath{40}\ensuremath{^\circ}) ÷ \ensuremath{2} = \ensuremath{70}\ensuremath{^\circ}。\allowbreak{}\par};
\end{scope}\end{tikzpicture}\newpage
\definecolor{answer}{rgb}{0.850980,0.047059,0.070588}
\null\begin{tikzpicture}[remember picture,overlay]\begin{scope}[shift={(current page.north west)},x=1bp,y=-1bp]
\node[anchor=north west,inner sep=0pt,text=answer,align=left,text width=484.0bp] at (15,15) {\fontsize{7.5bp}{9.5bp}\selectfont 析：\allowbreak{}看千分位四舍五入：\allowbreak{}A约\ensuremath{14.09},\allowbreak{} B约\ensuremath{14.08},\allowbreak{} C约\ensuremath{14.07},\allowbreak{}选B。\allowbreak{}\par};
\end{scope}\end{tikzpicture}\newpage
\definecolor{answer}{rgb}{0.850980,0.047059,0.070588}
\null\begin{tikzpicture}[remember picture,overlay]\begin{scope}[shift={(current page.north west)},x=1bp,y=-1bp]
\node[anchor=north west,inner sep=0pt,text=answer,align=left,text width=484.0bp] at (15,15) {\fontsize{7.5bp}{9.5bp}\selectfont 析：\allowbreak{}\ensuremath{12}与其余两边不能围成三角形；\allowbreak{}最长可选\ensuremath{4}、\allowbreak{}\ensuremath{5}、\allowbreak{}\ensuremath{7},\allowbreak{}且\ensuremath{4} + \ensuremath{5}＞\ensuremath{7},\allowbreak{}周长\ensuremath{16}最大，\allowbreak{}选C。\allowbreak{}\par};
\end{scope}\end{tikzpicture}\newpage
\definecolor{answer}{rgb}{0.850980,0.047059,0.070588}
\null\begin{tikzpicture}[remember picture,overlay]\begin{scope}[shift={(current page.north west)},x=1bp,y=-1bp]
\node[anchor=north west,inner sep=0pt,text=answer,align=left,text width=484.0bp] at (15,15) {\fontsize{7.5bp}{9.5bp}\selectfont 析：\allowbreak{}返程用时\ensuremath{6} \ensuremath{-}\ensuremath{1} = \ensuremath{5}时，\allowbreak{}同路程\ensuremath{300}千米；\allowbreak{}\ensuremath{300} ÷ \ensuremath{5} = \ensuremath{60}千米/时，\allowbreak{}选B。\allowbreak{}\par};
\end{scope}\end{tikzpicture}\newpage
\definecolor{answer}{rgb}{0.850980,0.047059,0.070588}
\null\begin{tikzpicture}[remember picture,overlay]\begin{scope}[shift={(current page.north west)},x=1bp,y=-1bp]
\node[anchor=north west,inner sep=0pt,text=answer,align=left,text width=484.0bp] at (15,15) {\fontsize{7.5bp}{9.5bp}\selectfont 析：\allowbreak{}\ensuremath{49} = \ensuremath{50} \ensuremath{-}\ensuremath{1},\allowbreak{}所以\ensuremath{37} × \ensuremath{49} = \ensuremath{37} × \ensuremath{50} \ensuremath{-}\ensuremath{37},\allowbreak{}选B。\allowbreak{}\par};
\end{scope}\end{tikzpicture}\newpage
\definecolor{answer}{rgb}{0.850980,0.047059,0.070588}
\null\begin{tikzpicture}[remember picture,overlay]\begin{scope}[shift={(current page.north west)},x=1bp,y=-1bp]
\node[anchor=north west,inner sep=0pt,text=answer,align=left,text width=484.0bp] at (15,15) {\fontsize{7.5bp}{9.5bp}\selectfont 析：\allowbreak{}平均每班须用\ensuremath{3}天总数除以\ensuremath{4}班：\allowbreak{}(\ensuremath{87} + \ensuremath{75} + \ensuremath{65}) ÷ \ensuremath{4},\allowbreak{}选B。\allowbreak{}\par};
\end{scope}\end{tikzpicture}\newpage
\definecolor{answer}{rgb}{0.850980,0.047059,0.070588}
\null\begin{tikzpicture}[remember picture,overlay]\begin{scope}[shift={(current page.north west)},x=1bp,y=-1bp]
\node[anchor=north west,inner sep=0pt,text=answer,align=left,text width=225.0bp] at (15,15) {\fontsize{7.5bp}{9.5bp}\selectfont 连续减去，\allowbreak{}先凑\ensuremath{20}\par};
\end{scope}\end{tikzpicture}\newpage
\definecolor{answer}{rgb}{0.850980,0.047059,0.070588}
\null\begin{tikzpicture}[remember picture,overlay]\begin{scope}[shift={(current page.north west)},x=1bp,y=-1bp]
\node[anchor=north west,inner sep=0pt,text=answer,align=left,text width=484.0bp] at (15,15) {\fontsize{7.5bp}{9.5bp}\selectfont 析：\allowbreak{}比较对应顶点，\allowbreak{}先向右\ensuremath{8}格，\allowbreak{}再向下\ensuremath{6}格；\allowbreak{}数格距，\allowbreak{}非经过的线数。\allowbreak{}\par};
\end{scope}\end{tikzpicture}\newpage
\definecolor{answer}{rgb}{0.850980,0.047059,0.070588}
\null\begin{tikzpicture}[remember picture,overlay]\begin{scope}[shift={(current page.north west)},x=1bp,y=-1bp]
\node[anchor=north west,inner sep=0pt,text=answer,align=left,text width=484.0bp] at (15,15) {\fontsize{8.5bp}{10.62bp}\selectfont \ensuremath{6}月份在\ensuremath{4}月份每升\ensuremath{7.06}元的基础上调\ensuremath{0.21}元，\allowbreak{}\ensuremath{9}月份再下调\ensuremath{0.38}元；\allowbreak{}涨、\allowbreak{}跌方向不能弄反。\allowbreak{}\par};
\end{scope}\end{tikzpicture}\newpage
\definecolor{answer}{rgb}{0.850980,0.047059,0.070588}
\null\begin{tikzpicture}[remember picture,overlay]\begin{scope}[shift={(current page.north west)},x=1bp,y=-1bp]
\node[anchor=north west,inner sep=0pt,text=answer,align=left,text width=484.0bp] at (15,15) {\fontsize{8.5bp}{10.62bp}\selectfont 芹菜比西红柿多\ensuremath{136}千克，\allowbreak{}先求西红柿\ensuremath{328} \ensuremath{-}\ensuremath{136},\allowbreak{}即\ensuremath{192}千克。\allowbreak{}辣椒是西红柿的\ensuremath{2}倍再少\ensuremath{17}千克。\allowbreak{}\par};
\end{scope}\end{tikzpicture}\newpage
\definecolor{answer}{rgb}{0.850980,0.047059,0.070588}
\null\begin{tikzpicture}[remember picture,overlay]\begin{scope}[shift={(current page.north west)},x=1bp,y=-1bp]
\node[anchor=north west,inner sep=0pt,text=answer,align=left,text width=484.0bp] at (15,15) {\fontsize{7.5bp}{9.5bp}\selectfont 析：\allowbreak{}依等级边界逐一计数：\allowbreak{}优秀\ensuremath{9}、\allowbreak{}良好\ensuremath{11}、\allowbreak{}合格\ensuremath{9}、\allowbreak{}不合格\ensuremath{6},\allowbreak{}共\ensuremath{35}人；\allowbreak{}\ensuremath{100}归合格，\allowbreak{}\ensuremath{120}归良好，\allowbreak{}\ensuremath{140}归优秀。\allowbreak{}\par};
\end{scope}\end{tikzpicture}\newpage
\definecolor{answer}{rgb}{0.850980,0.047059,0.070588}
\null\begin{tikzpicture}[remember picture,overlay]\begin{scope}[shift={(current page.north west)},x=1bp,y=-1bp]
\node[anchor=north west,inner sep=0pt,text=answer,align=left,text width=484.0bp] at (15,15) {\fontsize{8.5bp}{10.62bp}\selectfont 只需比较四个等级的人数\ensuremath{9}、\allowbreak{}\ensuremath{11}、\allowbreak{}\ensuremath{9}、\allowbreak{}\ensuremath{6},\allowbreak{}无须列脱式算式。\allowbreak{}\ensuremath{11}最大，\allowbreak{}\ensuremath{6}最小。\allowbreak{}\par};
\end{scope}\end{tikzpicture}\newpage
\definecolor{answer}{rgb}{0.850980,0.047059,0.070588}
\null\begin{tikzpicture}[remember picture,overlay]\begin{scope}[shift={(current page.north west)},x=1bp,y=-1bp]
\node[anchor=north west,inner sep=0pt,text=answer,align=left,text width=484.0bp] at (15,15) {\fontsize{7.5bp}{9.5bp}\selectfont 析：\allowbreak{}三角形边长确定后形状固定，\allowbreak{}利用其稳定性。\allowbreak{}\par};
\end{scope}\end{tikzpicture}\newpage
\definecolor{answer}{rgb}{0.850980,0.047059,0.070588}
\null\begin{tikzpicture}[remember picture,overlay]\begin{scope}[shift={(current page.north west)},x=1bp,y=-1bp]
\node[anchor=north west,inner sep=0pt,text=answer,align=left,text width=484.0bp] at (15,15) {\fontsize{7.5bp}{9.5bp}\selectfont 析：\allowbreak{}千克→吨除\ensuremath{1000};\allowbreak{}平方米→平方分米乘\ensuremath{100};\allowbreak{}元→分乘\ensuremath{100};\allowbreak{}米→千米除\ensuremath{1000},\allowbreak{}得\ensuremath{8.05}、\allowbreak{}\ensuremath{360}、\allowbreak{}\ensuremath{2}和\ensuremath{60}、\allowbreak{}\ensuremath{6.08}。\allowbreak{}\par};
\end{scope}\end{tikzpicture}\newpage
\definecolor{answer}{rgb}{0.850980,0.047059,0.070588}
\null\begin{tikzpicture}[remember picture,overlay]\begin{scope}[shift={(current page.north west)},x=1bp,y=-1bp]
\node[anchor=north west,inner sep=0pt,text=answer,align=left,text width=484.0bp] at (15,15) {\fontsize{7.5bp}{9.5bp}\selectfont 析：\allowbreak{}末尾添\ensuremath{0}大小不变，\allowbreak{}计数单位由\ensuremath{0.1}变\ensuremath{0.001};\allowbreak{}\ensuremath{0.7}为\ensuremath{7}个十分之一，\allowbreak{}\ensuremath{0.700}为\ensuremath{700}个千分之一。\allowbreak{}\par};
\end{scope}\end{tikzpicture}\newpage
\definecolor{answer}{rgb}{0.850980,0.047059,0.070588}
\null\begin{tikzpicture}[remember picture,overlay]\begin{scope}[shift={(current page.north west)},x=1bp,y=-1bp]
\node[anchor=north west,inner sep=0pt,text=answer,align=left,text width=484.0bp] at (15,15) {\fontsize{7.5bp}{9.5bp}\selectfont 析：\allowbreak{}\ensuremath{7}不能作腰（\ensuremath{7} + \ensuremath{7}＜\ensuremath{15}),\allowbreak{}故两腰\ensuremath{15}、\allowbreak{}底\ensuremath{7};\allowbreak{}周长\ensuremath{15} + \ensuremath{15} + \ensuremath{7} = \ensuremath{37}厘米。\allowbreak{}\par};
\end{scope}\end{tikzpicture}\newpage
\definecolor{answer}{rgb}{0.850980,0.047059,0.070588}
\null\begin{tikzpicture}[remember picture,overlay]\begin{scope}[shift={(current page.north west)},x=1bp,y=-1bp]
\node[anchor=north west,inner sep=0pt,text=answer,align=left,text width=484.0bp] at (15,15) {\fontsize{7.5bp}{9.5bp}\selectfont 析：\allowbreak{}垂线段最短，\allowbreak{}\ensuremath{216}、\allowbreak{}\ensuremath{178}、\allowbreak{}\ensuremath{243}中最短\ensuremath{178},\allowbreak{}垂直小路长\ensuremath{178}米。\allowbreak{}\par};
\end{scope}\end{tikzpicture}\newpage
\definecolor{answer}{rgb}{0.850980,0.047059,0.070588}
\null\begin{tikzpicture}[remember picture,overlay]\begin{scope}[shift={(current page.north west)},x=1bp,y=-1bp]
\node[anchor=north west,inner sep=0pt,text=answer,align=left,text width=484.0bp] at (15,15) {\fontsize{7.5bp}{9.5bp}\selectfont 析：\allowbreak{}小数不都属于整数，\allowbreak{}A关系错误；\allowbreak{}等边属等腰，\allowbreak{}平行四边形属四边形，\allowbreak{}B、\allowbreak{}C正确，\allowbreak{}选A。\allowbreak{}\par};
\end{scope}\end{tikzpicture}\newpage
\definecolor{answer}{rgb}{0.850980,0.047059,0.070588}
\null\begin{tikzpicture}[remember picture,overlay]\begin{scope}[shift={(current page.north west)},x=1bp,y=-1bp]
\node[anchor=north west,inner sep=0pt,text=answer,align=left,text width=484.0bp] at (15,15) {\fontsize{7.5bp}{9.5bp}\selectfont 析：\allowbreak{}若最大角＜\ensuremath{60}\ensuremath{^\circ},\allowbreak{}三角和就＜\ensuremath{180}\ensuremath{^\circ},\allowbreak{}不可能；\allowbreak{}至少\ensuremath{60}\ensuremath{^\circ},\allowbreak{}选B。\allowbreak{}\par};
\end{scope}\end{tikzpicture}\newpage
\definecolor{answer}{rgb}{0.850980,0.047059,0.070588}
\null\begin{tikzpicture}[remember picture,overlay]\begin{scope}[shift={(current page.north west)},x=1bp,y=-1bp]
\node[anchor=north west,inner sep=0pt,text=answer,align=left,text width=484.0bp] at (15,15) {\fontsize{7.5bp}{9.5bp}\selectfont 析：\allowbreak{}A沿垂线移动使高和顶角变化，\allowbreak{}可依次出现锐角、\allowbreak{}直角、\allowbreak{}钝角三角形，\allowbreak{}选C。\allowbreak{}\par};
\end{scope}\end{tikzpicture}\newpage
\definecolor{answer}{rgb}{0.850980,0.047059,0.070588}
\null\begin{tikzpicture}[remember picture,overlay]\begin{scope}[shift={(current page.north west)},x=1bp,y=-1bp]
\node[anchor=north west,inner sep=0pt,text=answer,align=left,text width=334.8704833984375bp] at (15,15) {\fontsize{7.5bp}{9.5bp}\selectfont 析：\allowbreak{}B在O正右方，\allowbreak{}绕O顺时针\ensuremath{90}\ensuremath{^\circ}到正下方C,\allowbreak{}选C。\allowbreak{}\par};
\end{scope}\end{tikzpicture}\newpage
\definecolor{answer}{rgb}{0.850980,0.047059,0.070588}
\null\begin{tikzpicture}[remember picture,overlay]\begin{scope}[shift={(current page.north west)},x=1bp,y=-1bp]
\node[anchor=north west,inner sep=0pt,text=answer,align=left,text width=432.4945983886719bp] at (15,15) {\fontsize{8.5bp}{10.62bp}\selectfont 三种费用为\ensuremath{400}、\allowbreak{}\ensuremath{336}、\allowbreak{}\ensuremath{330}元；\allowbreak{}简算：\allowbreak{}\ensuremath{400} + \ensuremath{336} + \ensuremath{330} = (\ensuremath{400} + \ensuremath{330}) + \ensuremath{336} = \ensuremath{730} + \ensuremath{336} = \ensuremath{1066}(元)。\allowbreak{}\par};
\end{scope}\end{tikzpicture}\newpage
\definecolor{answer}{rgb}{0.850980,0.047059,0.070588}
\null\begin{tikzpicture}[remember picture,overlay]\begin{scope}[shift={(current page.north west)},x=1bp,y=-1bp]
\node[anchor=north west,inner sep=0pt,text=answer,align=left,text width=225bp] at (15,15) {\fontsize{7.5bp}{9.5bp}\selectfont 析：\allowbreak{}按角三分：\allowbreak{}锐角、\allowbreak{}直角、\allowbreak{}钝角三角形互不重叠，\allowbreak{}合起来是全部三角形。\allowbreak{}\par};
\end{scope}\end{tikzpicture}\newpage
\definecolor{answer}{rgb}{0.850980,0.047059,0.070588}
\null\begin{tikzpicture}[remember picture,overlay]\begin{scope}[shift={(current page.north west)},x=1bp,y=-1bp]
\node[anchor=north west,inner sep=0pt,text=answer,align=left,text width=206bp] at (15,15) {\fontsize{7.5bp}{9.5bp}\selectfont 析：\allowbreak{}等边一定等腰，\allowbreak{}等腰属于三角形；\allowbreak{}外到内填三角形、\allowbreak{}等腰三角形、\allowbreak{}等边三角形。\allowbreak{}\par};
\end{scope}\end{tikzpicture}\newpage
\definecolor{answer}{rgb}{0.850980,0.047059,0.070588}
\null\begin{tikzpicture}[remember picture,overlay]\begin{scope}[shift={(current page.north west)},x=1bp,y=-1bp]
\node[anchor=north west,inner sep=0pt,text=answer,align=left,text width=264bp] at (15,15) {\fontsize{7.5bp}{9.5bp}\selectfont 析：\allowbreak{}扫描近似长\ensuremath{4}厘米、\allowbreak{}宽\ensuremath{2}厘米；\allowbreak{}两组对边平行、\allowbreak{}邻边垂直。\allowbreak{}测量以实际纸面为准。\allowbreak{}\par};
\end{scope}\end{tikzpicture}\newpage
\definecolor{answer}{rgb}{0.850980,0.047059,0.070588}
\null\begin{tikzpicture}[remember picture,overlay]\begin{scope}[shift={(current page.north west)},x=1bp,y=-1bp]
\node[anchor=north west,inner sep=0pt,text=answer,align=left,text width=191bp] at (15,15) {\fontsize{7.5bp}{9.5bp}\selectfont 析：\allowbreak{}给定边扫描近似\ensuremath{2}厘米，\allowbreak{}另三边等长且相邻垂直；\allowbreak{}实际以纸面测量为准。\allowbreak{}\par};
\end{scope}\end{tikzpicture}\newpage
\definecolor{answer}{rgb}{0.850980,0.047059,0.070588}
\null\begin{tikzpicture}[remember picture,overlay]\begin{scope}[shift={(current page.north west)},x=1bp,y=-1bp]
\node[anchor=north west,inner sep=0pt,text=answer,align=left,text width=244bp] at (15,15) {\fontsize{8.5bp}{10.62bp}\selectfont 表中单价单位应为元/千克。\allowbreak{}销量=营业额÷单价；\allowbreak{}分别求销量后比较。\allowbreak{}\par};
\end{scope}\end{tikzpicture}\newpage
\definecolor{answer}{rgb}{0.850980,0.047059,0.070588}
\null\begin{tikzpicture}[remember picture,overlay]\begin{scope}[shift={(current page.north west)},x=1bp,y=-1bp]
\node[anchor=north west,inner sep=0pt,text=answer,align=left,text width=501bp] at (15,15) {\fontsize{7.5bp}{9.5bp}\selectfont 析：\allowbreak{}（\ensuremath{2}）整边拼接共\ensuremath{6}种：\allowbreak{}沿\ensuremath{3}厘米边得等腰三角形、\allowbreak{}平行四边形，\allowbreak{}周长均\ensuremath{18}厘米；\allowbreak{}沿\ensuremath{4}厘米边同样两种，\allowbreak{}周长均\ensuremath{16}厘米；\allowbreak{}沿\ensuremath{5}厘米边得长方形和另一个四边形，\allowbreak{}周长均\ensuremath{14}厘米。\allowbreak{}仅改变摆放方向或翻转不重复计数。\allowbreak{}\par};
\end{scope}\end{tikzpicture}\newpage
\definecolor{answer}{rgb}{0.850980,0.047059,0.070588}
\null\begin{tikzpicture}[remember picture,overlay]\begin{scope}[shift={(current page.north west)},x=1bp,y=-1bp]
\node[anchor=north west,inner sep=0pt,text=answer,align=left,text width=501bp] at (15,15) {\fontsize{7.5bp}{9.5bp}\selectfont 析：\allowbreak{}梯形只有一组对边平行，\allowbreak{}不混入平行四边形。\allowbreak{}分类数：\allowbreak{}平行四边形\ensuremath{4} + \ensuremath{2} + \ensuremath{2} + \ensuremath{1} = \ensuremath{9}，\allowbreak{}梯形\ensuremath{2} + \ensuremath{3} + \ensuremath{2} = \ensuremath{7}，\allowbreak{}三角形\ensuremath{1} + \ensuremath{1} = \ensuremath{2}；\allowbreak{}也可按平行线组逐类数。\allowbreak{}\par};
\end{scope}\end{tikzpicture}\newpage
\definecolor{answer}{rgb}{0.850980,0.047059,0.070588}
\null\begin{tikzpicture}[remember picture,overlay]\begin{scope}[shift={(current page.north west)},x=1bp,y=-1bp]
\node[anchor=north west,inner sep=0pt,text=answer,align=left,text width=169.0bp] at (15,15) {\fontsize{7.5bp}{9.5bp}\selectfont 析：\allowbreak{}不对：\allowbreak{}只见一组对边平行，\allowbreak{}另一组被遮住，\allowbreak{}可能是平行四边形，\allowbreak{}也可能是梯形。\allowbreak{}\par};
\end{scope}\end{tikzpicture}\newpage
\definecolor{answer}{rgb}{0.850980,0.047059,0.070588}
\null\begin{tikzpicture}[remember picture,overlay]\begin{scope}[shift={(current page.north west)},x=1bp,y=-1bp]
\node[anchor=north west,inner sep=0pt,text=answer,align=left,text width=196.67000000000002bp] at (15,15) {\fontsize{9.33bp}{11.66bp}\selectfont 只能看出一组对边平行;\allowbreak{}另一组对边的情况被遮住,\allowbreak{}不能确定是平行四边形。\allowbreak{}\par};
\end{scope}\end{tikzpicture}\newpage
\definecolor{answer}{rgb}{0.850980,0.047059,0.070588}
\null\begin{tikzpicture}[remember picture,overlay]\begin{scope}[shift={(current page.north west)},x=1bp,y=-1bp]
\node[anchor=north west,inner sep=0pt,text=answer,align=left,text width=241bp] at (15,15) {\fontsize{7.5bp}{9.5bp}\selectfont 移动减去的\ensuremath{255}须保留减号，\allowbreak{}把\ensuremath{355} \ensuremath{-}\ensuremath{255}先凑成\ensuremath{100}，\allowbreak{}再与\ensuremath{172}、\allowbreak{}\ensuremath{1285}相加。\allowbreak{}\par};
\end{scope}\end{tikzpicture}\newpage
\definecolor{answer}{rgb}{0.850980,0.047059,0.070588}
\null\begin{tikzpicture}[remember picture,overlay]\begin{scope}[shift={(current page.north west)},x=1bp,y=-1bp]
\node[anchor=north west,inner sep=0pt,text=answer,align=left,text width=222bp] at (15,15) {\fontsize{7.5bp}{9.5bp}\selectfont 析：\allowbreak{}（\ensuremath{1})小数点右第\ensuremath{2}位是百分位，\allowbreak{}单位\ensuremath{0.01};\allowbreak{}整数右第\ensuremath{3}位是百位，\allowbreak{}单位百。\allowbreak{}\par};
\end{scope}\end{tikzpicture}\newpage
\definecolor{answer}{rgb}{0.850980,0.047059,0.070588}
\null\begin{tikzpicture}[remember picture,overlay]\begin{scope}[shift={(current page.north west)},x=1bp,y=-1bp]
\node[anchor=north west,inner sep=0pt,text=answer,align=left,text width=256bp] at (15,15) {\fontsize{7.5bp}{9.5bp}\selectfont 析：\allowbreak{}（\ensuremath{2}）\ensuremath{154.31}读一百五十四点三一，\allowbreak{}\ensuremath{5}在十位、\allowbreak{}表示\ensuremath{5}个十；\allowbreak{}\ensuremath{0.13} = \ensuremath{13}个\ensuremath{0.01} = \ensuremath{1}个\ensuremath{0.1} + \ensuremath{3}个\ensuremath{0.01}。\allowbreak{}\par};
\end{scope}\end{tikzpicture}\newpage
\definecolor{answer}{rgb}{0.850980,0.047059,0.070588}
\null\begin{tikzpicture}[remember picture,overlay]\begin{scope}[shift={(current page.north west)},x=1bp,y=-1bp]
\node[anchor=north west,inner sep=0pt,text=answer,align=left,text width=276bp] at (15,15) {\fontsize{7.5bp}{9.5bp}\selectfont 析：\allowbreak{}大于\ensuremath{3},\allowbreak{}整数部分只能取\ensuremath{5};\allowbreak{}\ensuremath{0}、\allowbreak{}\ensuremath{1}、\allowbreak{}\ensuremath{2}全排列，\allowbreak{}共\ensuremath{6}种，\allowbreak{}见原空。\allowbreak{}\par};
\end{scope}\end{tikzpicture}\newpage
\definecolor{answer}{rgb}{0.850980,0.047059,0.070588}
\null\begin{tikzpicture}[remember picture,overlay]\begin{scope}[shift={(current page.north west)},x=1bp,y=-1bp]
\node[anchor=north west,inner sep=0pt,text=answer,align=left,text width=276bp] at (15,15) {\fontsize{7.5bp}{9.5bp}\selectfont 析：\allowbreak{}小于\ensuremath{1},\allowbreak{}整数部分只能取\ensuremath{0};\allowbreak{}\ensuremath{1}、\allowbreak{}\ensuremath{2}、\allowbreak{}\ensuremath{5}全排列，\allowbreak{}共\ensuremath{6}种，\allowbreak{}见原空。\allowbreak{}\par};
\end{scope}\end{tikzpicture}\newpage
\definecolor{answer}{rgb}{0.850980,0.047059,0.070588}
\null\begin{tikzpicture}[remember picture,overlay]\begin{scope}[shift={(current page.north west)},x=1bp,y=-1bp]
\node[anchor=north west,inner sep=0pt,text=answer,align=left,text width=484bp] at (15,15) {\fontsize{7.5bp}{9.5bp}\selectfont \ensuremath{273.2}万是按万为单位保留一位小数得到的数，\allowbreak{}相邻两个刻度相差\ensuremath{0.1}万，\allowbreak{}取近似数的边界各相差\ensuremath{0.05}万。\allowbreak{}\ensuremath{273.15}万能向上进成\ensuremath{273.2}万，\allowbreak{}应包含；\allowbreak{}\ensuremath{273.25}万会进成\ensuremath{273.3}万，\allowbreak{}不包含。\allowbreak{}\par};
\end{scope}\end{tikzpicture}\newpage
\definecolor{answer}{rgb}{0.850980,0.047059,0.070588}
\null\begin{tikzpicture}[remember picture,overlay]\begin{scope}[shift={(current page.north west)},x=1bp,y=-1bp]
\node[anchor=north west,inner sep=0pt,text=answer,align=left,text width=152bp] at (15,15) {\fontsize{7.5bp}{9.5bp}\selectfont 析：\allowbreak{}整数只能\ensuremath{6};\allowbreak{}小数可\ensuremath{5.7}、\allowbreak{}\ensuremath{6.5}等，\allowbreak{}不限位数则无数个。\allowbreak{}\par};
\end{scope}\end{tikzpicture}\newpage
\definecolor{answer}{rgb}{0.850980,0.047059,0.070588}
\null\begin{tikzpicture}[remember picture,overlay]\begin{scope}[shift={(current page.north west)},x=1bp,y=-1bp]
\node[anchor=north west,inner sep=0pt,text=answer,align=left,text width=136bp] at (15,15) {\fontsize{7.5bp}{9.5bp}\selectfont 析：\allowbreak{}右框\ensuremath{1}至\ensuremath{8},\allowbreak{}左框取\ensuremath{0}至(右框-\ensuremath{1}),\allowbreak{}共\ensuremath{1} +\ensuremath{2} + · · · + \ensuremath{8} = \ensuremath{36}组。\allowbreak{}\par};
\end{scope}\end{tikzpicture}\newpage
\definecolor{answer}{rgb}{0.850980,0.047059,0.070588}
\null\begin{tikzpicture}[remember picture,overlay]\begin{scope}[shift={(current page.north west)},x=1bp,y=-1bp]
\node[anchor=north west,inner sep=0pt,text=answer,align=left,text width=162bp] at (15,15) {\fontsize{7.5bp}{9.5bp}\selectfont 析：\allowbreak{}全三组：\allowbreak{}(\ensuremath{1},\allowbreak{} \ensuremath{0})、\allowbreak{}(\ensuremath{2},\allowbreak{} \ensuremath{0})、\allowbreak{}(\ensuremath{2},\allowbreak{} \ensuremath{1})。\allowbreak{}\par};
\end{scope}\end{tikzpicture}\newpage
\definecolor{answer}{rgb}{0.850980,0.047059,0.070588}
\null\begin{tikzpicture}[remember picture,overlay]\begin{scope}[shift={(current page.north west)},x=1bp,y=-1bp]
\node[anchor=north west,inner sep=0pt,text=answer,align=left,text width=136bp] at (15,15) {\fontsize{7.5bp}{9.5bp}\selectfont 析：\allowbreak{}按\ensuremath{1}位数字只填\ensuremath{0};\allowbreak{}允多位时可写\ensuremath{01}等，\allowbreak{}组成\ensuremath{9.01}米。\allowbreak{}\par};
\end{scope}\end{tikzpicture}\newpage
\definecolor{answer}{rgb}{0.850980,0.047059,0.070588}
\null\begin{tikzpicture}[remember picture,overlay]\begin{scope}[shift={(current page.north west)},x=1bp,y=-1bp]
\node[anchor=north west,inner sep=0pt,text=answer,align=left,text width=162bp] at (15,15) {\fontsize{7.5bp}{9.5bp}\selectfont 分钟数<\ensuremath{30}，\allowbreak{}可填\ensuremath{00}至\ensuremath{29}。\allowbreak{}\par};
\end{scope}\end{tikzpicture}\newpage
\definecolor{answer}{rgb}{0.850980,0.047059,0.070588}
\null\begin{tikzpicture}[remember picture,overlay]\begin{scope}[shift={(current page.north west)},x=1bp,y=-1bp]
\node[anchor=north west,inner sep=0pt,text=answer,align=left,text width=222bp] at (15,15) {\fontsize{7.5bp}{9.5bp}\selectfont 析：\allowbreak{}环数越高越好：\allowbreak{}\ensuremath{9.7}＞\ensuremath{9.5}＞\ensuremath{9.1}＞\ensuremath{8.9};\allowbreak{}前三为李亮、\allowbreak{}赵娜、\allowbreak{}白雪。\allowbreak{}\par};
\end{scope}\end{tikzpicture}\newpage
\definecolor{answer}{rgb}{0.850980,0.047059,0.070588}
\null\begin{tikzpicture}[remember picture,overlay]\begin{scope}[shift={(current page.north west)},x=1bp,y=-1bp]
\node[anchor=north west,inner sep=0pt,text=answer,align=left,text width=256bp] at (15,15) {\fontsize{7.5bp}{9.5bp}\selectfont 析：\allowbreak{}用时越少越快：\allowbreak{}\ensuremath{14.3}＜\ensuremath{14.5}＜\ensuremath{14.8}＜\ensuremath{15.2};\allowbreak{}前三为赵娜、\allowbreak{}白雪、\allowbreak{}李亮。\allowbreak{}\par};
\end{scope}\end{tikzpicture}\newpage
\definecolor{answer}{rgb}{0.850980,0.047059,0.070588}
\null\begin{tikzpicture}[remember picture,overlay]\begin{scope}[shift={(current page.north west)},x=1bp,y=-1bp]
\node[anchor=north west,inner sep=0pt,text=answer,align=left,text width=222bp] at (15,15) {\fontsize{8.5bp}{10.62bp}\selectfont 两件原价合\ensuremath{157.8}元，\allowbreak{}达到满\ensuremath{150}元减\ensuremath{20}元的条件，\allowbreak{}只需付\ensuremath{137.8}元，\allowbreak{}再从钱包余额扣除。\allowbreak{}\par};
\end{scope}\end{tikzpicture}\newpage
\definecolor{answer}{rgb}{0.850980,0.047059,0.070588}
\null\begin{tikzpicture}[remember picture,overlay]\begin{scope}[shift={(current page.north west)},x=1bp,y=-1bp]
\node[anchor=north west,inner sep=0pt,text=answer,align=left,text width=272.375bp] at (15,15) {\fontsize{8.5bp}{10.62bp}\selectfont 门高标注待核。\allowbreak{}若门高为\ensuremath{1.7}米，\allowbreak{}则黑板顶比门顶高\ensuremath{0.35}米，\allowbreak{}底离地\ensuremath{0.8}米。\allowbreak{}\par};
\end{scope}\end{tikzpicture}\newpage
\definecolor{answer}{rgb}{0.850980,0.047059,0.070588}
\null\begin{tikzpicture}[remember picture,overlay]\begin{scope}[shift={(current page.north west)},x=1bp,y=-1bp]
\node[anchor=north west,inner sep=0pt,text=answer,align=left,text width=484bp] at (15,15) {\fontsize{8.5bp}{10.62bp}\selectfont 平均成绩\ensuremath{150}个表示二班\ensuremath{8}名参赛同学的整体水平。\allowbreak{}乐乐\ensuremath{145}个，\allowbreak{}比平均水平少\ensuremath{5}个。\allowbreak{}仅由平均数不能确定乐乐的具体名次，\allowbreak{}也不能确定有几个人比乐乐高。\allowbreak{}\par};
\end{scope}\end{tikzpicture}\newpage
\definecolor{answer}{rgb}{0.850980,0.047059,0.070588}
\null\begin{tikzpicture}[remember picture,overlay]\begin{scope}[shift={(current page.north west)},x=1bp,y=-1bp]
\node[anchor=north west,inner sep=0pt,text=answer,align=left,text width=484bp] at (15,15) {\fontsize{8.5bp}{10.62bp}\selectfont 两组各\ensuremath{20}人，\allowbreak{}可比较同一分数段的人数。\allowbreak{}男教师投中\ensuremath{11}个及以上共\ensuremath{16}人，\allowbreak{}女教师共\ensuremath{8}人。\allowbreak{}据成绩分段分布，\allowbreak{}可认为男教师整体成绩较好；\allowbreak{}表格未给组内每人的准确成绩，\allowbreak{}不能据此求精确平均数。\allowbreak{}\par};
\end{scope}\end{tikzpicture}\newpage
\definecolor{answer}{rgb}{0.850980,0.047059,0.070588}
\null\begin{tikzpicture}[remember picture,overlay]\begin{scope}[shift={(current page.north west)},x=1bp,y=-1bp]
\node[anchor=north west,inner sep=0pt,text=answer,align=left,text width=484bp] at (15,15) {\fontsize{8.5bp}{10.62bp}\selectfont 两船同时相向行驶，\allowbreak{}\ensuremath{3}小时共航行\ensuremath{126}千米；\allowbreak{}原距离\ensuremath{172}千米大于\ensuremath{126},\allowbreak{}尚未相遇，\allowbreak{}应减去已航行的路程。\allowbreak{}原页线段图表示两船分别走了\ensuremath{24} × \ensuremath{3}、\allowbreak{}\ensuremath{18} × \ensuremath{3}千米，\allowbreak{}中间剩余部分就是所求距离。\allowbreak{}\par};
\end{scope}\end{tikzpicture}\newpage
\definecolor{answer}{rgb}{0.850980,0.047059,0.070588}
\null\begin{tikzpicture}[remember picture,overlay]\begin{scope}[shift={(current page.north west)},x=1bp,y=-1bp]
\node[anchor=north west,inner sep=0pt,text=answer,align=left,text width=197.462bp] at (15,15) {\fontsize{9.23bp}{11.54bp}\selectfont 分别向两腰作垂线,\allowbreak{} 垂足D、\allowbreak{}E即固定点。\allowbreak{}点到直线的垂线段最短。\allowbreak{}\par};
\end{scope}\end{tikzpicture}\newpage
\definecolor{answer}{rgb}{0.850980,0.047059,0.070588}
\null\begin{tikzpicture}[remember picture,overlay]\begin{scope}[shift={(current page.north west)},x=1bp,y=-1bp]
\node[anchor=north west,inner sep=0pt,text=answer,align=left,text width=484bp] at (15,15) {\fontsize{8.5bp}{10.62bp}\selectfont \ensuremath{6}小时后还相差\ensuremath{160}千米，\allowbreak{}说明两车尚未相遇，\allowbreak{}合行路程为\ensuremath{1120} \ensuremath{-}\ensuremath{160},\allowbreak{}即\ensuremath{960}千米。\allowbreak{}把\ensuremath{960}平均分给\ensuremath{6}小时，\allowbreak{}求速度和，\allowbreak{}再减甲车\ensuremath{78}千米/时，\allowbreak{}得乙车速度。\allowbreak{}原页线段图已标两车路程与剩余\ensuremath{160}千米。\allowbreak{}\par};
\end{scope}\end{tikzpicture}\newpage
\definecolor{answer}{rgb}{0.850980,0.047059,0.070588}
\null\begin{tikzpicture}[remember picture,overlay]\begin{scope}[shift={(current page.north west)},x=1bp,y=-1bp]
\node[anchor=north west,inner sep=0pt,text=answer,align=left,text width=484bp] at (15,15) {\fontsize{7.5bp}{9.5bp}\selectfont 析：\allowbreak{}①\ensuremath{75} + \ensuremath{50} = \ensuremath{125}表示速度之和；\allowbreak{}②\ensuremath{75} × \ensuremath{4} = \ensuremath{300}为快车路程；\allowbreak{}③\ensuremath{50} × \ensuremath{4} = \ensuremath{200}为慢车路程；\allowbreak{}④(\ensuremath{75} + \ensuremath{50}) × \ensuremath{4} = \ensuremath{500}为两城距离（也可\ensuremath{75} × \ensuremath{4} + \ensuremath{50} × \ensuremath{4})。\allowbreak{}\par};
\end{scope}\end{tikzpicture}\newpage
\definecolor{answer}{rgb}{0.850980,0.047059,0.070588}
\null\begin{tikzpicture}[remember picture,overlay]\begin{scope}[shift={(current page.north west)},x=1bp,y=-1bp]
\node[anchor=north west,inner sep=0pt,text=answer,align=left,text width=483.0bp] at (15,15) {\fontsize{7.5bp}{9.5bp}\selectfont 析：\allowbreak{}例如\ensuremath{5}、\allowbreak{}\ensuremath{5}、\allowbreak{}\ensuremath{6},\allowbreak{}和\ensuremath{16}且\ensuremath{5} + \ensuremath{5} > \ensuremath{6}。\allowbreak{}若限整数，\allowbreak{}不计换序共有(\ensuremath{2},\allowbreak{} \ensuremath{7},\allowbreak{} \ensuremath{7})、\allowbreak{}(\ensuremath{3},\allowbreak{} \ensuremath{6},\allowbreak{} \ensuremath{7})、\allowbreak{}(\ensuremath{4},\allowbreak{} \ensuremath{5},\allowbreak{} \ensuremath{7})、\allowbreak{}(\ensuremath{4},\allowbreak{} \ensuremath{6},\allowbreak{} \ensuremath{6})、\allowbreak{}(\ensuremath{5},\allowbreak{} \ensuremath{5},\allowbreak{} \ensuremath{6});\allowbreak{}不限整数有无穷多组。\allowbreak{}\par};
\end{scope}\end{tikzpicture}\newpage
\definecolor{answer}{rgb}{0.850980,0.047059,0.070588}
\null\begin{tikzpicture}[remember picture,overlay]\begin{scope}[shift={(current page.north west)},x=1bp,y=-1bp]
\node[anchor=north west,inner sep=0pt,text=answer,align=left,text width=501bp] at (15,15) {\fontsize{8.5bp}{10.62bp}\selectfont 过\ensuremath{A}作平行线：\allowbreak{}将三角板的一边对齐下水道，\allowbreak{}沿直尺平移到\ensuremath{A}点，\allowbreak{}过\ensuremath{A}画直线。\allowbreak{}\par};
\end{scope}\end{tikzpicture}\newpage
\definecolor{answer}{rgb}{0.850980,0.047059,0.070588}
\null\begin{tikzpicture}[remember picture,overlay]\begin{scope}[shift={(current page.north west)},x=1bp,y=-1bp]
\node[anchor=north west,inner sep=0pt,text=answer,align=left,text width=501bp] at (15,15) {\fontsize{8.5bp}{10.62bp}\selectfont 沿垂线连接下水道：\allowbreak{}从\ensuremath{A}向下水道作垂线，\allowbreak{}取\ensuremath{A}与垂足之间的线段作为排水沟，\allowbreak{}垂足标直角。\allowbreak{}\par};
\end{scope}\end{tikzpicture}\newpage
\definecolor{answer}{rgb}{0.850980,0.047059,0.070588}
\null\begin{tikzpicture}[remember picture,overlay]\begin{scope}[shift={(current page.north west)},x=1bp,y=-1bp]
\node[anchor=north west,inner sep=0pt,text=answer,align=left,text width=483.0bp] at (15,15) {\fontsize{8.5bp}{10.62bp}\selectfont 两车同时从同地反向开出，\allowbreak{}每小时相隔距离为速度和，\allowbreak{}\ensuremath{3}小时的距离用速度和乘\ensuremath{3}。\allowbreak{}\par};
\end{scope}\end{tikzpicture}\newpage
\definecolor{answer}{rgb}{0.850980,0.047059,0.070588}
\null\begin{tikzpicture}[remember picture,overlay]\begin{scope}[shift={(current page.north west)},x=1bp,y=-1bp]
\node[anchor=north west,inner sep=0pt,text=answer,align=left,text width=483.0bp] at (15,15) {\fontsize{7.5bp}{9.5bp}\selectfont 析：\allowbreak{}沿对角线分成两全同三角形；\allowbreak{}连接一组对边的中点，\allowbreak{}分两全同平行四边形；\allowbreak{}过中心斜切且不过顶点、\allowbreak{}也不平行于边，\allowbreak{}可得两全同梯形。\allowbreak{}\par};
\end{scope}\end{tikzpicture}\newpage
\definecolor{answer}{rgb}{0.850980,0.047059,0.070588}
\null\begin{tikzpicture}[remember picture,overlay]\begin{scope}[shift={(current page.north west)},x=1bp,y=-1bp]
\node[anchor=north west,inner sep=0pt,text=answer,align=left,text width=501.0bp] at (15,15) {\fontsize{7.5bp}{9.5bp}\selectfont 析：\allowbreak{}按教材常见意图判√，\allowbreak{}说明平均数不等于每次值；\allowbreak{}严格说，\allowbreak{}仅凭成年人总体平均\ensuremath{75},\allowbreak{}不能推出妈妈必有高于、\allowbreak{}低于\ensuremath{75}的时刻，\allowbreak{}条件不足。\allowbreak{}\par};
\end{scope}\end{tikzpicture}\newpage
\definecolor{answer}{rgb}{0.850980,0.047059,0.070588}
\null\begin{tikzpicture}[remember picture,overlay]\begin{scope}[shift={(current page.north west)},x=1bp,y=-1bp]
\node[anchor=north west,inner sep=0pt,text=answer,align=left,text width=501.0bp] at (15,15) {\fontsize{7.5bp}{9.5bp}\selectfont 析：\allowbreak{}端点只归一段：\allowbreak{}\ensuremath{140}及以上、\allowbreak{}\ensuremath{130}～\ensuremath{139}、\allowbreak{}\ensuremath{120}～\ensuremath{129}、\allowbreak{}\ensuremath{119}及以下。\allowbreak{}两班频数\ensuremath{4}、\allowbreak{}\ensuremath{5}、\allowbreak{}\ensuremath{4}、\allowbreak{}\ensuremath{2}和\ensuremath{1}、\allowbreak{}\ensuremath{4}、\allowbreak{}\ensuremath{6}、\allowbreak{}\ensuremath{4}，\allowbreak{}各\ensuremath{15}人；\allowbreak{}纵向\ensuremath{5}、\allowbreak{}\ensuremath{9}、\allowbreak{}\ensuremath{10}、\allowbreak{}\ensuremath{6}，\allowbreak{}共\ensuremath{30}人。\allowbreak{}\par};
\end{scope}\end{tikzpicture}\newpage
\definecolor{answer}{rgb}{0.850980,0.047059,0.070588}
\null\begin{tikzpicture}[remember picture,overlay]\begin{scope}[shift={(current page.north west)},x=1bp,y=-1bp]
\node[anchor=north west,inner sep=0pt,text=answer,align=left,text width=501.0bp] at (15,15) {\fontsize{7.5bp}{9.5bp}\selectfont 析：\allowbreak{}A*：\allowbreak{}每分钟分别转\ensuremath{45}°、\allowbreak{}\ensuremath{36}°，\allowbreak{}(\ensuremath{45} + \ensuremath{36}) × \ensuremath{14} \ensuremath{-}\ensuremath{360} × \ensuremath{3} = \ensuremath{54}\ensuremath{^\circ}。\allowbreak{}按示意相邻位置选A；\allowbreak{}插图未按比例，\allowbreak{}若按精确角度三图均不准确。\allowbreak{}\par};
\end{scope}\end{tikzpicture}\newpage
\definecolor{answer}{rgb}{0.850980,0.047059,0.070588}
\null\begin{tikzpicture}[remember picture,overlay]\begin{scope}[shift={(current page.north west)},x=1bp,y=-1bp]
\node[anchor=north west,inner sep=0pt,text=answer,align=left,text width=501.0bp] at (15,15) {\fontsize{8.5bp}{10.62bp}\selectfont 可按销量判断需求，\allowbreak{}以同样销售条件为前提举例建议多进香蕉；\allowbreak{}不能仅由一天销售量保证今后必然相同。\allowbreak{}\par};
\end{scope}\end{tikzpicture}\newpage
\definecolor{answer}{rgb}{0.850980,0.047059,0.070588}
\null\begin{tikzpicture}[remember picture,overlay]\begin{scope}[shift={(current page.north west)},x=1bp,y=-1bp]
\node[anchor=north west,inner sep=0pt,text=answer,align=left,text width=222.0bp] at (15,15) {\fontsize{7.5bp}{9.5bp}\selectfont 析：\allowbreak{}①选同一顶点按箭头数格：\allowbreak{}先下\ensuremath{4}格，\allowbreak{}再右\ensuremath{3}格。\allowbreak{}\par};
\end{scope}\end{tikzpicture}\newpage
\definecolor{answer}{rgb}{0.850980,0.047059,0.070588}
\null\begin{tikzpicture}[remember picture,overlay]\begin{scope}[shift={(current page.north west)},x=1bp,y=-1bp]
\node[anchor=north west,inner sep=0pt,text=answer,align=left,text width=484.0bp] at (15,15) {\fontsize{8.5bp}{10.62bp}\selectfont 两人都走\ensuremath{6}分钟，\allowbreak{}小红每分钟比小明多走\ensuremath{5}米，\allowbreak{}路程就多\ensuremath{30}米。\allowbreak{}比较的是各自家离学校的距离，\allowbreak{}不是两个家的距离。\allowbreak{}\par};
\end{scope}\end{tikzpicture}\newpage
\definecolor{answer}{rgb}{0.850980,0.047059,0.070588}
\null\begin{tikzpicture}[remember picture,overlay]\begin{scope}[shift={(current page.north west)},x=1bp,y=-1bp]
\node[anchor=north west,inner sep=0pt,text=answer,align=left,text width=484.0bp] at (15,15) {\fontsize{7.5bp}{9.5bp}\selectfont 析：\allowbreak{}\ensuremath{140}～\ensuremath{149}厘米\ensuremath{10}人、\allowbreak{}\ensuremath{150}～\ensuremath{159}厘米\ensuremath{8}人、\allowbreak{}\ensuremath{160}厘米及以上\ensuremath{2}人，\allowbreak{}共\ensuremath{20}人；\allowbreak{}\ensuremath{149}、\allowbreak{}\ensuremath{150}、\allowbreak{}\ensuremath{160}分别归第一、\allowbreak{}第二、\allowbreak{}第三段。\allowbreak{}\par};
\end{scope}\end{tikzpicture}\newpage
\definecolor{answer}{rgb}{0.850980,0.047059,0.070588}
\null\begin{tikzpicture}[remember picture,overlay]\begin{scope}[shift={(current page.north west)},x=1bp,y=-1bp]
\node[anchor=north west,inner sep=0pt,text=answer,align=left,text width=261.3334045410156bp] at (15,15) {\fontsize{10.0bp}{12.5bp}\selectfont 平均分成\ensuremath{10}份,\allowbreak{}涂其中\ensuremath{6}份。\allowbreak{}\par};
\end{scope}\end{tikzpicture}\newpage
\definecolor{answer}{rgb}{0.850980,0.047059,0.070588}
\null\begin{tikzpicture}[remember picture,overlay]\begin{scope}[shift={(current page.north west)},x=1bp,y=-1bp]
\node[anchor=north west,inner sep=0pt,text=answer,align=left,text width=483.0bp] at (15,15) {\fontsize{8.5bp}{10.62bp}\selectfont 单价乘每捆棵数得每捆价，\allowbreak{}再乘捆数；\allowbreak{}也可先求总棵数。\allowbreak{}\par};
\end{scope}\end{tikzpicture}\newpage
\definecolor{answer}{rgb}{0.850980,0.047059,0.070588}
\null\begin{tikzpicture}[remember picture,overlay]\begin{scope}[shift={(current page.north west)},x=1bp,y=-1bp]
\node[anchor=north west,inner sep=0pt,text=answer,align=left,text width=483.0bp] at (15,15) {\fontsize{7.5bp}{9.5bp}\selectfont 析：\allowbreak{}总和\ensuremath{372} ÷ \ensuremath{4} = \ensuremath{93}分；\allowbreak{}也可以\ensuremath{92}为基准，\allowbreak{}多\ensuremath{3}与多\ensuremath{1}均分给\ensuremath{4}次，\allowbreak{}\ensuremath{92} + \ensuremath{4} ÷ \ensuremath{4} = \ensuremath{93}。\allowbreak{}\par};
\end{scope}\end{tikzpicture}\newpage
\definecolor{answer}{rgb}{0.850980,0.047059,0.070588}
\null\begin{tikzpicture}[remember picture,overlay]\begin{scope}[shift={(current page.north west)},x=1bp,y=-1bp]
\node[anchor=north west,inner sep=0pt,text=answer,align=left,text width=483.0bp] at (15,15) {\fontsize{7.5bp}{9.5bp}\selectfont 析：\allowbreak{}六数和\ensuremath{180},\allowbreak{} \ensuremath{180} ÷ \ensuremath{6} = \ensuremath{30};\allowbreak{}或以\ensuremath{30}为基准，\allowbreak{}三个\ensuremath{28}少\ensuremath{6},\allowbreak{} \ensuremath{32}、\allowbreak{}\ensuremath{34}多\ensuremath{6},\allowbreak{}正好抵消。\allowbreak{}旧红\ensuremath{29}应改\ensuremath{30}。\allowbreak{}\par};
\end{scope}\end{tikzpicture}\newpage
\definecolor{answer}{rgb}{0.850980,0.047059,0.070588}
\null\begin{tikzpicture}[remember picture,overlay]\begin{scope}[shift={(current page.north west)},x=1bp,y=-1bp]
\node[anchor=north west,inner sep=0pt,text=answer,align=left,text width=483.0bp] at (15,15) {\fontsize{7.5bp}{9.5bp}\selectfont 析：\allowbreak{}五个连续数中间数就是平均\ensuremath{20},\allowbreak{}最小\ensuremath{20} \ensuremath{-}\ensuremath{2} = \ensuremath{18},\allowbreak{}最大\ensuremath{20} + \ensuremath{2} = \ensuremath{22};\allowbreak{}也可(\ensuremath{20} × \ensuremath{5} \ensuremath{-}\ensuremath{1} \ensuremath{-}\ensuremath{2} \ensuremath{-}\ensuremath{3} \ensuremath{-}\ensuremath{4}) ÷ \ensuremath{5} = \ensuremath{18}。\allowbreak{}\par};
\end{scope}\end{tikzpicture}\newpage
\definecolor{answer}{rgb}{0.850980,0.047059,0.070588}
\null\begin{tikzpicture}[remember picture,overlay]\begin{scope}[shift={(current page.north west)},x=1bp,y=-1bp]
\node[anchor=north west,inner sep=0pt,text=answer,align=left,text width=501.0bp] at (15,15) {\fontsize{7.5bp}{9.5bp}\selectfont 析：\allowbreak{}[(\ensuremath{11} + \ensuremath{4}) × \ensuremath{2} \ensuremath{-}\ensuremath{12}] ÷ \ensuremath{2} = \ensuremath{9}厘米；\allowbreak{}或\ensuremath{11} + \ensuremath{4} \ensuremath{-}\ensuremath{12} ÷ \ensuremath{2} = \ensuremath{9}厘米。\allowbreak{}还须检验\ensuremath{9} + \ensuremath{9} > \ensuremath{12}，\allowbreak{}确能围成三角形。\allowbreak{}\par};
\end{scope}\end{tikzpicture}\newpage
\definecolor{answer}{rgb}{0.850980,0.047059,0.070588}
\null\begin{tikzpicture}[remember picture,overlay]\begin{scope}[shift={(current page.north west)},x=1bp,y=-1bp]
\node[anchor=north west,inner sep=0pt,text=answer,align=left,text width=484.0bp] at (15,15) {\fontsize{7.5bp}{9.5bp}\selectfont 析：\allowbreak{}每盒\ensuremath{630} ÷ \ensuremath{30} = \ensuremath{21}元；\allowbreak{}每支\ensuremath{21} ÷ \ensuremath{7} = \ensuremath{3}元（也可\ensuremath{630} ÷ (\ensuremath{7} × \ensuremath{30}) = \ensuremath{3}元）。\allowbreak{}\par};
\end{scope}\end{tikzpicture}\newpage
\definecolor{answer}{rgb}{0.850980,0.047059,0.070588}
\null\begin{tikzpicture}[remember picture,overlay]\begin{scope}[shift={(current page.north west)},x=1bp,y=-1bp]
\node[anchor=north west,inner sep=0pt,text=answer,align=left,text width=483.9998474121094bp] at (15,15) {\fontsize{7.5bp}{9.5bp}\selectfont 析：\allowbreak{}去掉最高\ensuremath{99}、\allowbreak{}最低\ensuremath{92}，\allowbreak{}以减小极端评分的影响；\allowbreak{}余\ensuremath{5}个评分除以\ensuremath{5}得\ensuremath{96}分；\allowbreak{}也可用\ensuremath{96}为基准补差，\allowbreak{}仍为\ensuremath{96}分。\allowbreak{}\par};
\end{scope}\end{tikzpicture}\newpage
\definecolor{answer}{rgb}{0.850980,0.047059,0.070588}
\null\begin{tikzpicture}[remember picture,overlay]\begin{scope}[shift={(current page.north west)},x=1bp,y=-1bp]
\node[anchor=north west,inner sep=0pt,text=answer,align=left,text width=484.0bp] at (15,15) {\fontsize{8.5bp}{10.62bp}\selectfont 同买\ensuremath{12}千克，\allowbreak{}先求每千克便宜多少，\allowbreak{}再乘\ensuremath{12};\allowbreak{}也可分别求两种鱼的总价后相减。\allowbreak{}\par};
\end{scope}\end{tikzpicture}\newpage
\definecolor{answer}{rgb}{0.850980,0.047059,0.070588}
\null\begin{tikzpicture}[remember picture,overlay]\begin{scope}[shift={(current page.north west)},x=1bp,y=-1bp]
\node[anchor=north west,inner sep=0pt,text=answer,align=left,text width=484.0bp] at (15,15) {\fontsize{8.5bp}{10.62bp}\selectfont 题目给出第\ensuremath{1}根比第\ensuremath{2}根长\ensuremath{1.8}米，\allowbreak{}所以第\ensuremath{2}根为\ensuremath{9.6} \ensuremath{-}\ensuremath{1.8},\allowbreak{}即\ensuremath{7.8}米。\allowbreak{}第\ensuremath{3}根比第\ensuremath{2}根长\ensuremath{3.2}米，\allowbreak{}所求第\ensuremath{3}根为\ensuremath{7.8} + \ensuremath{3.2},\allowbreak{}即\ensuremath{11}米。\allowbreak{}也可先比较第\ensuremath{3}根与第\ensuremath{1}根：\allowbreak{}它比第\ensuremath{1}根长\ensuremath{3.2} \ensuremath{-}\ensuremath{1.8},\allowbreak{}即\ensuremath{1.4}米。\allowbreak{}\par};
\end{scope}\end{tikzpicture}\newpage
\definecolor{answer}{rgb}{0.850980,0.039216,0.070588}
\null\begin{tikzpicture}[remember picture,overlay]\begin{scope}[shift={(current page.north west)},x=1bp,y=-1bp]
\node[anchor=north west,inner sep=0pt,text=answer,align=left,text width=241bp] at (15,15) {\fontsize{7.5bp}{9.5bp}\selectfont 析：\allowbreak{}\ensuremath{B}：\allowbreak{}乘除同级从左到右，\allowbreak{}\ensuremath{75}×\ensuremath{3}÷\ensuremath{75}×\ensuremath{3}=\ensuremath{225}÷\ensuremath{75}×\ensuremath{3}=\ensuremath{3}×\ensuremath{3}=\ensuremath{9}。\allowbreak{}\par};
\end{scope}\end{tikzpicture}\newpage
\definecolor{answer}{rgb}{0.850980,0.047059,0.070588}
\null\begin{tikzpicture}[remember picture,overlay]\begin{scope}[shift={(current page.north west)},x=1bp,y=-1bp]
\node[anchor=north west,inner sep=0pt,text=answer,align=left,text width=254.375bp] at (15,15) {\fontsize{7.5bp}{9.5bp}\selectfont 简算旁注：\allowbreak{}\ensuremath{75}×\ensuremath{3}÷\ensuremath{75}×\ensuremath{3}=(\ensuremath{75}÷\ensuremath{75})×(\ensuremath{3}×\ensuremath{3})=\ensuremath{1}×\ensuremath{9}=\ensuremath{9}。\allowbreak{}\par};
\end{scope}\end{tikzpicture}\newpage
\definecolor{answer}{rgb}{0.850980,0.039216,0.070588}
\null\begin{tikzpicture}[remember picture,overlay]\begin{scope}[shift={(current page.north west)},x=1bp,y=-1bp]
\node[anchor=north west,inner sep=0pt,text=answer,align=left,text width=55bp] at (15,15) {\fontsize{9.96bp}{13bp}\selectfont \ensuremath{5}厘米\par};
\end{scope}\end{tikzpicture}\newpage
\definecolor{answer}{rgb}{0.850980,0.039216,0.070588}
\null\begin{tikzpicture}[remember picture,overlay]\begin{scope}[shift={(current page.north west)},x=1bp,y=-1bp]
\node[anchor=north west,inner sep=0pt,text=answer,align=left,text width=55bp] at (15,15) {\fontsize{9.96bp}{13bp}\selectfont \ensuremath{5}厘米\par};
\end{scope}\end{tikzpicture}\newpage
\definecolor{answer}{rgb}{0.850980,0.039216,0.070588}
\null\begin{tikzpicture}[remember picture,overlay]\begin{scope}[shift={(current page.north west)},x=1bp,y=-1bp]
\node[anchor=north west,inner sep=0pt,text=answer,align=left,text width=55bp] at (15,15) {\fontsize{9.96bp}{13bp}\selectfont \ensuremath{6}厘米\par};
\end{scope}\end{tikzpicture}\newpage
\end{document}